\documentclass[11pt]{article}

\usepackage[margin=1.1in]{geometry}

\usepackage{amssymb,mathtools,natbib}

% Anchored formal environments for causalsmith papers.
% First mandatory argument = obj_id (machine anchor joining the paper to the
% Lean crosswalk; invisible in the rendered PDF). Optional argument = name.
% Each environment also defines \label{obj:<obj_id>} so prose can use
% \cref{obj:T-1}; cleveref derives the printed kind and number from the target.
\usepackage{amsmath}
\usepackage{mathrsfs}
\usepackage{amsthm}
\usepackage{booktabs}
% Long displays may break across pages; let TeX stretch a little harder before an
% overfull \hbox so wide formulas/tables are less likely to run off the page.
\allowdisplaybreaks
\setlength{\emergencystretch}{3em}
\newtheorem{theorem}{Theorem}
\newtheorem{lemma}{Lemma}
\newtheorem{assumption}{Assumption}
\newtheorem{definition}{Definition}
\newtheorem{proposition}{Proposition}
\newtheorem{remark}{Remark}
\newtheorem{citedresult}{Cited result}
% The amsthm counter is `algorithmbox` (not `algorithm`) so a later `\usepackage{algorithm}` cannot clash.
\newcounter{algorithmbox}
\renewcommand{\thealgorithmbox}{\arabic{algorithmbox}}
\NewDocumentEnvironment{theoremv}{m o s}{\begin{theorem}\label{obj:#1}{\normalfont[\texttt{\detokenize{#1}}]\IfValueT{#2}{\ (#2)}\IfBooleanT{#3}{\verificationfootnotemark}\ }}{\end{theorem}}
\NewDocumentEnvironment{lemmav}{m o s}{\begin{lemma}\label{obj:#1}{\normalfont[\texttt{\detokenize{#1}}]\IfValueT{#2}{\ (#2)}\IfBooleanT{#3}{\verificationfootnotemark}\ }}{\end{lemma}}
\NewDocumentEnvironment{assumptionv}{m o s}{\begin{assumption}\label{obj:#1}{\normalfont[\texttt{\detokenize{#1}}]\IfValueT{#2}{\ (#2)}\IfBooleanT{#3}{\verificationfootnotemark}\ }}{\end{assumption}}
\NewDocumentEnvironment{definitionv}{m o s}{\begin{definition}\label{obj:#1}{\normalfont[\texttt{\detokenize{#1}}]\IfValueT{#2}{\ (#2)}\IfBooleanT{#3}{\verificationfootnotemark}\ }}{\end{definition}}
% A `propositionv` is a proved result tiered as a Proposition (theorem-grade, machine-verified).
\NewDocumentEnvironment{propositionv}{m o s}{\begin{proposition}\label{obj:#1}{\normalfont[\texttt{\detokenize{#1}}]\IfValueT{#2}{\ (#2)}\IfBooleanT{#3}{\verificationfootnotemark}\ }}{\end{proposition}}
% A `remarkv` holds NON-load-bearing interpretive / open-question content (e.g. a causalsmith `oeq:`
% open question). It is neither proved nor assumed; no theorem may depend on it.
\NewDocumentEnvironment{remarkv}{m o s}{\begin{remark}\label{obj:#1}{\normalfont[\texttt{\detokenize{#1}}]\IfValueT{#2}{\ (#2)}\IfBooleanT{#3}{\verificationfootnotemark}\ }}{\end{remark}}
% An `algorithmv` is a DEFINITION-kind object presented as a boxed, numbered-step procedure (an
% estimator / selection rule built in stages). Same anchor contract as `definitionv`; its body is
% ordinary LaTeX whose steps are `\begin{enumerate}` items, so tex2html needs no extra machinery.
% Presented in the CONVENTIONAL algorithm style readers expect from `algorithm`+`algorithmic`:
% a rule above, an "Algorithm N (Title)" caption line, a rule under the caption, the numbered
% steps, and a closing rule. It is NOT a float — the anchor contract requires the object to stay
% where the outline places it, and floats drift. The obj-id tag is suppressed here (it is
% bookkeeping, and a caption line carrying it reads as debug output in a finished paper).
\NewDocumentEnvironment{algorithmv}{m o s}%
  {\par\addvspace{\medskipamount}\noindent\rule{\linewidth}{0.8pt}\par\nobreak\vspace{2pt}%
   \refstepcounter{algorithmbox}\label{obj:#1}%
   \noindent\textbf{Algorithm \thealgorithmbox}\IfValueT{#2}{ \textnormal{(#2)}}%
   \IfBooleanT{#3}{\verificationfootnotemark}\par\nobreak\vspace{2pt}%
   \noindent\rule{\linewidth}{0.4pt}\par\nobreak\vspace{2pt}\normalfont}%
  {\par\nobreak\vspace{2pt}\noindent\rule{\linewidth}{0.8pt}\par\addvspace{\medskipamount}}
% A `citedv` is an IMPORTED external result the paper relies on but does NOT prove
% (a causalsmith `gate_class:"cited"` node). It is machine-checked against the cited
% source at F2.5, never proved here; the fixed provenance note makes that explicit
% so the environment can never be read as a contribution of this paper. The \citep
% attribution is supplied by the surrounding prose (P2), keyed to references.bib.
\NewDocumentEnvironment{citedv}{m o s}{\begin{citedresult}\label{obj:#1}{\normalfont[\texttt{\detokenize{#1}}]\IfValueT{#2}{\ (#2)}\IfBooleanT{#3}{\verificationfootnotemark}\ \textit{(imported result; machine-checked against the cited source, not proved here.)}\ }}{\end{citedresult}}
% \leanref{obj_id}{display text}: an inline first-mention link to a from-note object's
% verified Lean code. In the PDF it renders the display text plainly (the Lean link is a
% web-only affordance); tex2html turns it into a drawer-opening span.
\newcommand{\leanref}[2]{#2}
% For a theorem/lemma whose Lean declaration accepts a source-matched published proposition as an
% input, the starred anchored environment puts the mark in its heading; the generated text remains
% outside the frozen mathematical environment. Keep the combined macro for older paper bundles.
\newcommand{\verificationfootnotemark}{\footnotemark}
\newcommand{\verificationfootnotetext}[1]{\footnotetext{#1}}
\newcommand{\verificationfootnote}[1]{\verificationfootnotemark\verificationfootnotetext{#1}}
% xcolor before hyperref (hyperref would otherwise pull in plain `color` for colorlinks, and a
% later xcolor load clashes). The page-1 identity stamp at the end of this file needs a grey.
\usepackage{xcolor}
\usepackage{graphicx}
\usepackage{eso-pic}
% hyperref follows natbib and the environment macros; cleveref must follow hyperref.
\usepackage[colorlinks=true,linkcolor=blue,citecolor=blue,urlcolor=blue]{hyperref}
\usepackage[capitalise,nameinlink,noabbrev]{cleveref}
% Pin the names explicitly: labels are placed inside the underlying amsthm environments, and these
% declarations make the PDF contract independent of cleveref's theorem-name inference. House style:
% reference names always print with a capital first letter ("Theorem 1", "Definition 2"), in any
% sentence position — hence `capitalise` above and capitalized \crefname forms below.
\crefname{theorem}{Theorem}{Theorems}
\Crefname{theorem}{Theorem}{Theorems}
\crefname{lemma}{Lemma}{Lemmas}
\Crefname{lemma}{Lemma}{Lemmas}
\crefname{assumption}{Assumption}{Assumptions}
\Crefname{assumption}{Assumption}{Assumptions}
\crefname{definition}{Definition}{Definitions}
\Crefname{definition}{Definition}{Definitions}
\crefname{proposition}{Proposition}{Propositions}
\Crefname{proposition}{Proposition}{Propositions}
\crefname{remark}{Remark}{Remarks}
\Crefname{remark}{Remark}{Remarks}
\crefname{citedresult}{Cited result}{Cited results}
\Crefname{citedresult}{Cited result}{Cited results}
\crefname{algorithmbox}{Algorithm}{Algorithms}
\Crefname{algorithmbox}{Algorithm}{Algorithms}
% ---------------------------------------------------------------------------
% Page-1 identity stamp (arXiv style) and the pinned title-block date.
%
% Split of responsibility: THIS file owns the FORMAT (placement, font, colour, punctuation),
% the generated `paper_stamp.tex` owns only the VALUES (number+version, area, revised date,
% DOI, link target). P4 refreshes this template on every emit, so a layout fix reaches every
% bundle from a recompile; and because `paper_stamp.tex` is not one of the P2 assembly
% digest's authored inputs, a DOI that arrives after publication needs only that one small
% file rewritten plus a recompile — never a redraft, never a reassembly.
%
% Defaults first, so a bundle WITHOUT a stamp file compiles exactly as it always did:
% `\cspaperdate` falls back to `\today` and no stamp is drawn.
\providecommand*{\cspaperdate}{\today}  % the \date{} target
\providecommand*{\csstampid}{}          % e.g. CSWP-2026-008v3 (empty ⇒ no stamp at all)
\providecommand*{\csstamparea}{}        % e.g. Stat
\providecommand*{\csstampdate}{}        % e.g. 27 Aug 2026
\providecommand*{\csstampdoi}{}         % e.g. 10.5281/zenodo.123 (empty ⇒ DOI part omitted)
\providecommand*{\csstamplink}{}        % href target (DOI url, else the paper's page; may be empty)
\IfFileExists{paper_stamp.tex}{% GENERATED by CausalSmith P4 — paper identity stamp values. Do not hand-edit.
%
% Field order on the stamp (owned by paper_macros.tex):
%   <number>v<version> · doi:<doi> · [<Area>] <date>   — the DOI is never the last token, so a
%   text extractor cannot run it into whatever follows. Without a DOI that part is omitted.
% Format lives in paper_macros.tex; only the values are here, and this file is NOT one of
% the P2 assembly digest's authored inputs. A DOI arriving after publication therefore needs
% this file rewritten plus a `latexmk` recompile — never a redraft or a reassembly.
\renewcommand*{\csstampid}{CSWP-2026-036v3}
\renewcommand*{\csstamparea}{Stat}
\renewcommand*{\csstampdate}{3 Oct 2026}
\renewcommand*{\csstampdoi}{}
\renewcommand*{\csstamplink}{https://causalsmith.org/papers/stat_ldp_optvalue_uniform_frontier_v1}
\renewcommand*{\cspaperdate}{3 October 2026}
}{}
\makeatletter
% Field order is deliberate: the DOI is NEVER the last token on the line. A trailing identifier
% runs into whatever a text extractor emits next, which makes it unsafe to match mechanically
% (the Zenodo stamp matcher reads exactly this line); the date is a safe terminator.
%   <number>v<version> · doi:<doi> · [<Area>] <date>      — with a DOI
%   <number>v<version> · [<Area>] <date>                  — without
\newcommand{\cs@stampsep}{\,\textperiodcentered\,}
\newcommand{\cs@stamptext}{%
  \csstampid
  \ifx\csstampdoi\@empty\else\cs@stampsep doi:\csstampdoi\fi
  \ifx\csstamparea\@empty\else\cs@stampsep[\csstamparea]\fi
  \ifx\csstampdate\@empty\else\quad\csstampdate\fi
}
\ifx\csstampid\@empty\else
  \definecolor{csstampgrey}{gray}{0.45}
  % Background shipout material: it occupies no space, so the text block and the \thanks
  % footnote are byte-identical to an unstamped compile. The starred form applies to the
  % NEXT page shipped only — i.e. page 1. Placed in the left margin (the text block starts
  % at 1.1in), reading bottom-to-top, in a Times-like face at 8pt.
  \AddToShipoutPictureBG*{%
    \AtPageLowerLeft{%
      \put(\LenToUnit{0.42in},\LenToUnit{1.1in}){%
        \rotatebox{90}{%
          \fontfamily{ptm}\fontsize{8}{10}\selectfont
          \ifx\csstamplink\@empty
            \textcolor{csstampgrey}{\cs@stamptext}%
          \else
            \expandafter\href\expandafter{\csstamplink}{\textcolor{csstampgrey}{\cs@stamptext}}%
          \fi
        }%
      }%
    }%
  }
\fi
\makeatother


\title{Optimal treatment welfare under local privacy: Minimax estimation and honest inference}

\author{CausalSmith\thanks{Machine-generated by the CausalSmith research pipeline.}}

\date{\cspaperdate}

\begin{document}

\maketitle

\begin{abstract}
We characterize finite-resource precision for estimating first-best treatment welfare in a randomized binary trial with public uniform stratum probabilities, fair assignment, consistency, and binary potential outcomes whose arm-specific means lie in \([1/4,3/4]\). Independent participants each release one message satisfying pure local differential privacy for the complete observed record. Conditional on the classical Bernstein limit for polynomial approximation of absolute value, we establish matching bounds, up to universal constants, for minimax squared error and worst-law expected length of connected confidence intervals with coverage at least \(0.90\) at every law in this causal class, including exact treatment ties. For participant count \(n\ge2\), stratum count \(d\ge2\), and privacy budget \(0<\varepsilon\le1\), let \(t=n\varepsilon^2\) be the effective private resource and \(L=\log(ed)\) the logarithmic stratum-count factor. Minimax squared error has order \(d^2/(tL)\) when \(t\ge d^2L\) and \(\min\{1,[\log(e+d^2L/t)/L]^2\}\) otherwise; honest length has the square-root order. An explicit noninteractive polynomial procedure attains both orders, and matching converses cover one-message sequential interaction with independent public randomness. The bounds characterize saturation, consistency, and the bounded-dimension boundary for inverse-resource squared error. A signed-experiment correspondence yields the same complete characterization for uniform-reference total-variation distance on a paired family, together with estimation and honest-length lower bounds on the full probability simplex.
\end{abstract}

\section{Introduction}

The population welfare attainable by assigning the better treatment in each stratum is a natural benchmark for a randomized trial. Reporting this benchmark requires estimating the larger of two conditional outcome means in every stratum and then averaging over the population. Treatment ties make this scalar target nonsmooth: a small change in a treatment contrast can change which arm supplies the maximum. When each participant reports through a confidential local channel, precision depends jointly on the number of strata, the participant count, the privacy budget, and the permitted interaction during collection.

This paper characterizes that dependence for a specified categorical trial model. Strata have public uniform probabilities, treatment is assigned fairly conditional on the stratum and both potential outcomes, observed outcomes satisfy consistency, and binary potential-outcome means belong to \([1/4,3/4]\) in every arm and stratum. Participant records are independent and identically distributed. Pure local differential privacy protects the complete observed record, including the realized stratum, assignment, and outcome. Each participant sends one message. Collection may use independent public randomness, participant-specific kernels, and standard-Borel message spaces; sequential collection additionally permits a participant's message rule to depend on earlier releases.

Under these conditions, and conditional on the classical Bernstein absolute-approximation limit, we establish matching finite-resource bounds for two reporting criteria: minimax squared-error risk and minimax worst-law expected length of connected closed intervals with coverage at least \(0.90\) at every law in the causal class. The bounds are uniform over arbitrary admissible treatment contrasts, including exact ties. An explicit noninteractive procedure attains both precision orders simultaneously, while the lower bounds apply to every admissible one-message sequentially interactive protocol.

The welfare target admits a useful decomposition. Write \(d\ge2\) for the stratum count. Let \(V(P)\), the first-best population welfare under a causal law \(P\), be the average of the larger arm mean across strata. Let \(B(P)\), the baseline observed-outcome mean under fair assignment, and \(\tau_j(P)\), the treatment contrast in stratum \(j\), denote its two components. Then \cref{obj:prop:signed-causal-bridge} gives
\[
V(P)=B(P)+\frac{1}{2d}\sum_{j=1}^d|\tau_j(P)|.
\]
The baseline is a bounded-mean target. The average absolute contrast captures the gain from choosing the better arm separately in each stratum and supplies the nonsmooth component governing the dimension-dependent precision bounds.

To state the characterization, write \(n\ge2\) for the participant count and \(0<\varepsilon\le1\) for the pure local privacy budget. Define \(t=n\varepsilon^2\), the effective private resource, and \(L=\log(ed)\), the logarithmic dimension factor. The common squared-error scale is
\[
\rho(t,d)=
\begin{cases}
\dfrac{d^2}{tL},&t\ge d^2L,\\[5pt]
\displaystyle\min\left\{1,
\left[\dfrac{\log(e+d^2L/t)}{L}\right]^2\right\},
&0<t<d^2L.
\end{cases}
\]
For both noninteractive and sequentially interactive collection, minimax squared error is comparable to \(\rho(t,d)\), and globally honest expected interval length is comparable to \(\sqrt{\rho(t,d)}\), with universal comparison constants \cref{obj:thm:uniform-private-value-frontiers}. Thus sequential adaptation improves either minimax criterion by at most a universal multiplicative factor within the declared protocol classes.

The evaluated rate gives precise resource implications. Its saturated range ends at
\[
t=\frac{d^2\log(ed)}{ed-e}.
\]
Along any admissible sequence of resource tuples, estimation error vanishes and honest intervals shrink exactly when effective private resource diverges and
\[
\frac{\log(1+d^2/t)}{\log(ed)}\longrightarrow0.
\]
Along sequences with diverging effective private resource, inverse-resource squared-error order and inverse-square-root-resource honest-length order are equivalent to eventual boundedness of dimension. When dimension diverges and \(t/d^2\) stays in a fixed positive finite interval, squared error instead has order \([\log\log(ed)/\log(ed)]^2\), with the corresponding unsquared order for honest length. The two-stratum result supplies a direct calibration across the entire resource range \cref{obj:thm:two-cell-calibration}.

The construction combines separate private estimates of the baseline and the absolute-contrast component. A sign-vector message law supplies scaled coordinate observations whose means equal the treatment contrasts. Products over distinct participants give unbiased estimates of contrast powers, and elementary-symmetric recursions compute those moments. Resource-dependent polynomial approximation estimates absolute value; in the dense branch, an independent pilot selects polynomial estimation near zero and a signed evaluation mean for larger contrasts \cref{obj:def:frontier-handle}. The analysis accounts explicitly for dependence among coordinates released in the same messages \cref{obj:lem:private-moment-and-dependence}. For the converse, moment-matched coordinate priors create separated welfare values, while higher-order transcript contraction controls their distinguishability after history-dependent privatization \cref{obj:lem:adaptive-moment-contraction,obj:lem:separated-value-mixtures}. The common approximation input is the Bernstein limit reported by \citet[Section~3.1, p.~8]{CaiLow2011Nonsmooth}.

These results complement statistical treatment choice and policy learning, which evaluate treatment rules through welfare and regret \citep{Manski2004,Kitagawa2018,Athey2021}. Here precision is evaluated for the scalar population optimum under the declared randomized-trial model. Optimal-value inference supplies closely related estimand comparisons: \citet{LuedtkeVanDerLaan2016OptimalValue} develop inference accommodating possibly nonunique optimal strategies under their conditions, and \citet[Theorem~3.3]{WhitehouseChenAusternSyrgkanis2026Softmax} establish asymptotic inference through controlled smoothing. Our characterization evaluates finite-resource minimax precision and global connected-interval honesty over the categorical causal class.

Methodologically, the paper implements the polynomial-approximation and moment-matching approach of \citet[Theorem~1 and Sections~3--5]{CaiLow2011Nonsmooth} through private participant messages. It also complements local-private functional estimation for positive power sums \citep[Theorems~2.4--2.7]{ButuceaIssartel2021Functionals} and quadratic density functionals \citep[Theorems~3.2--3.3 and~4.1--4.2]{ButuceaRohdeSteinberger2023Quadratic}. The resulting interaction comparison concerns the average magnitude of signed treatment contrasts, with its kink at equality of the arm means.

A further consequence isolates the statistical component of the welfare problem. The signed record induces a paired law with equal total mass in each pair, and its total-variation distance from the public uniform reference equals one half of the average absolute contrast. On a symmetric causal family, the protocol correspondence preserves both collection classes and the joint seed-transcript law \cref{obj:prop:signed-causal-bridge}. Conditional on the same cited Bernstein absolute-approximation limit, we obtain matching estimation and honest-length bounds on the paired family \cref{obj:thm:paired-uniform-tv-frontier}; embedding that family in the full simplex yields the corresponding lower bounds for arbitrary laws on the same alphabet \cref{obj:thm:full-simplex-tv-converse}.



The paper proceeds from related work to the trial model, privacy requirements, and reporting criteria. It then presents the welfare construction and uniform precision bounds, followed by the total-variation consequences, welfare interpretation, and limitations and future work. The appendices develop approximation and private-moment bounds, sequential contraction and decision lower bounds, a guide to the main proofs and verification scope, and proofs of the main results.

\section{Related work}

\paragraph{Welfare and optimal-value inference.}
The estimand is the population welfare attained by choosing the better treatment in each stratum. It corresponds to the unrestricted binary first-best benchmark in \citet[Section~2, equations~(3)--(7)]{Cerulli2026OptimalPolicy}. Statistical treatment choice and policy learning evaluate data-dependent treatment rules through their welfare and regret \citep{Manski2004,Kitagawa2018,Athey2021}. Here the decision criterion concerns precision in estimating the scalar population optimum. For the \leanref{S-1}{randomized binary trial with public categorical strata} studied here, the causal model combines uniform stratum probabilities, fair assignment, consistency, and interior arm means, with independent participant records \cref{obj:def:causal-model,obj:ass:iid-people}. The welfare decomposition in \cref{obj:prop:signed-causal-bridge} connects this optimum to the baseline outcome mean and the average absolute treatment contrast. Exact equality of the arm means is therefore part of the parameter space over which precision is evaluated.

Optimal-value inference has developed several approaches to the resulting nonsmoothness. \citet{LuedtkeVanDerLaan2016OptimalValue} characterize pathwise differentiability and construct asymptotic inference for a possibly non-unique optimal treatment strategy, including inference at exceptional laws under their stated conditions. Softmax smoothing yields asymptotic inference under controlled behavior of treatment gaps near zero, suitable smoothing sequences, and nuisance-estimation conditions \citep[Theorem~3.3]{WhitehouseChenAusternSyrgkanis2026Softmax}. For optimal welfare gains integrated against a known target density, \citet[Theorem~1, Assumptions~1, 2(b)(c), and~3]{Chen2025} obtain asymptotic normality under smooth regression functions, overlap, a bounded target-to-training density ratio, a regular zero level set, and the specified uniform first-stage convergence rate. \citet[Sections~3.2 and~4]{Feng2026} develop differentiability results for sorting operators and asymptotic inference for allocation values under their level-set regularity and estimation conditions. Adaptive smoothing and efficiency results in \citet[Theorems~2.2 and~3.1, Assumptions~1--8]{Xu2026} accommodate nonregularity under the stated overlap, moment, margin, nuisance-rate, bandwidth, and homoscedasticity conditions. These contributions establish condition-specific asymptotic guarantees, with several explicitly accommodating treatment ties. Our comparison concerns finite-resource precision uniformly over the declared categorical causal class, including exact ties, and connected confidence intervals with coverage at least \(0.90\) at every law \cref{obj:def:risk,obj:def:honest-length}.

\paragraph{Nonsmooth estimation and private moments.}
Polynomial approximation and moment matching provide the closest statistical precedents. \citet[Theorem~1 and Sections~3--5]{CaiLow2011Nonsmooth} combine a composite-testing lower-bound inequality with Gaussian absolute-mean estimation using polynomial approximation, unbiased Hermite-polynomial estimates, and a hybrid construction. For independent unit-variance Gaussian observations with a fixed coordinate bound, their minimax mean-squared error is asymptotic to \(\beta_*^2M^2[\log\log d/\log d]^2\). The same squared log-log over log form appears here when effective private resource \(t=n\varepsilon^2\) is proportional to \(d^2\), although private transcripts replace Gaussian coordinates and the present result is uniform over the full resource domain. In the discrete sampling model, \citet[Theorems~2--4]{JiaoHanWeissman2018L1} obtain worst-reference mean-squared error of order \(S/(n\log n)\) for raw-sample \(L_1\)-distance estimation under their large-alphabet comparability conditions. That setting observes categorical draws directly; the present paired problem estimates a population distance from locally private original-record messages.

The present construction implements this approximation strategy through participant-level private messages. Products over distinct participants furnish unbiased estimates of contrast powers \cref{obj:def:private-moments}. Their use in polynomial estimation requires bounds that account for dependence among coordinates released in the same messages \cref{obj:lem:private-moment-and-dependence,obj:lem:finite-private-polynomial-calibration}. On the lower-bound side, moment-matched priors are compared through the distribution of the entire private transcript. The sequential contraction result controls their distinguishability when each participant's release can depend on preceding messages \cref{obj:lem:adaptive-moment-contraction}. Thus the approximation comparison is attached to two features of the observation model: attainable moments arise from distinct-person private products, and the converse applies after history-dependent privatization.

\paragraph{Confidentiality and sequential interaction.}
Under \leanref{S-2}{pure local differential privacy}, the target functional and the permitted interaction jointly determine statistical precision. For positive power sums with exponent \(\gamma>1\), \citet[Theorems~2.4--2.7]{ButuceaIssartel2021Functionals} give sequential mean-squared error bounded at order \(t^{-((\gamma-1)\wedge1)}\), with matching \(t^{-1}\) converses for \(\gamma\ge2\). For integrated squared densities over their bounded Besov classes, \citet[Theorems~3.2--3.3 and~4.1--4.2]{ButuceaRohdeSteinberger2023Quadratic} obtain noninteractive power \(t^{-8s/(4s+3)}\) and sequential power \(t^{-4s/(2s+1)}\), up to their stated logarithmic factors; their parametric upper regimes begin at \(s>3/4\) and \(s>1/2\), respectively. The functional here averages absolute signed contrasts, whose kink occurs at equality of the treatment means. Conditional on the cited, unformalized Bernstein absolute-approximation limit, the matching noninteractive construction and adaptive sequential converse establish the same minimax rates in both protocol classes for the specified causal and paired models \cref{obj:thm:uniform-private-value-frontiers,obj:thm:paired-uniform-tv-frontier}. The extra work is to control dependence among private coordinate moments and to contract higher-order, moment-matched transcript mixtures uniformly over adaptive histories and resources.

Average treatment effects provide a complementary causal comparison. \citet[Lemma~4 and Theorem~7]{OhnishiAwan2025Randomized} give an order-\(t^{-1}\) known-propensity estimator based on a bounded inverse-probability score, a matching private minimax comparison, and asymptotic intervals. Their average treatment effect is a signed mean contrast; optimal welfare incorporates the average magnitude of the stratum-specific contrasts. The confidentiality condition here protects the complete observed record, including stratum membership, assignment, and outcome \cref{obj:ass:full-record-privacy}.

The interval criterion also connects to honest-confidence theory. \citet{Low1997} studies expected-length limits under uniform coverage, \citet{CaiLow2004AdaptationCI} develops modulus-based expected-length theory for linear functionals, and \citet{ArmstrongKolesar2020HonestCI} constructs bias-aware nonparametric intervals. Here honesty is finite-sample and global over the declared nonlinear causal or paired model, performance is worst-law expected connected length, and the upper interval is obtained from the uniform squared-error construction. The lower bound is interval-specific: connectedness forces a single reported set to span separated target neighborhoods when private transcript mixtures are difficult to distinguish.

Public randomness also matters for private identity testing: \citet{AcharyaCanonneFreitagTyagi2019TestWithoutTrust} and \citet{Acharya2021} establish its benefits for distinguishing a reference distribution from separated alternatives. Our paired-model criterion estimates the numerical distance from the public uniform reference and assesses globally honest connected intervals for that distance \cref{obj:def:tv-private-decision-values}. A further comparison concerns empirical distribution sketches. \citet[Theorems~4.6, 4.10, and~4.13]{CormodeTing2025FederatedShift} provide sketch accuracy guarantees and approximate item-level local privacy for empirical distribution shifts in the federated setting. The present population criteria use pure privacy for each complete participant record and allow public randomness within the admissible sequential protocols.

Conditional on the cited, unformalized Bernstein absolute-approximation limit, the contribution is the joint finite-resource characterization, up to universal constants, of minimax squared error and minimax worst-law expected length of globally honest connected intervals for the specified causal and paired signed models. Noninteractive procedures attain these bounds, while the converses cover every admissible one-message sequentially interactive protocol \cref{obj:thm:uniform-private-value-frontiers,obj:thm:paired-uniform-tv-frontier}. Under the same premise, embedding the paired family in the full probability simplex also yields the corresponding full-simplex estimation and honest-length lower bounds \cref{obj:thm:full-simplex-tv-converse}.


\section{Randomized trials, local privacy, and welfare}

Throughout, \leanref{sym:n}{\ensuremath{n\ge2}} denotes the participant count and \leanref{sym:d}{\ensuremath{d\ge2}} the stratum count. Stratum indices range over \leanref{sym:[d]}{\ensuremath{[d]=\{1,\ldots,d\}}}, and participant indices range from \(1\) to \(n\). Generic positive numerical constants are denoted by \leanref{sym:c}{\ensuremath{c}} and \leanref{sym:C_0}{\ensuremath{C_0}}; their values may differ across comparisons. For a finite vector, the \(\ell_1\) norm is the sum of its absolute coordinates. The notation \(\asymp\) denotes comparison within fixed positive multiplicative constants, with its sequence convention specified in \cref{obj:synth_3}. We first describe the trial population and welfare target, then the private observation scheme and the statistical criteria.

\subsection{Trial population and welfare}

Each participant has a stratum, a treatment assignment, an observed outcome, and two potential outcomes. We first state the public uniform-stratum condition; the complete record space and observed-record marginal are collected explicitly below.

\begin{assumptionv}{ass:uniform-covariates}[Uniform participant strata]
The participant stratum \(X\) is uniformly distributed on \([d]\):
\[
P(X=j)=\frac{1}{d}\qquad(j\in[d]).
\]
\end{assumptionv}

Thus the stratum \leanref{sym:X}{\ensuremath{X}}, assignment \leanref{sym:A}{\ensuremath{A}}, and observed outcome \leanref{sym:Y}{\ensuremath{Y}} form the observed record \leanref{sym:O}{\ensuremath{O}} in the alphabet \leanref{sym:\mathcal O}{\ensuremath{\mathcal O}}. The potential outcomes \leanref{sym:Y^0}{\ensuremath{Y^0}} and \leanref{sym:Y^1}{\ensuremath{Y^1}} enter the full-data law, while \leanref{sym:Y^{\mathrm{pot}}}{\ensuremath{Y^{\mathrm{pot}}_A}} denotes the outcome selected by assignment. The observed-record distribution is \leanref{sym:P_O}{\ensuremath{P_O}}, as defined in \cref{obj:synth_1}.

For laws assigning positive probability to every stratum, define the arm-specific conditional mean \leanref{sym:\mu}{\ensuremath{\mu_{aj}(P)}} and treatment contrast \leanref{sym:\tau}{\ensuremath{\tau_j(P)}} by
\[
\mu_{aj}(P)=\mathbb E_P[Y^a\mid X=j],
\qquad
\tau_j(P)=\mu_{1j}(P)-\mu_{0j}(P),
\qquad
\tau(P)=(\tau_1(P),\ldots,\tau_d(P)).
\]
Here the arm index \leanref{sym:a}{\ensuremath{a\in\{0,1\}}} distinguishes control from treatment, and \leanref{sym:j}{\ensuremath{j\in[d]}} indexes the stratum. A cellwise treatment rule selects an arm in each stratum. Under uniform strata, its first-best population welfare is obtained by choosing the arm with the larger conditional mean. We write the first-best value \leanref{sym:V}{\ensuremath{V(P)}} and the observed-outcome baseline \leanref{sym:B}{\ensuremath{B(P)}} as
\[
V(P)=\frac1d\sum_{j=1}^d\max\{\mu_{0j}(P),\mu_{1j}(P)\},
\qquad
B(P)=\mathbb E_P[Y].
\]
This target evaluates the welfare attainable by an optimal stratum-specific treatment choice.

The remaining single-participant restrictions specify randomized assignment, the relation between observed and potential outcomes, and the admissible conditional means.

\begin{assumptionv}{ass:fair-randomization}[Fair conditional treatment assignment]
The binary treatment assignment \(A\), conditional on the participant stratum \(X\) and potential outcomes \(Y^0,Y^1\), satisfies
\[
\mathcal L_P(A\mid X,Y^0,Y^1)=\operatorname{Bernoulli}(1/2).
\]
\end{assumptionv}

Uniform categorical sampling fixes the stratum probabilities at a public reference distribution, specializing the multinomial sampling framework used for distributional functionals \citep{JiaoHanWeissman2018L1}. Equal population weights make each cell's contribution to welfare explicit, while the realized participant stratum remains part of the private record.

\begin{assumptionv}{ass:consistency}[Observed outcome consistency]
The observed outcome \(Y\) equals \(Y^{\mathrm{pot}}_A\), the potential outcome selected by the realized assignment, almost surely:
\[
P(Y=Y^{\mathrm{pot}}_A)=1.
\]
\end{assumptionv}

Known randomized treatment assignment is maintained here in its fair-assignment form \citep{OhnishiAwan2025Randomized}. Conditioning on both potential outcomes ensures that assignment supplies experimental variation within each stratum and gives the two arms equal weight in the baseline mean.

\begin{assumptionv}{ass:interior-means}[Interior potential-outcome means]
The conditional potential-outcome means \(\mu_{aj}(P)\) satisfy
\[
\mu_{aj}(P)\in[1/4,3/4]
\qquad(a\in\{0,1\},\ j\in[d]).
\]
\end{assumptionv}

Potential-outcome consistency connects the recorded response to the outcome under the assigned arm, giving the observed data their treatment-welfare interpretation \citep{LuedtkeVanDerLaan2016OptimalValue}. Together with fair assignment, it identifies each conditional arm mean from the corresponding observed treatment group.

\begin{assumptionv}{ass:iid-people}[Independent identically distributed participants]
The vector of observed participant records \(\mathbf O\) has the product law
\[
\mathcal L_P(\mathbf O)=P_O^{\otimes n},
\]
where \(P_O\) is the observed-record marginal law induced by \(P\).
\end{assumptionv}

Fixed interior Bernoulli means strengthen the bounded-outcome framework for optimal treatment value \citep{LuedtkeVanDerLaan2016OptimalValue}. The maintained interval keeps every arm mean away from the outcome boundaries, places treatment contrasts in \([-1/2,1/2]\), and places first-best welfare in \([1/4,3/4]\).

We now define the complete record space before collecting the four single-participant restrictions into the causal class.

\begin{definitionv}{synth_1}[Full-data causal law space]
For an integer \(d\ge2\), write \([d]=\{1,\ldots,d\}\). Define \(\mathcal P_d\) as the class of all probability laws on the finite record space
\[
[d]\times\{0,1\}^4,
\]
equipped with its discrete sigma-algebra. The coordinate variables are ordered as \((X,A,Y,Y^0,Y^1)\): \(X\in[d]\) is the covariate stratum, \(A\in\{0,1\}\) is the treatment assignment, \(Y\in\{0,1\}\) is the observed outcome, and \(Y^0,Y^1\in\{0,1\}\) are the potential outcomes under control and treatment, respectively. Write \(Y^{\mathrm{pot}}_0=Y^0\) and \(Y^{\mathrm{pot}}_1=Y^1\). For \(P\in\mathcal P_d\), the observed record is \(O=(X,A,Y)\in\mathcal O=[d]\times\{0,1\}^2\), and \(P_O\) denotes its marginal law under \(P\).
\end{definitionv}

The causal class \leanref{sym:\mathcal V_d}{\ensuremath{\mathcal V_d}} in \cref{obj:def:causal-model} leaves the dependence between the two potential outcomes unrestricted subject to these conditions. Under this class, fair assignment and consistency give
\[
B(P)=\frac1{2d}\sum_{j=1}^d
\bigl(\mu_{0j}(P)+\mu_{1j}(P)\bigr).
\]
For a contrast vector, let the average absolute contrast \leanref{sym:F}{\ensuremath{F(\theta)}} be defined by
\[
F(\theta)=\frac1d\sum_{j=1}^d|\theta_j|.
\]
The identity \(\max\{u,v\}=(u+v+|u-v|)/2\) then yields
\[
V(P)=B(P)+\frac12F(\tau(P)).
\]
The baseline records welfare under fair assignment; the absolute-contrast component records the gain from selecting the better arm separately in each stratum.

The following definition collects the four single-participant restrictions. Sampling is specified separately through the iid condition above, with observed sample \leanref{sym:\mathbf O}{\ensuremath{\mathbf O=(O_1,\ldots,O_n)}}.

\begin{definitionv}{def:causal-model}[Causal law class $\mathcal V_d$]
\(\mathcal V_d\) is the class of full-data probability laws \(P\in\mathcal P_d\) satisfying all four conditions:
\begin{itemize}
\item \textbf{(Uniform strata.)} \Cref{obj:ass:uniform-covariates}.
\item \textbf{(Fair assignment.)} \Cref{obj:ass:fair-randomization}.
\item \textbf{(Consistency.)} \Cref{obj:ass:consistency}.
\item \textbf{(Interior means.)} \Cref{obj:ass:interior-means}.
\end{itemize}
\end{definitionv}

Independent identically distributed one-record participants provide the standard sampling structure for local-private minimax analysis \citep{DuchiJordanWainwright2018Minimax}. Every participant contributes one draw from the same trial population, and the collection protocol acts on these separate records.

\subsection{Private collection and statistical criteria}

Fix the privacy budget \leanref{sym:\varepsilon}{\ensuremath{0<\varepsilon\le1}}. Collection may use a shared public seed \leanref{sym:R}{\ensuremath{R}} taking values in a standard-Borel seed space \leanref{sym:\mathcal R}{\ensuremath{\mathcal R}}, with probability law \leanref{sym:\lambda}{\ensuremath{\lambda}}. Participant \leanref{sym:i}{\ensuremath{i}} sends one message in a standard-Borel message space \leanref{sym:\mathcal Z}{\ensuremath{\mathcal Z_i}}. The analyst also has a uniform randomizer \leanref{sym:U}{\ensuremath{U}} for randomized statistical decisions. We first state the independence condition on public and analyst randomness.

\begin{assumptionv}{ass:independent-randomness}[Independent public and analyst randomness]
The shared public seed \(R\), with probability law \(\lambda\), and the analyst randomizer \(U\) have conditional joint law
\[
\mathcal L(R,U\mid\mathbf O)
=\lambda\otimes\operatorname{Uniform}[0,1],
\]
where \(\mathbf O\) is the vector of observed participant records.
\end{assumptionv}

The history space \leanref{sym:\mathcal H}{\ensuremath{\mathcal H_i}} in \cref{obj:synth_2} contains the messages preceding participant \(i\). A history \leanref{sym:\eta}{\ensuremath{\eta}} and a seed realization \leanref{sym:r}{\ensuremath{r}} therefore index the participant's conditional message distribution \leanref{sym:Q}{\ensuremath{Q_i(\cdot\mid o,\eta,r)}}, where \leanref{sym:o}{\ensuremath{o}} denotes an observed-record input. Sequential collection generates
\[
Z_i\mid \mathbf O,R,Z_1,\ldots,Z_{i-1}
\sim Q_i(\cdot\mid O_i,(Z_1,\ldots,Z_{i-1}),R).
\]
The private transcript is the message vector \leanref{sym:\mathbf Z}{\ensuremath{\mathbf Z=(Z_1,\ldots,Z_n)}}; the analyst observes this vector and the public seed. The next condition imposes privacy on every history-indexed message kernel.

\begin{assumptionv}{ass:full-record-privacy}[Full-record local privacy]
For every participant index \(i\), observed records \(o,o'\), history \(\eta\), public seed \(r\), and measurable message set \(E\) in their declared domains,
\[
Q_i(E\mid o,\eta,r)\le e^\varepsilon Q_i(E\mid o',\eta,r).
\]
\end{assumptionv}

Independent public coins and independent randomized decision rules allow coordination and randomized analysis within the observation scheme \citep{AcharyaCanonneFreitagTyagi2019TestWithoutTrust}. The product conditional law makes their joint distribution the same for every observed sample.

Privacy compares the message distributions induced by any two complete observed records, including their strata, assignments, and outcomes. Let \leanref{sym:o'}{\ensuremath{o'}} denote the second record and \leanref{sym:E}{\ensuremath{E}} a measurable message event. The following definition makes the history domain used in that inequality explicit.

\begin{definitionv}{synth_2}[Participant message histories]
For a protocol with \(n\) participants, let \((\mathcal Z_i)_{i=1}^n\) be its standard-Borel participant message spaces. For each \(i\in\{1,\ldots,n\}\), define the history space
\[
\mathcal H_i=\prod_{\ell=1}^{i-1}\mathcal Z_\ell
\]
with the product sigma-algebra; the empty product \(\mathcal H_1\) is a singleton. A history \(\eta\in\mathcal H_i\) records the messages of participants \(1,\ldots,i-1\). Thus, with public-seed space \(\mathcal R\), the participant kernel \(Q_i\) is a Markov kernel from \(\mathcal O\times\mathcal H_i\times\mathcal R\) to \(\mathcal Z_i\), and the realized history preceding participant \(i\) is \((Z_1,\ldots,Z_{i-1})\).
\end{definitionv}

Pure sequential local differential privacy controls the change in each participant's message distribution when the complete input record changes \citep{DuchiJordanWainwright2018Minimax}. Requiring the bound at every seed and history applies the same protection to every conditional collection rule, including histories that arise under other population laws.

The sequential class permits later collection rules to depend on earlier messages. Its noninteractive subclass imposes a further condition, with \leanref{sym:\eta\prime}{\ensuremath{\eta'}} denoting a second history.

\begin{assumptionv}{ass:noninteraction}[Independence from message history]
For every participant index \(i\), observed record \(o\), public seed \(r\), and histories \(\eta,\eta'\in\mathcal H_i\),
\[
Q_i(\cdot\mid o,\eta,r)=Q_i(\cdot\mid o,\eta',r).
\]
\end{assumptionv}

Noninteractive channels use rules fixed independently of the preceding messages, while permitting participant-specific rules and shared public coins \citep{AcharyaCanonneFreitagTyagi2019TestWithoutTrust}. Thus coordination through the independent seed remains available within this subclass.

We collect these requirements into two protocol classes, both of which give each participant one message.

\begin{definitionv}{def:sequential-class}[Sequential protocol class $\mathfrak Q_{\mathrm{SI}}(n,d,\varepsilon)$]
\(\mathfrak Q_{\mathrm{SI}}(n,d,\varepsilon)\) is the class of sequences
\[
Q=(Q_i)_{i=1}^n
\]
of the declared Markov kernels satisfying \cref{obj:ass:full-record-privacy}.
\end{definitionv}

The one-message sequential class \leanref{sym:\mathfrak Q_{\mathrm{SI}}}{\ensuremath{\mathfrak Q_{\mathrm{SI}}(n,d,\varepsilon)}} in \cref{obj:def:sequential-class} allows the conditional message distribution to respond to earlier releases. The next definition restricts this dependence.

\begin{definitionv}{def:noninteractive-class}[Noninteractive protocol class $\mathfrak Q_{\mathrm{NI}}(n,d,\varepsilon)$]
\(\mathfrak Q_{\mathrm{NI}}(n,d,\varepsilon)\) is the subclass of \(\mathfrak Q_{\mathrm{SI}}(n,d,\varepsilon)\) satisfying both \cref{obj:ass:full-record-privacy,obj:ass:noninteraction}.
\end{definitionv}

The noninteractive class \leanref{sym:\mathfrak Q_{\mathrm{NI}}}{\ensuremath{\mathfrak Q_{\mathrm{NI}}(n,d,\varepsilon)}} in \cref{obj:def:noninteractive-class} is contained in the sequential class. Let the class index \leanref{sym:\mathcal C}{\ensuremath{\mathcal C\in\{\mathrm{NI},\mathrm{SI}\}}} select the admissible protocol class \leanref{sym:\mathfrak Q}{\ensuremath{\mathfrak Q_{\mathcal C}(n,d,\varepsilon)}}. A welfare estimate \leanref{sym:T}{\ensuremath{T(\mathbf Z,R,U)}} is a Borel real-valued function of the messages, public seed, and analyst randomizer. Expectations and probabilities below include participant sampling, message generation, and the independent randomizers. Our precision criterion optimizes over both the collection protocol and the estimate.

\begin{definitionv}{def:risk}[Minimax welfare risk $\mathfrak R_{\mathcal C}(n,d,\varepsilon)$]
\[
\mathfrak R_{\mathcal C}(n,d,\varepsilon)
=
\inf_{Q\in\mathfrak Q_{\mathcal C}(n,d,\varepsilon)}
\inf_T
\sup_{P\in\mathcal V_d}
\mathbb E_{P,Q}
\left[
\bigl(T(\mathbf Z,R,U)-V(P)\bigr)^2
\right].
\]
Here \(\mathcal C\in\{\mathrm{NI},\mathrm{SI}\}\) selects the corresponding declared protocol class \(\mathfrak Q_{\mathcal C}\), and the second infimum ranges over Borel real-valued welfare estimates \(T(\mathbf Z,R,U)\) based on the private message vector \(\mathbf Z\), the public seed \(R\), and the independent uniform analyst randomizer \(U\). The target \(V(P)\) is the average of the larger arm mean in each stratum under \(P\).
\end{definitionv}

The minimax squared-error risk \leanref{sym:\mathfrak R}{\ensuremath{\mathfrak R_{\mathcal C}(n,d,\varepsilon)}} in \cref{obj:def:risk} measures the best worst-law precision available within the selected protocol class. Its population supremum permits arbitrary treatment-effect patterns satisfying the causal restrictions.

For inference, a welfare interval \leanref{sym:I}{\ensuremath{I(\mathbf Z,R,U)}} is connected and closed, with ordered Borel endpoints in the welfare range \([1/4,3/4]\). Its length is the upper endpoint minus the lower endpoint. Global honesty requires coverage of at least \(0.90\) at every law in the causal class; the corresponding criterion minimizes worst-law expected length.

\begin{definitionv}{def:honest-length}[Honest interval length $\mathfrak H_{\mathcal C}(n,d,\varepsilon)$]
The minimax expected length of a globally honest welfare interval is
\[
\mathfrak H_{\mathcal C}(n,d,\varepsilon)
=
\inf_{Q\in\mathfrak Q_{\mathcal C}(n,d,\varepsilon)}
\ \inf_{\substack{I:\ \inf_{P\in\mathcal V_d}
\Pr_{P,Q}\{V(P)\in I(\mathbf Z,R,U)\}\ge0.90}}
\ \sup_{P\in\mathcal V_d}
\mathbb E_{P,Q}\!\left[
\operatorname{length} I(\mathbf Z,R,U)
\right].
\]
The second infimum ranges over connected closed intervals whose ordered endpoints are Borel functions of the private messages \(\mathbf Z\), public seed \(R\), and analyst randomizer \(U\), and take values in \([1/4,3/4]\).
\end{definitionv}

The honest-length criterion \leanref{sym:\mathfrak H}{\ensuremath{\mathfrak H_{\mathcal C}(n,d,\varepsilon)}} in \cref{obj:def:honest-length} couples interval precision to uniform coverage over the same population class used for squared-error risk. Connectedness makes the reported uncertainty a single welfare range, and the endpoint restriction respects the range implied by the interior means.

\subsection{The signed experiment}

The welfare decomposition separates a population mean from an absolute-value functional of treatment contrasts. To connect the latter to a finite-alphabet experiment, define the signed treatment-outcome statistic \leanref{sym:S}{\ensuremath{S}} by
\[
S=(2A-1)(2Y-1).
\]
The signed input consists of the stratum and this sign, taking values in the paired alphabet \leanref{sym:\mathcal A_d}{\ensuremath{\mathcal A_d=[d]\times\{-1,1\}}}. For a contrast parameter \leanref{sym:\theta}{\ensuremath{\theta\in[-1/2,1/2]^d}}, define the paired law \leanref{sym:P_\theta}{\ensuremath{P_\theta}} and the uniform reference law \leanref{sym:\mathsf U_{2d}}{\ensuremath{\mathsf U_{2d}}} by
\[
P_\theta(j,s)=\frac{1+s\theta_j}{2d},
\qquad
\mathsf U_{2d}(j,s)=\frac1{2d}.
\]
Here \leanref{sym:s}{\ensuremath{s\in\{-1,1\}}} is the sign coordinate. Total variation on this finite alphabet is one half the sum of the absolute differences in probability masses, so
\[
\operatorname{TV}(P_\theta,\mathsf U_{2d})
=\frac1{2d}\sum_{j=1}^d|\theta_j|
=\frac12F(\theta).
\]
These definitions connect the absolute-contrast welfare component to distance from a public reference distribution.

A symmetric causal family embeds every allowed contrast vector into the trial model while keeping the baseline fixed. Its construction specifies the dependence of the potential outcomes as well as their means.

\begin{definitionv}{synth_3}[Private resource rates and their regimes]
For \(n,d\in\{2,3,\ldots\}\) and \(0<\varepsilon\le1\), put \(t=n\varepsilon^2\) and \(L=\log(ed)\). Define
\[
\rho(t,d)=
\begin{cases}
\dfrac{d^2}{tL},&t\ge d^2L,\\[4pt]
\displaystyle\min\left\{1,
\left[\dfrac{\log(e+d^2L/t)}{L}\right]^2\right\},
&0<t<d^2L,
\end{cases}
\qquad
r_{\mathcal C}(n,d,\varepsilon)=\rho(t,d),
\qquad
h_{\mathcal C}(n,d,\varepsilon)=\sqrt{\rho(t,d)}
\]
for each \(\mathcal C\in\{\mathrm{NI},\mathrm{SI}\}\).

For positive numerical sequences \(f_\ell,g_\ell\), write \(f_\ell\asymp g_\ell\) when there exist constants \(0<c\le C_0<\infty\) and an integer \(\ell_0\) such that
\[
c g_\ell\le f_\ell\le C_0g_\ell
\qquad(\ell\ge\ell_0).
\]
We say that these rates satisfy the resource conclusions when all the following properties hold.

For every integer \(d\ge2\) and every \(0<t<d^2\log(ed)\),
\[
\rho(t,d)=1
\quad\Longleftrightarrow\quad
t\le\frac{d^2\log(ed)}{ed-e}.
\]
Thus the saturated regime ends at \(d^2\log(ed)/(ed-e)\), the intermediate regime occupies
\[
\frac{d^2\log(ed)}{ed-e}<t<d^2\log(ed),
\]
and the dense regime starts at \(t=d^2\log(ed)\).

For every sequence \((n_\ell,d_\ell,\varepsilon_\ell)_{\ell\ge1}\) satisfying
\[
n_\ell,d_\ell\in\{2,3,\ldots\},
\qquad 0<\varepsilon_\ell\le1,
\]
put \(t_\ell=n_\ell\varepsilon_\ell^2\) and \(L_\ell=\log(ed_\ell)\). Then
\[
\rho(t_\ell,d_\ell)\longrightarrow0
\quad\Longleftrightarrow\quad
t_\ell\longrightarrow\infty
\ \text{and}\
\frac{\log(1+d_\ell^2/t_\ell)}{L_\ell}\longrightarrow0,
\]
and
\[
\rho(t_\ell,d_\ell)\asymp t_\ell^{-1}
\quad\Longleftrightarrow\quad
\left(\exists a>0,\ \exists\ell_0,\ \forall\ell\ge\ell_0:
t_\ell\ge a\right)
\ \text{and}\
\left(\exists M<\infty,\ \exists\ell_1,\ \forall\ell\ge\ell_1:
\min\{t_\ell,d_\ell\}\le M\right).
\]
For every such sequence with \(t_\ell\to\infty\),
\[
\rho(t_\ell,d_\ell)\asymp t_\ell^{-1}
\quad\Longleftrightarrow\quad
\exists M\in\mathbb N,\ \exists\ell_0,\ \forall\ell\ge\ell_0:
d_\ell\le M.
\]

Finally, for every \(0<a\le b<\infty\), there exist constants \(0<c\le C_0<\infty\) and an integer \(D_0\) such that, for every integer \(d\ge\max\{2,D_0\}\) and every \(t>0\) with \(a\le t/d^2\le b\),
\[
c\left[\frac{\log\log(ed)}{\log(ed)}\right]^2
\le \rho(t,d)\le
C_0\left[\frac{\log\log(ed)}{\log(ed)}\right]^2.
\]
The constants and dimension threshold in this last comparison may depend on \(a,b\).
\end{definitionv}

The symmetric law \leanref{sym:G_\theta}{\ensuremath{G_\theta}} in \cref{obj:def:symmetric-law} pairs arm means equally around \(1/2\). The following result states its membership in the causal class and establishes an exact correspondence between private signed observations and private trial records. The correspondence preserves both collection classes and the joint law of the public seed and messages.

\begin{definitionv}{def:symmetric-law}[Symmetric causal law $G_\theta$]
The full-data law \(G_\theta\) is generated by
\[
X\sim\operatorname{Uniform}[d],
\qquad
Y^0\mid X=j\sim\operatorname{Bernoulli}\!\left(\frac{1-\theta_j}{2}\right),
\qquad
Y^1\mid X=j\sim\operatorname{Bernoulli}\!\left(\frac{1+\theta_j}{2}\right),
\]
with \(Y^0\) and \(Y^1\) conditionally independent given \(X\). Independently of \((X,Y^0,Y^1)\), generate the treatment assignment and set the observed outcome by
\[
A\sim\operatorname{Bernoulli}(1/2),
\qquad
Y=Y^{\mathrm{pot}}_A.
\]
\end{definitionv}

\Cref{obj:prop:signed-causal-bridge} identifies the absolute-contrast component with total variation from the uniform reference for every law in the causal class. On the symmetric family, the baseline equals \(1/2\), so variation in welfare is precisely variation in this distance. The assignment bit is fair conditional on the signed input, and the observed outcome is determined by that input and the assignment. Averaging over the assignment therefore reproduces the private transcript of any admissible trial-record protocol on this family. Conversely, a signed-input protocol can be applied to trial records by computing the sign locally. These transformations preserve full-record privacy and history dependence, allowing estimation procedures and decision lower bounds for the signed experiment to inform welfare analysis in both collection classes.

The absolute value also explains the statistical focus on small treatment contrasts. Selecting the better arm contributes \(|\tau_j(P)|/2\) beyond the cell's fair-assignment mean, and this contribution changes slope at a zero contrast. The signed experiment retains the entire contrast vector in its input probabilities, while the target aggregates coordinate magnitudes. Thus the welfare problem combines estimation of the linear baseline with estimation of this nonlinear aggregate.

\subsection{Resource scales and paired procedures}

To state the finite-resource comparisons uniformly across sample sizes, dimensions, and privacy budgets, we use the effective resource \(t=n\varepsilon^2\) and the logarithmic dimension factor \(L=\log(ed)\). The following definition records the sequence conventions used later for the paired minimax values; the paired model and decision criteria are introduced immediately afterward.

\begin{definitionv}{synth_6}[Sequence conclusions for paired minimax values]
For each \(\mathcal C\in\{\mathrm{NI},\mathrm{SI}\}\), let
\[
\mathcal R_{\mathcal C}(n,d,\varepsilon)
=\mathfrak R^{\mathrm{TV}}_{\mathcal C}
(\mathcal P_d^{\mathrm{pair}};n,d,\varepsilon),
\qquad
\mathcal H_{\mathcal C}(n,d,\varepsilon)
=\mathfrak H^{\mathrm{TV}}_{\mathcal C}
(\mathcal P_d^{\mathrm{pair}};n,d,\varepsilon),
\]
where the minimax risk and globally \(0.90\)-honest connected-interval length are defined in \cref{obj:def:tv-private-decision-values}. We say that the sequence-quantified value conclusions hold when the following properties hold for each class.

For every sequence \((n_\ell,d_\ell,\varepsilon_\ell)_{\ell\ge1}\) with
\[
n_\ell,d_\ell\in\{2,3,\ldots\},
\qquad 0<\varepsilon_\ell\le1,
\]
write
\[
t_\ell=n_\ell\varepsilon_\ell^2,\qquad
L_\ell=\log(ed_\ell),\qquad
R_\ell=\mathcal R_{\mathcal C}(n_\ell,d_\ell,\varepsilon_\ell),
\qquad
H_\ell=\mathcal H_{\mathcal C}(n_\ell,d_\ell,\varepsilon_\ell).
\]
Then
\[
R_\ell\longrightarrow0
\quad\Longleftrightarrow\quad
t_\ell\longrightarrow\infty
\ \text{and}\
\frac{\log(1+d_\ell^2/t_\ell)}{L_\ell}\longrightarrow0,
\]
and, separately,
\[
H_\ell\longrightarrow0
\quad\Longleftrightarrow\quad
t_\ell\longrightarrow\infty
\ \text{and}\
\frac{\log(1+d_\ell^2/t_\ell)}{L_\ell}\longrightarrow0.
\]

Define \(R_\ell\asymp t_\ell^{-1}\) to mean that there exist \(0<c\le C_0<\infty\) and \(\ell_0\) such that
\[
c/t_\ell\le R_\ell\le C_0/t_\ell
\qquad(\ell\ge\ell_0).
\]
For every allowed sequence,
\[
R_\ell\asymp t_\ell^{-1}
\quad\Longleftrightarrow\quad
\left(\exists a>0,\ \exists\ell_0,\ \forall\ell\ge\ell_0:
t_\ell\ge a\right)
\ \text{and}\
\left(\exists M<\infty,\ \exists\ell_1,\ \forall\ell\ge\ell_1:
\min\{t_\ell,d_\ell\}\le M\right).
\]
For every allowed sequence satisfying \(t_\ell\to\infty\),
\[
R_\ell\asymp t_\ell^{-1}
\quad\Longleftrightarrow\quad
\exists M\in\mathbb N,\ \exists\ell_0,\ \forall\ell\ge\ell_0:
d_\ell\le M.
\]

Finally, for every \(0<a\le b<\infty\), there exist \(0<c\le C_0<\infty\) and an integer \(D_0\) such that, for every allowed tuple \((n,d,\varepsilon)\) with \(d\ge D_0\) and
\[
a\le\frac{n\varepsilon^2}{d^2}\le b,
\]
both comparisons hold:
\[
c\left[\frac{\log\log(ed)}{\log(ed)}\right]^2
\le\mathcal R_{\mathcal C}(n,d,\varepsilon)
\le C_0\left[\frac{\log\log(ed)}{\log(ed)}\right]^2,
\]
\[
\sqrt c\left|\frac{\log\log(ed)}{\log(ed)}\right|
\le\mathcal H_{\mathcal C}(n,d,\varepsilon)
\le\sqrt{C_0}\left|\frac{\log\log(ed)}{\log(ed)}\right|.
\]
These constants and the dimension threshold may depend on \(a,b\) and the fixed class, and are uniform over the displayed resource domain.
\end{definitionv}

The squared-error scale \leanref{sym:r_{\mathcal C}}{\ensuremath{r_{\mathcal C}(n,d,\varepsilon)}} and length scale \leanref{sym:h_{\mathcal C}}{\ensuremath{h_{\mathcal C}(n,d,\varepsilon)}} in \cref{obj:synth_3} use the same resource function for both collection classes. The defined resource conclusions distinguish a saturated range, an intermediate range, and a dense range, and express consistency and inverse-resource comparisons along arbitrary admissible sequences. They provide the notation for the statistical comparisons made in \cref{obj:thm:uniform-private-value-frontiers,obj:thm:paired-uniform-tv-frontier}.

We also specify the attainment guarantees for the paired experiment. Its input family \leanref{sym:\mathcal P_d^{\mathrm{pair}}}{\ensuremath{\mathcal P_d^{\mathrm{pair}}}} consists of the paired signed laws, and its target \leanref{sym:D_{\mathrm{TV}}}{\ensuremath{D_{\mathrm{TV}}(P_\theta)}} is their total-variation distance from the uniform reference, as defined in \cref{obj:def:tv-corollary-models}. The paired-input protocol class \leanref{sym:\mathfrak K}{\ensuremath{\mathfrak K_{\mathcal C}(n,d,\varepsilon)}} uses the same public randomness, message histories, and one-message privacy requirement on the paired alphabet; its full definition is in \cref{obj:def:tv-private-decision-values}. A paired protocol \leanref{sym:\mathsf K}{\ensuremath{\mathsf K}} generates the private message vector \leanref{sym:\mathbf Z^{\mathsf K}}{\ensuremath{\mathbf Z^{\mathsf K}}}. The estimate \leanref{sym:\widehat D}{\ensuremath{\widehat D}} and interval \leanref{sym:J_{\mathrm{TV}}}{\ensuremath{J_{\mathrm{TV}}}} target the numerical distance. The next definition packages the finite-message attainment properties that the subsequent results establish.

\begin{definitionv}{synth_5}[Resource-selected paired polynomial procedure]
Fix \(n,d\in\{2,3,\ldots\}\), \(0<\varepsilon\le1\), and \(C_0>0\). Put
\[
t=n\varepsilon^2,\qquad L=\log(ed),\qquad
\delta=\tanh(\varepsilon/2),\qquad b=d/\delta,
\qquad m=\lfloor n/3\rfloor.
\]
Use a singleton public seed. For participants \(i=1,\ldots,2m\), release a sign vector \(Z_i\in\{-1,1\}^d\) with conditional probabilities
\[
\mathsf K_i(z\mid(j,s))
=2^{-d}(1+\delta s z_j).
\]
The remaining participants release the constant vector \((-1,\ldots,-1)\). All releases use independent local randomness and are independent of histories.

If \(n=2\), set \(\widehat D=1/8\). If \(n>2\), use the first \(m\) active participants as the pilot block and the next \(m\) as the evaluation block. Define
\[
W^{\mathrm{pil}}_{ij}=bZ_{i,j},
\qquad
W^{\mathrm{ev}}_{ij}=bZ_{m+i,j},
\qquad 1\le i\le m,\quad j\in[d],
\]
and
\[
\sigma^2=\frac{b^2}{m},\qquad
\sigma=\sqrt{\sigma^2},\qquad
H=\log(e+\sigma^2L),\qquad
M_j=\frac1m\sum_{i=1}^mW^{\mathrm{pil}}_{ij},
\qquad
\overline W_j=\frac1m\sum_{i=1}^mW^{\mathrm{ev}}_{ij}.
\]
For the evaluation block, define the distinct-person moments
\[
\mathsf U_{0j}=1,\qquad
\mathsf U_{vj}
=\binom mv^{-1}
\sum_{\substack{J\subseteq\{1,\ldots,m\}\\ |J|=v}}
\prod_{i\in J}W^{\mathrm{ev}}_{ij},
\qquad 1\le v\le m.
\]
Their finite computation uses elementary-symmetric recursions.

For each even degree \(D\ge2\) selected below, define
\[
\mathcal T_0(x)=1,\qquad
\mathcal T_1(x)=x,\qquad
\mathcal T_{v+1}(x)=2x\mathcal T_v(x)-\mathcal T_{v-1}(x)
\quad(v\ge1),
\]
\[
p_D^{\mathrm{abs}}(x)
=\frac2\pi+\frac4\pi
\sum_{v=1}^{D/2}
\frac{(-1)^{v+1}\mathcal T_{2v}(x)}{4v^2-1}
=\sum_{v=0}^Dc_{D,v}x^v,
\]
and, for a positive radius \(a_{\mathrm{rad}}\),
\[
\widehat p_{a_{\mathrm{rad}},D,j}
=\sum_{v=0}^D
c_{D,v}a_{\mathrm{rad}}^{\,1-v}\mathsf U_{vj}.
\]

For \(n>2\), define \(\widehat F\) by the following ordered branches, taking the first applicable branch:
\[
\widehat F=
\begin{cases}
\displaystyle\frac1d\sum_{j=1}^d|\overline W_j|,
&L<4096,\\[6pt]
\displaystyle\frac1d\sum_{j=1}^d
\left[
\widehat p_{a_{\mathrm{rad}},D,j}
\mathbf1_{\{|M_j|\le T_*\}}
+\operatorname{sign}(M_j)\overline W_j
\mathbf1_{\{|M_j|>T_*\}}
\right],
&L\ge4096,\ t\ge d^2L,\\[6pt]
1/4,
&L\ge4096,\ t<d^2L,\ L/(1024H)<4,\\[4pt]
\displaystyle\frac1d\sum_{j=1}^d
\widehat p_{a_{\mathrm{rad}},D,j},
&L\ge4096,\ t<d^2L,\ L/(1024H)\ge4.
\end{cases}
\]
In the second branch, set
\[
D=2\left\lfloor\frac{L}{2048}\right\rfloor,
\qquad T_*=4\sigma\sqrt L,\qquad
a_{\mathrm{rad}}=8\sigma\sqrt L,
\qquad \operatorname{sign}(0)=0.
\]
In the fourth branch, set
\[
a_{\mathrm{rad}}=\frac12,
\qquad D=2\left\lfloor\frac{L}{2048H}\right\rfloor.
\]
For \(n>2\), set
\[
\widehat D=\max\left\{0,\min\left\{\frac14,\frac{\widehat F}{2}\right\}\right\}.
\]
For every \(n\ge2\), report
\[
q=\sqrt{10C_0\rho(t,d)},\qquad
J_{\mathrm{TV}}
=\left[
\max\{0,\widehat D-q\},
\min\{1/4,\widehat D+q\}
\right].
\]

We say that the concrete procedure guarantees hold with constant \(C_0\) when this fixed protocol belongs to \(\mathfrak K_{\mathrm{NI}}(n,d,\varepsilon)\) and, for every \(\theta\in[-1/2,1/2]^d\),
\[
\begin{aligned}
\mathbb E_{P_\theta,\mathsf K}
\left[(\widehat D-D_{\mathrm{TV}}(P_\theta))^2\right]
&\le C_0\rho(t,d),\\
\Pr_{P_\theta,\mathsf K}
\{D_{\mathrm{TV}}(P_\theta)\in J_{\mathrm{TV}}\}
&\ge0.90,\\
\mathbb E_{P_\theta,\mathsf K}
\left[\operatorname{length}J_{\mathrm{TV}}\right]
&\le C_0\sqrt{\rho(t,d)}.
\end{aligned}
\]
Here \(P_\theta(j,s)=(1+s\theta_j)/(2d)\) and
\(D_{\mathrm{TV}}(P_\theta)=(2d)^{-1}\sum_j|\theta_j|\).
Both decisions are Borel functions of the transcript, and the reported interval is nonempty, connected, closed, and contained in \([0,1/4]\).
\end{definitionv}

The attainment properties in \cref{obj:synth_4} require a finite sign-vector message from each participant and place both decisions on the paired target's natural range. They combine squared-error control, global coverage, and expected connected-interval length under a common protocol. A more explicit specification selects its polynomial estimate from public resource parameters. In that construction, the privacy attenuation \leanref{sym:\delta}{\ensuremath{\delta}} determines the scaled-message magnitude \leanref{sym:b}{\ensuremath{b}}, and the block size \leanref{sym:m}{\ensuremath{m}} separates pilot information from evaluation information.

\begin{definitionv}{synth_4}[Finite sign-vector attainment]
Fix \(n,d\in\{2,3,\ldots\}\), \(0<\varepsilon\le1\), and \(C_0>0\). Let
\[
P_\theta(j,s)=\frac{1+s\theta_j}{2d},
\qquad
\theta\in[-1/2,1/2]^d,\quad j\in[d],\quad s\in\{-1,1\}.
\]
The paired family and its total-variation target are
\[
\mathcal P_d^{\mathrm{pair}}
=\{P_\theta:\theta\in[-1/2,1/2]^d\},
\qquad
D_{\mathrm{TV}}(P_\theta)
=\frac1{2d}\sum_{j=1}^d|\theta_j|.
\]
We say that the general attainment guarantees hold with constant \(C_0\) when there exist a protocol
\[
\mathsf K\in\mathfrak K_{\mathrm{NI}}(n,d,\varepsilon),
\]
a Borel estimator \(\widehat D\), and a Borel interval decision \(J_{\mathrm{TV}}\) with the following properties.

Each participant's message space admits a measurable bijection with \(\{-1,1\}^d\). The protocol is a one-person-one-message pure-local protocol on the original alphabet \(\mathcal A_d=[d]\times\{-1,1\}\), as specified in \cref{obj:def:tv-private-decision-values}. In particular, its kernels are constant in the history and satisfy
\[
\mathsf K_i(E\mid(j,s),\eta,r)
\le e^\varepsilon
\mathsf K_i(E\mid(j',s'),\eta,r)
\]
for every participant, public seed, history, Borel message event, and pair of original inputs.

The estimator depends on the transcript alone: there is a Borel function \(f\) such that
\[
\widehat D(\mathbf Z^{\mathsf K},R,U)=f(\mathbf Z^{\mathsf K}).
\]
The interval \(J_{\mathrm{TV}}(\mathbf Z^{\mathsf K},R,U)\) is nonempty, connected and closed, has ordered Borel endpoints, and is contained in \([0,1/4]\).

For every \(\theta\in[-1/2,1/2]^d\), under \(n\) iid inputs from \(P_\theta\),
\[
\begin{aligned}
\mathbb E_{P_\theta,\mathsf K}
\left[(\widehat D-D_{\mathrm{TV}}(P_\theta))^2\right]
&\le C_0\,\rho(n\varepsilon^2,d),\\
\Pr_{P_\theta,\mathsf K}
\{D_{\mathrm{TV}}(P_\theta)\in J_{\mathrm{TV}}\}
&\ge0.90,\\
\mathbb E_{P_\theta,\mathsf K}
\left[\operatorname{length}J_{\mathrm{TV}}\right]
&\le C_0\sqrt{\rho(n\varepsilon^2,d)}.
\end{aligned}
\]
Probabilities and expectations include input sampling, channel randomness, the independent public seed, and independent analyst randomization. The same protocol belongs to \(\mathfrak K_{\mathrm{SI}}(n,d,\varepsilon)\).
\end{definitionv}

The construction in \cref{obj:synth_5} fixes collection independently of histories and uses resource-dependent branches for the statistical decision. Its scaled messages \leanref{sym:W}{\ensuremath{W^{\mathrm{pil}}}} and \(W^{\mathrm{ev}}\) support pilot averages and distinct-person moments \leanref{sym:\mathsf U}{\ensuremath{\mathsf U_{vj}}}. In the dense, large-dimension branch, the pilot selects between a polynomial estimate near zero and a signed evaluation mean away from zero. The other branches use an absolute empirical mean, a constant, or a polynomial across the full contrast range, according to the displayed resource conditions. This deliberately conservative construction serves as an explicit uniform rate-attainment witness. The concrete guarantees are the three uniform decision bounds specified in the definition; their establishment is part of \cref{obj:thm:paired-uniform-tv-frontier}.

Finally, the paired minimax squared-error criterion \leanref{sym:\mathfrak R^{\mathrm{TV}}}{\ensuremath{\mathfrak R^{\mathrm{TV}}_{\mathcal C}}} and honest-length criterion \leanref{sym:\mathfrak H^{\mathrm{TV}}}{\ensuremath{\mathfrak H^{\mathrm{TV}}_{\mathcal C}}} are defined in \cref{obj:def:tv-private-decision-values}. The next definition records the sequence comparisons for these values that accompany the finite-resource bounds. It keeps the quantifiers explicit when dimension and privacy vary with sample size.

\begin{propositionv}{prop:signed-causal-bridge}[Signed causal experiment equivalence]
Fix a sample size \(n\in\mathbb N\), a stratum count \(d\in\mathbb N\), and a sampling scheme satisfying:
\begin{itemize}
\item \textbf{(Dimension.)} \(d\ge 2\).
\item \textbf{(Independent participants.)} \cref{obj:ass:iid-people}.
\item \textbf{(Independent randomization.)} \cref{obj:ass:independent-randomness}.
\end{itemize}
Let \(\mathcal V_d\) be the class of causal probability laws in \cref{obj:def:causal-model}. For the stratum \(X\), binary assignment \(A\), observed outcome \(Y\), and potential outcomes \(Y^0,Y^1\), write
\[
S=(2A-1)(2Y-1),\qquad
\mu_{aj}(P)=\mathbb E_P[Y^a\mid X=j],\qquad
\tau_j(P)=\mu_{1j}(P)-\mu_{0j}(P).
\]
Write \(\tau(P)\) for the vector of these contrasts, and define the baseline mean, first-best value, and average absolute contrast by
\[
B(P)=\mathbb E_P[Y],\qquad
V(P)=\frac1d\sum_{j=1}^d\max\{\mu_{0j}(P),\mu_{1j}(P)\},
\qquad
F(\theta)=\frac1d\sum_{j=1}^d|\theta_j|.
\]
The paired signed law \(P_\theta\) has masses
\[
P_\theta(j,s)=\frac{1+s\theta_j}{2d},
\qquad j\in[d],\quad s\in\{-1,1\},
\]
and \(\mathsf U_{2d}\) is the uniform law on this paired alphabet.

For every \(P\in\mathcal V_d\) and every \(j\in[d]\),
\[
\mathbb E_P[S\mid X=j]=\tau_j(P),
\qquad
\mathcal L_P(X,S)=P_{\tau(P)},
\]
and
\[
V(P)
=B(P)+\frac{F(\tau(P))}{2}
=B(P)+\operatorname{TV}(P_{\tau(P)},\mathsf U_{2d}).
\]

For every \(\theta\in[-1/2,1/2]^d\), the symmetric law \(G_\theta\) of \cref{obj:def:symmetric-law} is a probability law in \(\mathcal V_d\), and
\[
V(G_\theta)=\frac12+\frac{F(\theta)}2.
\]
Moreover, for every \(j\in[d]\), \(s\in\{-1,1\}\), and \(a\in\{0,1\}\),
\[
G_\theta(X=j,S=s,A=a)=\frac{P_\theta(j,s)}2,
\]
and
\[
2Y-1=S(2A-1)
\qquad G_\theta\text{-almost surely}.
\]

For every privacy parameter \(\varepsilon\in\mathbb R\), the following two protocol transformations preserve both the sequential and noninteractive classes of \cref{obj:def:sequential-class,obj:def:noninteractive-class}. Let \(\mathcal H_i\) denote the message-history space preceding participant \(i\), and let \(\mathcal R\) denote the declared space of the public-seed random variable \(R\). For every original-record protocol \(Q=(Q_i)_{i=1}^n\), define its paired-alphabet protocol \(\mathsf K\) by averaging over the fair assignment bit:
\[
\mathsf K_i(\,\cdot\mid(j,s),h,r)
=
\frac12\sum_{a\in\{0,1\}}
Q_i\!\left(
\,\cdot\mid
\left(j,a,\frac{1+s(2a-1)}2\right),h,r
\right),
\qquad
h\in\mathcal H_i,\quad r\in\mathcal R.
\]
Sequential membership of \(Q\) at \(\varepsilon\) implies sequential membership of \(\mathsf K\) at \(\varepsilon\), and noninteractive membership of \(Q\) implies noninteractive membership of \(\mathsf K\) at \(\varepsilon\). For every \(\theta\in[-1/2,1/2]^d\), the joint law of the public seed \(R\) and message vector \(\mathbf Z\) generated by \(Q\) under \(G_\theta\) and the fixed sampling scheme equals the joint law generated by \(\mathsf K\) from \(n\) independent \(P_\theta\) inputs and an independent public seed with the same law.

Conversely, for every paired-alphabet protocol \(\mathsf K\), define its original-record protocol \(Q\) by
\[
Q_i(\,\cdot\mid(j,a,y),h,r)
=
\mathsf K_i(\,\cdot\mid(j,(2a-1)(2y-1)),h,r),
\qquad
h\in\mathcal H_i,\quad r\in\mathcal R.
\]
Sequential membership of \(\mathsf K\) at \(\varepsilon\) implies sequential membership of \(Q\) at \(\varepsilon\), and noninteractive membership of \(\mathsf K\) implies noninteractive membership of \(Q\) at \(\varepsilon\). For every \(\theta\in[-1/2,1/2]^d\), the joint public-seed and message law generated by this \(Q\) under \(G_\theta\) and the fixed sampling scheme equals that generated by \(\mathsf K\) from \(n\) independent \(P_\theta\) inputs and an independent public seed with the same law.
\end{propositionv}

The comparisons collected in \cref{obj:synth_6} distinguish vanishing estimation error, shrinking honest intervals, and inverse-resource squared error. Their uniform resource-domain formulation also specifies the dimension-dependent scale when effective resource is proportional to \(d^2\). Together with the exact signed-experiment correspondence in \cref{obj:prop:signed-causal-bridge}, these definitions provide the population, observation, and decision framework for the welfare and total-variation results.


\section{Minimax precision and an attaining private procedure}

The signed-experiment correspondence in \cref{obj:prop:signed-causal-bridge} separates optimal welfare into the baseline outcome mean and one half of the average absolute treatment contrast:
\[
V(P)=B(P)+\frac12F(\tau(P)).
\]
This decomposition suggests collecting separate private messages for the two components. Estimating the baseline is a bounded-mean problem; estimating the absolute-contrast component requires accounting for the behavior of absolute value near treatment ties. The procedure below combines private mean estimation with polynomial approximation and selects its statistical decision using the public sample size, dimension, and privacy budget.

\subsection{Private messages and the welfare procedure}

We first specify the sign-vector message law and the distinct-person moments used to estimate polynomials of the treatment contrasts. Under the \leanref{S-4}{scaled-message convention}, the privacy attenuation is \(\delta=\tanh(\varepsilon/2)\), the message magnitude is \(b=d/\delta\), and a sign-vector coordinate is scaled as \(W_{ij}=bZ_{ij}\). The following message law places the record-dependent information in the participant's realized stratum while drawing the remaining coordinates as fair signs.

\begin{definitionv}{def:vector-kernel}[Sign-vector kernel $Q^{\mathrm{vec}}$]
The sign-vector message kernel is
\[
Q^{\mathrm{vec}}(z\mid O=(j,A,Y))
=
2^{-d}\bigl(1+\delta(2A-1)(2Y-1)z_j\bigr),
\qquad z\in\{-1,1\}^d,
\]
where \(\delta=\tanh(\varepsilon/2)\) is the privacy attenuation. A sampler draws the coordinates independently: every coordinate other than \(j\) is a fair sign, and coordinate \(j\) has mean \(\delta S\), where
\[
S=(2A-1)(2Y-1).
\]
\end{definitionv}

The message law \leanref{sym:Q^{\mathrm{vec}}}{\ensuremath{Q^{\mathrm{vec}}}} in \cref{obj:def:vector-kernel} protects the full observed record: its likelihood ratio across any two records is bounded by \((1+\delta)/(1-\delta)=e^\varepsilon\). Uniform strata and fair assignment make the scaled coordinate mean equal to the corresponding treatment contrast. Products formed from distinct participants therefore supply unbiased estimates of powers of that contrast. The next definition collects these products into symmetric averages.

\begin{definitionv}{def:private-moments}[Private moments $\mathsf U_{kj}$]
For a block of \(m\) participants with scaled coordinate messages \(W_{ij}\), define
\[
\mathsf U_{kj}
=
\binom{m}{k}^{-1}
\sum_{\substack{J\subseteq\{1,\ldots,m\}\\ |J|=k}}
\prod_{i\in J}W_{ij}
\quad(1\le k\le m),
\qquad
\mathsf U_{0j}=1.
\]
These quantities are computed by elementary-symmetric recursions.
\end{definitionv}

At moment degree \leanref{sym:k}{\ensuremath{k}}, the distinct-person construction in \cref{obj:def:private-moments} uses independence across participants to estimate \(\tau_j(P)^k\). To compute all moments through a selected degree \(D\), initialize \(E_0=1\) and \(E_v=0\) for \(1\le v\le D\). For each observation \(W_{ij}\), update in descending order
\[
E_v\leftarrow E_v+W_{ij}E_{v-1}
\qquad(v=\min\{D,i\},\ldots,1),
\]
and then set \(\mathsf U_{vj}=E_v/\binom mv\). This costs \(O(mD)\) arithmetic operations per coordinate and avoids enumerating participant subsets. It is the classical symmetric-statistic structure of \citet{Hoeffding1948}, applied here to scaled private messages.

The complete procedure uses disjoint baseline, pilot, and evaluation blocks. Binary messages estimate the baseline; the pilot determines which treatment contrasts are close to zero; and evaluation messages supply both coordinate means and polynomial moments. All message laws are fixed before collection. Items 1--7 and 9 below specify collection and the reported estimate and interval. Item 8 records converse tuning packaged in the same verified declaration; it is an analytical choice and is not an operation performed by participants or the analyst.

\begin{algorithmv}{def:frontier-handle}[Public-resource welfare construction $\mathscr H$]
The construction \(\mathscr H\) takes the participant records, sample size \(n\ge2\), dimension \(d\), and privacy parameter \(\varepsilon\) as inputs; its tuning parameters are fixed below, and all auxiliary quantities and summation indices are local to the construction.
\begin{enumerate}
\item Fix the numerical constants, effective private resource, and logarithmic scales:
\[
\eta_*=\frac1{1024},
\qquad L_0=4096,
\qquad t=n\varepsilon^2,
\qquad L=\log(ed),
\qquad z_*=\frac{d^2L}{t},
\qquad w=\log(e+z_*).
\]
Define the squared-error rate function by
\[
\rho(t,d)=
\begin{cases}
\displaystyle\frac{d^2}{tL},&t\ge d^2L,\\[3pt]
\displaystyle\min\left\{1,
\left[\frac{\log(e+d^2L/t)}{L}\right]^2\right\},
&0<t<d^2L.
\end{cases}
\]
Use a singleton public seed and ignore analyst randomization.

\item If \(n=2\), each participant sends a constant message, and set
\[
\widehat V=\frac12.
\]
For this branch, proceed directly to the converse tuning and interval output in the final two steps. For \(n>2\), define the block sizes by
\[
m=\lfloor n/3\rfloor,
\qquad m_0=n-2m.
\]
In the fixed participant order, allocate a baseline block of size \(m_0\), a disjoint sign-vector pilot block of size \(m\), and a disjoint sign-vector evaluation block of size \(m\).

\item For \(n>2\), define the privacy attenuation and scaled-message magnitude by
\[
\delta=\tanh(\varepsilon/2),
\qquad b=\frac d\delta.
\]
Each baseline participant sends a binary message with law
\[
\Pr(\mathsf H_i=h\mid O)
=\frac{1+h\delta(2Y-1)}2,
\qquad h\in\{-1,1\},
\qquad 1\le i\le m_0.
\]
Define the baseline mean estimate by
\[
\widehat B
=\frac12+\frac1{2\delta m_0}
\sum_{i=1}^{m_0}\mathsf H_i.
\]

\item For \(n>2\), each participant in either vector block sends one message with law
\[
Q^{\mathrm{vec}}(z\mid O=(j,A,Y))
=
2^{-d}\bigl(1+\delta(2A-1)(2Y-1)z_j\bigr),
\qquad z\in\{-1,1\}^d,
\]
using independent local randomness. Write \(Z^{\mathrm{pil}}_{ij}\) and \(Z^{\mathrm{ev}}_{ij}\) for coordinate \(j\) of participant \(i\)'s message in the two blocks, and define the scaled-message matrices by
\[
W^{\mathrm{pil}}_{ij}=bZ^{\mathrm{pil}}_{ij},
\qquad
W^{\mathrm{ev}}_{ij}=bZ^{\mathrm{ev}}_{ij},
\qquad 1\le i\le m,\quad 1\le j\le d.
\]
The participant order, binary and sign-vector message alphabets, and kernels are predetermined. This is one fixed noninteractive protocol on the original participant inputs for both protocol classes \(\mathfrak Q_{\mathrm{NI}}\) and \(\mathfrak Q_{\mathrm{SI}}\).

\item For \(n>2\), define the noise scale, polynomial calibration factor, and coordinate means by
\[
\sigma^2=\frac{b^2}{m},
\qquad \sigma=\sqrt{b^2/m},
\qquad H=\log(e+\sigma^2L),
\]
\[
M_j=\frac1m\sum_{i=1}^m W^{\mathrm{pil}}_{ij},
\qquad
\overline W_j=\frac1m\sum_{i=1}^m W^{\mathrm{ev}}_{ij},
\qquad 1\le j\le d.
\]
Define the evaluation-block moments by
\[
\mathsf U_{vj}
=
\binom{m}{v}^{-1}
\sum_{\substack{J\subseteq\{1,\ldots,m\}\\ |J|=v}}
\prod_{i\in J}W^{\mathrm{ev}}_{ij}
\quad(1\le v\le m),
\qquad
\mathsf U_{0j}=1.
\]
Compute these moments by elementary-symmetric recursions.

\item For each even degree \(D\ge2\) selected below, define the Chebyshev polynomials and the absolute-value approximation by
\[
\mathcal T_0(x)=1,
\qquad \mathcal T_1(x)=x,
\qquad
\mathcal T_{v+1}(x)
=2x\mathcal T_v(x)-\mathcal T_{v-1}(x)
\quad(v\ge1),
\]
\[
p_D^{\mathrm{abs}}(x)
=
\frac2\pi+\frac4\pi\sum_{v=1}^{D/2}
\frac{(-1)^{v+1}\mathcal T_{2v}(x)}{4v^2-1}
=
\sum_{v=0}^D c_{D,v}x^v.
\]
For the positive radius \(a_{\mathrm{rad}}\) selected below, define the coordinate polynomial estimate by
\[
\widehat p_{a_{\mathrm{rad}},D,j}
=
\sum_{v=0}^D
c_{D,v}a_{\mathrm{rad}}^{\,1-v}\mathsf U_{vj}.
\]

\item For \(n>2\), define the estimate \(\widehat F\) of the average absolute treatment contrast \(F(\tau(P))\) by taking the first applicable branch below.

If \(L<L_0\), set
\[
\widehat F=\frac1d\sum_{j=1}^d|\overline W_j|.
\]
Otherwise, if \(d^2L\le t\), define
\[
D=2\left\lfloor\frac{L}{2048}\right\rfloor,
\qquad T_*=4\sigma\sqrt L,
\qquad a_{\mathrm{rad}}=8\sigma\sqrt L,
\]
and set
\[
\widehat F
=
\frac1d\sum_{j=1}^d
\left[
\widehat p_{a_{\mathrm{rad}},D,j}
\mathbf1_{\{|M_j|\le T_*\}}
+
\operatorname{sign}(M_j)\overline W_j
\mathbf1_{\{|M_j|>T_*\}}
\right],
\qquad \operatorname{sign}(0)=0.
\]
Otherwise, if \(L/(1024H)<4\), set
\[
\widehat F=\frac14.
\]
Otherwise, define
\[
a_{\mathrm{rad}}=\frac12,
\qquad
D=2\left\lfloor\frac{L}{2048H}\right\rfloor,
\]
and set
\[
\widehat F
=\frac1d\sum_{j=1}^d
\widehat p_{a_{\mathrm{rad}},D,j}.
\]

\item For \(n>2\), combine the baseline and absolute-contrast estimates and clip the resulting welfare estimate:
\[
\widetilde V=\widehat B+\frac12\widehat F,
\qquad
\widehat V=\max\{1/4,\min\{3/4,\widetilde V\}\}.
\]

\item For every allowed resource tuple, define the converse amplitude and integer degrees by
\[
(\alpha,k)=
\begin{cases}
\left(\displaystyle\frac1{100}\sqrt{z_*},
\,2\lceil16L\rceil+1\right),&z_*\le1,\\[3pt]
\left(\displaystyle\frac1{100},
\,2\lceil16L/w\rceil+1\right),&z_*>1,
\end{cases}
\qquad K=k-1.
\]
Use these values in the scaled symmetric moment-matching measures, their independent coordinate product priors, and the symmetric causal laws \(G_\theta\).

\item For every \(n\ge2\), define the interval radius by
\[
q_*=\sqrt{10^{17}\rho(t,d)}.
\]
Output the welfare estimate and connected closed interval
\[
\widehat V=
\begin{cases}
1/2,&n=2,\\
\max\{1/4,\min\{3/4,\widehat B+\widehat F/2\}\},&n>2,
\end{cases}
\qquad
\left[
\max\{1/4,\widehat V-q_*\},
\min\{3/4,\widehat V+q_*\}
\right].
\]
\end{enumerate}
\end{algorithmv}

The welfare construction \leanref{sym:\mathscr H}{\ensuremath{\mathscr H}} in \cref{obj:def:frontier-handle} uses pilot information solely in the statistical decision, so collection remains noninteractive. Independence between the pilot and evaluation blocks separates coordinate selection from evaluation. Near zero, the polynomial estimates absolute value through distinct-person moments; beyond the pilot threshold, the signed evaluation mean uses the locally linear form of absolute value. The remaining branches calibrate the estimate to bounded dimension or lower effective resource. Clipping places the welfare estimate and interval endpoints in the model's welfare range.

For clarity, the converse tuning recorded in item 8 of \cref{obj:def:frontier-handle} belongs to the lower-bound analysis rather than the reporting algorithm. It specifies the moment-matching priors and resource-dependent symmetric causal laws used later in the converse. The operational construction consists of message collection, moment computation, branch selection, clipping, and the interval output; its interval radius is calibrated to the uniform squared-error bound.

\subsection{Uniform precision across resource regimes}

The next result characterizes minimax squared error and globally honest expected interval length for the causal model in \cref{obj:def:causal-model}. Its approximation input is the Bernstein limit for best polynomial approximation to absolute value, as presented by \citet[Section 3.1, p.~8]{CaiLow2011Nonsmooth}. The comparison is uniform over the stated resources and covers both noninteractive and sequentially interactive collection.

\begin{theoremv}{thm:uniform-private-value-frontiers}[Uniform private welfare frontiers]*
This result is conditional on the cited, unformalized Bernstein absolute-approximation limit. If \(e_K=\inf_{\deg p\le K}\sup_{|x|\le1}\lvert |x|-p(x)\rvert\) is the best uniform error for approximating \(|x|\) on \([-1,1]\), then \(2k e_{2k}\to\beta_*\) for a constant \(0.278<\beta_*<0.286\). See \citet[Section 3.1]{CaiLow2011Nonsmooth} and \citet{Bernstein1913}.

There are universal constants \(0<c\le C_0<\infty\) with the following properties. For every sample size \(n\), dimension \(d\), privacy budget \(\varepsilon\), and sampling scheme satisfying
\begin{itemize}
\item \textbf{(Resources.)} \(n\ge2\), \(d\ge2\), and \(0<\varepsilon\le1\);
\item \textbf{(Sampling.)} \cref{obj:ass:iid-people};
\item \textbf{(Randomization.)} \cref{obj:ass:independent-randomness};
\end{itemize}
let \(t=n\varepsilon^2\) be the effective private resource and \(L=\log(ed)\) the log-dimension factor. For either protocol class \(\mathcal C\in\{\mathrm{NI},\mathrm{SI}\}\), define
\[
r_{\mathcal C}(n,d,\varepsilon)=\rho(t,d)=
\begin{cases}
\dfrac{d^2}{tL},&t\ge d^2L,\\[4pt]
\displaystyle\min\left\{1,
\left[\frac{\log(e+d^2L/t)}{L}\right]^2\right\},
&0<t<d^2L,
\end{cases}
\qquad
h_{\mathcal C}(n,d,\varepsilon)=\sqrt{\rho(t,d)}.
\]
The minimax squared-error risk \(\mathfrak R_{\mathcal C}\) and globally \(0.90\)-honest connected interval length \(\mathfrak H_{\mathcal C}\), defined in \cref{obj:def:risk,obj:def:honest-length} for this sampling scheme, satisfy
\[
c\,r_{\mathcal C}(n,d,\varepsilon)
\le \mathfrak R_{\mathcal C}(n,d,\varepsilon)
\le C_0\,r_{\mathcal C}(n,d,\varepsilon),
\qquad
c\,h_{\mathcal C}(n,d,\varepsilon)
\le \mathfrak H_{\mathcal C}(n,d,\varepsilon)
\le C_0\,h_{\mathcal C}(n,d,\varepsilon).
\]

For each class, a noninteractive \(\varepsilon\)-private protocol \(Q\), a transcript-only estimator \(T\), and a connected closed interval decision \(I\) with endpoints in \([1/4,3/4]\) attain these upper bounds simultaneously: for every \(P\in\mathcal V_d\), the declared causal model,
\[
\mathbb E_P\!\left[(T(\mathbf Z,R,U)-V(P))^2\right]
\le C_0\,r_{\mathcal C}(n,d,\varepsilon),
\qquad
\mathbb P_P\!\left\{V(P)\in I(\mathbf Z,R,U)\right\}\ge0.90,
\]
\[
\mathbb E_P\!\left[\operatorname{length}
 I(\mathbf Z,R,U)\right]
\le C_0\,h_{\mathcal C}(n,d,\varepsilon).
\]
Here expectations and probabilities use the specified sampling scheme and protocol. The resource-selected explicit noninteractive welfare procedure also satisfies these three guarantees with the same \(C_0\) and the NI rates.

For both classes, risk and squared honest length obey the quantitative comparisons
\[
\frac{c^2}{C_0}\,
\mathfrak R_{\mathcal C}(n,d,\varepsilon)
\le
\mathfrak H_{\mathcal C}(n,d,\varepsilon)^2
\le
\frac{C_0^2}{c}\,
\mathfrak R_{\mathcal C}(n,d,\varepsilon).
\]
The class comparisons are
\[
1\le
\frac{\mathfrak R_{\mathrm{NI}}(n,d,\varepsilon)}
     {\mathfrak R_{\mathrm{SI}}(n,d,\varepsilon)}
\le\frac{C_0}{c},
\qquad
1\le
\frac{\mathfrak H_{\mathrm{NI}}(n,d,\varepsilon)}
     {\mathfrak H_{\mathrm{SI}}(n,d,\varepsilon)}
\le\frac{C_0}{c}.
\]

For every \(d\ge2\) and \(0<t<d^2\log(ed)\), the saturation condition is exactly
\[
\rho(t,d)=1
\quad\Longleftrightarrow\quad
t\le\frac{d^2\log(ed)}{ed-e}.
\]

Along every sequence of allowed resource tuples
\((n_k,d_k,\varepsilon_k)\), write \(t_k=n_k\varepsilon_k^2\). Then
\[
\rho(t_k,d_k)\longrightarrow0
\quad\Longleftrightarrow\quad
t_k\longrightarrow\infty
\ \text{and}\
\frac{\log(1+d_k^2/t_k)}{\log(ed_k)}\longrightarrow0.
\]
Moreover,
\[
\rho(t_k,d_k)\asymp t_k^{-1}
\quad\Longleftrightarrow\quad
\left(\exists a>0:\ t_k\ge a\ \text{eventually}\right)
\ \text{and}\
\left(\exists M\in\mathbb R:\ \min\{t_k,d_k\}\le M\
\text{eventually}\right),
\]
where \(\asymp\) denotes eventual comparability by positive finite constants. In particular, when \(t_k\to\infty\),
\[
\rho(t_k,d_k)\asymp t_k^{-1}
\quad\Longleftrightarrow\quad
\exists M\in\mathbb N:\ d_k\le M\ \text{eventually}.
\]
For every family of sampling schemes satisfying
\cref{obj:ass:iid-people,obj:ass:independent-randomness} at each \(n,d\), and for either fixed class \(\mathcal C\), these consistency and private-parametric characterizations also hold for the corresponding minimax risk. Honest length tends to zero under the same consistency condition, and its private-parametric order is \(t_k^{-1/2}\) under the same comparability conditions.

Finally, for every \(0<a\le b<\infty\), there are constants \(0<c\le C_0<\infty\) and an integer \(D_0\), possibly depending on \(a,b\), such that for every \(d\ge\max\{D_0,2\}\) and every \(t\) with \(a\le t/d^2\le b\),
\[
c\left[\frac{\log\log(ed)}{\log(ed)}\right]^2
\le \rho(t,d)\le
C_0\left[\frac{\log\log(ed)}{\log(ed)}\right]^2.
\]
Consequently, along allowed sequences with \(d_k\to\infty\) and \(t_k/d_k^2\) in a fixed positive finite interval, either class has
\[
\mathfrak R_{\mathcal C}(n_k,d_k,\varepsilon_k)
\asymp
\left[\frac{\log\log(ed_k)}{\log(ed_k)}\right]^2,
\qquad
\mathfrak H_{\mathcal C}(n_k,d_k,\varepsilon_k)
\asymp
\frac{\log\log(ed_k)}{\log(ed_k)}.
\]
\end{theoremv}
% CAUSALSMITH-CITED-SCOPE-BEGIN thm:uniform-private-value-frontiers
\verificationfootnotetext{\textbf{Formalization scope.} This result uses the cited conclusion \citep{CaiLow2011Nonsmooth} (Section 3.1, definition of the Bernstein constant and immediately following paragraph, page 8). The cited source proof is not formalized here: Lean verifies its use and all remaining steps. If that cited conclusion is false, the portions of this result that depend on it are not certified.}
% CAUSALSMITH-CITED-SCOPE-END thm:uniform-private-value-frontiers

The intuition behind \cref{obj:thm:uniform-private-value-frontiers} is a balance between approximating absolute value and estimating its polynomial approximation from noisy messages. Higher degrees reduce approximation bias while increasing the variability of private moments. The resource-dependent choices in \cref{obj:def:frontier-handle} balance these effects, with the pilot localizing the approximation when resources are large. The squared-error rate \leanref{sym:r_{\mathcal C}}{\ensuremath{r_{\mathcal C}(n,d,\varepsilon)}} and honest-length rate \leanref{sym:h_{\mathcal C}}{\ensuremath{h_{\mathcal C}(n,d,\varepsilon)}} express the same precision scale: honest length has the order of the square root of minimax risk.

For \(t\ge d^2L\), squared error has order \(d^2/(tL)\), and honest length has order \(d/\sqrt{tL}\). Below this threshold, precision depends logarithmically on the resource ratio \(d^2L/t\). At and below the stated saturation threshold, risk and honest length both have constant order. These regimes translate participant count and privacy budget into uncertainty for reporting population welfare, uniformly over treatment contrasts that include exact ties.

Consistency requires both diverging effective private resource and a vanishing normalized logarithmic resource deficit:
\[
t_k\longrightarrow\infty,
\qquad
\frac{\log(1+d_k^2/t_k)}{\log(ed_k)}\longrightarrow0.
\]
The first requirement controls the overall information available after privatization; the second describes how dimension may grow relative to that information. Under these conditions, minimax squared error vanishes and globally honest intervals shrink in either protocol class.

The private-parametric boundary is especially sharp along sequences with \(t_k\to\infty\). Eventual boundedness of dimension is equivalent to squared-error order \(t_k^{-1}\) and honest-length order \(t_k^{-1/2}\). The theorem also gives the full comparability criterion for sequences whose effective resource stays bounded away from zero. When dimension diverges and \(t_k/d_k^2\) remains in a fixed positive finite interval, the logarithmic regime yields squared-error order \([\log\log(ed_k)/\log(ed_k)]^2\) and interval-length order \(\log\log(ed_k)/\log(ed_k)\).

The protocol comparisons establish equality of statistical orders between noninteractive and sequentially interactive collection. A fixed noninteractive procedure attains the common orders, while the lower bounds apply throughout the sequential class. The constants in these comparisons quantify uniform comparability; equality of orders leaves room for differences in leading constants. The explicit numerical calibration in \cref{obj:def:frontier-handle} serves simultaneous risk, coverage, and length guarantees.

\subsection{Two-cell calibration}

A two-stratum trial provides a direct calibration of the bounded-dimension characterization. The following result states the saturated and private-parametric orders across the entire two-cell causal model, together with simultaneous attainment by a transcript-only noninteractive decision.

\begin{theoremv}{thm:two-cell-calibration}[Two-cell risk and length calibration]
There exist universal constants \(0<c\le C_0<\infty\) such that, for every sample size \(n\), privacy budget \(\varepsilon\), and sampling scheme for two-cell observed records satisfying
\begin{itemize}
\item \textbf{(Resources.)} \(n\ge2\) and \(0<\varepsilon\le1\).
\item \textbf{(Sampling.)} \cref{obj:ass:iid-people}.
\item \textbf{(Randomness.)} \cref{obj:ass:independent-randomness}.
\end{itemize}
the minimax squared-error risk and globally \(0.90\)-honest expected interval length defined in \cref{obj:def:risk,obj:def:honest-length}, evaluated under this sampling scheme, satisfy, for every \(\mathcal C\in\{\mathrm{NI},\mathrm{SI}\}\),
\[
c\min\{1,(n\varepsilon^2)^{-1}\}
\le \mathfrak R_{\mathcal C}(n,2,\varepsilon)
\le C_0\min\{1,(n\varepsilon^2)^{-1}\},
\]
\[
c\min\{1,(\sqrt n\,\varepsilon)^{-1}\}
\le \mathfrak H_{\mathcal C}(n,2,\varepsilon)
\le C_0\min\{1,(\sqrt n\,\varepsilon)^{-1}\}.
\]

Moreover, for each such sampling scheme, \(n\), and \(\varepsilon\), there exist a total noninteractive protocol \(Q\in\mathfrak Q_{\mathrm{NI}}(n,2,\varepsilon)\) on the original observed records, an estimator \(T(\mathbf Z,R,U)\) equal to a Borel function of the message transcript \(\mathbf Z\) alone, and a connected closed interval \(I(\mathbf Z,R,U)\) with ordered Borel endpoints in \([1/4,3/4]\), such that, simultaneously for every \(P\in\mathcal V_2\), the causal model in \cref{obj:def:causal-model},
\[
\mathbb E\!\left[(T(\mathbf Z,R,U)-V(P))^2\right]
\le C_0\min\{1,(n\varepsilon^2)^{-1}\},
\]
\[
\mathbb P\!\left\{V(P)\in I(\mathbf Z,R,U)\right\}\ge0.90,
\qquad
\mathbb E\!\left[\operatorname{length}(I(\mathbf Z,R,U))\right]
\le C_0\min\{1,(\sqrt n\,\varepsilon)^{-1}\}.
\]
The expectations and probability are taken under the decision law induced by the sampling scheme, \(Q\), and \(P\).
\end{theoremv}

With two strata, the dimension factors are fixed, so effective private resource governs the precision orders in \cref{obj:thm:two-cell-calibration}. Diverging \(n\varepsilon^2\) gives inverse-resource squared error and inverse-square-root-resource honest length, while bounded resource retains the saturated scale. The guarantees hold uniformly over the two-cell model, including exact treatment ties. This calibration connects the general characterization to the familiar privacy order for bounded-mean estimation.

\subsection{Analytical ingredients}

The upper-bound analysis separates approximation bias from aggregate private-moment variance. The Chebyshev approximation and exact moment-duality conclusions are established in \cref{obj:lem:cai-low-chebyshev-approximation,obj:lem:cai-low-moment-duality}. The cited Bernstein limit in \citet[Section 3.1, p.~8]{CaiLow2011Nonsmooth} supplies the quantitative inverse-degree approximation scale used in the general minimax converse. Private-moment variability and dependence across coordinates are treated in \cref{obj:lem:private-moment-and-dependence}, and \cref{obj:lem:finite-private-polynomial-calibration} connects those bounds to the resource-selected polynomial estimates. The bounded-mean concentration bound in \cref{obj:lem:bounded-mean-concentration}, associated with \citet{Hoeffding1963Bounded}, controls pilot selection. These ingredients assign separate roles to approximation, evaluation noise, and selection error.

For the lower bounds, the \leanref{S-3}{sequential converse construction} combines moment-matching coordinate priors with separated welfare values under the symmetric causal laws in \cref{obj:def:symmetric-law}. The contraction bound in \cref{obj:lem:adaptive-moment-contraction} controls the distinguishability of the resulting sequential transcript mixtures. Their welfare separation is quantified in \cref{obj:lem:separated-value-mixtures}, while \cref{obj:lem:dimension-free-private-two-point} supplies the complementary private two-point bound. The same indistinguishability argument supports squared-error and honest-length lower bounds: estimating separated welfare values requires distinguishing the mixtures, and global coverage requires intervals long enough to accommodate their separation. Detailed calculations and the proofs of the main results are assigned to the appendices.


\section{Total-variation estimation and interpretation}

The signed experiment in \cref{obj:prop:signed-causal-bridge} connects the welfare analysis to estimation of total-variation distance from a public uniform reference. The resulting statistical characterization concerns a family with equal mass across pairs and an unknown signed contrast within each pair. A separate converse applies to the full probability simplex on the same alphabet. To state these conclusions, we first specify the two statistical models and their common estimand.

\begin{definitionv}{def:tv-corollary-models}[Total-variation models $\mathcal P_d^{\mathrm{pair}}$ and $\Delta_{2d}$]
Define the finite paired alphabet and its probability simplex by
\[
\mathcal A_d=[d]\times\{-1,1\},
\qquad
\Delta_{2d}
=
\left\{
\mathsf p:\mathcal A_d\to[0,1]:
\sum_{j=1}^d\sum_{s\in\{-1,1\}}\mathsf p(j,s)=1
\right\}.
\]
Here \(\mathsf p\) is a probability law on the paired alphabet. Define the paired family and its uniform-reference total-variation functional by
\[
\mathcal P_d^{\mathrm{pair}}
=
\{P_\theta:\theta\in[-1/2,1/2]^d\}
\subseteq\Delta_{2d},
\]
\[
D_{\mathrm{TV}}(\mathsf p)
=
\frac12\sum_{j=1}^d\sum_{s\in\{-1,1\}}
\left|\mathsf p(j,s)-\frac1{2d}\right|.
\]
The reference law \(\mathsf U_{2d}\) is public and uniform on \(\mathcal A_d\). In either statistical model
\[
\mathcal M=\mathcal P_d^{\mathrm{pair}}
\quad\text{or}\quad
\mathcal M=\Delta_{2d},
\]
the data are \(n\) independent, identically distributed draws from one unknown law \(\mathsf p\in\mathcal M\), and the estimand is \(D_{\mathrm{TV}}(\mathsf p)\). The causal model \(\mathcal V_d\) retains its full-data laws and welfare estimand \(V(P)\).
\end{definitionv}

The paired family \leanref{sym:\mathcal P_d^{\mathrm{pair}}}{\ensuremath{\mathcal P_d^{\mathrm{pair}}}} fixes each pair's total probability at \(1/d\), whereas the simplex \leanref{sym:\Delta_{2d}}{\ensuremath{\Delta_{2d}}} permits arbitrary probability masses on the alphabet. In both models, the scalar target \leanref{sym:D_{\mathrm{TV}}}{\ensuremath{D_{\mathrm{TV}}(\mathsf p)}} measures departure from the known reference law \leanref{sym:\mathsf U_{2d}}{\ensuremath{\mathsf U_{2d}}}. This common reference makes the model comparison precise: the input alphabet and target remain fixed while the set of possible sampling laws expands.

The estimation and interval criteria use pure local privacy across every pair of input symbols. The following specification also fixes the information available to the analyst and the meaning of sequential interaction.

\begin{definitionv}{def:tv-private-decision-values}[Private total-variation decision values $\mathfrak R^{\mathrm{TV}}_{\mathcal C}$ and $\mathfrak H^{\mathrm{TV}}_{\mathcal C}$]
The class \(\mathfrak K_{\mathcal C}(n,d,\varepsilon)\), for \(\mathcal C\in\{\mathrm{NI},\mathrm{SI}\}\), consists of one-message-per-participant protocols on the paired input alphabet \(\mathcal A_d=[d]\times\{-1,1\}\). A protocol \(\mathsf K\) satisfies the following conditions:
\begin{itemize}
\item \textbf{(Spaces and kernels.)} The protocol has an arbitrary standard-Borel public-seed space with probability law \(\lambda\), standard-Borel participant message spaces, and participant-specific message kernels \(\mathsf K_i\) conditional on the input, preceding message history, and public seed.
\item \textbf{(Local privacy.)} For every participant \(i\), public seed \(r\), preceding message history \(\eta\), Borel message event \(E\), \(j,j'\in[d]\), and \(s,s'\in\{-1,1\}\),
\[
\mathsf K_i(E\mid(j,s),\eta,r)
\le e^\varepsilon\mathsf K_i(E\mid(j',s'),\eta,r).
\]
\item \textbf{(Public specification.)} All protocol choices depend only on \(n,d,\varepsilon\), the specified model, and public bounds.
\item \textbf{(Interaction.)} For \(\mathcal C=\mathrm{NI}\), every kernel is constant in its message history; participant-specific kernels and shared public seeds are permitted. For \(\mathcal C=\mathrm{SI}\), kernels may depend on preceding messages, and each participant is accessed once.
\item \textbf{(Independent randomizers.)} The public seed \(R\) and analyst randomizer \(U\) have laws
\[
R\sim\lambda,
\qquad
U\sim\operatorname{Uniform}[0,1],
\]
and are mutually independent and independent of the independent, identically distributed inputs.
\end{itemize}

Write
\[
\mathbf Z^{\mathsf K}=(Z_i^{\mathsf K})_{i=1}^n
\]
for the resulting message vector. An estimate \(\widehat D\) is any Borel real-valued function of \((\mathbf Z^{\mathsf K},R,U)\). An interval \(J_{\mathrm{TV}}\) is any nonempty connected closed interval contained in \([0,1]\), with ordered Borel endpoints that depend only on the same information.

For either statistical model
\[
\mathcal M\in\{\mathcal P_d^{\mathrm{pair}},\Delta_{2d}\},
\]
the minimax squared-error risk and globally honest expected interval length for the uniform-reference total-variation distance \(D_{\mathrm{TV}}(\mathsf p)\) are
\[
\begin{aligned}
\mathfrak R^{\mathrm{TV}}_{\mathcal C}(\mathcal M;n,d,\varepsilon)
&=
\inf_{\mathsf K\in\mathfrak K_{\mathcal C}(n,d,\varepsilon)}
\inf_{\widehat D}
\sup_{\mathsf p\in\mathcal M}
\mathbb E_{\mathsf p,\mathsf K}
\bigl[(\widehat D-D_{\mathrm{TV}}(\mathsf p))^2\bigr],\\
\mathfrak H^{\mathrm{TV}}_{\mathcal C}(\mathcal M;n,d,\varepsilon)
&=
\inf_{\mathsf K\in\mathfrak K_{\mathcal C}(n,d,\varepsilon)}
\inf_{\substack{J_{\mathrm{TV}}:\,
\inf_{\mathsf p\in\mathcal M}
\Pr_{\mathsf p,\mathsf K}
\{D_{\mathrm{TV}}(\mathsf p)\in J_{\mathrm{TV}}\}\ge0.90}}
\sup_{\mathsf p\in\mathcal M}
\mathbb E_{\mathsf p,\mathsf K}
\operatorname{length}J_{\mathrm{TV}}.
\end{aligned}
\]
Here \(\operatorname{length}J_{\mathrm{TV}}\) is the difference between its ordered endpoints. Probabilities and expectations include input sampling, channel randomness, public randomness, and analyst randomization. Singleton intervals are permitted. For the paired model, the attaining intervals specified below are contained in \([0,1/4]\).
\end{definitionv}

Global honesty requires coverage at every sampling law in the selected model, and expected length is evaluated at the worst such law. The connected interval \leanref{sym:J_{\mathrm{TV}}}{\ensuremath{J_{\mathrm{TV}}}} therefore measures the precision of a single reported range for the scalar target. The noninteractive and sequential classes both allow shared public randomness and participant-specific message distributions; their distinction is whether a participant's message distribution may depend on preceding messages.

\subsection{Matching precision on the paired family}

For the paired family, total variation is one half of the average absolute contrast in the signed experiment. This identity carries the approximation and private-moment analysis into the statistical problem. The published Bernstein approximation limit used in the comparison is given by \citet[Section 3.1, p.~8]{CaiLow2011Nonsmooth}. Under this input, the following result supplies matching squared-error and honest-length bounds, together with a concrete noninteractive procedure attaining their orders.

\begin{theoremv}{thm:paired-uniform-tv-frontier}[Paired-uniform total-variation frontier]*
This result is conditional on the cited, unformalized Bernstein absolute-approximation limit. If \(e_K=\inf_{\deg p\le K}\sup_{|x|\le1}\lvert |x|-p(x)\rvert\) is the best uniform error for approximating \(|x|\) on \([-1,1]\), then \(2k e_{2k}\to\beta_*\) for a constant \(0.278<\beta_*<0.286\). See \citet[Section 3.1]{CaiLow2011Nonsmooth} and \citet{Bernstein1913}.

For the paired family in \cref{obj:def:tv-corollary-models}, every integer \(d\ge2\) and every \(\theta\in[-1/2,1/2]^d\) satisfy
\[
 P_\theta(j,s)=\frac{1+s\theta_j}{2d},
 \qquad
 D_{\mathrm{TV}}(P_\theta)
 =\operatorname{TV}(P_\theta,\mathsf U_{2d})
 =\frac12F(\theta)
 =\frac{1}{2d}\sum_{j=1}^{d}|\theta_j|.
\]
Here \(\mathsf U_{2d}\) is the uniform law on the paired alphabet and \(F(\theta)\) is the average absolute contrast.

There are universal real constants \(0<c\le C_0<\infty\) such that the following conclusions hold for every resource tuple satisfying
\begin{itemize}
\item \textbf{(Sample size.)} \(n\) is an integer with \(n\ge2\).
\item \textbf{(Dimension.)} \(d\) is an integer with \(d\ge2\).
\item \textbf{(Privacy.)} \(0<\varepsilon\le1\).
\end{itemize}
For either \(\mathcal C\in\{\mathrm{NI},\mathrm{SI}\}\), let
\[
 t=n\varepsilon^2,\qquad
 r_{\mathcal C}(n,d,\varepsilon)=\rho(t,d),\qquad
 h_{\mathcal C}(n,d,\varepsilon)=\sqrt{\rho(t,d)},
\]
where \(\rho\) is the evaluated rate defined in \cref{obj:def:frontier-handle}. The decision values of \cref{obj:def:tv-private-decision-values} obey
\[
\begin{aligned}
 c\,r_{\mathcal C}(n,d,\varepsilon)
 &\le
 \mathfrak R^{\mathrm{TV}}_{\mathcal C}
   (\mathcal P_d^{\mathrm{pair}};n,d,\varepsilon)
 \le C_0\,r_{\mathcal C}(n,d,\varepsilon),\\
 c\,h_{\mathcal C}(n,d,\varepsilon)
 &\le
 \mathfrak H^{\mathrm{TV}}_{\mathcal C}
   (\mathcal P_d^{\mathrm{pair}};n,d,\varepsilon)
 \le C_0\,h_{\mathcal C}(n,d,\varepsilon).
\end{aligned}
\]

For each such resource tuple there exists a pure \(\varepsilon\)-locally private noninteractive paired-input protocol \(\mathsf K\), with each participant's message space in bijection with \(\{-1,1\}^d\), together with a transcript-only estimator \(\widehat D\) and a connected closed interval \(J_{\mathrm{TV}}\subseteq[0,1/4]\), such that, for every \(\mathsf p\in\mathcal P_d^{\mathrm{pair}}\), under independent sampling from \(\mathsf p\),
\[
\begin{aligned}
 \mathbb E_{\mathsf p,\mathsf K}
   [(\widehat D-D_{\mathrm{TV}}(\mathsf p))^2]
 &\le C_0\,r_{\mathcal C}(n,d,\varepsilon),\\
 \mathbb P_{\mathsf p,\mathsf K}
   \{D_{\mathrm{TV}}(\mathsf p)\in J_{\mathrm{TV}}\}
 &\ge0.90,\\
 \mathbb E_{\mathsf p,\mathsf K}
   [\operatorname{length}(J_{\mathrm{TV}})]
 &\le C_0\,h_{\mathcal C}(n,d,\varepsilon).
\end{aligned}
\]
These guarantees also hold for the concrete resource-selected sign-vector polynomial protocol and its total-variation estimator and interval, with the construction specified in \cref{obj:def:frontier-handle,obj:thm:uniform-private-value-frontiers}: this protocol is noninteractive and pure \(\varepsilon\)-locally private, and its squared risk, coverage, and expected length satisfy the three displayed bounds with \(\mathcal C=\mathrm{NI}\).

For the same universal constants and every admissible resource tuple,
\[
\begin{aligned}
1
&\le
\frac{\mathfrak R^{\mathrm{TV}}_{\mathrm{NI}}
  (\mathcal P_d^{\mathrm{pair}};n,d,\varepsilon)}
 {\mathfrak R^{\mathrm{TV}}_{\mathrm{SI}}
  (\mathcal P_d^{\mathrm{pair}};n,d,\varepsilon)}
\le \frac{C_0}{c},\\
1
&\le
\frac{\mathfrak H^{\mathrm{TV}}_{\mathrm{NI}}
  (\mathcal P_d^{\mathrm{pair}};n,d,\varepsilon)}
 {\mathfrak H^{\mathrm{TV}}_{\mathrm{SI}}
  (\mathcal P_d^{\mathrm{pair}};n,d,\varepsilon)}
\le \frac{C_0}{c}.
\end{aligned}
\]

The rate has the following exact resource properties. Write \(L=\log(ed)\). For every integer \(d\ge2\) and every \(0<t<d^2L\),
\[
 \rho(t,d)=1
 \quad\Longleftrightarrow\quad
 t\le\frac{d^2L}{ed-e}.
\]
For every sequence of admissible resource tuples
\((n_k,d_k,\varepsilon_k)\), write
\[
 t_k=n_k\varepsilon_k^2,\qquad L_k=\log(ed_k).
\]
Then
\[
 \rho(t_k,d_k)\longrightarrow0
 \quad\Longleftrightarrow\quad
 t_k\longrightarrow\infty
 \ \text{and}\
 \frac{\log(1+d_k^2/t_k)}{L_k}\longrightarrow0.
\]
Moreover,
\[
 \rho(t_k,d_k)\asymp t_k^{-1}
 \quad\Longleftrightarrow\quad
 \bigl(\exists a>0:\ t_k\ge a\ \text{eventually}\bigr)
 \ \text{and}\
 \bigl(\exists M\in\mathbb R:\ \min\{t_k,d_k\}\le M\
       \text{eventually}\bigr).
\]
Here \(\asymp\) means two-sided comparison by fixed positive finite constants for all sufficiently large \(k\). If \(t_k\to\infty\), this equivalence becomes
\[
 \rho(t_k,d_k)\asymp t_k^{-1}
 \quad\Longleftrightarrow\quad
 \exists M\in\mathbb N:\ d_k\le M\ \text{eventually}.
\]
For every pair of real numbers \(0<a\le b\), there exist real constants \(0<c\le C_0<\infty\) and an integer \(D_0\) such that, for every integer \(d\ge D_0\) with \(d\ge2\) and every real \(t\) satisfying \(a\le t/d^2\le b\),
\[
 c\left(\frac{\log L}{L}\right)^2
 \le \rho(t,d)
 \le C_0\left(\frac{\log L}{L}\right)^2.
\]

Finally, these sequence properties transfer to the actual minimax values for each fixed \(\mathcal C\in\{\mathrm{NI},\mathrm{SI}\}\). Along every admissible sequence,
\[
\begin{aligned}
 \mathfrak R^{\mathrm{TV}}_{\mathcal C}
   (\mathcal P_{d_k}^{\mathrm{pair}};n_k,d_k,\varepsilon_k)
   \longrightarrow0
 &\quad\Longleftrightarrow\quad
 t_k\longrightarrow\infty
 \ \text{and}\\
 \frac{\log(1+d_k^2/t_k)}{L_k}\longrightarrow0,\\
 \mathfrak H^{\mathrm{TV}}_{\mathcal C}
   (\mathcal P_{d_k}^{\mathrm{pair}};n_k,d_k,\varepsilon_k)
   \longrightarrow0
 &\quad\Longleftrightarrow\quad
 t_k\longrightarrow\infty
 \ \text{and}\\
 \frac{\log(1+d_k^2/t_k)}{L_k}\longrightarrow0.
\end{aligned}
\]
The squared-error value satisfies
\[
 \mathfrak R^{\mathrm{TV}}_{\mathcal C}
   (\mathcal P_{d_k}^{\mathrm{pair}};n_k,d_k,\varepsilon_k)
   \asymp t_k^{-1}
 \quad\Longleftrightarrow\quad
 \bigl(\exists a>0:\ t_k\ge a\ \text{eventually}\bigr)
 \ \text{and}\
 \bigl(\exists M\in\mathbb R:\ \min\{t_k,d_k\}\le M\
       \text{eventually}\bigr).
\]
Under \(t_k\to\infty\), equivalently,
\[
 \mathfrak R^{\mathrm{TV}}_{\mathcal C}
   (\mathcal P_{d_k}^{\mathrm{pair}};n_k,d_k,\varepsilon_k)
   \asymp t_k^{-1}
 \quad\Longleftrightarrow\quad
 \exists M\in\mathbb N:\ d_k\le M\ \text{eventually}.
\]
For each fixed \(\mathcal C\) and every \(0<a\le b\), there exist real constants \(0<c\le C_0<\infty\) and an integer \(D_0\) such that every admissible tuple with \(d\ge D_0\) and \(a\le n\varepsilon^2/d^2\le b\) satisfies
\[
\begin{aligned}
 c\left(\frac{\log L}{L}\right)^2
 &\le
 \mathfrak R^{\mathrm{TV}}_{\mathcal C}
   (\mathcal P_d^{\mathrm{pair}};n,d,\varepsilon)
 \le C_0\left(\frac{\log L}{L}\right)^2,\\
 \sqrt c\left|\frac{\log L}{L}\right|
 &\le
 \mathfrak H^{\mathrm{TV}}_{\mathcal C}
   (\mathcal P_d^{\mathrm{pair}};n,d,\varepsilon)
 \le \sqrt{C_0}\left|\frac{\log L}{L}\right|.
\end{aligned}
\]
\end{theoremv}
% CAUSALSMITH-CITED-SCOPE-BEGIN thm:paired-uniform-tv-frontier
\verificationfootnotetext{\textbf{Formalization scope.} This result uses the cited conclusion \citep{CaiLow2011Nonsmooth} (Section 3.1, definition of the Bernstein constant and immediately following paragraph, page 8). The cited source proof is not formalized here: Lean verifies its use and all remaining steps. If that cited conclusion is false, the portions of this result that depend on it are not certified.}
% CAUSALSMITH-CITED-SCOPE-END thm:paired-uniform-tv-frontier

The identity in \cref{obj:thm:paired-uniform-tv-frontier} explains the common precision benchmark: departure from uniformity within a pair contributes the absolute value of its signed contrast. Thus the scalar total-variation target retains the same absolute-value structure that governs welfare in \cref{obj:prop:signed-causal-bridge}. The concrete paired protocol, estimator, and interval in \cref{obj:synth_5} use the same private-moment and polynomial-approximation ingredients as the welfare construction in \cref{obj:def:frontier-handle}. They attain both the squared-error order and the square-root order for globally honest expected interval length.

The resource characterizations describe when these guarantees yield increasing precision. Both decision values vanish exactly when effective private resource diverges and the normalized logarithmic resource deficit vanishes. Along sequences with diverging effective private resource, the usual inverse-resource squared-error order is equivalent to an eventually bounded dimension. When effective private resource is proportional to the square of the dimension, the theorem instead identifies squared-error order \((\log L/L)^2\) and honest-length order \(|\log L/L|\), where \(L=\log(ed)\) is the log-dimension factor.

The interaction comparison is uniform over admissible sample sizes, dimensions, and privacy budgets. Within the one-person-one-message classes of \cref{obj:def:tv-private-decision-values}, sequential adaptation can improve the minimax criteria by at most a universal multiplicative factor. A noninteractive procedure already attains their common orders. This comparison concerns adaptation to preceding participants' messages under the same pure local privacy requirement and with each participant accessed once.

\subsection{A converse on the full simplex}

The paired family is a subset of the full simplex in \cref{obj:def:tv-corollary-models}. Consequently, estimation over the simplex includes the paired-family problem, and global simplex coverage includes coverage over the paired family. With the same alphabet and protocol restrictions, this inclusion transfers the paired lower bounds to arbitrary probability laws. The approximation input remains the Bernstein limit cited in \citet[Section 3.1, p.~8]{CaiLow2011Nonsmooth}.

\begin{theoremv}{thm:full-simplex-tv-converse}[Full-simplex total-variation lower bounds]*
This result is conditional on the cited, unformalized Bernstein absolute-approximation limit. If \(e_K=\inf_{\deg p\le K}\sup_{|x|\le1}\lvert |x|-p(x)\rvert\) is the best uniform error for approximating \(|x|\) on \([-1,1]\), then \(2k e_{2k}\to\beta_*\) for a constant \(0.278<\beta_*<0.286\). See \citet[Section 3.1]{CaiLow2011Nonsmooth} and \citet{Bernstein1913}.

Let \(\Delta_{2d}\) be the probability simplex on the paired alphabet \(\mathcal A_d=[d]\times\{-1,1\}\), and let
\[
D_{\mathrm{TV}}(\mathsf p)=\operatorname{TV}(\mathsf p,\mathsf U_{2d}),
\]
where \(\mathsf U_{2d}\) is the uniform reference law. There is a universal constant \(c>0\) such that, for every sample size \(n\), dimension \(d\), and privacy budget \(\varepsilon\) satisfying
\begin{itemize}
\item \textbf{(Sample size.)} \(n\ge2\).
\item \textbf{(Dimension.)} \(d\ge2\).
\item \textbf{(Privacy budget.)} \(0<\varepsilon\le1\).
\end{itemize}
and every protocol-class label \(\mathcal C\in\{\mathrm{NI},\mathrm{SI}\}\), the private decision values of \cref{obj:def:tv-private-decision-values}, with iid inputs drawn from \(\mathsf p\in\Delta_{2d}\), satisfy
\[
\begin{aligned}
\mathfrak R^{\mathrm{TV}}_{\mathcal C}(\Delta_{2d};n,d,\varepsilon)
&\ge c\,r_{\mathcal C}(n,d,\varepsilon),\\
\mathfrak H^{\mathrm{TV}}_{\mathcal C}(\Delta_{2d};n,d,\varepsilon)
&\ge c\,h_{\mathcal C}(n,d,\varepsilon),
\end{aligned}
\]
where the squared-error and connected-length rates are
\[
r_{\mathcal C}(n,d,\varepsilon)=\rho(n\varepsilon^2,d),
\qquad
h_{\mathcal C}(n,d,\varepsilon)=\sqrt{r_{\mathcal C}(n,d,\varepsilon)}.
\]
The interval value uses connected closed intervals in \([0,1]\) with coverage at least \(0.90\) at every law in \(\Delta_{2d}\). The protocol classes permit arbitrary public seeds, participant-specific kernels, standard-Borel message spaces, and, for the sequential class, dependence on preceding messages, with one message per participant and pure local privacy.
\end{theoremv}
% CAUSALSMITH-CITED-SCOPE-BEGIN thm:full-simplex-tv-converse
\verificationfootnotetext{\textbf{Formalization scope.} This result uses the cited conclusion \citep{CaiLow2011Nonsmooth} (Section 3.1, definition of the Bernstein constant and immediately following paragraph, page 8). The cited source proof is not formalized here: Lean verifies its use and all remaining steps. If that cited conclusion is false, the portions of this result that depend on it are not certified.}
% CAUSALSMITH-CITED-SCOPE-END thm:full-simplex-tv-converse

The intuition behind \cref{obj:thm:full-simplex-tv-converse} is that a larger model retains the difficult sampling laws contained in its paired subfamily. Allowing arbitrary pair masses therefore preserves the lower bounds for the minimax squared-error value \leanref{sym:\mathfrak R^{\mathrm{TV}}}{\ensuremath{\mathfrak R^{\mathrm{TV}}_{\mathcal C}}} and the globally honest expected-length value \leanref{sym:\mathfrak H^{\mathrm{TV}}}{\ensuremath{\mathfrak H^{\mathrm{TV}}_{\mathcal C}}}. The result supplies necessary precision bounds on the simplex, while \cref{obj:thm:paired-uniform-tv-frontier} supplies a matching characterization on the paired family.

\subsection{Interpretation for welfare reporting}

The causal result in \cref{obj:thm:uniform-private-value-frontiers} gives a precision benchmark for reporting a fixed population's first-best welfare under the model in \cref{obj:def:causal-model}. The population law determines the target, and independent participant sampling and private messages determine the uncertainty in its estimate. Squared-error risk measures the accuracy of the reported welfare value; global honesty requires coverage across the declared population laws, with expected connected interval length measuring reporting precision.

The signed bridge in \cref{obj:prop:signed-causal-bridge} explains why the paired statistical problem is informative for this interpretation. Both analyses involve the average absolute stratum contrast, and \cref{obj:thm:paired-uniform-tv-frontier} isolates that component as a distance from a known uniform law. The welfare and paired-family interaction comparisons in \cref{obj:thm:uniform-private-value-frontiers,obj:thm:paired-uniform-tv-frontier} establish common minimax orders for noninteractive and sequential protocols within their respective declared classes. These conclusions give a sample-size and privacy benchmark for scalar reporting under independent sampling and one message per participant.

\subsection{Limitations and future work}

The causal model uses public uniform stratum probabilities and fair treatment assignment. Unknown stratum probabilities or treatment propensities would require an analysis of their estimation and their contribution to welfare precision under local privacy. Repeated records per participant would likewise require a participant-level privacy specification and a sampling model that accounts for dependence within participants. The present interaction comparison concerns protocols that access each participant once.

The displayed procedure is a conservative witness to uniform rate attainment rather than a calibrated recommendation for applied reporting. Its polynomial branches require \(L=\log(ed)\ge4096\), hence \(d\ge e^{4095}\), and each active participant releases a \(d\)-coordinate sign vector. Its welfare interval uses radius \(\sqrt{10^{17}\rho(t,d)}\); whenever \(\rho(t,d)\ge2.5\times10^{-18}\), clipping makes that interval span the entire welfare range regardless of its center. Sharper constants and communication-efficient procedures are therefore important practical questions even though these conservative choices establish the stated uniform rates.

The matching bounds characterize minimax orders up to universal constants. Sharp leading constants and procedures optimized for those constants remain topics for further work. Learned-policy regret is a separate target from estimation of the population's first-best welfare: it would require evaluating the welfare achieved by a policy selected from private data. Design-based inference for a finite population would require uncertainty and coverage statements based on the treatment-assignment mechanism with the population held fixed, rather than the independent population sampling used here.

Finally, \cref{obj:thm:full-simplex-tv-converse} supplies lower bounds for total-variation estimation and honest intervals on the full simplex. Matching full-simplex upper bounds remain an open question; the attaining guarantees in \cref{obj:thm:paired-uniform-tv-frontier} apply to the paired family.


\appendix

\section*{Appendices}

\section{Approximation and private-moment bounds}

The attaining procedure in \cref{obj:def:frontier-handle} combines polynomial estimation near zero with a signed mean estimate for larger treatment contrasts. Its analysis separates three sources of error: approximation of absolute value, variance of the estimated polynomial, and mistakes in selecting a branch from the pilot sample. The results below supply the corresponding bounds, together with the moment-matching construction used in the converse.

\subsection{Concentration and approximation inputs}

Independent bounded observations provide the concentration control needed for pilot selection. The following formulation records both tails and includes the degenerate bound explicitly. For a positive bound on the summands, it is the bounded-variable inequality of \citet[Theorem~2]{Hoeffding1963Bounded}.

\begin{theoremv}{lem:bounded-mean-concentration}[Bounded mean concentration]
Let \(m\ge1\) be an integer, let \(b\in\mathbb R\), and let \(H_1,\ldots,H_m\) be mutually independent, measurable real random variables on a common probability space. Suppose that
\[
-b\le H_i\le b
\]
at every sample point, for each \(i=1,\ldots,m\). Then, for every real \(x\ge0\),
\[
\begin{aligned}
\Pr\left\{m^{-1}\sum_{i=1}^m(H_i-\mathbb E H_i)\ge x\right\}
&\le
\begin{cases}
\exp\!\left(-\dfrac{mx^2}{2b^2}\right),& b\ne0,\\
1,& b=0,
\end{cases}\\
\Pr\left\{m^{-1}\sum_{i=1}^m(H_i-\mathbb E H_i)\le -x\right\}
&\le
\begin{cases}
\exp\!\left(-\dfrac{mx^2}{2b^2}\right),& b\ne0,\\
1,& b=0.
\end{cases}
\end{aligned}
\]
\end{theoremv}

The exponential bounds quantify the precision of a block mean through its sample size and the magnitude of its summands. In the private construction, those summands are scaled sign messages, so their magnitude incorporates the privacy attenuation. Both tails matter because pilot selection uses the absolute value of a coordinate mean.

Polynomial estimation addresses the nondifferentiability of absolute value at zero. The next result specifies an explicit approximation and controls its monomial coefficients. The approximation error governs bias, while the coefficient bound enters the variance calibration when powers are replaced by estimated moments.

\begin{theoremv}{lem:cai-low-chebyshev-approximation}[Chebyshev approximation of absolute value]
Let \(D\ge 2\) be an even integer. Define the Chebyshev polynomials by
\[
\mathcal T_0(x)=1,\qquad \mathcal T_1(x)=x,\qquad
\mathcal T_{v+1}(x)=2x\mathcal T_v(x)-\mathcal T_{v-1}(x)
\quad (v\ge 1),
\]
and define the polynomial \(q_D\) and its monomial coefficients \(a_{Dv}\) by
\[
q_D(x)=\frac{2}{\pi}
+\frac{4}{\pi}\sum_{v=1}^{D/2}
\frac{(-1)^{v+1}\mathcal T_{2v}(x)}{4v^2-1}
=\sum_{v=0}^{D}a_{Dv}x^v.
\]
Then, for every \(x\in[-1,1]\),
\[
\bigl|q_D(x)-|x|\bigr|\le \frac{2}{\pi(D+1)},
\]
and, for every integer \(v\) with \(0\le v\le D\),
\[
|a_{Dv}|\le 2^{3D/2}.
\]
\end{theoremv}

Increasing the degree reduces the uniform approximation error at an inverse-degree rate. The accompanying coefficient bound makes the statistical cost of that choice explicit: estimating a higher-degree polynomial requires controlling a larger collection of weighted moments. The finite calibration below balances these two effects.

The converse uses the complementary approximation principle. Two symmetric distributions can agree on every moment through a prescribed degree while retaining a separation in their expected absolute values. The following statement expresses that separation in terms of the best uniform approximation error.

\begin{theoremv}{lem:cai-low-moment-duality}[Symmetric moment matching duality]
For every positive even integer \(K\), define the best uniform polynomial approximation error
\[
e_K=\inf_{\deg p\le K}\sup_{|x|\le 1}\bigl|p(x)-|x|\bigr|,
\]
where the infimum is over real polynomials \(p\). There exist probability measures \(\xi_0,\xi_1\) on \(\mathbb R\), each supported on \([-1,1]\) and invariant under reflection \(x\mapsto -x\), such that
\[
\int x^\ell\,d\xi_1(x)=\int x^\ell\,d\xi_0(x)
\qquad\text{for every integer }0\le \ell\le K,
\]
and
\[
\int |x|\,d\xi_1(x)-\int |x|\,d\xi_0(x)=2e_K.
\]
\end{theoremv}

Moment agreement makes polynomial expectations through the stated degree identical under the two measures, whereas the absolute-value expectations differ by twice the approximation error. Thus the same approximation problem supplies both an estimation device and a quantitative separation for moment-matched alternatives. Reflection invariance aligns this construction with the symmetric contrast distributions used in the converse.

\subsection{Private moments and dependence across strata}

The signed-vector channel \leanref{sym:Q^{\mathrm{vec}}}{\ensuremath{Q^{\mathrm{vec}}}} in \cref{obj:def:vector-kernel} releases one vector per participant. Independent participant sampling supplies the distinct observations needed for the moment estimates in \cref{obj:def:private-moments}; coordinates within a released vector share the same participant record. The next result establishes privacy, the required moment identities, and a dependence bound for arbitrary measurable column statistics with the displayed variances.

\begin{lemmav}{lem:private-moment-and-dependence}[Private moments and column dependence]
Let \(n,d,m\) be nonnegative integers and let \(\varepsilon\) be the privacy budget. Suppose that:
\begin{itemize}
\item \textbf{(Sampling.)} The sampling scheme satisfies \cref{obj:ass:iid-people}.
\item \textbf{(Population.)} The full-data law \(P\) belongs to \(\mathcal V_d\), as defined in \cref{obj:def:causal-model}.
\item \textbf{(Privacy calibration.)} \(0<\varepsilon\le1\).
\item \textbf{(Dimension and block size.)} \(d\ge2\) and \(0<m\le n\).
\end{itemize}
Apply the vector channel \(Q^{\mathrm{vec}}\) once per participant, with a singleton public seed. Define the privacy attenuation and scaled-message magnitude by
\[
\delta=\tanh(\varepsilon/2),
\qquad
b=\frac d\delta.
\]
For an \(m\)-participant message block, write \(Z_{ij}\in\{-1,1\}\) for coordinate \(j\) of participant \(i\)'s sign-vector message and define the scaled messages by
\[
W_{ij}=bZ_{ij},
\qquad 1\le i\le m,\quad j\in[d].
\]
Write \(\tau_j(P)\) for the cellwise treatment contrast,
\[
\tau_j(P)=\mu_{1j}(P)-\mu_{0j}(P).
\]
All expectations, variances, and covariances below are taken under the block message law induced by \(P\) and \(Q^{\mathrm{vec}}\).

The protocol satisfies
\[
Q^{\mathrm{vec}}\in\mathfrak Q_{\mathrm{NI}}(n,d,\varepsilon),
\]
where membership has the full-record privacy and noninteraction requirements of \cref{obj:def:noninteractive-class}. For every \(1\le i\le m\) and \(j\in[d]\),
\[
\mathbb E[W_{ij}]=\tau_j(P),
\qquad
\mathbb E[W_{ij}^{2}]=b^{2}.
\]
For every \(1\le i\le m\) and distinct \(j,j'\in[d]\),
\[
\mathbb E[W_{ij}W_{ij'}]=0.
\]
The private moments \(\mathsf U_{kj}\) of \(W\), defined in \cref{obj:def:private-moments}, satisfy, for every integer \(0\le k\le m\) and every \(j\in[d]\),
\[
\mathbb E[\mathsf U_{kj}]=\tau_j(P)^k.
\]
For every nonnegative integer \(k\) with \(2k\le m\) and every \(j\in[d]\),
\[
\mathbb E[\mathsf U_{kj}^{2}]
\le
\left(\tau_j(P)^2+\frac{2kb^2}{m}\right)^k.
\]

For any deterministic measurable column statistics \(T_j^{\mathrm{col}}\), each a function of \((W_{1j},\ldots,W_{mj})\), every pair of distinct coordinates \(j,j'\in[d]\) satisfies
\[
\left|\operatorname{Cov}
\left(T_j^{\mathrm{col}},T_{j'}^{\mathrm{col}}\right)\right|
\le
\frac{|\tau_j(P)\tau_{j'}(P)|}
{\sqrt{(b^2-\tau_j(P)^2)(b^2-\tau_{j'}(P)^2)}}
\sqrt{\operatorname{Var}(T_j^{\mathrm{col}})
      \operatorname{Var}(T_{j'}^{\mathrm{col}})}.
\]
The same collection satisfies
\[
\operatorname{Var}\left(\frac1d\sum_{j=1}^{d}T_j^{\mathrm{col}}\right)
\le
\frac2d\max_{j\in[d]}\operatorname{Var}(T_j^{\mathrm{col}}).
\]
\end{lemmav}

\begin{proof}[Proof of \cref{obj:lem:private-moment-and-dependence}]
Write
\[
[m]=\{1,\ldots,m\},
\qquad
S=(2A-1)(2Y-1).
\]
We use the indicator convention
\[
\mathbf1_{\{\text{condition}\}}
=
\begin{cases}
1,&\text{if the condition holds},\\
0,&\text{otherwise},
\end{cases}
\]
and take empty products and zeroth powers to equal one, including \(0^0=1\).
By \cref{obj:ass:iid-people}, the singleton seed and independent local draws give independent, identically distributed message rows in the block. All the statistics considered below have finite moments because the block message alphabet is finite.

\begin{enumerate}
\item \textit{Single-coordinate means.}
By \cref{obj:ass:consistency},
\[
\mathbf1_{\{X=j\}}S
=
2\mathbf1_{\{X=j,A=1,Y^1=1\}}
-\mathbf1_{\{X=j,A=1\}}
-2\mathbf1_{\{X=j,A=0,Y^0=1\}}
+\mathbf1_{\{X=j,A=0\}}
\quad P\text{-almost surely}.
\]
Fair assignment in \cref{obj:ass:fair-randomization} supplies, for \(a\in\{0,1\}\),
\[
P(X=j,A=a,Y^a=1)=\frac12P(X=j,Y^a=1),
\qquad
P(X=j,A=a)=\frac12P(X=j).
\]
Taking expectations in the preceding identity and using
\cref{obj:ass:uniform-covariates} therefore gives
\[
\mathbb E_P[\mathbf1_{\{X=j\}}S]
=
P(X=j,Y^1=1)-P(X=j,Y^0=1)
=
\frac{\mu_{1j}(P)-\mu_{0j}(P)}d
=
\frac{\tau_j(P)}d.
\]
The sampler in \cref{obj:def:vector-kernel} gives
\[
\mathbb E[Z_{ij}\mid O]
=
\begin{cases}
\delta S,&X=j,\\
0,&X\ne j.
\end{cases}
\]
Consequently, for every \(i\in[m]\) and \(j\in[d]\),
\[
\mathbb E[W_{ij}]
=
b\delta\,\mathbb E_P[\mathbf1_{\{X=j\}}S]
=
\frac d\delta\,\delta\,\frac{\tau_j(P)}d
=
\tau_j(P).
\]
% lean: CausalSmith.Stat.LdpOptvalueUniformFrontier.vectorMessageLaw_scaled_sign_integral

\item \textit{Protocol membership.}
The exponential formula for the hyperbolic tangent yields
\[
\delta=\frac{e^\varepsilon-1}{e^\varepsilon+1},
\qquad
0<\delta<1,
\qquad
1+\delta=e^\varepsilon(1-\delta).
\]
For any observed records \(o,o'\) and any
\(z\in\{-1,1\}^d\), the masses in
\cref{obj:def:vector-kernel} satisfy
\[
2^{-d}(1-\delta)
\le Q^{\mathrm{vec}}(z\mid o)
\le 2^{-d}(1+\delta).
\]
It follows that
\[
Q^{\mathrm{vec}}(z\mid o)
\le
2^{-d}(1+\delta)
=
e^\varepsilon2^{-d}(1-\delta)
\le
e^\varepsilon Q^{\mathrm{vec}}(z\mid o').
\]
Summing over any message event \(E\subseteq\{-1,1\}^d\) proves
\[
Q^{\mathrm{vec}}(E\mid o)
\le e^\varepsilon Q^{\mathrm{vec}}(E\mid o').
\]
The participant kernels are constant in message history. Thus they satisfy
\cref{obj:ass:full-record-privacy,obj:ass:noninteraction}, and
\cref{obj:def:noninteractive-class} gives
\[
Q^{\mathrm{vec}}\in\mathfrak Q_{\mathrm{NI}}(n,d,\varepsilon).
\]
% lean: CausalSmith.Stat.LdpOptvalueUniformFrontier.vectorProtocol_noninteractive

\item \textit{Second and cross-coordinate row moments.}
Since \(Z_{ij}\in\{-1,1\}\),
\[
W_{ij}^2=b^2,
\qquad
\mathbb E[W_{ij}^2]=b^2.
\]
For distinct \(j,j'\in[d]\), flipping coordinate \(j\) pairs opposite terms in the first sum below. Flipping every coordinate pairs opposite terms in the second sum, for every \(\ell\in[d]\):
\[
\sum_{z\in\{-1,1\}^d}z_jz_{j'}=0,
\qquad
\sum_{z\in\{-1,1\}^d}z_jz_{j'}z_\ell=0.
\]
Expanding the kernel mass from \cref{obj:def:vector-kernel} hence gives
\[
\begin{aligned}
\mathbb E[Z_{ij}Z_{ij'}\mid O]
&=
2^{-d}\sum_{z\in\{-1,1\}^d}z_jz_{j'}
+
2^{-d}\delta S
\sum_{z\in\{-1,1\}^d}z_jz_{j'}z_X\\
&=0.
\end{aligned}
\]
Averaging over the record and scaling proves
\[
\mathbb E[W_{ij}W_{ij'}]=0
\qquad(j\ne j').
\]
% lean: CausalSmith.Stat.LdpOptvalueUniformFrontier.blockMean_scaledMessages_cross

\item \textit{Unbiased private moments.}
For each integer \(0\le k\le m\), define the family of participant subsets
\[
\mathcal J_k=\{J\subseteq[m]:|J|=k\}.
\]
Independent rows imply, for every \(J\subseteq[m]\),
\[
\mathbb E\left[\prod_{i\in J}W_{ij}\right]
=
\prod_{i\in J}\mathbb E[W_{ij}]
=
\tau_j(P)^{|J|}.
\]
There are \(\binom mk>0\) members of \(\mathcal J_k\). Applying
\cref{obj:def:private-moments} therefore gives
\[
\mathbb E[\mathsf U_{kj}]
=
\binom mk^{-1}\sum_{J\in\mathcal J_k}\tau_j(P)^k
=
\tau_j(P)^k.
\]
For \(k=0\), the family consists of the empty subset and this is the identity
\(\mathbb E[\mathsf U_{0j}]=1\).
% lean: CausalSmith.Stat.LdpOptvalueUniformFrontier.blockMean_privateMoment_eq_rowMean_pow

\item \textit{The second-moment bound at degrees zero and one.}
At \(k=0\), both sides of the claimed second-moment bound equal one.
At \(k=1\), \cref{obj:def:private-moments} gives
\[
\mathsf U_{1j}=\frac1m\sum_{i=1}^mW_{ij}.
\]
The diagonal row terms have expectation \(b^2\), while independence gives
\(\mathbb E[W_{ij}W_{i'j}]=\tau_j(P)^2\) for \(i\ne i'\).
Thus
\[
\begin{aligned}
\mathbb E[\mathsf U_{1j}^2]
&=
\frac{mb^2+m(m-1)\tau_j(P)^2}{m^2}\\
&=
\tau_j(P)^2+\frac{b^2-\tau_j(P)^2}{m}\\
&\le
\tau_j(P)^2+\frac{2b^2}{m},
\end{aligned}
\]
where the last inequality follows from nonnegativity of the squares.
This establishes the required bound in the degree-one case.
% lean: CausalSmith.Stat.LdpOptvalueUniformFrontier.blockMean_privateMoment_sq_le

\item \textit{The second-moment bound at higher degrees.}
Suppose now that \(k\ge2\) and \(2k\le m\). For
\(J,J'\in\mathcal J_k\), each row in \(J\cap J'\) contributes a second moment,
each row in either set difference contributes a first moment, and every
remaining row contributes one. Independence therefore gives
\[
\begin{aligned}
\mathbb E\left[
\left(\prod_{i\in J}W_{ij}\right)
\left(\prod_{i\in J'}W_{ij}\right)\right]
&=
b^{2|J\cap J'|}
\tau_j(P)^{|J\setminus J'|+|J'\setminus J|}\\
&=
b^{2|J\cap J'|}
\bigl(\tau_j(P)^2\bigr)^{k-|J\cap J'|}.
\end{aligned}
\]
Expanding the square of the subset average gives the exact identity
\[
\mathbb E[\mathsf U_{kj}^2]
=
\binom mk^{-2}
\sum_{J\in\mathcal J_k}\sum_{J'\in\mathcal J_k}
b^{2|J\cap J'|}
\bigl(\tau_j(P)^2\bigr)^{k-|J\cap J'|}.
\]

We first establish the inclusion count bound
\[
\binom{m-v}{k-v}
\le
\binom mk\left(\frac{2k}{m}\right)^v
\qquad(0\le v\le k).
\]
At \(v=0\) this is equality. For \(0\le v<k\), the assumptions give
\(m-v\ge m/2>0\), and hence
\[
\frac{k-v}{m-v}\le\frac{2k}{m}.
\]
The binomial recurrence and induction then yield
\[
\begin{aligned}
\binom{m-(v+1)}{k-(v+1)}
&=
\binom{m-v}{k-v}\frac{k-v}{m-v}\\
&\le
\binom mk\left(\frac{2k}{m}\right)^{v+1}.
\end{aligned}
\]
In particular, for any \(J_{\mathrm{sub}}\subseteq[m]\) with
\(|J_{\mathrm{sub}}|\le k\),
\[
\#\{J'\in\mathcal J_k:J_{\mathrm{sub}}\subseteq J'\}
=
\binom{m-|J_{\mathrm{sub}}|}{k-|J_{\mathrm{sub}}|}
\le
\binom mk
\left(\frac{2k}{m}\right)^{|J_{\mathrm{sub}}|}.
\]

Fix \(J\in\mathcal J_k\). All moment weights are nonnegative, and the
subset \(J_{\mathrm{sub}}=J\cap J'\) is one of the terms in
\[
b^{2|J\cap J'|}
\bigl(\tau_j(P)^2\bigr)^{k-|J\cap J'|}
\le
\sum_{J_{\mathrm{sub}}\subseteq J\cap J'}
b^{2|J_{\mathrm{sub}}|}
\bigl(\tau_j(P)^2\bigr)^{k-|J_{\mathrm{sub}}|}.
\]
Summing over \(J'\), interchanging the finite sums, and applying the
inclusion count bound gives
\[
\begin{aligned}
&\sum_{J'\in\mathcal J_k}
b^{2|J\cap J'|}
\bigl(\tau_j(P)^2\bigr)^{k-|J\cap J'|}\\
&\quad\le
\sum_{J_{\mathrm{sub}}\subseteq J}
b^{2|J_{\mathrm{sub}}|}
\bigl(\tau_j(P)^2\bigr)^{k-|J_{\mathrm{sub}}|}
\#\{J'\in\mathcal J_k:J_{\mathrm{sub}}\subseteq J'\}\\
&\quad\le
\binom mk
\sum_{J_{\mathrm{sub}}\subseteq J}
\left(\frac{2kb^2}{m}\right)^{|J_{\mathrm{sub}}|}
\bigl(\tau_j(P)^2\bigr)^{k-|J_{\mathrm{sub}}|}\\
&\quad=
\binom mk
\sum_{v=0}^k\binom kv
\left(\frac{2kb^2}{m}\right)^v
\bigl(\tau_j(P)^2\bigr)^{k-v}\\
&\quad=
\binom mk
\left(\tau_j(P)^2+\frac{2kb^2}{m}\right)^k.
\end{aligned}
\]
The penultimate equality counts the \(v\)-subsets of \(J\), and the final
equality is the binomial expansion. Averaging this row bound over
\(J\in\mathcal J_k\) proves
\[
\begin{aligned}
\mathbb E[\mathsf U_{kj}^2]
&\le
\binom mk^{-2}\binom mk\binom mk
\left(\tau_j(P)^2+\frac{2kb^2}{m}\right)^k\\
&=
\left(\tau_j(P)^2+\frac{2kb^2}{m}\right)^k.
\end{aligned}
\]
% lean: CausalSmith.Stat.LdpOptvalueUniformFrontier.blockMean_privateMoment_sq_le

\item \textit{Finite orthonormal column bases.}
By \cref{obj:ass:interior-means},
\[
-\frac12\le\tau_j(P)\le\frac12,
\qquad
\tau_j(P)^2\le\frac14.
\]
Since \(0<\delta<1\) and \(d\ge2\),
\[
b=\frac d\delta>d\ge2,
\qquad
b^2-\tau_j(P)^2>0.
\]
For \(i\in[m]\), \(j\in[d]\), and \(J\subseteq[m]\), define
\[
\psi^{\mathrm{row}}_{ij}
=
\frac{W_{ij}-\tau_j(P)}
{\sqrt{b^2-\tau_j(P)^2}},
\qquad
\psi^{\mathrm{col}}_{Jj}
=
\prod_{i\in J}\psi^{\mathrm{row}}_{ij}.
\]
In particular, \(\psi^{\mathrm{col}}_{\varnothing j}=1\).
The row moments already established give
\[
\mathbb E[W_{ij}-\tau_j(P)]=0,
\]
\[
\mathbb E[(W_{ij}-\tau_j(P))^2]
=
b^2-2\tau_j(P)\mathbb E[W_{ij}]+\tau_j(P)^2
=
b^2-\tau_j(P)^2,
\]
and, for \(j\ne j'\),
\[
\begin{aligned}
&\mathbb E[(W_{ij}-\tau_j(P))(W_{ij'}-\tau_{j'}(P))]\\
&\quad=
\mathbb E[W_{ij}W_{ij'}]
-\tau_{j'}(P)\mathbb E[W_{ij}]
-\tau_j(P)\mathbb E[W_{ij'}]
+\tau_j(P)\tau_{j'}(P)\\
&\quad=-\tau_j(P)\tau_{j'}(P).
\end{aligned}
\]
Thus
\[
\mathbb E[\psi^{\mathrm{row}}_{ij}]=0,
\qquad
\mathbb E[(\psi^{\mathrm{row}}_{ij})^2]=1.
\]

If \(J\ne J'\), a row in \(J\setminus J'\) or \(J'\setminus J\)
contributes a single centered factor. Factoring over independent rows
therefore gives, for any \(j,j'\in[d]\),
\[
\mathbb E[
\psi^{\mathrm{col}}_{Jj}\psi^{\mathrm{col}}_{J'j'}]=0.
\]
For equal subsets in the same column, row independence instead gives
\[
\mathbb E[(\psi^{\mathrm{col}}_{Jj})^2]
=
\prod_{i\in J}\mathbb E[(\psi^{\mathrm{row}}_{ij})^2]
=
1.
\]
Consequently,
\[
\mathbb E[
\psi^{\mathrm{col}}_{Jj}\psi^{\mathrm{col}}_{J'j}]
=
\begin{cases}
1,&J=J',\\
0,&J\ne J',
\end{cases}
\qquad
\mathbb E[\psi^{\mathrm{col}}_{Jj}]
=
\begin{cases}
1,&J=\varnothing,\\
0,&J\ne\varnothing.
\end{cases}
\]
For each \(j\), these functions of the column are linearly independent by
orthonormality. Their number equals the dimension of the real function
space on the column patterns:
\[
\#\{J:J\subseteq[m]\}
=
2^m
=
\dim_{\mathbb R}\mathbb R^{\{-b,b\}^m}.
\]
They therefore form a basis of that function space.
% lean: CausalSmith.Stat.LdpOptvalueUniformFrontier.booleanColumnWalshBasis

\item \textit{Centered column expansions.}
Fix the given collection of deterministic measurable statistics.
For every \(j\in[d]\) and \(J\subseteq[m]\), define
\[
\kappa^{\mathrm{raw}}_{Jj}
=
\mathbb E[T_j^{\mathrm{col}}\psi^{\mathrm{col}}_{Jj}],
\qquad
\kappa_{Jj}
=
\begin{cases}
0,&J=\varnothing,\\
\kappa^{\mathrm{raw}}_{Jj},&J\ne\varnothing.
\end{cases}
\]
The finite basis and its orthonormality give the pointwise expansion
\[
T_j^{\mathrm{col}}
=
\sum_{J\subseteq[m]}
\kappa^{\mathrm{raw}}_{Jj}\psi^{\mathrm{col}}_{Jj}.
\]
Indeed, multiplying a basis expansion by
\(\psi^{\mathrm{col}}_{Jj}\) and taking expectations identifies its
coefficient with the displayed definition.
Taking expectations of the expansion shows that
\[
\mathbb E[T_j^{\mathrm{col}}]
=
\kappa^{\mathrm{raw}}_{\varnothing j}.
\]
Subtracting this constant deletes precisely the empty-subset coefficient:
\[
T_j^{\mathrm{col}}-\mathbb E[T_j^{\mathrm{col}}]
=
\sum_{J\subseteq[m]}
\kappa_{Jj}\psi^{\mathrm{col}}_{Jj},
\qquad
\kappa_{\varnothing j}=0.
\]
% lean: CausalSmith.Stat.LdpOptvalueUniformFrontier.columnStatistic_centered_walsh_expansion

\item \textit{The dimension bound for row correlation.}
For each distinct pair \(j,j'\in[d]\), define
\[
\gamma_{jj'}
=
-\frac{\tau_j(P)\tau_{j'}(P)}
{\sqrt{(b^2-\tau_j(P)^2)(b^2-\tau_{j'}(P)^2)}}.
\]
Its denominator is positive, and
\[
|\gamma_{jj'}|
=
\frac{|\tau_j(P)\tau_{j'}(P)|}
{\sqrt{(b^2-\tau_j(P)^2)(b^2-\tau_{j'}(P)^2)}}.
\]
For every \(\ell\in[d]\), the bounds above and \(d\ge2\) imply
\[
b^2-\tau_\ell(P)^2
\ge d^2-\frac14
\ge d.
\]
Multiplying the two nonnegative lower bounds and taking square roots gives
\[
(b^2-\tau_j(P)^2)(b^2-\tau_{j'}(P)^2)\ge d^2,
\qquad
\sqrt{(b^2-\tau_j(P)^2)(b^2-\tau_{j'}(P)^2)}\ge d.
\]
Also, \(|\tau_j(P)\tau_{j'}(P)|\le1\), since both contrasts have
absolute value at most \(1/2\). Hence
\[
|\gamma_{jj'}|
\le\frac1d
\le1.
\]
% lean: CausalSmith.Stat.LdpOptvalueUniformFrontier.column_correlation_le_inv_dimension

\item \textit{Covariance of arbitrary column statistics.}
For \(j\ne j'\), the centered row cross moment and positivity of the
variances give
\[
\mathbb E[
\psi^{\mathrm{row}}_{ij}\psi^{\mathrm{row}}_{ij'}]
=
\frac{-\tau_j(P)\tau_{j'}(P)}
{\sqrt{b^2-\tau_j(P)^2}\sqrt{b^2-\tau_{j'}(P)^2}}
=
\gamma_{jj'}.
\]
Factoring over independent rows therefore yields
\[
\mathbb E[
\psi^{\mathrm{col}}_{Jj}\psi^{\mathrm{col}}_{Jj'}]
=
\gamma_{jj'}^{|J|}
\qquad(J\subseteq[m]).
\]
Expanding the centered statistics, discarding unequal-subset terms by
their zero mixed moments, and using the within-column unit moments gives
\[
\operatorname{Cov}(T_j^{\mathrm{col}},T_{j'}^{\mathrm{col}})
=
\sum_{J\subseteq[m]}
\kappa_{Jj}\kappa_{Jj'}\gamma_{jj'}^{|J|},
\qquad
\operatorname{Var}(T_j^{\mathrm{col}})
=
\sum_{J\subseteq[m]}\kappa_{Jj}^2.
\]

Because \(|\gamma_{jj'}|\le1\), induction gives
\[
|\gamma_{jj'}|^{v+1}\le|\gamma_{jj'}|
\qquad(v\ge0):
\]
the degree-one case is equality, and the induction step uses
\[
|\gamma_{jj'}|^{v+2}
\le|\gamma_{jj'}|^2
\le|\gamma_{jj'}|.
\]
For the empty subset the covariance summand vanishes because
\(\kappa_{\varnothing j}=0\). For every nonempty subset the preceding
power bound applies. The triangle inequality and Cauchy--Schwarz thus give
\[
\begin{aligned}
\left|
\operatorname{Cov}(T_j^{\mathrm{col}},T_{j'}^{\mathrm{col}})
\right|
&\le
\sum_{J\subseteq[m]}
\left|\kappa_{Jj}\kappa_{Jj'}\gamma_{jj'}^{|J|}\right|\\
&\le
|\gamma_{jj'}|
\sum_{J\subseteq[m]}|\kappa_{Jj}|\,|\kappa_{Jj'}|\\
&\le
|\gamma_{jj'}|
\sqrt{\left(\sum_{J\subseteq[m]}\kappa_{Jj}^2\right)
      \left(\sum_{J\subseteq[m]}\kappa_{Jj'}^2\right)}\\
&=
\frac{|\tau_j(P)\tau_{j'}(P)|}
{\sqrt{(b^2-\tau_j(P)^2)(b^2-\tau_{j'}(P)^2)}}
\sqrt{\operatorname{Var}(T_j^{\mathrm{col}})
      \operatorname{Var}(T_{j'}^{\mathrm{col}})}.
\end{aligned}
\]
% lean: CausalSmith.Stat.LdpOptvalueUniformFrontier.column_covariance_bound

\item \textit{Variance of the column average.}
Define the finite maximum
\[
v_{\max}
=
\max_{j\in[d]}\operatorname{Var}(T_j^{\mathrm{col}}).
\]
Since \(d\ge2\), the index set is nonempty, and
\[
0\le\operatorname{Var}(T_j^{\mathrm{col}})\le v_{\max},
\qquad
0\le v_{\max}.
\]
It follows that
\[
\sqrt{\operatorname{Var}(T_j^{\mathrm{col}})
      \operatorname{Var}(T_{j'}^{\mathrm{col}})}
\le\sqrt{v_{\max}^2}=v_{\max}.
\]
Combining the covariance bound with
\(|\gamma_{jj'}|\le1/d\) gives, for distinct coordinates,
\[
\operatorname{Cov}(T_j^{\mathrm{col}},T_{j'}^{\mathrm{col}})
\le
\left|\operatorname{Cov}(T_j^{\mathrm{col}},T_{j'}^{\mathrm{col}})\right|
\le
\frac1d
\sqrt{\operatorname{Var}(T_j^{\mathrm{col}})
      \operatorname{Var}(T_{j'}^{\mathrm{col}})}
\le
\frac{v_{\max}}d.
\]
On the diagonal,
\[
\operatorname{Cov}(T_j^{\mathrm{col}},T_j^{\mathrm{col}})
=
\operatorname{Var}(T_j^{\mathrm{col}})
\le v_{\max}+\frac{v_{\max}}d.
\]
Thus each covariance row satisfies
\[
\sum_{j'=1}^d
\operatorname{Cov}(T_j^{\mathrm{col}},T_{j'}^{\mathrm{col}})
\le
v_{\max}+\frac{v_{\max}}d
+(d-1)\frac{v_{\max}}d
=
2v_{\max}.
\]
Finally, the finite covariance expansion of the average gives
\[
\begin{aligned}
\operatorname{Var}\left(\frac1d\sum_{j=1}^dT_j^{\mathrm{col}}\right)
&=
\frac1{d^2}\sum_{j=1}^d\sum_{j'=1}^d
\operatorname{Cov}(T_j^{\mathrm{col}},T_{j'}^{\mathrm{col}})\\
&\le
\frac1{d^2}\sum_{j=1}^d2v_{\max}\\
&=
\frac2d\,v_{\max}.
\end{aligned}
\]
% lean: CausalSmith.Stat.LdpOptvalueUniformFrontier.blockVar_average_le_of_covariance
\end{enumerate}
\end{proof}

The sampling and population conditions connect the released coordinate means to the treatment contrasts. Privacy calibration determines the scaled-message magnitude, while the positive block size supplies the observations used by the distinct-person moments. Their unbiasedness permits a polynomial in a contrast to be estimated by replacing each power with its corresponding private moment. The restriction \(2k\le m\) accompanies the second-moment bound needed to control the variance of this substitution.

For the column statistics \leanref{sym:T^{\mathrm{col}}}{\ensuremath{T_j^{\mathrm{col}}}}, the covariance inequality explicitly accounts for dependence across strata. Its aggregate consequence gives a factor of \(1/d\) in the variance of their average, up to the stated factor of two. This control applies to nonlinear column corrections and supplies the averaging bound required by the polynomial procedure.

\subsection{Global and pilot-selected calibration}

The final result combines the preceding inputs into two finite-sample calibrations. A fixed approximation radius covers the entire contrast range supplied by \cref{obj:def:causal-model}. A second calibration uses independent pilot and evaluation blocks: the pilot selects between polynomial estimation near zero and a signed evaluation mean for larger contrasts. The concentration and approximation conditions are stated explicitly so that each calibration's resources and error bounds remain visible.

\begin{lemmav}{lem:finite-private-polynomial-calibration}[Finite private polynomial calibration]
Let \(n,d\) be sample size and dimension, \(\varepsilon\) the privacy budget, and \(m\) the evaluation-block size. Suppose that:
\begin{itemize}
\item \textbf{(Resources.)} \(n\ge2\), \(d\ge2\), \(0<\varepsilon\le1\), and \(0<m\le n\).
\item \textbf{(Sampling.)} The sampling scheme satisfies \cref{obj:ass:iid-people}.
\item \textbf{(Causal model.)} The full-data law \(P\) belongs to \(\mathcal V_d\), as defined in \cref{obj:def:causal-model}.
\item \textbf{(Bounded mean concentration.)} For every positive integer \(m\), every \(b>0\), and every collection of independent measurable real random variables \(H_1,\ldots,H_m\) taking values in \([-b,b]\), the following bounds hold for every \(x\ge0\):
\[
 \Pr\!\left\{\frac1m\sum_{i=1}^m(H_i-\mathbb EH_i)\ge x\right\}
 \le \exp\!\left(-\frac{mx^2}{2b^2}\right),
 \qquad
 \Pr\!\left\{\frac1m\sum_{i=1}^m(H_i-\mathbb EH_i)\le -x\right\}
 \le \exp\!\left(-\frac{mx^2}{2b^2}\right).
\]
\item \textbf{(Chebyshev approximation.)} For every even integer \(D\ge2\), the explicit Chebyshev absolute-value approximation, with monomial coefficients \(a_{Dv}\),
\[
 q_D(x)=\sum_{v=0}^D a_{Dv}x^v,
\]
satisfies
\[
 \sup_{x\in[-1,1]}|q_D(x)-|x||\le\frac{2}{\pi(D+1)},
 \qquad
 |a_{Dv}|\le 2^{3D/2}\quad(0\le v\le D).
\]
\end{itemize}

Set the privacy attenuation, scaled-message magnitude, logarithmic dimension, calibration constant, and block noise scale to
\[
 \delta=\tanh(\varepsilon/2),\qquad
 b=\frac d\delta,\qquad
 L=\log(ed),\qquad
 \eta_*=\frac1{1024},\qquad
 \sigma^2=\frac{b^2}{m},\qquad
 \sigma=\sqrt{\sigma^2}.
\]
Use independent participant releases from the signed-vector kernel
\[
 S=(2A-1)(2Y-1),\qquad
 Q^{\mathrm{vec}}(z\mid O=(j,A,Y))
 =2^{-d}(1+\delta S z_j),
 \qquad z\in\{-1,1\}^d,
\]
and let \(W_{ij}\) be coordinate \(j\) of release \(i\), multiplied by \(b\). On an evaluation block of size \(m\), define the distinct-person moments and polynomial estimates by
\[
 \mathsf U_{0j}=1,\qquad
 \mathsf U_{vj}
 =\binom mv^{-1}
   \sum_{\substack{I\subseteq\{1,\ldots,m\}\\|I|=v}}
      \prod_{i\in I}W_{ij}
 \quad(1\le v\le m),
\]
\[
 p_{a,D}(x)=a q_D(x/a),\qquad
 \widehat p_{a,D,j}
 =\sum_{v=0}^D a_{Dv}a^{1-v}\mathsf U_{vj}.
\]
The target coordinate and its average absolute magnitude are
\[
 \tau_j(P)=\mu_{1j}(P)-\mu_{0j}(P),\qquad
 F(\tau(P))=\frac1d\sum_{j=1}^d|\tau_j(P)|.
\]

For every even integer \(D\ge2\) satisfying
\[
 2D\le m,\qquad
 H=\log(e+\sigma^2L),\qquad
 D\le\frac{\eta_*L}{H},
\]
take \(a=1/2\). Then, for every \(j\in[d]\),
\[
 \left|\mathbb E\widehat p_{1/2,D,j}-|\tau_j(P)|\right|
 \le\frac1{D+1},
 \qquad
 \operatorname{Var}\!\left(\frac1d\sum_{j=1}^d
                    \widehat p_{1/2,D,j}\right)
 \le\frac{\exp(9DH)}{2d}.
\]

For every integer \(D\ge0\) satisfying
\[
 L\ge4096,\qquad
 D=2\left\lfloor\frac{\eta_*L}{2}\right\rfloor,\qquad
 2D\le m,\qquad
 2m\le n,
\]
use two independent signed-vector release blocks of size \(m\), with scaled messages \(W^{\mathrm{pil}}\) and \(W^{\mathrm{ev}}\). Define
\[
 M_j=\frac1m\sum_{i=1}^mW^{\mathrm{pil}}_{ij},\qquad
 \overline W_j=\frac1m\sum_{i=1}^mW^{\mathrm{ev}}_{ij},\qquad
 T_*=4\sigma\sqrt L,\qquad a=2T_*,
\]
and form \(\widehat p_{a,D,j}\) from the evaluation-block moments. With \(\operatorname{sign}(0)=0\), put
\[
 H_j=\widehat p_{a,D,j}\mathbf1_{\{|M_j|\le T_*\}}
       +\operatorname{sign}(M_j)\overline W_j
          \mathbf1_{\{|M_j|>T_*\}}.
\]
Then, for every \(j\in[d]\),
\[
 \left|\mathbb EH_j-|\tau_j(P)|\right|
 \le\frac{20000\sigma}{\sqrt L},
 \qquad
 \operatorname{Var}(H_j)
 \le400\sigma^2L\exp(L/100),
\]
and
\[
 \mathbb E\!\left[
   \left(\frac1d\sum_{j=1}^dH_j-F(\tau(P))\right)^2
 \right]
 \le500000000\,\frac{\sigma^2}{L}.
\]
\end{lemmav}

\begin{proof}[Proof of \cref{obj:lem:finite-private-polynomial-calibration}]
The resource assumptions give
\[
\delta=\frac{e^\varepsilon-1}{e^\varepsilon+1}>0,
\qquad b>0,
\qquad \sigma^2=\frac{b^2}{m}>0,
\qquad \sigma>0,
\qquad L\ge1.
\]
All statistics below are measurable and have finite moments, since their message spaces are finite.

\begin{enumerate}
\item \textit{Global polynomial bias.}
Fix an even degree \(D\ge2\) satisfying the first set of degree conditions. In particular, \(D\le m\). By the moment unbiasedness in \cref{obj:lem:private-moment-and-dependence}, for every positive radius \(a\),
\[
\mathbb E\widehat p_{a,D,j}
=\sum_{v=0}^D a_{Dv}a^{1-v}\tau_j(P)^v
=a q_D\!\left(\frac{\tau_j(P)}a\right)
=p_{a,D}(\tau_j(P)).
\]
The interior means in \cref{obj:ass:interior-means} imply
\[
-\frac12\le\tau_j(P)\le\frac12.
\]
Consequently, at \(a=1/2\), the assumed approximation bound gives
\[
\left|\mathbb E\widehat p_{1/2,D,j}-|\tau_j(P)|\right|
=\frac12\left|q_D(2\tau_j(P))-|2\tau_j(P)|\right|
\le\frac1{\pi(D+1)}
\le\frac1{D+1}.
\]
% lean: CausalSmith.Stat.LdpOptvalueUniformFrontier.global_polynomial_bias_of_gate

\item \textit{Global polynomial variance.}
For \(1\le v\le D\), the moment second-moment bound in \cref{obj:lem:private-moment-and-dependence} applies because \(2v\le m\). It yields
\[
\begin{aligned}
\mathbb E\left[\left((1/2)^{1-v}\mathsf U_{vj}\right)^2\right]
&=\frac14\,4^v\mathbb E[\mathsf U_{vj}^2]\\
&\le\frac14\left(4\tau_j(P)^2+8v\sigma^2\right)^v\\
&\le\frac14(1+8D\sigma^2)^D.
\end{aligned}
\]
For \(v=0\), the same final bound follows from \(\mathsf U_{0j}=1\). Here we used \(\tau_j(P)^2\le1/4\) and \(1+8D\sigma^2\ge1\).

The assumed coefficient bound and \(\log2\le1\) give, for \(0\le v\le D\),
\[
|a_{Dv}|
\le 2^{3D/2}
=\exp\!\left(\frac{3D}{2}\log2\right)
\le e^{3D/2},
\qquad
a_{Dv}^2\le e^{3D}.
\]
Finite Cauchy--Schwarz therefore gives
\[
\begin{aligned}
\mathbb E[\widehat p_{1/2,D,j}^{\,2}]
&\le(D+1)\sum_{v=0}^D
 a_{Dv}^2
 \mathbb E\left[\left((1/2)^{1-v}\mathsf U_{vj}\right)^2\right]\\
&\le\frac14(D+1)^2e^{3D}(1+8D\sigma^2)^D\\
&\le\frac14e^{5D}(1+8D\sigma^2)^D,
\end{aligned}
\]
where the last inequality uses
\[
D+1\le e^D,
\qquad (D+1)^2\le e^{2D}.
\]

For the stipulated \(H=\log(e+\sigma^2L)\), we have
\[
H\ge1,
\qquad
DH\le\frac L{1024},
\qquad D\le L,
\qquad e^H=e+\sigma^2L.
\]
Since \(8\le e^3\), these inequalities imply
\[
1+8D\sigma^2
\le8(e+\sigma^2L)
=8e^H
\le e^{4H}.
\]
Thus
\[
e^{5D}(1+8D\sigma^2)^D
\le e^{5D+4DH}
\le e^{9DH},
\]
and hence
\[
\operatorname{Var}(\widehat p_{1/2,D,j})
\le\mathbb E[\widehat p_{1/2,D,j}^{\,2}]
\le\frac14e^{9DH}.
\]
Each polynomial estimate is a deterministic measurable statistic of its own column. Applying the average-variance conclusion of \cref{obj:lem:private-moment-and-dependence}, with \(0<m\le n\), proves
\[
\operatorname{Var}\left(\frac1d\sum_{j=1}^d\widehat p_{1/2,D,j}\right)
\le\frac2d\max_{j\in[d]}
       \operatorname{Var}(\widehat p_{1/2,D,j})
\le\frac{e^{9DH}}{2d}.
\]
% lean: CausalSmith.Stat.LdpOptvalueUniformFrontier.global_polynomial_calibration_of_gate

\item \textit{Hybrid bias inside the approximation window.}
Now fix a degree satisfying the second set of hypotheses. The floor inequalities give
\[
D\ \text{is even},
\qquad
D>\frac L{1024}-2,
\qquad
2\le D,
\qquad
\frac L{2048}\le D\le\frac L{1024},
\]
where \(L\ge4096\) supplies the lower bounds. Also \(D\le m\), and
\[
T_*=4\sigma\sqrt L>0,
\qquad
a=2T_*=8\sigma\sqrt L>0,
\qquad
a^2=64\sigma^2L.
\]

The two release blocks are independent by \cref{obj:ass:iid-people} and the independent participant releases. Within either block, the rows of a fixed column are independent. The row moments in \cref{obj:lem:private-moment-and-dependence} therefore give
\[
\mathbb EM_j=\mathbb E\overline W_j=\tau_j(P),
\]
and
\[
\begin{aligned}
\mathbb E[\overline W_j^{\,2}]
&=\frac{mb^2+m(m-1)\tau_j(P)^2}{m^2}
=\tau_j(P)^2+\frac{b^2-\tau_j(P)^2}{m},\\
\operatorname{Var}(\overline W_j)
&=\frac{b^2-\tau_j(P)^2}{m}
\le\sigma^2.
\end{aligned}
\]

First suppose that \(\tau_j(P)\ge0\). Define the pilot probabilities
\[
h_j=\Pr\{|M_j|\le T_*\},
\qquad
w_j=\Pr\{M_j<-T_*\}.
\]
They satisfy
\[
0\le h_j\le1,
\qquad w_j\ge0,
\qquad
h_j+\Pr\{M_j>T_*\}+w_j=1.
\]
The bounded mean concentration assumption applies to the independent pilot messages, each bounded by \(b\) in absolute value. For every real \(u\ge0\), it gives
\[
\Pr\{M_j-\tau_j(P)\ge u\}
\le e^{-u^2/(2\sigma^2)},
\qquad
\Pr\{M_j-\tau_j(P)\le-u\}
\le e^{-u^2/(2\sigma^2)}.
\]
In particular,
\[
\begin{aligned}
w_j
&\le\exp\!\left(-\frac{(\tau_j(P)+T_*)^2}{2\sigma^2}\right)\\
&\le\exp\!\left(-\frac{\tau_j(P)^2}{2\sigma^2}\right)e^{-8L},
\end{aligned}
\]
because \(T_*^2=16\sigma^2L\) and \(\tau_j(P)T_*\ge0\).

The elementary exponential envelopes give
\[
\tau_j(P)e^{-\tau_j(P)^2/(2\sigma^2)}\le\sigma,
\qquad
\tau_j(P)^2e^{-\tau_j(P)^2/(2\sigma^2)}\le\sigma^2.
\]
For the first inequality, multiply
\[
\frac{\tau_j(P)}{\sigma}
\le1+\frac{\tau_j(P)^2}{2\sigma^2}
\le e^{\tau_j(P)^2/(2\sigma^2)}
\]
by \(\sigma e^{-\tau_j(P)^2/(2\sigma^2)}\). For the second, use the elementary maximum
\[
u e^{-u}\le e^{-1}\qquad(u\ge0)
\]
at \(u=\tau_j(P)^2/(2\sigma^2)\), together with \(2/e\le1\). Consequently,
\[
\tau_j(P)w_j\le\sigma e^{-8L},
\qquad
\tau_j(P)^2w_j\le\sigma^2e^{-8L}\le\sigma^2.
\]
Moreover,
\[
\sqrt L\le L\le e^{8L},
\qquad
e^{-8L}\le\frac1{\sqrt L},
\qquad
\tau_j(P)w_j\le\frac{\sigma}{\sqrt L}.
\]

Suppose in addition that \(\tau_j(P)\le a\). Approximation and the degree lower bound imply
\[
\begin{aligned}
|p_{a,D}(\tau_j(P))-\tau_j(P)|
&\le\frac{2a}{\pi(D+1)}
\le\frac a{D+1}\\
&\le16384\,\frac{\sigma}{\sqrt L}.
\end{aligned}
\]
To compute the hybrid mean, its three disjoint pilot branches give
\[
H_j
=\widehat p_{a,D,j}\mathbf1_{\{|M_j|\le T_*\}}
+\overline W_j\mathbf1_{\{M_j>T_*\}}
-\overline W_j\mathbf1_{\{M_j<-T_*\}}.
\]
Block independence, polynomial unbiasedness, and the probability partition yield
\[
\begin{aligned}
\mathbb EH_j
&=h_jp_{a,D}(\tau_j(P))
 +(1-h_j-2w_j)\tau_j(P),\\
\mathbb EH_j-\tau_j(P)
&=h_j\bigl(p_{a,D}(\tau_j(P))-\tau_j(P)\bigr)
 -2\tau_j(P)w_j.
\end{aligned}
\]
Hence, throughout \(0\le\tau_j(P)\le a\),
\[
\begin{aligned}
|\mathbb EH_j-\tau_j(P)|
&\le h_j|p_{a,D}(\tau_j(P))-\tau_j(P)|
       +2\tau_j(P)w_j\\
&\le(16384+2)\frac{\sigma}{\sqrt L}
\le20000\,\frac{\sigma}{\sqrt L}.
\end{aligned}
\]
% lean: CausalSmith.Stat.LdpOptvalueUniformFrontier.hybrid_inside_nonneg_bias

\item \textit{Hybrid bias outside the approximation window.}
Consider the complementary case \(\tau_j(P)>a\), and define
\[
y_j=\frac{\tau_j(P)}a>1.
\]
On the central pilot event,
\[
M_j-\tau_j(P)
\le T_*-\tau_j(P)
\le-\frac{\tau_j(P)}2.
\]
The pilot concentration bound therefore gives
\[
h_j
\le\exp\!\left(-\frac{\tau_j(P)^2}{8\sigma^2}\right)
=e^{-8Ly_j^2}.
\]
The coefficient bound gives the polynomial growth estimate
\[
\begin{aligned}
|q_D(y_j)|
&\le\sum_{v=0}^D|a_{Dv}|y_j^v\\
&\le(D+1)e^{3D/2}y_j^D
\le e^{3D}y_j^D.
\end{aligned}
\]

For \(y_j\ge1\), the inequality \(\log y_j\le y_j-1\) and the nonnegativity of \((y_j-1)^2\) imply
\[
\log y_j\le\frac{y_j^2}{2}.
\]
For every nonnegative integer \(v\le L/1024\), it follows that
\[
3v+v\log y_j
\le\frac72v y_j^2
\le\frac7{2048}Ly_j^2
\le Ly_j^2,
\]
so
\[
e^{-8Ly_j^2}e^{3v}y_j^v\le e^{-7Ly_j^2}.
\]
Applying this with \(v=D\) proves
\[
h_j|p_{a,D}(\tau_j(P))|
\le a e^{-8Ly_j^2}e^{3D}y_j^D
\le a e^{-7Ly_j^2}.
\]
Since \(1,2\le L/1024\), the same bound with \(v=1,2\) gives
\[
h_jy_j^v
\le e^{-8Ly_j^2}y_j^v
\le e^{-8Ly_j^2}e^{3v}y_j^v
\le e^{-7Ly_j^2}
\qquad(v=1,2).
\]
Thus
\[
h_j\tau_j(P)\le a e^{-7Ly_j^2},
\qquad
h_j\tau_j(P)^2\le a^2e^{-7Ly_j^2}.
\]

The exponential lower bound \(e^{7L}\ge1+7L\ge L\) implies
\[
e^{-7Ly_j^2}\le e^{-7L}\le\frac1L,
\qquad
\frac aL=\frac{8\sigma}{\sqrt L}.
\]
Using the hybrid mean identity and the wrong-sign bound consequently gives
\[
\begin{aligned}
|\mathbb EH_j-\tau_j(P)|
&\le h_j|p_{a,D}(\tau_j(P))|
      +h_j\tau_j(P)+2\tau_j(P)w_j\\
&\le(8+8+2)\frac{\sigma}{\sqrt L}
\le20000\,\frac{\sigma}{\sqrt L}.
\end{aligned}
\]
Together with the inside case, this proves the bias bound for every nonnegative contrast, including zero and the boundary \(\tau_j(P)=a\).
% lean: CausalSmith.Stat.LdpOptvalueUniformFrontier.hybrid_nonneg_bias

\item \textit{Reflection for the hybrid bias.}
For a coordinate with \(\tau_j(P)<0\), define the reflected causal law
\[
P^{\mathrm{ref}}=G_{-\tau(P)},
\]
using \cref{obj:def:symmetric-law}. The parameter vector belongs to \([-1/2,1/2]^d\). The construction gives a probability law with uniform strata, fair assignment, and consistency, and
\[
\mu_{1j}(P^{\mathrm{ref}})=\frac{1-\tau_j(P)}2,
\qquad
\mu_{0j}(P^{\mathrm{ref}})=\frac{1+\tau_j(P)}2.
\]
Both means lie in \([1/4,3/4]\). Thus
\[
P^{\mathrm{ref}}\in\mathcal V_d,
\qquad
\tau_j(P^{\mathrm{ref}})=-\tau_j(P)
\quad(j\in[d]).
\]

We establish the reflection identity for the hybrid mean. The Chebyshev polynomials appearing in the explicit approximation are characterized, for \(\zeta\in\mathbb R\), by
\[
\mathcal T_0(\zeta)=1,
\qquad
\mathcal T_1(\zeta)=\zeta,
\qquad
\mathcal T_{v+1}(\zeta)
=2\zeta\mathcal T_v(\zeta)-\mathcal T_{v-1}(\zeta)
\quad(v\ge1).
\]
Induction gives
\[
\mathcal T_v(-\zeta)=(-1)^v\mathcal T_v(\zeta).
\]
The explicit approximation has the expansion
\[
q_D(\zeta)
=\frac2\pi+\frac4\pi
 \sum_{v=1}^{\lfloor D/2\rfloor}
 \frac{(-1)^{v+1}}{4v^2-1}\mathcal T_{2v}(\zeta).
\]
Consequently,
\[
q_D(-\zeta)=q_D(\zeta),
\qquad
(-1)^v a_{Dv}=a_{Dv}\quad(0\le v\le D),
\]
where the coefficient identity follows by comparing the two polynomials.

Viewing moments and estimates as functions of their message arrays, negation gives
\[
\mathsf U_{vj}(-W)=(-1)^v\mathsf U_{vj}(W),
\qquad
\widehat p_{a,D,j}(-W^{\mathrm{ev}})
=\widehat p_{a,D,j}(W^{\mathrm{ev}}).
\]
The pilot and evaluation means change sign, while
\[
|-M_j|=|M_j|,
\qquad
\operatorname{sign}(-M_j)(-\overline W_j)
=\operatorname{sign}(M_j)\overline W_j.
\]
These identities include \(M_j=0\) under the stipulated sign convention. They prove
\[
H_j(-W^{\mathrm{pil}},-W^{\mathrm{ev}})
=H_j(W^{\mathrm{pil}},W^{\mathrm{ev}}),
\]
with threshold ties preserved.

It remains to identify the reflected message law. Fair assignment, consistency, and uniform strata, as specified in
\cref{obj:ass:fair-randomization,obj:ass:consistency,obj:ass:uniform-covariates}, give
\[
\begin{aligned}
\mathbb E_P[S\mathbf1_{\{X=j\}}]
&=\frac12\left(2P(X=j,Y^1=1)-P(X=j)\right)\\
&\quad-\frac12\left(2P(X=j,Y^0=1)-P(X=j)\right)\\
&=\frac{\tau_j(P)}d.
\end{aligned}
\]
Averaging the vector kernel over a participant record therefore gives, for every row \(i\) in either block and every \(z\in\{-1,1\}^d\),
\[
\begin{aligned}
\Pr_P\!\left\{\left(\frac{W_{ij}}b\right)_{j=1}^d=z\right\}
&=2^{-d}\left(1+\delta\sum_{j=1}^d
 z_j\mathbb E_P[S\mathbf1_{\{X=j\}}]\right)\\
&=2^{-d}\left(1+\frac\delta d\sum_{j=1}^d\tau_j(P)z_j\right).
\end{aligned}
\]
Replacing \(z\) by \(-z\) gives precisely the row law under \(P^{\mathrm{ref}}\). Independence of participant rows and of the two blocks then yields
\[
\mathcal L_{P^{\mathrm{ref}}}(W^{\mathrm{pil}},W^{\mathrm{ev}})
=\mathcal L_P(-W^{\mathrm{pil}},-W^{\mathrm{ev}}).
\]
Combining this law identity with the invariance of \(H_j\) proves
\[
\mathbb E_{P^{\mathrm{ref}}}H_j=\mathbb E_PH_j.
\]
The nonnegative-contrast bias bound applied under \(P^{\mathrm{ref}}\) now gives
\[
\begin{aligned}
\left|\mathbb E_PH_j-|\tau_j(P)|\right|
&=\left|\mathbb E_{P^{\mathrm{ref}}}H_j
       -\tau_j(P^{\mathrm{ref}})\right|\\
&\le20000\,\frac{\sigma}{\sqrt L}.
\end{aligned}
\]
This proves the coordinate bias bound for every \(j\).
% lean: CausalSmith.Stat.LdpOptvalueUniformFrontier.hybrid_bias_bound

\item \textit{Polynomial second moments at the hybrid radius.}
For \(1\le v\le D\), \cref{obj:lem:private-moment-and-dependence} and \(2D\le m\) give
\[
\begin{aligned}
\mathbb E[(a^{1-v}\mathsf U_{vj})^2]
&\le a^2
 \left(\frac{\tau_j(P)^2+2v\sigma^2}{a^2}\right)^v\\
&\le a^2
 \left(\max\left\{1,\frac{\tau_j(P)^2+2D\sigma^2}{a^2}\right\}\right)^D.
\end{aligned}
\]
For \(v=0\), the left side is \(a^2\), so the final bound also holds. Finite Cauchy--Schwarz, \(a_{Dv}^2\le e^{3D}\), and \((D+1)^2\le e^{2D}\) yield
\[
\begin{aligned}
\mathbb E[\widehat p_{a,D,j}^{\,2}]
&\le(D+1)\sum_{v=0}^D
 a_{Dv}^2\mathbb E[(a^{1-v}\mathsf U_{vj})^2]\\
&\le a^2(D+1)^2e^{3D}
 \left(\max\left\{1,\frac{\tau_j(P)^2+2D\sigma^2}{a^2}\right\}\right)^D\\
&\le a^2e^{5D}
 \left(\max\left\{1,\frac{\tau_j(P)^2+2D\sigma^2}{a^2}\right\}\right)^D.
\end{aligned}
\]
% lean: CausalSmith.Stat.LdpOptvalueUniformFrontier.radius_polynomial_second_bound

\item \textit{Hybrid variance inside the approximation window.}
First suppose \(0\le\tau_j(P)\le a\). Since
\[
2D\sigma^2\le64\sigma^2L=a^2,
\]
we have
\[
\max\left\{1,\frac{\tau_j(P)^2+2D\sigma^2}{a^2}\right\}\le2.
\]
Using \(2\le e\) and \(D\le L/1024\), the polynomial second-moment bound becomes
\[
\mathbb E[\widehat p_{a,D,j}^{\,2}]
\le a^2e^{5D}2^D
\le a^2e^{6D}
\le a^2e^{L/100}.
\]

The three disjoint pilot branches give the pointwise identity
\[
\begin{aligned}
(H_j-\tau_j(P))^2
&=\mathbf1_{\{|M_j|\le T_*\}}
       (\widehat p_{a,D,j}-\tau_j(P))^2\\
&\quad+\mathbf1_{\{M_j>T_*\}}
       (\overline W_j-\tau_j(P))^2\\
&\quad+\mathbf1_{\{M_j<-T_*\}}
       (-\overline W_j-\tau_j(P))^2.
\end{aligned}
\]
Unbiasedness of the evaluation mean implies
\[
\begin{aligned}
\mathbb E[(-\overline W_j-\tau_j(P))^2]
&=\mathbb E[(\overline W_j-\tau_j(P))^2]
   +4\tau_j(P)\mathbb E\overline W_j\\
&=\operatorname{Var}(\overline W_j)+4\tau_j(P)^2.
\end{aligned}
\]
Block independence and the pilot probability partition therefore yield
\[
\begin{aligned}
\mathbb E[(H_j-\tau_j(P))^2]
&=h_j\mathbb E[(\widehat p_{a,D,j}-\tau_j(P))^2]\\
&\quad+(1-h_j)\operatorname{Var}(\overline W_j)
       +4\tau_j(P)^2w_j.
\end{aligned}
\]
The pointwise inequality
\[
(\widehat p_{a,D,j}-\tau_j(P))^2
\le2\widehat p_{a,D,j}^{\,2}+2\tau_j(P)^2
\]
and \(\operatorname{Var}(\overline W_j)\le\sigma^2\) give
\[
\begin{aligned}
\operatorname{Var}(H_j)
&\le\mathbb E[(H_j-\tau_j(P))^2]\\
&\le2h_j\mathbb E[\widehat p_{a,D,j}^{\,2}]
     +2h_j\tau_j(P)^2+\sigma^2+4\tau_j(P)^2w_j.
\end{aligned}
\]
This centered-error bound also applies to the outside case.

Inside the window, \(h_j\le1\), \(\tau_j(P)^2\le a^2\), and the wrong-sign estimate gives \(\tau_j(P)^2w_j\le\sigma^2\). Hence
\[
\begin{aligned}
\operatorname{Var}(H_j)
&\le2a^2e^{L/100}+2a^2+5\sigma^2\\
&=128\sigma^2Le^{L/100}+128\sigma^2L+5\sigma^2\\
&\le400\sigma^2Le^{L/100},
\end{aligned}
\]
using \(L\ge1\) and \(e^{L/100}\ge1\).
% lean: CausalSmith.Stat.LdpOptvalueUniformFrontier.hybrid_inside_nonneg_variance

\item \textit{Hybrid variance outside the approximation window.}
Suppose \(\tau_j(P)>a\), and again use \(y_j=\tau_j(P)/a>1\). Then
\[
\frac{\tau_j(P)^2+2D\sigma^2}{a^2}
=y_j^2+\frac{2D\sigma^2}{a^2}
\le y_j^2+1
\le2y_j^2,
\]
so
\[
\max\left\{1,\frac{\tau_j(P)^2+2D\sigma^2}{a^2}\right\}
\le2y_j^2.
\]
Combining the pilot tail and polynomial second-moment bounds gives
\[
h_j\mathbb E[\widehat p_{a,D,j}^{\,2}]
\le a^2e^{-8Ly_j^2}e^{5D}(2y_j^2)^D.
\]
Since \(2\le e\),
\[
(2y_j^2)^D\le e^D e^{2D\log y_j}.
\]
Moreover, \(\log y_j\le y_j^2/2\), \(y_j^2\ge1\), and \(D\le L/1024\) imply
\[
6D+2D\log y_j
\le7Dy_j^2
\le\frac7{1024}Ly_j^2
\le Ly_j^2.
\]
It follows that
\[
h_j\mathbb E[\widehat p_{a,D,j}^{\,2}]
\le a^2e^{-7Ly_j^2}\le a^2.
\]
The contrast-weighted pilot bound also gives
\[
h_j\tau_j(P)^2\le a^2e^{-7Ly_j^2}\le a^2.
\]
Using these inequalities and \(\tau_j(P)^2w_j\le\sigma^2\) in the centered-error bound yields
\[
\begin{aligned}
\operatorname{Var}(H_j)
&\le4a^2+5\sigma^2\\
&=256\sigma^2L+5\sigma^2\\
&\le400\sigma^2Le^{L/100}.
\end{aligned}
\]
The inside and outside cases establish the variance bound for every nonnegative contrast.
% lean: CausalSmith.Stat.LdpOptvalueUniformFrontier.hybrid_nonneg_variance

\item \textit{Reflection for the hybrid variance.}
If \(\tau_j(P)<0\), use the same law \(P^{\mathrm{ref}}=G_{-\tau(P)}\). The reflected message-law identity and invariance of \(H_j\), together with equality of its means, imply
\[
\begin{aligned}
\operatorname{Var}_{P^{\mathrm{ref}}}(H_j)
&=\mathbb E_{P^{\mathrm{ref}}}
       [(H_j-\mathbb E_{P^{\mathrm{ref}}}H_j)^2]\\
&=\mathbb E_P[(H_j-\mathbb E_PH_j)^2]
=\operatorname{Var}_P(H_j).
\end{aligned}
\]
Since \(\tau_j(P^{\mathrm{ref}})>0\), the nonnegative-contrast variance bound gives
\[
\operatorname{Var}_P(H_j)
=\operatorname{Var}_{P^{\mathrm{ref}}}(H_j)
\le400\sigma^2Le^{L/100}.
\]
Thus the asserted coordinate variance bound holds for every \(j\).
% lean: CausalSmith.Stat.LdpOptvalueUniformFrontier.hybrid_variance_bound

\item \textit{Variance of the hybrid average.}
Stack the pilot rows followed by the evaluation rows. Their joint law is the \(2m\)-participant product release law, and each \(H_j\) is a deterministic measurable function of the combined column
\[
(W^{\mathrm{pil}}_{1j},\ldots,W^{\mathrm{pil}}_{mj},
 W^{\mathrm{ev}}_{1j},\ldots,W^{\mathrm{ev}}_{mj})
\in\mathbb R^{2m}.
\]
Because \(0<2m\le n\), \cref{obj:lem:private-moment-and-dependence} applies to this combined block. The coordinate variance bounds therefore give
\[
\operatorname{Var}\left(\frac1d\sum_{j=1}^dH_j\right)
\le\frac2d\max_{j\in[d]}\operatorname{Var}(H_j)
\le\frac{800\sigma^2Le^{L/100}}d.
\]

We absorb the dimension factor using \(L\ge4096\). The elementary exponential lower bound gives
\[
\frac L4\le e^{L/4},
\qquad
L^2\le16e^{L/2},
\]
and
\[
4\le\frac{49L}{100}
=1+\left(\frac{49L}{100}-1\right)
\le e^{49L/100-1}.
\]
Since \(e^{L-1}=d\), it follows that
\[
\begin{aligned}
L^2e^{L/100}
&\le16e^{L/2}e^{L/100}\\
&=16e^{51L/100}\\
&\le4e^{51L/100}e^{49L/100-1}
=4d.
\end{aligned}
\]
As \(d,L>0\), this proves
\[
\frac{800\sigma^2Le^{L/100}}d
\le\frac{3200\sigma^2}{L},
\qquad
\operatorname{Var}\left(\frac1d\sum_{j=1}^dH_j\right)
\le\frac{3200\sigma^2}{L}.
\]
% lean: CausalSmith.Stat.LdpOptvalueUniformFrontier.hybrid_average_variance_envelope_le

\item \textit{Mean squared error of the hybrid average.}
Linearity of expectation, the definition of \(F\), and the coordinate bias bounds give
\[
\begin{aligned}
\left|\mathbb E\left[\frac1d\sum_{j=1}^dH_j\right]
      -F(\tau(P))\right|
&=\left|\frac1d\sum_{j=1}^d
       \bigl(\mathbb EH_j-|\tau_j(P)|\bigr)\right|\\
&\le\frac1d\sum_{j=1}^d
       \left|\mathbb EH_j-|\tau_j(P)|\right|\\
&\le20000\,\frac{\sigma}{\sqrt L}.
\end{aligned}
\]
Squaring, with \(\sigma^2>0\) and \(L>0\), yields
\[
\left(\mathbb E\left[\frac1d\sum_{j=1}^dH_j\right]
      -F(\tau(P))\right)^2
\le400000000\,\frac{\sigma^2}{L}.
\]
Finally, expansion around the expectation gives the variance-plus-squared-bias identity
\[
\begin{aligned}
\mathbb E\left[
 \left(\frac1d\sum_{j=1}^dH_j-F(\tau(P))\right)^2\right]
&=\operatorname{Var}\left(\frac1d\sum_{j=1}^dH_j\right)\\
&\quad+\left(\mathbb E\left[\frac1d\sum_{j=1}^dH_j\right]
             -F(\tau(P))\right)^2\\
&\le(3200+400000000)\frac{\sigma^2}{L}\\
&=400003200\,\frac{\sigma^2}{L}\\
&\le500000000\,\frac{\sigma^2}{L}.
\end{aligned}
\]
% lean: CausalSmith.Stat.LdpOptvalueUniformFrontier.hybrid_average_mse_bound
\end{enumerate}
% lean: CausalSmith.Stat.LdpOptvalueUniformFrontier.finite_private_polynomial_calibration
\end{proof}

The concentration and approximation hypotheses are supplied by \cref{obj:lem:bounded-mean-concentration,obj:lem:cai-low-chebyshev-approximation}. The sampling and causal conditions give the moment identities in \cref{obj:lem:private-moment-and-dependence}. Together, these conditions connect deterministic approximation error to bias and estimated-moment variability to polynomial variance.

In the global calibration, the radius \(1/2\) covers every admissible treatment contrast. The bias bound decreases with the degree, while the variance bound grows exponentially in the degree times the noise-dependent calibration factor. The stated upper bound on the degree balances this growth against averaging across strata, and the block-size condition supplies the second moments required for all polynomial terms.

For the pilot-selected calibration, the threshold \(T_*\) is measured in units of the block noise scale, enlarged by the square root of the logarithmic dimension. The approximation radius is twice that threshold. Independent blocks separate branch selection from evaluation, while concentration controls threshold and sign mistakes. Under the stated logarithmic-dimension and block-size conditions, the coordinate bias and variance bounds account for these selections, and the final bound controls the squared error of the average absolute-contrast estimate. This is the nonlinear component of the welfare procedure in \cref{obj:def:frontier-handle}.


\section{Sequential contraction and decision lower bounds}

The approximation and private-moment bounds describe the precision attainable by the welfare procedure. The complementary lower bounds connect moment matching to the information retained by a sequential private experiment, and then connect that information to welfare estimation and honest interval length. The argument accommodates independent public randomness and message rules that depend on the preceding private messages, as specified in \cref{obj:def:sequential-class}.

\subsection{Measurable densities for conditional experiments}

Sequential experiments require density versions that are measurable jointly in the conditioning variables and the next message. For the auxiliary result below, \(H\) denotes a measurable parameter space and \(Z\) a standard-Borel output space. The finite kernels \(K\) and \(L\) assign output measures to each parameter value, with \(L\) serving as the reference kernel. A jointly measurable Radon--Nikodym density \(f\) represents \(K\) relative to \(L\) throughout this family of conditional experiments. The following result supplies that representation under pointwise absolute continuity.

\begin{theoremv}{lem:measurable-kernel-radon-nikodym}[Jointly measurable Radon–Nikodym density]
Let \(H\) and \(Z\) be measurable spaces, and let \(K,L:H\to Z\) be kernels satisfying:
\begin{itemize}
\item \textbf{(Standard Borel target.)} \(Z\) is a standard-Borel space.
\item \textbf{(Finiteness.)} Both \(K\) and \(L\) are finite kernels.
\item \textbf{(Pointwise domination.)} \(K(h,\cdot)\ll L(h,\cdot)\) for every \(h\in H\).
\end{itemize}
Then there exists a jointly measurable function \(f:H\times Z\to[0,\infty]\) such that, for every \(h\in H\) and every measurable set \(E\subseteq Z\),
\[
K(h,E)=\int_E f(h,z)\,L(h,dz).
\]
\end{theoremv}

The standard-Borel restriction applies to the output space, while the parameter space remains an arbitrary measurable space. This distinction allows the conditioning parameter to collect the public seed, the preceding messages, and the contrast vector. Joint measurability makes the density suitable for integration over both public randomness and parameter priors. The representation holds at every parameter value, matching the all-history formulation of privacy in \cref{obj:def:sequential-class}.

\subsection{Higher-order contraction under sequential privacy}

Fix a sequential protocol \(Q\) from \cref{obj:def:sequential-class}. For a public-seed realization \(r\), the reference experiment is the conditional transcript law under the zero-contrast causal law \(G_0\). The symmetric causal family \(G_\theta\) is defined in \cref{obj:def:symmetric-law}. Write \leanref{sym:p^Q}{\ensuremath{p^Q(\theta,r)}} for its conditional transcript density relative to that reference. The privacy contraction scale is \leanref{sym:\beta}{\ensuremath{\beta}}, equal to \((e^\varepsilon-1)/(2d)\). Conditioning on the seed preserves the protocol's dependence on the entire preceding message history.

To compare experiments generated by different contrast priors, let \leanref{sym:\alpha}{\ensuremath{\alpha}} be a positive amplitude at most \(1/2\). The coordinate priors \leanref{sym:\nu_0}{\ensuremath{\nu_0}} and \leanref{sym:\nu_1(same)}{\ensuremath{\nu_1}} are probability laws on \([-\alpha,\alpha]\), chosen independently of the public seed. Drawing the coordinates independently from either prior and then running the protocol gives the joint seed-transcript mixtures \leanref{sym:\mathsf M_0^Q}{\ensuremath{\mathsf M_0^Q}} and \leanref{sym:\mathsf M_1^Q}{\ensuremath{\mathsf M_1^Q}}. Moment matching means that the two coordinate priors have equal moments below a specified order. The next result bounds the resulting transcript proximity through the corresponding higher-order density derivative.

\begin{lemmav}{lem:adaptive-moment-contraction}[Adaptive higher-order transcript contraction]
Let \(n,d\) be integers and let \(\varepsilon\) be a privacy budget. Suppose that:
\begin{itemize}
\item \textbf{(Parameter regime.)} \(n\ge2\), \(d\ge2\), and \(0<\varepsilon\le1\).
\item \textbf{(Sampling and randomness.)} The sampling scheme satisfies \cref{obj:ass:iid-people,obj:ass:independent-randomness}.
\item \textbf{(Sequential privacy.)} \(Q\in\mathfrak Q_{\mathrm{SI}}(n,d,\varepsilon)\), as specified in \cref{obj:def:sequential-class}.
\item \textbf{(Measurable kernel densities.)} For every measurable parameter space \(H\), standard-Borel output space \(Z\), and finite kernels \(K,L\) from \(H\) to \(Z\) satisfying \(K(h,\cdot)\ll L(h,\cdot)\) for every \(h\), there exists a jointly measurable \(f:H\times Z\to[0,\infty]\) such that, for every \(h\) and measurable \(E\subseteq Z\),
\[
K(h,E)=\int_E f(h,z)\,L(h,dz).
\]
\end{itemize}
Set the privacy contraction scale to
\[
\beta=\frac{e^\varepsilon-1}{2d}.
\]
There exists a jointly measurable real-valued density version \(p^Q(\theta,r)(\mathbf z)\), defined for all \(\theta\in\mathbb R^d\), public seeds \(r\), and transcripts \(\mathbf z\), with the following properties. For every \(\theta\in[-1/2,1/2]^d\) and every \(r\), it is nonnegative at every transcript and gives the conditional transcript law under \(G_\theta\) relative to the zero-contrast reference:
\[
\Pr_Q(d\mathbf z\mid r,G_\theta)
=
p^Q(\theta,r)(\mathbf z)\,
\Pr_Q(d\mathbf z\mid r,G_0).
\]
For every \(\theta,r,\mathbf z,j\), its dependence on coordinate \(\theta_j\), with the other coordinates fixed, is a polynomial on \(\mathbb R\) of degree at most \(n\). For every \(\theta\in[-1/2,1/2]^d\), every \(r,j\), and every integer \(1\le k\le n\),
\[
\int
\left|\partial_{\theta_j}^{\,k}p^Q(\theta,r)(\mathbf z)\right|
\,\Pr_Q(d\mathbf z\mid r,G_0)
\le
k!\sqrt{\binom nk}\,\beta^k.
\]
For every \(\theta\in\mathbb R^d\), every \(r,\mathbf z,j\), and every integer \(k>n\),
\[
\partial_{\theta_j}^{\,k}p^Q(\theta,r)(\mathbf z)=0.
\]

For every amplitude \(0<\alpha\le1/2\), let \(\nu_0,\nu_1\) be probability laws supported on \([-\alpha,\alpha]\), fixed independently of the public seed. Define their seed-transcript mixtures by
\[
\mathsf M_\ell^Q(dr,d\mathbf z)
=
\lambda(dr)
\int_{[-\alpha,\alpha]^d}
\Pr_Q(d\mathbf z\mid r,G_\theta)\,
\nu_\ell^{\otimes d}(d\theta),
\qquad \ell\in\{0,1\}.
\]
For every integer \(k\ge1\) such that
\[
\int u^q\,\nu_0(du)=\int u^q\,\nu_1(du),
\qquad q=0,\ldots,k-1,
\]
the following conclusions hold:
\[
\operatorname{TV}(\mathsf M_0^Q,\mathsf M_1^Q)
\le
\min\left\{1,\,
d\alpha^k\sqrt{\binom nk}\,\beta^k\right\},
\qquad 1\le k\le n,
\]
and
\[
\mathsf M_0^Q=\mathsf M_1^Q,
\qquad k>n.
\]
Moreover, every Borel measurable seed-transcript decision \(f(r,\mathbf z)\in[0,1]\), interpreted as a randomized test's acceptance probability, satisfies
\[
\left|
\int f(r,\mathbf z)\,\mathsf M_0^Q(dr,d\mathbf z)
-
\int f(r,\mathbf z)\,\mathsf M_1^Q(dr,d\mathbf z)
\right|
\le
\operatorname{TV}(\mathsf M_0^Q,\mathsf M_1^Q).
\]
\end{lemmav}

\begin{proof}[Proof of \cref{obj:lem:adaptive-moment-contraction}]
All products over empty index sets equal one. We use the fixed-seed protocol experiment as the conditional-law version, including at seeds of zero \(\lambda\)-probability.

\begin{enumerate}
\item \textbf{Construction of the density.}
Define the observed-record and transcript spaces by
\[
\mathcal O_d=[d]\times\{0,1\}^2,
\qquad
\mathcal Z=\prod_{i=1}^n\mathcal Z_i.
\]
For \(o=(j_o,a_o,y_o)\in\mathcal O_d\), and for
\(\theta\in\mathbb R^d\) and
\(\mathbf o=(o_1,\ldots,o_n)\in\mathcal O_d^n\), set
\[
\mathsf{sgn}_{\mathrm{obs}}(o)
=(2a_o-1)(2y_o-1),
\qquad
\mathsf{mass}^{(1)}_\theta(o)
=\frac{1+\mathsf{sgn}_{\mathrm{obs}}(o)\theta_{j_o}}{4d},
\qquad
\mathsf{mass}^{(n)}_\theta(\mathbf o)
=\prod_{i=1}^n\mathsf{mass}^{(1)}_\theta(o_i).
\]
The reconstruction identities in \cref{obj:prop:signed-causal-bridge} give
\[
G_\theta(O=o)=\mathsf{mass}^{(1)}_\theta(o)
\qquad
\bigl(\theta\in[-1/2,1/2]^d\bigr).
\]
For every \(\mathbf o\in\mathcal O_d^n\) and \(r\in\mathcal R\), define the fixed-record transcript law by successive kernel integration:
\[
\mathsf{Law}^Q_{\mathbf o,r}(d\mathbf z)
=\prod_{i=1}^n Q_i(dz_i\mid o_i,\mathbf z_{<i},r).
\]
For \(\theta\in[-1/2,1/2]^d\), define the conditional transcript law and its zero-contrast reference by
\[
\mathsf{Prob}^Q_{\theta,r}
=\sum_{\mathbf o\in\mathcal O_d^n}
  \mathsf{mass}^{(n)}_\theta(\mathbf o)\,
  \mathsf{Law}^Q_{\mathbf o,r},
\qquad
\Lambda_r^Q=\mathsf{Prob}^Q_{0,r}.
\]
In particular,
\[
\Lambda_r^Q
=(4d)^{-n}\sum_{\mathbf o\in\mathcal O_d^n}
  \mathsf{Law}^Q_{\mathbf o,r}.
\]
Since \(d\ge2\), every coefficient in this last sum is positive. Consequently,
\[
\mathsf{Law}^Q_{\mathbf o,r}\ll\Lambda_r^Q
\qquad
\text{for every \(\mathbf o,r\)}.
\]

Apply the measurable-kernel-density hypothesis to these two probability-kernel families, parametrized by \((\mathbf o,r)\). It supplies jointly measurable functions
\[
\varphi^{\mathrm{ext}}_{\mathbf o}:
\mathcal R\times\mathcal Z\longrightarrow[0,\infty],
\qquad
\mathsf{Law}^Q_{\mathbf o,r}(E)
=\int_E\varphi^{\mathrm{ext}}_{\mathbf o}(r,\mathbf z)\,
  \Lambda_r^Q(d\mathbf z).
\]
Their integrals over \(\mathcal Z\) equal one, so their infinite-value sets are \(\Lambda_r^Q\)-null for each \(\mathbf o,r\). Define everywhere finite versions by
\[
\varphi_{\mathbf o}(r,\mathbf z)
=
\begin{cases}
\varphi^{\mathrm{ext}}_{\mathbf o}(r,\mathbf z),
 &\varphi^{\mathrm{ext}}_{\mathbf o}(r,\mathbf z)<\infty,\\
0,&\varphi^{\mathrm{ext}}_{\mathbf o}(r,\mathbf z)=\infty.
\end{cases}
\]
These versions are jointly measurable, nonnegative, and satisfy
\[
\mathsf{Law}^Q_{\mathbf o,r}(d\mathbf z)
=\varphi_{\mathbf o}(r,\mathbf z)\Lambda_r^Q(d\mathbf z),
\qquad
\int\varphi_{\mathbf o}(r,\mathbf z)\Lambda_r^Q(d\mathbf z)=1.
\]

Now define, on the full domain
\(\mathbb R^d\times\mathcal R\times\mathcal Z\),
\[
p^Q(\theta,r)(\mathbf z)
=
\sum_{\mathbf o\in\mathcal O_d^n}
 \mathsf{mass}^{(n)}_\theta(\mathbf o)\,
 \varphi_{\mathbf o}(r,\mathbf z).
\]
This is jointly measurable and real valued. On the parameter cube its summands are nonnegative, and finite summation of the density identities gives
\[
\mathsf{Prob}^Q_{\theta,r}(d\mathbf z)
=p^Q(\theta,r)(\mathbf z)\Lambda_r^Q(d\mathbf z).
\]

For \(j\in[d]\) and \(u\in\mathbb R\), define coordinate replacement by
\[
\theta^{[j:u]}_x
=
\begin{cases}
u,&x=j,\\
\theta_x,&x\ne j.
\end{cases}
\]
Each observed-record mass obeys
\[
\mathsf{mass}^{(1)}_{\theta^{[j:u]}}(o)
=
\mathsf{mass}^{(1)}_{\theta^{[j:0]}}(o)
+\frac{\mathbf1_{\{j_o=j\}}\mathsf{sgn}_{\mathrm{obs}}(o)}{4d}\,u.
\]
Thus each sample mass, and hence \(p^Q\), is a polynomial of degree at most \(n\) in each coordinate on all of \(\mathbb R\). Its coordinate coefficients are finite linear combinations of the integrable functions \(\varphi_{\mathbf o}(r,\cdot)\), so those coefficients and all coordinate derivatives are integrable against \(\Lambda_r^Q\).
% lean: CausalSmith.Stat.LdpOptvalueUniformFrontier.canonicalTranscriptDensity_integrable_coordinate_polynomial

\item \textbf{Pointwise normalized stage densities.}
Use the paired-input protocol \(\mathsf K\) supplied by
\cref{obj:prop:signed-causal-bridge}. It satisfies the same sequential privacy inequality at every history and seed. For
\(i\in[n]\), \(r\in\mathcal R\), and \(\eta\in\mathcal H_i\), define
\[
K_i^0(\,\cdot\mid r,\eta)
=\frac1{2d}\sum_{x=1}^d\sum_{s\in\{-1,1\}}
 \mathsf K_i(\,\cdot\mid(x,s),\eta,r).
\]
This is a probability kernel, and each summand is absolutely continuous with respect to it. The density hypothesis, with the same finite-value convention as above, therefore gives jointly measurable nonnegative real functions
\[
\psi^{\mathrm{raw}}_{i,x,s}:
\mathcal R\times\mathcal H_i\times\mathcal Z_i\longrightarrow[0,\infty)
\]
such that
\[
\mathsf K_i(d\zeta\mid(x,s),\eta,r)
=\psi^{\mathrm{raw}}_{i,x,s}(r,\eta,\zeta)
 K_i^0(d\zeta\mid r,\eta).
\]

For each \(r,\eta\), the definition of \(K_i^0\) and uniqueness of densities imply
\[
\sum_{x,s}\psi^{\mathrm{raw}}_{i,x,s}(r,\eta,\zeta)=2d
\quad
K_i^0(\cdot\mid r,\eta)\text{-almost everywhere}.
\]
Privacy gives, for every measurable \(E\subseteq\mathcal Z_i\),
\[
\int_E\psi^{\mathrm{raw}}_{i,x,s}(r,\eta,\zeta)
 K_i^0(d\zeta\mid r,\eta)
\le
e^\varepsilon
\int_E\psi^{\mathrm{raw}}_{i,x',s'}(r,\eta,\zeta)
 K_i^0(d\zeta\mid r,\eta).
\]
The comparison of these integrals over all measurable sets implies the corresponding density inequality almost everywhere.

Define the measurable good set
\[
\begin{split}
\mathcal G_i
=\biggl\{(r,\eta,\zeta):\;&
 \sum_{x,s}\psi^{\mathrm{raw}}_{i,x,s}(r,\eta,\zeta)=2d,\\
&\psi^{\mathrm{raw}}_{i,x,s}(r,\eta,\zeta)
 \le e^\varepsilon\psi^{\mathrm{raw}}_{i,x',s'}(r,\eta,\zeta)
 \ \text{for all }x,s,x',s'\biggr\},
\end{split}
\]
and set, on their full domains,
\[
\psi_{i,x,s}(r,\eta,\zeta)
=
\begin{cases}
\psi^{\mathrm{raw}}_{i,x,s}(r,\eta,\zeta),
 &(r,\eta,\zeta)\in\mathcal G_i,\\
1,&(r,\eta,\zeta)\notin\mathcal G_i.
\end{cases}
\]
For each \(r,\eta\), the good set has full reference measure. Hence this repair preserves every represented kernel and gives, everywhere,
\[
\sum_{x,s}\psi_{i,x,s}=2d,
\qquad
\psi_{i,x,s}\le e^\varepsilon\psi_{i,x',s'}.
\]
Summing the second inequality first over the left-hand input and then over the right-hand input yields
\[
2d\le 2d\,e^\varepsilon\psi_{i,x,s},
\qquad
2d\,\psi_{i,x,s}\le e^\varepsilon\sum_{x',s'}\psi_{i,x',s'}
=2d\,e^\varepsilon.
\]
Therefore
\[
e^{-\varepsilon}\le\psi_{i,x,s}(r,\eta,\zeta)\le e^\varepsilon
\qquad\text{everywhere}.
\]
% lean: CausalSmith.Stat.LdpOptvalueUniformFrontier.stage_density_repaired_of_gate

\item \textbf{The stage product and coordinate identification.}
For all \(\theta\in\mathbb R^d\) and all declared stage arguments, define
\[
\mathsf{dens}_i(\theta,r,\eta,\zeta)
=
\frac1{2d}\sum_{x,s}(1+s\theta_x)\psi_{i,x,s}(r,\eta,\zeta),
\qquad
\mathsf{slope}_{ij}(r,\eta,\zeta)
=
\frac{\psi_{i,j,1}(r,\eta,\zeta)
      -\psi_{i,j,-1}(r,\eta,\zeta)}{2d}.
\]
Normalization gives the identities
\[
\mathsf{dens}_i(\theta,r,\eta,\zeta)
=1+\sum_{j=1}^d\theta_j\mathsf{slope}_{ij}(r,\eta,\zeta),
\]
\[
\mathsf{dens}_i(\theta^{[j:u]},r,\eta,\zeta)
=
\mathsf{dens}_i(\theta,r,\eta,\zeta)
+(u-\theta_j)\mathsf{slope}_{ij}(r,\eta,\zeta).
\]
Define the reference stage-chain law and the stage product by
\[
\Lambda_r^{Q,\mathrm{stage}}(d\mathbf z)
=\prod_{i=1}^n K_i^0(dz_i\mid r,\mathbf z_{<i}),
\qquad
\mathsf{prod}^Q(\theta,r,\mathbf z)
=\prod_{i=1}^n
 \mathsf{dens}_i(\theta,r,\mathbf z_{<i},z_i).
\]

For a fixed paired-input vector
\(\boldsymbol\chi=((j_i,s_i))_{i=1}^n
 \in([d]\times\{-1,1\})^n\), define
\[
\mathsf{Law}^{\mathsf K}_{\boldsymbol\chi,r}(d\mathbf z)
=\prod_{i=1}^n
 \mathsf K_i(dz_i\mid(j_i,s_i),\mathbf z_{<i},r).
\]
Successive integration of the stage density identities gives
\[
\mathsf{Law}^{\mathsf K}_{\boldsymbol\chi,r}(d\mathbf z)
=
\left\{\prod_{i=1}^n
 \psi_{i,j_i,s_i}(r,\mathbf z_{<i},z_i)\right\}
 \Lambda_r^{Q,\mathrm{stage}}(d\mathbf z).
\]
The fair ancillary-bit reconstruction in
\cref{obj:prop:signed-causal-bridge}, followed by expansion of the averaged kernels at successive stages, gives the fixed-seed identity
\[
\mathsf{Prob}^Q_{\theta,r}
=
\sum_{\boldsymbol\chi}
 \left\{\prod_{i=1}^nP_\theta(j_i,s_i)\right\}
 \mathsf{Law}^{\mathsf K}_{\boldsymbol\chi,r}
\qquad
\bigl(\theta\in[-1/2,1/2]^d\bigr).
\]
Indeed, each ancillary assignment vector has probability \(2^{-n}\), and its recovered observed-record vector is obtained coordinatewise by
\[
o_i
=\left(j_i,a_i,\frac{1+s_i(2a_i-1)}2\right),
\qquad
(a_1,\ldots,a_n)\in\{0,1\}^n.
\]
Finite distribution of sums over products now gives
\[
\begin{split}
&\sum_{\boldsymbol\chi}
 \prod_{i=1}^n
 \left\{
 \frac{1+s_i\theta_{j_i}}{2d}
 \psi_{i,j_i,s_i}(r,\mathbf z_{<i},z_i)
 \right\}\\
&\hspace{3em}
=\prod_{i=1}^n
 \left\{\frac1{2d}\sum_{x,s}
 (1+s\theta_x)\psi_{i,x,s}(r,\mathbf z_{<i},z_i)\right\}
=\mathsf{prod}^Q(\theta,r,\mathbf z).
\end{split}
\]
Consequently,
\[
\mathsf{Prob}^Q_{\theta,r}(d\mathbf z)
=\mathsf{prod}^Q(\theta,r,\mathbf z)
 \Lambda_r^{Q,\mathrm{stage}}(d\mathbf z).
\]
At \(\theta=0\), every stage density equals one. Thus
\[
\Lambda_r^{Q,\mathrm{stage}}=\Lambda_r^Q,
\qquad
\mathsf{Prob}^Q_{\theta,r}(d\mathbf z)
=\mathsf{prod}^Q(\theta,r,\mathbf z)\Lambda_r^Q(d\mathbf z).
\]

Fix \(\theta\in[-1/2,1/2]^d\), \(r\), and \(j\). Both
\[
u\longmapsto p^Q(\theta^{[j:u]},r)(\mathbf z),
\qquad
u\longmapsto\mathsf{prod}^Q(\theta^{[j:u]},r,\mathbf z)
\]
are polynomials of degree at most \(n\). Define distinct interpolation points by
\[
\upsilon_\iota=\frac1{\iota+2},
\qquad \iota\in\mathbb N.
\]
Every \(\theta^{[j:\upsilon_\iota]}\) lies in the cube. Since the two nonnegative functions represent the same law there, uniqueness of densities gives equality almost everywhere at each \(\upsilon_\iota\). Taking the countable intersection of these full-measure sets, the two polynomials agree at infinitely many distinct points. Polynomial uniqueness therefore gives
\[
\text{for \(\Lambda_r^Q\)-almost every \(\mathbf z\), for every \(u\in\mathbb R\),}\qquad
p^Q(\theta^{[j:u]},r)(\mathbf z)
=\mathsf{prod}^Q(\theta^{[j:u]},r,\mathbf z).
\]
This identity also identifies every order of coordinate differentiation on the same full-measure set.
% lean: CausalSmith.Stat.LdpOptvalueUniformFrontier.transcript_coordinate_ae_eq_stageProduct

\item \textbf{Scores and their moments.}
Fix \(\theta\in[-1/2,1/2]^d\), \(r\), and \(j\). For each stage define
\[
\psi_{i,\min}(r,\eta,\zeta)
=\min_{x,s}\psi_{i,x,s}(r,\eta,\zeta).
\]
The weights \((1+s\theta_x)/(2d)\) are nonnegative and sum to one. Privacy and the minimum therefore imply
\[
\mathsf{dens}_i(\theta,r,\eta,\zeta)
\ge\psi_{i,\min}(r,\eta,\zeta)>0,
\]
\[
\left|\psi_{i,j,1}(r,\eta,\zeta)
      -\psi_{i,j,-1}(r,\eta,\zeta)\right|
\le(e^\varepsilon-1)\psi_{i,\min}(r,\eta,\zeta)
\le(e^\varepsilon-1)\mathsf{dens}_i(\theta,r,\eta,\zeta).
\]
Hence
\[
\left|
\frac{\mathsf{slope}_{ij}(r,\eta,\zeta)}
     {\mathsf{dens}_i(\theta,r,\eta,\zeta)}
\right|
\le\beta.
\]

Under \(\mathsf{Prob}^Q_{\theta,r}\), define expectation, the history filtration, and the scores by
\[
\mathbb E^Q_{\theta,r}[\cdot]
=\int_{\mathcal Z}(\cdot)\,\mathsf{Prob}^Q_{\theta,r}(d\mathbf z),
\qquad
\mathcal F_v=\sigma(Z_1,\ldots,Z_{\min\{v,n\}}),
\quad v\in\mathbb N,
\]
\[
\mathsf{score}^{\theta,r,j}_i(\mathbf z)
=
\begin{cases}
\dfrac{\mathsf{slope}_{ij}(r,\mathbf z_{<i},z_i)}
      {\mathsf{dens}_i(\theta,r,\mathbf z_{<i},z_i)},
 &1\le i\le n,\\[6pt]
0,&i>n,
\end{cases}
\qquad i\in\{1,2,\ldots\}.
\]
The parameter-specific stage kernel is
\[
\mathsf{dens}_i(\theta,r,\eta,\zeta)
 K_i^0(d\zeta\mid r,\eta).
\]
Its successive kernel integrations give
\(\mathsf{Prob}^Q_{\theta,r}\), by the product identity above. Each input-row density integrates to one, so
\[
\int\mathsf{slope}_{ij}(r,\eta,\zeta)
 K_i^0(d\zeta\mid r,\eta)
=\frac{1-1}{2d}=0.
\]
Cancellation of the positive stage density in the score integral gives
\[
\mathbb E^Q_{\theta,r}
 [\mathsf{score}^{\theta,r,j}_i\mid\mathcal F_{i-1}]
=0,
\qquad
|\mathsf{score}^{\theta,r,j}_i|\le\beta.
\]
The zero extension satisfies these identities after the horizon as well.

For \(k,v\in\mathbb N\), define the elementary score polynomials recursively by
\[
\mathsf{elem}_{0,v}=1,
\qquad
\mathsf{elem}_{k,0}=0\quad(k\ge1),
\]
\[
\mathsf{elem}_{k,v+1}
=\mathsf{elem}_{k,v}
+\mathsf{score}^{\theta,r,j}_{v+1}\mathsf{elem}_{k-1,v}
\quad(k\ge1).
\]
Here \(\mathsf{elem}_{k,v}\) is a function of the transcript, with the fixed
\(\theta,r,j\) understood. The recursion gives
\(\mathsf{elem}_{k,v}=0\) when \(k>v\), and each prefix polynomial is
\(\mathcal F_v\)-measurable. Conditional mean zero thus implies
\[
\mathbb E^Q_{\theta,r}
 \left[
 \mathsf{elem}_{k,v}\,
 \mathsf{score}^{\theta,r,j}_{v+1}\,
 \mathsf{elem}_{k-1,v}
 \right]=0
\qquad(k\ge1).
\]
Squaring the recursion and using the score bound yields
\[
\mathbb E^Q_{\theta,r}[\mathsf{elem}_{k,v+1}^2]
\le
\mathbb E^Q_{\theta,r}[\mathsf{elem}_{k,v}^2]
+\beta^2\mathbb E^Q_{\theta,r}[\mathsf{elem}_{k-1,v}^2].
\]
Induction from the displayed initial values, using
\(\binom{v+1}{k}=\binom vk+\binom v{k-1}\), gives
\[
\mathbb E^Q_{\theta,r}[\mathsf{elem}_{k,v}^2]
\le\binom vk\beta^{2k}.
\]
This includes \(k=0\) and \(k>v\).

The nonnegative centered square gives
\[
\begin{split}
0
&\le
\mathbb E^Q_{\theta,r}
 \left[
 \left(
 |\mathsf{elem}_{k,v}|
 -\mathbb E^Q_{\theta,r}|\mathsf{elem}_{k,v}|
 \right)^2
 \right]\\
&=
\mathbb E^Q_{\theta,r}[\mathsf{elem}_{k,v}^2]
-\left(\mathbb E^Q_{\theta,r}|\mathsf{elem}_{k,v}|\right)^2.
\end{split}
\]
Since \(\beta\ge0\), it follows that
\[
\mathbb E^Q_{\theta,r}|\mathsf{elem}_{k,v}|
\le\sqrt{\binom vk}\,\beta^k.
\]
% lean: CausalSmith.Stat.LdpOptvalueUniformFrontier.originalTranscript_scoreElementary_firstMoment

\item \textbf{The derivative certificate.}
Multiplying the affine coordinate-shift identities recursively gives
\[
\mathsf{prod}^Q(\theta^{[j:\theta_j+u]},r,\mathbf z)
=
\mathsf{prod}^Q(\theta,r,\mathbf z)
\sum_{q=0}^n u^q\mathsf{elem}_{q,n}(\mathbf z),
\qquad u\in\mathbb R.
\]
The simultaneous coordinate identity from the preceding construction therefore yields
\[
\partial_{\theta_j}^{\,k}p^Q(\theta,r)(\mathbf z)
=
k!\,p^Q(\theta,r)(\mathbf z)\,
\mathsf{elem}_{k,n}(\mathbf z)
\quad
\Lambda_r^Q\text{-almost everywhere}.
\]
Changing measure by the nonnegative density \(p^Q\), and applying the first-moment estimate, gives
\[
\begin{split}
\int
\left|\partial_{\theta_j}^{\,k}p^Q(\theta,r)(\mathbf z)\right|
 \Lambda_r^Q(d\mathbf z)
&=
k!\,\mathbb E^Q_{\theta,r}|\mathsf{elem}_{k,n}|\\
&\le k!\sqrt{\binom nk}\,\beta^k,
\qquad 1\le k\le n.
\end{split}
\]
Finally, the coordinate polynomial representation of the finite-sum density gives, pointwise on its full domain,
\[
\partial_{\theta_j}^{\,k}p^Q(\theta,r)(\mathbf z)=0
\qquad
\bigl(\theta\in\mathbb R^d,\ r\in\mathcal R,\
\mathbf z\in\mathcal Z,\ j\in[d],\ k>n\bigr).
\]
Together with the construction and representation already proved, these are all the asserted density properties.
% lean: CausalSmith.Stat.LdpOptvalueUniformFrontier.canonical_density_certificate_rows_of_gate

\item \textbf{A single-coordinate moment bound.}
Fix \(0<\alpha\le1/2\), the two supported probability priors, and
\(1\le k\le n\) with the asserted matching moments. For
\(\theta\in[-1/2,1/2]^d\), \(r\in\mathcal R\), and \(j\in[d]\), define
\[
\mathsf{diff}^{\theta,r}_j(\mathbf z)
=
\int p^Q(\theta^{[j:u]},r)(\mathbf z)\,\nu_0(du)
-\int p^Q(\theta^{[j:u]},r)(\mathbf z)\,\nu_1(du),
\]
\[
\mathsf{sign}^{\theta,r}_j(\mathbf z)
=
\begin{cases}
\dfrac{\mathsf{diff}^{\theta,r}_j(\mathbf z)}
      {|\mathsf{diff}^{\theta,r}_j(\mathbf z)|},
 &\mathsf{diff}^{\theta,r}_j(\mathbf z)\ne0,\\[6pt]
0,&\mathsf{diff}^{\theta,r}_j(\mathbf z)=0,
\end{cases}
\]
and the scalar polynomial, defined for every \(u\in\mathbb R\), by
\[
\mathsf{poly}^{\theta,r}_j(u)
=
\int
 \mathsf{sign}^{\theta,r}_j(\mathbf z)\,
 p^Q(\theta^{[j:u]},r)(\mathbf z)\,
 \Lambda_r^Q(d\mathbf z).
\]
The finite polynomial representation and integrable coefficients justify all these integrals, coefficientwise differentiation, and the exchanges of integrals below. In particular,
\[
|\mathsf{sign}^{\theta,r}_j|\le1,
\qquad
\mathsf{sign}^{\theta,r}_j\,\mathsf{diff}^{\theta,r}_j
=|\mathsf{diff}^{\theta,r}_j|,
\]
and
\[
\left(\mathsf{poly}^{\theta,r}_j\right)^{(q)}(u)
=
\int
 \mathsf{sign}^{\theta,r}_j(\mathbf z)\,
 \partial_{\theta_j}^{\,q}p^Q(\theta^{[j:u]},r)(\mathbf z)\,
 \Lambda_r^Q(d\mathbf z).
\]
For \(|u|\le\alpha\), the replacement parameter remains in the cube. Thus
\[
\left|
 \left(\mathsf{poly}^{\theta,r}_j\right)^{(k)}(u)
\right|
\le
\int
 \left|\partial_{\theta_j}^{\,k}
 p^Q(\theta^{[j:u]},r)(\mathbf z)\right|
 \Lambda_r^Q(d\mathbf z)
\le k!\sqrt{\binom nk}\,\beta^k.
\]

Define its Taylor polynomial by
\[
\mathsf{taylor}^{\theta,r}_{j,k}(u)
=
\sum_{q=0}^{k-1}
 \frac{\left(\mathsf{poly}^{\theta,r}_j\right)^{(q)}(0)}{q!}\,u^q,
\qquad u\in\mathbb R.
\]
At \(u=0\), the remainder equals zero. For \(u\ne0\), Taylor's theorem with Lagrange remainder on the segment joining \(0\) and \(u\) gives
\[
\begin{split}
\left|
 \mathsf{poly}^{\theta,r}_j(u)
 -\mathsf{taylor}^{\theta,r}_{j,k}(u)
\right|
&\le
\frac{|u|^k}{k!}
\sup_{u_*\in[\min\{0,u\},\,\max\{0,u\}]}
 \left|
 \left(\mathsf{poly}^{\theta,r}_j\right)^{(k)}(u_*)
 \right|\\
&\le |u|^k\sqrt{\binom nk}\,\beta^k
\qquad(|u|\le\alpha).
\end{split}
\]
Matching moments cancel its Taylor polynomial:
\[
\int\mathsf{taylor}^{\theta,r}_{j,k}(u)\,\nu_0(du)
=
\int\mathsf{taylor}^{\theta,r}_{j,k}(u)\,\nu_1(du).
\]
Each prior is supported on \([-\alpha,\alpha]\) and has mass one, so
\[
\left|
\int
 \left(\mathsf{poly}^{\theta,r}_j(u)
       -\mathsf{taylor}^{\theta,r}_{j,k}(u)\right)\nu_\ell(du)
\right|
\le\alpha^k\sqrt{\binom nk}\,\beta^k,
\qquad \ell\in\{0,1\}.
\]
Finally, exchanging the integrals in the scalar polynomial gives
\[
\begin{split}
\int|\mathsf{diff}^{\theta,r}_j(\mathbf z)|\,\Lambda_r^Q(d\mathbf z)
&=
\int\mathsf{poly}^{\theta,r}_j(u)\,\nu_0(du)
-\int\mathsf{poly}^{\theta,r}_j(u)\,\nu_1(du)\\
&\le 2\alpha^k\sqrt{\binom nk}\,\beta^k.
\end{split}
\]
This is uniform in the remaining coordinates, the seed, and the coordinate being replaced.
% lean: CausalSmith.Stat.LdpOptvalueUniformFrontier.canonical_density_coordinate_matching_L1

\item \textbf{Averaging and coordinate telescoping.}
Define the endpoint product priors and the intervening product priors by
\[
\Pi_\ell=\nu_\ell^{\otimes d},
\qquad \ell\in\{0,1\},
\]
\[
\Pi_v^{\mathrm{hyb}}
=\nu_1^{\otimes v}\otimes\nu_0^{\otimes(d-v)},
\qquad v=0,\ldots,d.
\]
Their coordinates are ordered from \(1\) to \(d\). They are probability laws supported on the parameter cube, with
\[
\Pi_0^{\mathrm{hyb}}=\Pi_0,
\qquad
\Pi_d^{\mathrm{hyb}}=\Pi_1.
\]
For each \(r\) and \(v\), define the averaged density
\[
\mathsf{avg}_v^r(\mathbf z)
=\int p^Q(\theta,r)(\mathbf z)\,\Pi_v^{\mathrm{hyb}}(d\theta).
\]
These functions are integrable against \(\Lambda_r^Q\), by their finite-row representation.

For \(v<d\), set \(j=v+1\). Resampling coordinate \(j\) from \(\nu_0\) preserves \(\Pi_v^{\mathrm{hyb}}\); resampling it from \(\nu_1\) gives \(\Pi_{v+1}^{\mathrm{hyb}}\). Explicitly, with \(\#\) denoting pushforward,
\[
(\theta,u\mapsto\theta^{[j:u]})_\#
 (\Pi_v^{\mathrm{hyb}}\otimes\nu_0)=\Pi_v^{\mathrm{hyb}},
\qquad
(\theta,u\mapsto\theta^{[j:u]})_\#
 (\Pi_v^{\mathrm{hyb}}\otimes\nu_1)=\Pi_{v+1}^{\mathrm{hyb}}.
\]
Fubini therefore gives
\[
\mathsf{avg}_v^r(\mathbf z)-\mathsf{avg}_{v+1}^r(\mathbf z)
=
\int\mathsf{diff}^{\theta,r}_j(\mathbf z)\,
 \Pi_v^{\mathrm{hyb}}(d\theta).
\]
The triangle inequality under the integral, followed by the single-coordinate estimate, yields
\[
\begin{split}
\int
 |\mathsf{avg}_v^r-\mathsf{avg}_{v+1}^r|\,d\Lambda_r^Q
&\le
\int
 \left\{\int|\mathsf{diff}^{\theta,r}_j(\mathbf z)|
              \Lambda_r^Q(d\mathbf z)\right\}
 \Pi_v^{\mathrm{hyb}}(d\theta)\\
&\le2\alpha^k\sqrt{\binom nk}\,\beta^k.
\end{split}
\]
Telescoping the \(d\) replacements now gives
\[
\begin{split}
&\int
 \left|
 \int p^Q(\theta,r)(\mathbf z)\,\Pi_0(d\theta)
 -\int p^Q(\theta,r)(\mathbf z)\,\Pi_1(d\theta)
 \right|
 \Lambda_r^Q(d\mathbf z)\\
&\qquad
=\int|\mathsf{avg}_0^r-\mathsf{avg}_d^r|\,d\Lambda_r^Q\\
&\qquad
\le\sum_{v=0}^{d-1}
 \int|\mathsf{avg}_v^r-\mathsf{avg}_{v+1}^r|\,d\Lambda_r^Q\\
&\qquad
\le2d\alpha^k\sqrt{\binom nk}\,\beta^k.
\end{split}
\]
% lean: CausalSmith.Stat.LdpOptvalueUniformFrontier.density_product_matching_of_gate

\item \textbf{Fixed-seed total variation.}
To define probability kernels on the full parameter domain, set
\[
\theta^{\mathrm{cl}}_j
=\max\left\{-\frac12,\min\left\{\frac12,\theta_j\right\}\right\},
\qquad
\theta\in\mathbb R^d,\ j\in[d].
\]
This parameter belongs to the cube for every \(\theta\), and equals \(\theta\) on the cube. Define, for \(\ell\in\{0,1\}\) and \(r\in\mathcal R\),
\[
\mathsf{condmix}_{\ell,r}(d\mathbf z)
=\int\mathsf{Prob}^Q_{\theta^{\mathrm{cl}},r}(d\mathbf z)\,
 \Pi_\ell(d\theta),
\]
\[
\mathsf{seedrow}_{\ell,r}
=(\mathbf z\mapsto(r,\mathbf z))_\#
 \mathsf{condmix}_{\ell,r},
\qquad
\mathsf{avden}_\ell(r,\mathbf z)
=\int p^Q(\theta,r)(\mathbf z)\,\Pi_\ell(d\theta).
\]
All these measure-valued families are measurable; the first two consist of probability measures. Mixing commutes with the displayed measurable map, so
\[
\mathsf{seedrow}_{\ell,r}
=
\int
 (\mathbf z\mapsto(r,\mathbf z))_\#
 \mathsf{Prob}^Q_{\theta^{\mathrm{cl}},r}\,
 \Pi_\ell(d\theta).
\]
Support on the cube, the density representation, and Fubini give
\[
\mathsf{condmix}_{\ell,r}(d\mathbf z)
=\mathsf{avden}_\ell(r,\mathbf z)\Lambda_r^Q(d\mathbf z),
\qquad
\mathsf{avden}_\ell(r,\mathbf z)\ge0,
\qquad
\int\mathsf{avden}_\ell(r,\mathbf z)\Lambda_r^Q(d\mathbf z)=1.
\]

For completeness, the last normalization gives
\[
\int\bigl(\mathsf{avden}_0(r,\mathbf z)
          -\mathsf{avden}_1(r,\mathbf z)\bigr)
 \Lambda_r^Q(d\mathbf z)=0.
\]
For every measurable \(E\subseteq\mathcal Z\), the integral of this difference over \(E\) is the negative of its integral over \(E^c\). Thus
\[
\begin{split}
&2\left|
 \mathsf{condmix}_{0,r}(E)-\mathsf{condmix}_{1,r}(E)
 \right|\\
&\qquad\le
 \int_E|\mathsf{avden}_0-\mathsf{avden}_1|\,d\Lambda_r^Q
 +\int_{E^c}|\mathsf{avden}_0-\mathsf{avden}_1|\,d\Lambda_r^Q\\
&\qquad=
 \int|\mathsf{avden}_0-\mathsf{avden}_1|\,d\Lambda_r^Q.
\end{split}
\]
Taking the supremum over measurable events gives the half-\(L^1\) bound. A measurable pushforward decreases total variation, since its event differences are differences over measurable preimages. Consequently, at every seed,
\[
\begin{split}
\operatorname{TV}(\mathsf{seedrow}_{0,r},\mathsf{seedrow}_{1,r})
&\le
\operatorname{TV}(\mathsf{condmix}_{0,r},\mathsf{condmix}_{1,r})\\
&\le
\frac12\int
 |\mathsf{avden}_0(r,\mathbf z)-\mathsf{avden}_1(r,\mathbf z)|
 \Lambda_r^Q(d\mathbf z)\\
&\le d\alpha^k\sqrt{\binom nk}\,\beta^k.
\end{split}
\]
% lean: CausalSmith.Stat.LdpOptvalueUniformFrontier.conditionalMixture_tv_le_half_L1

\item \textbf{Mixing over the common public seed.}
The sampling and randomness assumptions ensure that, at each cube parameter, the observed-record vector has its independent product law and is independent of the public seed and analyst randomizer. Successive kernel integration therefore gives the seed-transcript law
\[
\lambda(dr)\,\mathsf{Prob}^Q_{\theta,r}(d\mathbf z).
\]
The priors are supported on the cube and are fixed independently of the seed. Hence both \(\mathsf M_\ell^Q\) are probability measures, and
\[
\operatorname{TV}(\mathsf M_0^Q,\mathsf M_1^Q)\le1.
\]
Interchanging the independent prior and seed integrations gives, for every measurable seed-transcript event \(E\),
\[
\begin{split}
\mathsf M_\ell^Q(E)
&=
\int_{\mathbb R^d}\int_{\mathcal R}
 \left[
 (\mathbf z\mapsto(r,\mathbf z))_\#
 \mathsf{Prob}^Q_{\theta^{\mathrm{cl}},r}
 \right](E)\,\lambda(dr)\,\Pi_\ell(d\theta)\\
&=
\int_{\mathcal R}\mathsf{seedrow}_{\ell,r}(E)\,\lambda(dr).
\end{split}
\]
The measurability established above justifies these mixture identities. Applying the fixed-seed bound to each event gives
\[
\begin{split}
|\mathsf M_0^Q(E)-\mathsf M_1^Q(E)|
&\le
\int
 |\mathsf{seedrow}_{0,r}(E)-\mathsf{seedrow}_{1,r}(E)|
 \,\lambda(dr)\\
&\le
\int d\alpha^k\sqrt{\binom nk}\,\beta^k\,\lambda(dr)\\
&=d\alpha^k\sqrt{\binom nk}\,\beta^k.
\end{split}
\]
Taking the supremum and combining with the probability-measure bound proves
\[
\operatorname{TV}(\mathsf M_0^Q,\mathsf M_1^Q)
\le
\min\left\{1,\,
d\alpha^k\sqrt{\binom nk}\,\beta^k\right\},
\qquad 1\le k\le n.
\]
% lean: CausalSmith.Stat.LdpOptvalueUniformFrontier.probability_bind_tv_le_constant

\item \textbf{Exact cancellation above the sample degree.}
Now suppose \(k>n\). Define the prior-predictive input laws on
\(\mathcal O_d^n\) by
\[
\mathsf{Input}_\ell
=
\int
 \left(\sum_{o\in\mathcal O_d}
  \mathsf{mass}^{(1)}_{\theta^{\mathrm{cl}}}(o)\,\delta_o
 \right)^{\otimes n}
 \Pi_\ell(d\theta),
\qquad \ell\in\{0,1\}.
\]
Since clipping equals the identity on the prior support, their atom masses are
\[
\mathsf{Input}_\ell(\{\mathbf o\})
=\int\mathsf{mass}^{(n)}_\theta(\mathbf o)\,\Pi_\ell(d\theta).
\]
For \(\mathcal J\subseteq[n]\), \(\mathbf o\in\mathcal O_d^n\), and \(j\in[d]\), define the coordinate counts
\[
\mathsf{count}_{j,\mathcal J}(\mathbf o)
=\#\{i\in\mathcal J:j_{o_i}=j\}.
\]
Expanding the affine sample mass gives
\[
\mathsf{mass}^{(n)}_\theta(\mathbf o)
=
(4d)^{-n}
\sum_{\mathcal J\subseteq[n]}
 \left\{\prod_{i\in\mathcal J}\mathsf{sgn}_{\mathrm{obs}}(o_i)\right\}
 \prod_{i\in\mathcal J}\theta_{j_{o_i}}.
\]
Regrouping factors by coordinate and integrating under the product prior yields
\[
\begin{split}
\int\prod_{i\in\mathcal J}\theta_{j_{o_i}}\,\Pi_\ell(d\theta)
&=
\int\prod_{j=1}^d
 \theta_j^{\mathsf{count}_{j,\mathcal J}(\mathbf o)}
 \,\Pi_\ell(d\theta)\\
&=
\prod_{j=1}^d
 \int u^{\mathsf{count}_{j,\mathcal J}(\mathbf o)}\,\nu_\ell(du).
\end{split}
\]
Every count lies between \(0\) and \(n<k\), so every factor is the same for the two priors by moment matching. This also covers the empty subset and zero counts. Finite expansion therefore gives
\[
\mathsf{Input}_0(\{\mathbf o\})
=\mathsf{Input}_1(\{\mathbf o\})
\qquad
\text{for every \(\mathbf o\in\mathcal O_d^n\)},
\]
and hence
\[
\mathsf{Input}_0=\mathsf{Input}_1.
\]

Define the probability channel from a fixed input vector to the seed-transcript space by
\[
\mathsf{channel}^Q_{\mathbf o}(dr,d\mathbf z)
=\lambda(dr)\,\mathsf{Law}^Q_{\mathbf o,r}(d\mathbf z).
\]
The sampling and randomness assumptions give
\[
\mathsf M_\ell^Q
=
\sum_{\mathbf o\in\mathcal O_d^n}
 \mathsf{Input}_\ell(\{\mathbf o\})\,
 \mathsf{channel}^Q_{\mathbf o}.
\]
The independent analyst randomizer integrates to mass one in this representation. Applying the same channel to the identical input mixtures proves
\[
\mathsf M_0^Q=\mathsf M_1^Q
\qquad(k>n).
\]
% lean: CausalSmith.Stat.LdpOptvalueUniformFrontier.mixtureLaw_eq_of_matching_above_sample

\item \textbf{Bounded transcript decisions.}
Both mixtures are probability measures for every supported pair of priors. The standard bounded-range total-variation inequality, applied to a Borel measurable
\(f:\mathcal R\times\mathcal Z\to[0,1]\), gives
\[
\left|
\int f(r,\mathbf z)\,\mathsf M_0^Q(dr,d\mathbf z)
-
\int f(r,\mathbf z)\,\mathsf M_1^Q(dr,d\mathbf z)
\right|
\le
\operatorname{TV}(\mathsf M_0^Q,\mathsf M_1^Q).
\]
This applies for every \(k\ge1\) in the statement.
% lean: CausalSmith.Stat.LdpOptvalueUniformFrontier.mixtureLaw_bounded_test
\end{enumerate}
\end{proof}

The measurable-density condition is supplied by \cref{obj:lem:measurable-kernel-radon-nikodym}. The density has a probability interpretation on the admissible contrast cube and a coordinatewise polynomial extension to all real contrast vectors. Its degree bound reflects the finite participant count, while its integrated derivative bound quantifies the sensitivity of the private transcript at each order.

The contraction calculation retains the adaptive history within each conditional experiment. Both the density representation and the derivative bound hold for every public-seed realization, so the mixture conclusion accommodates any independent seed law \(\lambda\). The prior independence condition fixes the same comparison across those seed realizations. For matched moments below order \(k\), the displayed bound combines the amplitude factor \(\alpha^k\), the privacy factor \(\beta^k\), the participant factor \(\sqrt{\binom nk}\), and the dimension factor \(d\). Matching beyond the coordinatewise polynomial degree gives exact equality of the mixtures.

The final clause translates transcript proximity into a bound on every Borel randomized testing decision. Thus the result controls the information available to an analyst after observing both the public seed and all private messages. Together with the moment-matching measures in \cref{obj:lem:cai-low-moment-duality}, it supplies the experiment-comparison component of the welfare lower bounds.

\subsection{From welfare separation to decision lower bounds}

Transcript proximity becomes a decision lower bound when the two priors place their welfare targets near different centers. For two priors supported on the causal model \(\mathcal V_d\) from \cref{obj:def:causal-model}, let \(N_0\) and \(N_1\) denote their joint seed-transcript mixture laws. Write \(v_0<v_1\) for the welfare centers and \(\Delta=v_1-v_0\) for their separation. The quantity \(\eta_0\) bounds the probability, under each prior, that welfare lies farther than \(\Delta/8\) from its center; \(\omega\) bounds the total-variation distance between \(N_0\) and \(N_1\).

The following result separates these two requirements: concentration concerns the random target induced by a prior, whereas proximity concerns the distribution of the observable private experiment. Its decision laws also include the independent analyst randomizer \(U\) from \cref{obj:ass:independent-randomness}. This permits randomized estimators and intervals within the same comparison.

\begin{lemmav}{lem:separated-value-mixtures}[Separated welfare mixture bounds]
Fix a sample size \(n\), a dimension \(d\), a sampling scheme satisfying \cref{obj:ass:iid-people,obj:ass:independent-randomness}, and a law-independent one-message sequential protocol \(Q\) with probability message kernels and standard-Borel seed and message spaces. Write \(\mathcal V_d\) for the full-data probability laws satisfying the declared causal model of uniform strata, fair assignment, consistency, and interior arm means. The welfare target is
\[
V(P)=\frac{1}{d}\sum_{j=1}^{d}
\max\{\mu_{0j}(P),\mu_{1j}(P)\},
\]
where \(\mu_{aj}(P)\) is the conditional mean of potential outcome \(Y^a\) in stratum \(j\).

Consider two probability priors on full-data laws, called the first and second priors. Let \(v_0<v_1\) be their respective welfare centers, and define their separation by
\[
\Delta=v_1-v_0.
\]
Let \(\eta_0,\omega\in\mathbb R\). Suppose:
\begin{itemize}
\item \textbf{(Causal support.)} Each prior assigns probability one to \(\mathcal V_d\).
\item \textbf{(Decision mixtures.)} Mixing the induced joint law of \((\mathbf Z,R,U)\) over either prior yields a probability measure. Let \(N_0\) and \(N_1\) be the respective joint seed-transcript marginals, namely the mixture laws of \((R,\mathbf Z)\).
\item \textbf{(Target concentration.)} Under the first and second priors, respectively,
\[
\Pr\{|V(P)-v_0|>\Delta/8\}\le\eta_0,
\qquad
\Pr\{|V(P)-v_1|>\Delta/8\}\le\eta_0.
\]
\item \textbf{(Transcript proximity.)} The joint seed-transcript mixtures satisfy
\[
\operatorname{TV}(N_0,N_1)\le\omega.
\]
\end{itemize}

Then every Borel real-valued estimator \(T(\mathbf Z,R,U)\) satisfies
\[
\sup_{P\in\mathcal V_d}
\mathbb E_{P,Q}\bigl[(T(\mathbf Z,R,U)-V(P))^2\bigr]
\ge
\frac{9\Delta^2}{128}(1-\omega-2\eta_0)_+,
\]
with squared-error expectations interpreted in the extended nonnegative reals.

Moreover, every connected closed interval \(I(\mathbf Z,R,U)\) with ordered Borel endpoints in \([1/4,3/4]\) and global coverage
\[
\Pr_{P,Q}\{V(P)\in I(\mathbf Z,R,U)\}\ge0.90
\qquad\text{for every }P\in\mathcal V_d
\]
satisfies
\[
\sup_{P\in\mathcal V_d}
\mathbb E_{P,Q}\operatorname{length}I(\mathbf Z,R,U)
\ge
\frac{3\Delta}{4}(0.80-2\eta_0-\omega)_+.
\]
Here \(\operatorname{length}I\) is the upper endpoint minus the lower endpoint, and \(x_+=\max\{x,0\}\).
\end{lemmav}

\begin{proof}[Proof of \cref{obj:lem:separated-value-mixtures}]
\begin{enumerate}
\item \textit{Measurable decision laws.}
Denote the two priors by
\[
\pi_0:=\text{the first prior},
\qquad
\pi_1:=\text{the second prior}.
\]
Write
\[
\mathcal O_d:=[d]\times\{0,1\}^2,
\qquad
\mathcal Z:=\prod_{i=1}^n\mathcal Z_i,
\qquad
\mathcal R:=\text{the public seed space}.
\]
A decision triple has the form
\[
\mathsf w_{\mathrm{dec}}
=(\mathbf z,\mathsf r,\mathsf u)
\in\mathcal Z\times\mathcal R\times\mathbb R.
\]
For an observed record vector \(\mathbf o\in\mathcal O_d^n\) and a seed value \(\mathsf r\in\mathcal R\), define the probability transcript law by iterating the protocol kernels:
\[
\mathsf L_{\mathrm{tr}}(\mathbf o,\mathsf r)(d\mathbf z)
:=
\prod_{i=1}^n
Q_i(dz_i\mid o_i,\mathbf z_{<i},\mathsf r).
\]
Here the product denotes successive kernel integration in participant order. Define a probability measure for each fixed record vector by
\[
\mathsf D_{\mathbf o}(E)
:=
\int_{\mathcal R}\int_{[0,1]}\int_{\mathcal Z}
\mathbf 1_E(\mathbf z,\mathsf r,\mathsf u)\,
\mathsf L_{\mathrm{tr}}(\mathbf o,\mathsf r)(d\mathbf z)\,
d\mathsf u\,\lambda(d\mathsf r),
\]
for every Borel decision event \(E\). Each measure is a probability measure because all the message kernels and the seed law are probability measures.

For every nonnegative measure \(P\) on the full-record space, let \(P_O\) be its observed-record marginal and define
\[
\mathsf L_{\mathrm{dec}}(P)
:=
\sum_{\mathbf o\in\mathcal O_d^n}
\left(\prod_{i=1}^n P_O(\{o_i\})\right)\mathsf D_{\mathbf o}.
\]
This map is measurable: each coefficient is a finite product of measurable evaluations of the observed marginal, and the measures \(\mathsf D_{\mathbf o}\) are fixed.

The auxiliary target on all nonnegative measures uses the same welfare formula, with the following conventions:
\[
P_{\mathrm{real}}(E):=
\begin{cases}
P(E),&P(E)<\infty,\\
0,&P(E)=\infty,
\end{cases}
\qquad
\mu_{aj}(P):=
\frac{P_{\mathrm{real}}\{X=j,Y^a=1\}}
     {P_{\mathrm{real}}\{X=j\}},
\qquad
V(P):=
\frac1d\sum_{j=1}^d\max\{\mu_{0j}(P),\mu_{1j}(P)\}.
\]
Division by zero, including \(1/d\) when \(d=0\), is assigned value zero; empty sums and products have values zero and one, respectively. These formulas make the target measurable and agree with the stated conditional means and welfare on \(\mathcal V_d\).

For a probability law \(P\), the finite atomic decomposition is
\[
P_O^{\otimes n}
=
\sum_{\mathbf o\in\mathcal O_d^n}
\left(\prod_{i=1}^nP_O(\{o_i\})\right)\delta_{\mathbf o}.
\]
Consequently, integrating the protocol against the canonical product sampling law gives exactly \(\mathsf L_{\mathrm{dec}}(P)\). Moreover, \cref{obj:ass:iid-people,obj:ass:independent-randomness} give, for every \(P\in\mathcal V_d\),
\[
\mathcal L_P(\mathbf O,R,U)
=
P_O^{\otimes n}\otimes\lambda
\otimes\operatorname{Uniform}[0,1],
\qquad
\mathsf L_{\mathrm{dec}}(P)
=
\mathcal L_{P,Q}(\mathbf Z,R,U).
\]
Define the completed mixtures, for \(\ell\in\{0,1\}\), by
\[
\mathsf L_\ell(E)
:=
\int \mathsf L_{\mathrm{dec}}(P)(E)\,\pi_\ell(dP).
\]
Causal support and the preceding equality imply
\[
\mathsf L_\ell(E)
=
\int_{\mathcal V_d}
\Pr_{P,Q}\{(\mathbf Z,R,U)\in E\}\,\pi_\ell(dP).
\]
Thus these are the decision mixtures in the statement, and their seed-transcript marginals are \(N_\ell\).
% lean: CausalSmith.Stat.LdpOptvalueUniformFrontier.separated_value_mixtures

\item \textit{Independence of the analyst randomizer.}
For \(P\in\mathcal V_d\), the preceding construction gives
\[
\mathsf L_{\mathrm{dec}}(P)
(\mathcal Z\times\mathcal R\times\mathbb R)=1.
\]
For Borel sets \(E_Z\subseteq\mathcal Z\), \(E_R\subseteq\mathcal R\), and \(E_U\subseteq\mathbb R\), successive integration yields
\[
\begin{aligned}
&\mathsf L_{\mathrm{dec}}(P)(E_Z\times E_R\times E_U)\\
&\quad=
\left[
\int_{\mathcal O_d^n}\int_{\mathcal R}
\mathbf1_{E_R}(\mathsf r)\,
\mathsf L_{\mathrm{tr}}(\mathbf o,\mathsf r)(E_Z)\,
\lambda(d\mathsf r)\,P_O^{\otimes n}(d\mathbf o)
\right]
\operatorname{Uniform}[0,1](E_U).
\end{aligned}
\]
Indeed, for fixed \(\mathbf o,\mathsf r,\mathsf u\), the rectangle indicator contributes the transcript probability multiplied by
\(\mathbf1_{E_R}(\mathsf r)\mathbf1_{E_U}(\mathsf u)\); integrating in \(\mathsf u\) gives the displayed factor. Taking \(E_U=\mathbb R\) identifies the bracket as the transcript-seed marginal. More precisely, define
\[
\mathsf L_{\mathrm{dec}}^{ZR}(P)(E_{ZR})
:=
\mathsf L_{\mathrm{dec}}(P)
\{(\mathbf z,\mathsf r,\mathsf u):
(\mathbf z,\mathsf r)\in E_{ZR}\}.
\]
Equality on measurable rectangles therefore proves, with the product identified with decision triples,
\[
\mathsf L_{\mathrm{dec}}(P)
=
\mathsf L_{\mathrm{dec}}^{ZR}(P)
\otimes\operatorname{Uniform}[0,1],
\qquad P\in\mathcal V_d.
\]
% lean: CausalSmith.Stat.LdpOptvalueUniformFrontier.sampling_decisionLaw_randomizer

\item \textit{Total variation of the decision mixtures.}
Define a common probability kernel that appends the analyst randomizer:
\[
\mathsf C_{\mathrm{coin}}((\mathsf r,\mathbf z),E)
:=
\int_{[0,1]}
\mathbf1_E(\mathbf z,\mathsf r,\mathsf u)\,d\mathsf u.
\]
The rowwise product identity from the preceding step implies
\[
\mathsf L_{\mathrm{dec}}(P)(E)
=
\int_{\mathcal Z\times\mathcal R}
\mathsf C_{\mathrm{coin}}((\mathsf r,\mathbf z),E)\,
\mathsf L_{\mathrm{dec}}^{ZR}(P)(d\mathbf z,d\mathsf r),
\qquad P\in\mathcal V_d.
\]
Integrating this equality under either supported prior and interchanging the nonnegative integrals gives
\[
\mathsf L_\ell(E)
=
\int_{\mathcal R\times\mathcal Z}
\mathsf C_{\mathrm{coin}}((\mathsf r,\mathbf z),E)\,
N_\ell(d\mathsf r,d\mathbf z).
\]

For a Borel event \(E\), put
\[
f_E(\mathsf r,\mathbf z)
:=
\mathsf C_{\mathrm{coin}}((\mathsf r,\mathbf z),E),
\qquad
0\le f_E\le1.
\]
The function is measurable. Its layer-set representation gives
\[
\begin{aligned}
|\mathsf L_0(E)-\mathsf L_1(E)|
&=
\left|
\int_0^1
\left[
N_0\{f_E\ge t_{\mathrm{lev}}\}
-
N_1\{f_E\ge t_{\mathrm{lev}}\}
\right]dt_{\mathrm{lev}}
\right|\\
&\le
\int_0^1
\left|
N_0\{f_E\ge t_{\mathrm{lev}}\}
-
N_1\{f_E\ge t_{\mathrm{lev}}\}
\right|dt_{\mathrm{lev}}\\
&\le \operatorname{TV}(N_0,N_1).
\end{aligned}
\]
Taking the supremum over Borel decision events proves
\[
\operatorname{TV}(\mathsf L_0,\mathsf L_1)
\le \operatorname{TV}(N_0,N_1)
\le\omega.
\]
% lean: CausalSmith.Stat.LdpOptvalueUniformFrontier.supported_mixture_randomizer_tv

\item \textit{Probability rows for the testing argument.}
Since a probability prior is supported on \(\mathcal V_d\), that class is nonempty. Choose
\[
P_\star\in\mathcal V_d.
\]
Complete the measurable decision laws to probability rows on the entire auxiliary parameter space by setting
\[
\mathsf C_{\mathrm{dec}}(P):=
\begin{cases}
\mathsf L_{\mathrm{dec}}(P),
&\mathsf L_{\mathrm{dec}}(P)
(\mathcal Z\times\mathcal R\times\mathbb R)=1,\\
\mathsf L_{\mathrm{dec}}(P_\star),
&\text{otherwise}.
\end{cases}
\]
The condition selecting the first branch is measurable. The resulting kernel has probability rows, and
\[
\mathsf C_{\mathrm{dec}}(P)
=\mathsf L_{\mathrm{dec}}(P)
\quad(P\in\mathcal V_d),
\qquad
\int\mathsf C_{\mathrm{dec}}(P)\,\pi_\ell(dP)
=\mathsf L_\ell.
\]
Thus the completion preserves both prior-predictive laws and their total-variation bound. The concentration and proximity hypotheses also imply
\[
\eta_0\ge0,
\qquad
\omega\ge0,
\]
because probabilities and total variation are nonnegative.
% lean: CausalSmith.Stat.LdpOptvalueUniformFrontier.separated_prior_squared_risk

\item \textit{Squared-error bound.}
Fix a Borel estimator \(T\). Define
\[
\begin{gathered}
v_{\mathrm{mid}}:=\frac{v_0+v_1}{2},
\qquad
c_{\mathrm{err}}:=\left(\frac{3\Delta}{8}\right)^2,\\
E_T:=\{\mathsf w_{\mathrm{dec}}:
v_{\mathrm{mid}}\le T(\mathsf w_{\mathrm{dec}})\},
\qquad
E_{\mathrm{err},0}:=E_T,
\qquad
E_{\mathrm{err},1}:=E_T^c,\\
E_{\mathrm{esc},\ell}
:=\{P:\Delta/8<|V(P)-v_\ell|\},
\qquad \ell\in\{0,1\},\\
\mathcal E_T(P)
:=
\int
(T(\mathsf w_{\mathrm{dec}})-V(P))^2\,
\mathsf C_{\mathrm{dec}}(P)(d\mathsf w_{\mathrm{dec}}),
\qquad
\mathcal B_{\ell,T}
:=
\int\mathcal E_T(P)\,\pi_\ell(dP).
\end{gathered}
\]
The last two integrals are nonnegative integrals with values in
\([0,\infty]\).

For each prior index, the following pointwise inequality holds over the entire auxiliary parameter space:
\[
c_{\mathrm{err}}\,
\mathsf C_{\mathrm{dec}}(P)(E_{\mathrm{err},\ell})
\le
\mathcal E_T(P)
+
c_{\mathrm{err}}\mathbf1_{E_{\mathrm{esc},\ell}}(P).
\]
To see this, first suppose \(P\in E_{\mathrm{esc},\ell}\). The left side is at most \(c_{\mathrm{err}}\), since the row is a probability measure. Otherwise,
\(|V(P)-v_\ell|\le\Delta/8\). For \(\ell=0\), on \(E_T\),
\[
T(\mathsf w_{\mathrm{dec}})-V(P)
\ge v_{\mathrm{mid}}-(v_0+\Delta/8)
=\frac{3\Delta}{8}.
\]
For \(\ell=1\), on \(E_T^c\),
\[
V(P)-T(\mathsf w_{\mathrm{dec}})
>
(v_1-\Delta/8)-v_{\mathrm{mid}}
=\frac{3\Delta}{8}.
\]
In either case the squared loss dominates
\(c_{\mathrm{err}}\mathbf1_{E_{\mathrm{err},\ell}}\), proving the pointwise inequality by integration.

Averaging under each prior gives
\[
c_{\mathrm{err}}\mathsf L_0(E_T)
\le\mathcal B_{0,T}+c_{\mathrm{err}}\pi_0(E_{\mathrm{esc},0}),
\qquad
c_{\mathrm{err}}\mathsf L_1(E_T^c)
\le\mathcal B_{1,T}+c_{\mathrm{err}}\pi_1(E_{\mathrm{esc},1}).
\]
The testing inequality follows directly from total variation:
\[
\mathsf L_0(E_T)+\mathsf L_1(E_T^c)
=
1+\mathsf L_0(E_T)-\mathsf L_1(E_T)
\ge1-\omega.
\]
If \(\max\{\mathcal B_{0,T},\mathcal B_{1,T}\}=\infty\), the Bayes-risk lower bound below follows immediately. Otherwise, both Bayes risks are finite, so summing the preceding inequalities and using the escape bounds yields
\[
\begin{aligned}
c_{\mathrm{err}}(1-\omega-2\eta_0)
&\le \mathcal B_{0,T}+\mathcal B_{1,T}\\
&\le2\max\{\mathcal B_{0,T},\mathcal B_{1,T}\}.
\end{aligned}
\]
Consequently, in both cases,
\[
\frac{9\Delta^2}{128}(1-\omega-2\eta_0)
\le\max\{\mathcal B_{0,T},\mathcal B_{1,T}\}.
\]

Each supported prior average is bounded by the supremum on the causal class:
\[
\mathcal B_{\ell,T}
=
\int_{\mathcal V_d}\mathcal E_T(P)\,\pi_\ell(dP)
\le
\sup_{P\in\mathcal V_d}
\int
(T(\mathsf w_{\mathrm{dec}})-V(P))^2\,
\mathsf L_{\mathrm{dec}}(P)(d\mathsf w_{\mathrm{dec}}).
\]
If \(1-\omega-2\eta_0\ge0\), combining these inequalities gives the stated positive-part bound. If \(1-\omega-2\eta_0<0\), its positive part is zero and the bound follows from nonnegativity of squared risk. Finally, equality of the decision laws on \(\mathcal V_d\) identifies the supremum:
\[
\begin{aligned}
&\sup_{P\in\mathcal V_d}
\int
(T(\mathsf w_{\mathrm{dec}})-V(P))^2\,
\mathsf L_{\mathrm{dec}}(P)(d\mathsf w_{\mathrm{dec}})\\
&\qquad=
\sup_{P\in\mathcal V_d}
\mathbb E_{P,Q}\bigl[(T(\mathbf Z,R,U)-V(P))^2\bigr].
\end{aligned}
\]
This proves the risk assertion with its extended-nonnegative interpretation.
% lean: CausalSmith.Stat.LdpOptvalueUniformFrontier.separated_prior_squared_risk

\item \textit{Coverage of the target neighborhoods.}
Fix an interval satisfying the stated endpoint and coverage conditions. Write
\[
I(\mathsf w_{\mathrm{dec}})
=
[\underline I(\mathsf w_{\mathrm{dec}}),
 \overline I(\mathsf w_{\mathrm{dec}})],
\qquad
\operatorname{length}I(\mathsf w_{\mathrm{dec}})
=
\overline I(\mathsf w_{\mathrm{dec}})
-\underline I(\mathsf w_{\mathrm{dec}}).
\]
For \(\ell\in\{0,1\}\), define
\[
J_\ell:=[v_\ell-\Delta/8,v_\ell+\Delta/8],
\qquad
E_{I,\ell}:=
\{\mathsf w_{\mathrm{dec}}:
\underline I(\mathsf w_{\mathrm{dec}})\le v_\ell+\Delta/8,\
v_\ell-\Delta/8\le\overline I(\mathsf w_{\mathrm{dec}})\}.
\]
These are Borel events, and ordered endpoints make them exactly the events
\(I\cap J_\ell\ne\varnothing\).

For \(P\in\mathcal V_d\setminus E_{\mathrm{esc},\ell}\),
\[
\{\,\mathsf w_{\mathrm{dec}}:
V(P)\in I(\mathsf w_{\mathrm{dec}})\,\}
\subseteq E_{I,\ell},
\]
since \(V(P)\in J_\ell\). Global coverage therefore gives
\[
0.90
\le
\mathsf L_{\mathrm{dec}}(P)(E_{I,\ell})
+\mathbf1_{E_{\mathrm{esc},\ell}}(P),
\qquad P\in\mathcal V_d.
\]
For a law in the escape event, this inequality follows from the indicator being one and the event probability being nonnegative; for a law outside it, it follows from the displayed inclusion and coverage. Integrating under the supported prior proves
\[
0.90
\le
\mathsf L_\ell(E_{I,\ell})
+\pi_\ell(E_{\mathrm{esc},\ell}),
\qquad
\mathsf L_\ell(E_{I,\ell})\ge0.90-\eta_0.
\]
% lean: CausalSmith.Stat.LdpOptvalueUniformFrontier.supported_prior_neighborhood_mass

\item \textit{Intersection and expected length.}
The total-variation bound transfers the second neighborhood event to the first mixture:
\[
\mathsf L_0(E_{I,1})
\ge\mathsf L_1(E_{I,1})-\omega
\ge0.90-\eta_0-\omega.
\]
Under the first mixture,
\[
\begin{aligned}
\mathsf L_0(E_{I,0}\cap E_{I,1})
&=
\mathsf L_0(E_{I,0})
-\mathsf L_0(E_{I,0}\setminus E_{I,1})\\
&\ge
\mathsf L_0(E_{I,0})-\mathsf L_0(E_{I,1}^c)\\
&=
\mathsf L_0(E_{I,0})+\mathsf L_0(E_{I,1})-1\\
&\ge0.80-2\eta_0-\omega.
\end{aligned}
\]
Combining this with nonnegativity of probability gives
\[
\mathsf L_0(E_{I,0}\cap E_{I,1})
\ge(0.80-2\eta_0-\omega)_+.
\]

On the intersection event the endpoints satisfy
\[
\underline I(\mathsf w_{\mathrm{dec}})\le v_0+\Delta/8,
\qquad
\overline I(\mathsf w_{\mathrm{dec}})\ge v_1-\Delta/8.
\]
Hence
\[
\operatorname{length}I(\mathsf w_{\mathrm{dec}})
\ge
(v_1-\Delta/8)-(v_0+\Delta/8)
=\frac{3\Delta}{4}.
\]
Outside that event, ordered endpoints give nonnegative length. Thus the pointwise indicator bound and its integrated form are
\[
\operatorname{length}I(\mathsf w_{\mathrm{dec}})
\ge
\frac{3\Delta}{4}\,
\mathbf1_{E_{I,0}\cap E_{I,1}}(\mathsf w_{\mathrm{dec}}),
\]
\[
\int\operatorname{length}I\,d\mathsf L_0
\ge
\frac{3\Delta}{4}\,
\mathsf L_0(E_{I,0}\cap E_{I,1})
\ge
\frac{3\Delta}{4}(0.80-2\eta_0-\omega)_+.
\]
Finally, nonnegative integration through the mixture and causal support yield
\[
\begin{aligned}
\int\operatorname{length}I\,d\mathsf L_0
&=
\int_{\mathcal V_d}
\left[
\int\operatorname{length}I\,d\mathsf L_{\mathrm{dec}}(P)
\right]\pi_0(dP)\\
&\le
\sup_{P\in\mathcal V_d}
\int\operatorname{length}I\,d\mathsf L_{\mathrm{dec}}(P)\\
&=
\sup_{P\in\mathcal V_d}
\mathbb E_{P,Q}\operatorname{length}I(\mathbf Z,R,U).
\end{aligned}
\]
The same equality of decision laws transfers the global coverage condition to every completed row on \(\mathcal V_d\), so this establishes the interval assertion.
% lean: CausalSmith.Stat.LdpOptvalueUniformFrontier.separated_prior_interval_length
\end{enumerate}
\end{proof}

The squared-error bound has the scale of the squared welfare separation. Its positive-part factor records the remaining testing separation after accounting for transcript proximity and the two target-concentration failures. Causal support ensures that averaging over either prior compares decisions within the same model used by the worst-law risk. Extended nonnegative expectations include every Borel real-valued estimator in this reduction.

For intervals, honesty is a finite-sample coverage requirement at every law in \(\mathcal V_d\). Consequently, the coverage requirement applies throughout the support of both priors. The factor \(0.80\) reflects the two \(0.90\) coverage requirements, with target-concentration and transcript-proximity penalties appearing explicitly. Connectedness gives the geometric content of the length bound: an interval containing target values within \(\Delta/8\) of both centers must span at least \(3\Delta/4\). Ordered Borel endpoints make coverage and length measurable, and their range agrees with the welfare range in \cref{obj:def:causal-model}.

Applied to independent-coordinate priors on the symmetric causal family, \cref{obj:lem:adaptive-moment-contraction} supplies the transcript-proximity bound. The bounded-mean concentration result in \cref{obj:lem:bounded-mean-concentration} supplies the corresponding control of average coordinate contributions to welfare. The separation supplied by approximation duality can therefore enter both decision criteria through \cref{obj:lem:separated-value-mixtures}: squared-error risk scales with \(\Delta^2\), and expected connected interval length scales with \(\Delta\).

\subsection{A dimension-free calibration}

A further comparison supplies a lower bound uniform in the number of strata. It uses the same measurable-density property and the sampling and randomness conditions already required for sequential contraction. The decision criteria are the minimax squared-error risk and globally honest connected interval length defined in \cref{obj:def:risk,obj:def:honest-length}. The next result states their dimension-free dependence on the finite privacy resources.

\begin{lemmav}{lem:dimension-free-private-two-point}[Dimension-free private two-point lower bounds]
Assume the following measurable Radon--Nikodym property: for every measurable parameter space \(H\), every standard-Borel output space \(Z\), and every pair of finite kernels \(K,L\) from \(H\) to \(Z\) satisfying \(K(h,\cdot)\ll L(h,\cdot)\) for every \(h\in H\), there exists a jointly measurable density \(f:H\times Z\to[0,\infty]\) such that, for every \(h\in H\) and every measurable \(E\subseteq Z\),
\[
K(h,E)=\int_E f(h,z)\,L(h,dz).
\]
Then there exists a universal constant \(c_1>0\) such that, for every sample size \(n\in\mathbb N\), dimension \(d\in\mathbb N\), privacy budget \(\varepsilon\in\mathbb R\), and sampling scheme satisfying
\begin{itemize}
\item \textbf{(Resources.)} \(n\ge2\), \(d\ge2\), and \(0<\varepsilon\le1\).
\item \textbf{(Participant sampling.)} \cref{obj:ass:iid-people}.
\item \textbf{(Randomization.)} \cref{obj:ass:independent-randomness}.
\end{itemize}
the minimax squared-error risk \(\mathfrak R_{\mathcal C}(n,d,\varepsilon)\) and globally \(0.90\)-honest connected interval length \(\mathfrak H_{\mathcal C}(n,d,\varepsilon)\), defined in \cref{obj:def:risk,obj:def:honest-length} for that sampling scheme, satisfy, for every \(\mathcal C\in\{\mathrm{NI},\mathrm{SI}\}\),
\[
\mathfrak R_{\mathcal C}(n,d,\varepsilon)
\ge c_1\min\left\{1,\frac{1}{n\varepsilon^2}\right\},
\qquad
\mathfrak H_{\mathcal C}(n,d,\varepsilon)
\ge c_1\min\left\{1,\frac{1}{\sqrt n\,\varepsilon}\right\}.
\]
\end{lemmav}

\begin{proof}[Proof of \cref{obj:lem:dimension-free-private-two-point}]
Set
\[
c_1=\frac{1}{8192}>0.
\]
Fix admissible \(n,d,\varepsilon\) and a sampling scheme satisfying
\cref{obj:ass:iid-people,obj:ass:independent-randomness}.

\begin{enumerate}
\item Fix \(\mathcal C\in\{\mathrm{NI},\mathrm{SI}\}\) and an admissible protocol
\[
Q\in\mathfrak Q_{\mathcal C}(n,d,\varepsilon).
\]
If \(\mathcal C=\mathrm{NI}\), the inclusion in
\cref{obj:def:noninteractive-class} gives
\(Q\in\mathfrak Q_{\mathrm{SI}}(n,d,\varepsilon)\).
If \(\mathcal C=\mathrm{SI}\), this membership is the stipulated admissibility condition. Thus in both cases,
\[
Q\in\mathfrak Q_{\mathrm{SI}}(n,d,\varepsilon).
\]
We establish bounds uniform over all decisions based on this protocol.
% lean: CausalSmith.Stat.LdpOptvalueUniformFrontier.dimension_free_private_two_point

\item Define the positive scale and two-point amplitude by
\[
\kappa_{\mathrm{tp}}=\sqrt n\,\varepsilon>0,
\qquad
\alpha_{\mathrm{tp}}
=\frac18\min\left\{1,\frac1{\kappa_{\mathrm{tp}}}\right\}.
\]
Since \(n\ge2\) and \(\varepsilon>0\),
\[
0<\alpha_{\mathrm{tp}}\le\frac18\le\frac12,
\qquad
\alpha_{\mathrm{tp}}\sqrt n\,\varepsilon
=\frac{\kappa_{\mathrm{tp}}}{8}
  \min\left\{1,\frac1{\kappa_{\mathrm{tp}}}\right\}
\le\frac18.
\]
For the squared-rate calibration, splitting according to whether
\(\kappa_{\mathrm{tp}}\le1\) or \(\kappa_{\mathrm{tp}}>1\) gives
\[
\left(\min\left\{1,\frac1{\kappa_{\mathrm{tp}}}\right\}\right)^2
=
\begin{cases}
1,&\kappa_{\mathrm{tp}}\le1,\\
\kappa_{\mathrm{tp}}^{-2},&\kappa_{\mathrm{tp}}>1
\end{cases}
=
\min\left\{1,\frac1{\kappa_{\mathrm{tp}}^2}\right\}.
\]
Here the first case uses
\(1/\kappa_{\mathrm{tp}}\ge1\) and
\(1/\kappa_{\mathrm{tp}}^2\ge1\), while the second uses the reverse inequalities.
As
\[
\kappa_{\mathrm{tp}}^2=n\varepsilon^2,
\]
we obtain the two numerical comparisons
\[
\frac{63\alpha_{\mathrm{tp}}^2}{4096}
=
\frac{63}{262144}\min\left\{1,\frac1{n\varepsilon^2}\right\}
\ge
c_1\min\left\{1,\frac1{n\varepsilon^2}\right\},
\]
and
\[
\frac{81\alpha_{\mathrm{tp}}}{320}
=
\frac{81}{2560}\min\left\{1,\frac1{\sqrt n\,\varepsilon}\right\}
\ge
c_1\min\left\{1,\frac1{\sqrt n\,\varepsilon}\right\}.
\]
% lean: CausalSmith.Stat.LdpOptvalueUniformFrontier.two_point_amplitude_calibration

\item Define the two contrast vectors, coordinate priors, and privacy contraction scale by
\[
\theta^{(0)}_j=0,\qquad
\theta^{(1)}_j=\alpha_{\mathrm{tp}}
\quad(j\in[d]),
\qquad
\nu_0=\operatorname{Dirac}_0,\qquad
\nu_1=\operatorname{Dirac}_{\alpha_{\mathrm{tp}}},
\qquad
\beta=\frac{e^\varepsilon-1}{2d},
\]
where \(\operatorname{Dirac}_x\) denotes the unit point mass at \(x\).
These priors satisfy
\[
\nu_0([-\alpha_{\mathrm{tp}},\alpha_{\mathrm{tp}}])
=\nu_1([-\alpha_{\mathrm{tp}},\alpha_{\mathrm{tp}}])=1,
\qquad
\int u^0\,\nu_0(du)=\int u^0\,\nu_1(du)=1.
\]
Thus their moments match through order zero. Each coordinate under the corresponding product prior equals its point-mass location almost surely; intersecting these finitely many probability-one events yields
\[
\nu_\ell^{\otimes d}\bigl(\{\theta^{(\ell)}\}\bigr)=1
\qquad(\ell\in\{0,1\}).
\]
Consequently, writing the two joint seed-transcript laws as
\[
N_\ell=\mathcal L_{G_{\theta^{(\ell)}},Q}(R,\mathbf Z)
\qquad(\ell\in\{0,1\}),
\]
the mixtures in \cref{obj:lem:adaptive-moment-contraction} reduce to
\[
\mathsf M_\ell^Q=N_\ell
\qquad(\ell\in\{0,1\}).
\]
Apply that result with amplitude \(\alpha_{\mathrm{tp}}\), these coordinate priors, and \(k=1\). Its hypotheses hold by the assumed measurable Radon--Nikodym property, the sampling conditions, sequential membership of \(Q\), and \(1\le n\). It gives
\[
\begin{aligned}
\operatorname{TV}(N_0,N_1)
&\le
\min\left\{1,\,
d\alpha_{\mathrm{tp}}\sqrt{\binom n1}\,\beta\right\}\\
&\le
\frac{\alpha_{\mathrm{tp}}\sqrt n\,(e^\varepsilon-1)}2.
\end{aligned}
\]
For \(0<\varepsilon\le1\), the standard exponential estimate gives
\[
e^\varepsilon-1
\le |e^\varepsilon-1|
\le2|\varepsilon|
=2\varepsilon.
\]
Therefore,
\[
\operatorname{TV}(N_0,N_1)
\le\alpha_{\mathrm{tp}}\sqrt n\,\varepsilon
\le\frac18.
\]
% lean: CausalSmith.Stat.LdpOptvalueUniformFrontier.constant_contrast_tv_bound

\item Both contrast vectors lie in \([-1/2,1/2]^d\). By
\cref{obj:prop:signed-causal-bridge}, the symmetric laws of
\cref{obj:def:symmetric-law} satisfy
\[
G_{\theta^{(0)}},G_{\theta^{(1)}}\in\mathcal V_d.
\]
Since \(d>0\), their average absolute contrasts are
\[
F(\theta^{(0)})=\frac1d\sum_{j=1}^d0=0,
\qquad
F(\theta^{(1)})
=\frac1d\sum_{j=1}^d\alpha_{\mathrm{tp}}
=\alpha_{\mathrm{tp}}.
\]
The same cited result therefore gives
\[
V(G_{\theta^{(0)}})=\frac12,
\qquad
V(G_{\theta^{(1)}})=\frac12+\frac{\alpha_{\mathrm{tp}}}{2}.
\]
Define the full-data priors and their welfare centers by
\[
\Pi_0=\operatorname{Dirac}_{G_{\theta^{(0)}}},
\qquad
\Pi_1=\operatorname{Dirac}_{G_{\theta^{(1)}}},
\qquad
v_0=\frac12,
\qquad
v_1=\frac12+\frac{\alpha_{\mathrm{tp}}}{2},
\qquad
\Delta=v_1-v_0=\frac{\alpha_{\mathrm{tp}}}{2}>0.
\]
Their causal support and concentration properties are
\[
\Pi_\ell(\mathcal V_d)=1,
\qquad
\Pi_\ell\left(\left\{P:
|V(P)-v_\ell|>\frac{\Delta}{8}\right\}\right)=0
\qquad(\ell\in\{0,1\}).
\]
Mixing over a unit point mass returns the corresponding decision law. These decision laws are probability measures because the participant and randomization laws are probability measures and the message kernels are Markov kernels. Their seed-transcript marginals are precisely \(N_0,N_1\).

Set
\[
\eta_0=0,\qquad \omega=\frac18.
\]
The preceding displays discharge the hypotheses of
\cref{obj:lem:separated-value-mixtures}. Hence every admissible estimator \(T\) satisfies
\[
\sup_{P\in\mathcal V_d}
\mathbb E_{P,Q}\bigl[(T(\mathbf Z,R,U)-V(P))^2\bigr]
\ge
\frac{9\Delta^2}{128}\cdot\frac78
=
\frac{63\alpha_{\mathrm{tp}}^2}{4096}.
\]
Every admissible connected interval \(I\) with global coverage \(0.90\) satisfies
\[
\sup_{P\in\mathcal V_d}
\mathbb E_{P,Q}\operatorname{length}I(\mathbf Z,R,U)
\ge
\frac{3\Delta}{4}\cdot\frac{27}{40}
=
\frac{81\alpha_{\mathrm{tp}}}{320}.
\]
Here \(7/8=1-\omega-2\eta_0\) and
\(27/40=0.80-2\eta_0-\omega\), so both positive-part factors equal the displayed positive numbers.
% lean: CausalSmith.Stat.LdpOptvalueUniformFrontier.constant_contrast_decision_lower

\item Combining these bounds with the amplitude calibration yields
\[
\sup_{P\in\mathcal V_d}
\mathbb E_{P,Q}\bigl[(T(\mathbf Z,R,U)-V(P))^2\bigr]
\ge
c_1\min\left\{1,\frac1{n\varepsilon^2}\right\}
\]
for every admissible estimator, and
\[
\sup_{P\in\mathcal V_d}
\mathbb E_{P,Q}\operatorname{length}I(\mathbf Z,R,U)
\ge
c_1\min\left\{1,\frac1{\sqrt n\,\varepsilon}\right\}
\]
for every admissible globally honest interval.
% lean: CausalSmith.Stat.LdpOptvalueUniformFrontier.dimension_free_private_two_point_decisions

\item These inequalities hold for every admissible protocol in either selected class. Taking the infima over protocols and decisions in
\cref{obj:def:risk,obj:def:honest-length} therefore gives
\[
\mathfrak R_{\mathcal C}(n,d,\varepsilon)
\ge
c_1\min\left\{1,\frac1{n\varepsilon^2}\right\},
\qquad
\mathfrak H_{\mathcal C}(n,d,\varepsilon)
\ge
c_1\min\left\{1,\frac1{\sqrt n\,\varepsilon}\right\}.
\]
The infima are interpreted in the extended nonnegative reals, whose infimum over an empty decision class is \(+\infty\). Thus the same lower bounds cover that case as well.
% lean: CausalSmith.Stat.LdpOptvalueUniformFrontier.dimension_free_private_two_point
\end{enumerate}
\end{proof}

The two bounds depend on the participant count and privacy budget through \(n\varepsilon^2\), uniformly over every admissible dimension. When this quantity is at most one, they give constant lower bounds; above one, they give inverse-resource squared-error and inverse-square-root-resource length bounds. Both conclusions cover the noninteractive and sequential protocol classes, including the independent public and analyst randomness permitted by their definitions.

This comparison supplies a finite-resource baseline alongside the dimension-dependent moment-matching bounds. In particular, it applies at fixed small dimensions and throughout the stated sample-size range. Its universal constant allows the lower-bound calibration to retain the same dependence on privacy resources across dimensions, complementing the finite-sample calibrations in \cref{obj:lem:finite-private-polynomial-calibration} and the two-cell result in \cref{obj:thm:two-cell-calibration}.


\section{Guide to the main proofs and verification scope}

The welfare analysis begins with the signed causal representation in \cref{obj:prop:signed-causal-bridge}, which separates the baseline mean from the average absolute treatment contrast. For the attaining procedure in \cref{obj:def:frontier-handle}, the approximation bound in \cref{obj:lem:cai-low-chebyshev-approximation} controls polynomial bias. The moment identities and dependence bound in \cref{obj:lem:private-moment-and-dependence} control evaluation variability, while \cref{obj:lem:bounded-mean-concentration} controls pilot selection. These ingredients enter the finite calibration in \cref{obj:lem:finite-private-polynomial-calibration}, connecting the component bounds to the specified procedure.

The converse follows the complementary dependency chain. Approximation duality in \cref{obj:lem:cai-low-moment-duality} supplies moment-matched priors with separated absolute-value expectations. The measurable-density representation in \cref{obj:lem:measurable-kernel-radon-nikodym} supplies the density condition used by \cref{obj:lem:adaptive-moment-contraction}, which controls proximity of the resulting sequential private experiments. Concentration of the prior-induced welfare targets then permits \cref{obj:lem:separated-value-mixtures} to translate experiment proximity into bounds for estimation and honest interval length. The comparison in \cref{obj:lem:dimension-free-private-two-point} supplies the dimension-uniform component of the converse. The signed causal correspondence also carries the paired experiment into \cref{obj:thm:paired-uniform-tv-frontier}; inclusion of that family in the full simplex supplies the transfer in \cref{obj:thm:full-simplex-tv-converse}.

The causal conclusions concern independent population sampling under \cref{obj:def:causal-model,obj:ass:iid-people}. The model specifies binary outcomes, public uniform stratum probabilities, fair treatment assignment conditional on the stratum and both potential outcomes, observed-outcome consistency, and conditional arm means in \([1/4,3/4]\). These conditions give the welfare target its randomized-trial interpretation. The total-variation conclusions instead use independent draws from the paired family or full simplex specified in \cref{obj:def:tv-corollary-models}, with the public uniform distribution as reference.

Collection follows the one-message protocol classes in \cref{obj:def:sequential-class,obj:def:noninteractive-class}, with the independent public seed and analyst randomizer specified in \cref{obj:ass:independent-randomness}. Pure local privacy compares every pair of complete observed records at every public-seed realization and preceding message history, as required by \cref{obj:ass:full-record-privacy}. The seed and message spaces are standard Borel. Sequential collection permits dependence on preceding messages; the noninteractive subclass imposes history independence. The paired-input formulation in \cref{obj:def:tv-private-decision-values} imposes the corresponding privacy comparison across all paired symbols. Estimation and coverage are evaluated under the sampling, message-generation, and randomization laws specified by these definitions.

\paragraph{Verification scope.}
The theorem statements and proofs are machine-checked in Lean 4 under their stated hypotheses. For \cref{obj:thm:uniform-private-value-frontiers,obj:thm:paired-uniform-tv-frontier,obj:thm:full-simplex-tv-converse}, checking covers use of the published Bernstein premise and all remaining steps. Writing \(e_K\) for the best uniform approximation error defined in \cref{obj:lem:cai-low-moment-duality}, that premise is \(2k e_{2k}\to\beta_*\), where the Bernstein constant \(\beta_*\) satisfies \(0.278<\beta_*<0.286\) \citep[Section 3.1, definition of the Bernstein constant and immediately following paragraph, p.~8]{CaiLow2011Nonsmooth}. The theorem-local verification disclosures identify the source-proof boundary: the published proof of this premise remains external to the Lean checking scope, and the dependent conclusions are certified conditional on it.

\section{Proofs of the main results}\label{sec:deferred-proofs}

\begin{proof}[Proof of \cref{obj:prop:signed-causal-bridge}]
Throughout, \(d\ge2\), so every stratum has positive probability under a law in \(\mathcal V_d\).

\begin{enumerate}
\item \emph{Identification of the signed mean.}
Fix \(P\in\mathcal V_d\) and \(j\in[d]\). By \cref{obj:ass:fair-randomization}, summing over the potential outcome for the other arm gives, for \(a\in\{0,1\}\),
\[
\begin{aligned}
&P(X=j,A=a,Y^a=1)\\
&\qquad=\sum_{y\in\{0,1\}}
 P(X=j,A=a,Y^a=1,Y^{1-a}=y)\\
&\qquad=\frac12\sum_{y\in\{0,1\}}
 P(X=j,Y^a=1,Y^{1-a}=y)
 =\frac12P(X=j,Y^a=1).
\end{aligned}
\]
Summing the analogous identity over both values of \(Y^a\) yields
\[
P(X=j,A=a)=\frac12P(X=j)=\frac1{2d},
\]
where the last equality uses \cref{obj:ass:uniform-covariates}.

By \cref{obj:ass:consistency}, the following indicator identity holds almost surely:
\[
\begin{aligned}
\mathbf1_{\{X=j\}}S
={}&2\mathbf1_{\{X=j,A=1,Y^1=1\}}
-\mathbf1_{\{X=j,A=1\}}\\
&-2\mathbf1_{\{X=j,A=0,Y^0=1\}}
+\mathbf1_{\{X=j,A=0\}}.
\end{aligned}
\]
Integrating and using the fair-assignment identities gives
\[
\mathbb E_P[\mathbf1_{\{X=j\}}S]
=P(X=j,Y^1=1)-P(X=j,Y^0=1).
\]
Dividing by \(P(X=j)=1/d>0\), we obtain
\[
\mathbb E_P[S\mid X=j]
=\mu_{1j}(P)-\mu_{0j}(P)
=\tau_j(P).
\]
% lean: CausalSmith.Stat.LdpOptvalueUniformFrontier.signedCellMean_eq_contrast

\item \emph{Identification of the paired law.}
Because \(S\) takes values in \(\{-1,1\}\), its two signed atoms satisfy
\[
\begin{aligned}
P(X=j,S=1)+P(X=j,S=-1)&=\frac1d,\\
P(X=j,S=1)-P(X=j,S=-1)
&=\mathbb E_P[\mathbf1_{\{X=j\}}S]
=\frac{\tau_j(P)}d.
\end{aligned}
\]
Solving these two equations gives, for \(s\in\{-1,1\}\),
\[
P(X=j,S=s)=\frac{1+s\tau_j(P)}{2d}.
\]
Moreover, \cref{obj:ass:interior-means} gives
\[
-\frac12
=\frac14-\frac34
\le \tau_j(P)
\le \frac34-\frac14
=\frac12.
\]
For any \(\theta\in[-1/2,1/2]^d\), the paired masses satisfy
\[
\frac{1+s\theta_j}{2d}\ge0,
\qquad
\sum_{j=1}^d\sum_{s\in\{-1,1\}}
\frac{1+s\theta_j}{2d}
=\sum_{j=1}^d\frac1d=1.
\]
Thus \(P_{\tau(P)}\) is a probability law. Equality of the displayed atom probabilities on the finite alphabet proves
\[
\mathcal L_P(X,S)=P_{\tau(P)}.
\]
% lean: CausalSmith.Stat.LdpOptvalueUniformFrontier.causal_signed_law

\item \emph{Decomposition of the causal value.}
For each stratum, the two cases
\(\mu_{0j}(P)\le\mu_{1j}(P)\) and
\(\mu_{1j}(P)\le\mu_{0j}(P)\) give
\[
\max\{\mu_{0j}(P),\mu_{1j}(P)\}
=\frac{\mu_{0j}(P)+\mu_{1j}(P)}2
+\frac{|\tau_j(P)|}2.
\]
Consequently,
\[
V(P)
=\frac1d\sum_{j=1}^d
 \frac{\mu_{0j}(P)+\mu_{1j}(P)}2
+\frac{F(\tau(P))}2.
\]

To identify the first term, consistency gives the almost-sure identity
\[
\mathbf1_{\{X=j\}}Y
=\mathbf1_{\{X=j,A=0,Y^0=1\}}
+\mathbf1_{\{X=j,A=1,Y^1=1\}}.
\]
The fair-assignment identity
\(P(X=j,A=a,Y^a=1)=\frac12P(X=j,Y^a=1)\)
therefore yields
\[
\begin{aligned}
\mathbb E_P[\mathbf1_{\{X=j\}}Y]
&=\frac{P(X=j,Y^0=1)+P(X=j,Y^1=1)}2\\
&=\frac1d\,\frac{\mu_{0j}(P)+\mu_{1j}(P)}2.
\end{aligned}
\]
Since the strata partition the sample space,
\[
B(P)=\mathbb E_P[Y]
=\sum_{j=1}^d\mathbb E_P[\mathbf1_{\{X=j\}}Y]
=\frac1d\sum_{j=1}^d
 \frac{\mu_{0j}(P)+\mu_{1j}(P)}2.
\]
Substitution proves
\[
V(P)=B(P)+\frac{F(\tau(P))}2.
\]
% lean: CausalSmith.Stat.LdpOptvalueUniformFrontier.causal_value_decomposition

\item \emph{Total variation from the uniform law.}
Fix \(\theta\in[-1/2,1/2]^d\). Both \(P_\theta\) and
\(\mathsf U_{2d}=P_0\) are probability laws by the mass calculation above.
Define the event of nonnegative atom differences by
\[
\mathcal E_{\theta,+}
:=\left\{(j,s)\in[d]\times\{-1,1\}:
P_\theta(j,s)-\frac1{2d}\ge0\right\}.
\]
The atom differences sum to zero:
\[
\sum_{j=1}^d\sum_{s\in\{-1,1\}}
\left(P_\theta(j,s)-\frac1{2d}\right)=0.
\]
Splitting this sum into its nonnegative and negative terms gives
\[
P_\theta(\mathcal E_{\theta,+})
-\mathsf U_{2d}(\mathcal E_{\theta,+})
=\frac12\sum_{j=1}^d\sum_{s\in\{-1,1\}}
\left|P_\theta(j,s)-\frac1{2d}\right|.
\]
For every event \(\mathcal E\subseteq[d]\times\{-1,1\}\), its sum of atom
differences lies between the sum of all negative differences and the sum
of all nonnegative differences. Hence
\[
\left|P_\theta(\mathcal E)-\mathsf U_{2d}(\mathcal E)\right|
\le\frac12\sum_{j=1}^d\sum_{s\in\{-1,1\}}
\left|P_\theta(j,s)-\frac1{2d}\right|.
\]
The event \(\mathcal E_{\theta,+}\) attains this bound, proving the finite
atom formula for total variation. For each \(j\), its two summands reduce to
\[
\begin{aligned}
\sum_{s\in\{-1,1\}}
\left|\frac{1+s\theta_j}{2d}-\frac1{2d}\right|
&=\left|\frac{\theta_j}{2d}\right|
+\left|-\frac{\theta_j}{2d}\right|\\
&=\frac{|\theta_j|}{d}.
\end{aligned}
\]
Therefore
\[
\operatorname{TV}(P_\theta,\mathsf U_{2d})
=\frac1{2d}\sum_{j=1}^d|\theta_j|
=\frac{F(\theta)}2.
\]
Applying this identity at \(\theta=\tau(P)\) proves
\[
V(P)=B(P)+\operatorname{TV}(P_{\tau(P)},\mathsf U_{2d}).
\]
% lean: CausalSmith.Stat.LdpOptvalueUniformFrontier.tvFromUniform_pairedLaw

\item \emph{The symmetric causal construction.}
Fix \(\theta\in[-1/2,1/2]^d\). For
\(a,y\in\{0,1\}\) and \(j\in[d]\), define the Bernoulli masses
\[
\pi_{aj}(y)
:=
\left(\frac{1+(2a-1)\theta_j}{2}\right)^y
\left(\frac{1-(2a-1)\theta_j}{2}\right)^{1-y}.
\]
Their success probabilities belong to \([1/4,3/4]\), so
\[
\pi_{aj}(y)\ge0,
\qquad
\sum_{y\in\{0,1\}}\pi_{aj}(y)=1.
\]
By \cref{obj:def:symmetric-law}, the full-data atom masses are
\[
\begin{aligned}
&G_\theta(X=j,A=a,Y=y,Y^0=y^0,Y^1=y^1)\\
&\qquad=\frac1{2d}\,
\pi_{0j}(y^0)\pi_{1j}(y^1)
\mathbf1_{\{y=(1-a)y^0+ay^1\}},
\qquad a,y,y^0,y^1\in\{0,1\}.
\end{aligned}
\]
Summing first over \(y\), then over the two potential outcomes and
assignment, gives total mass
\[
\sum_{j=1}^d\sum_{a\in\{0,1\}}\frac1{2d}
\left(\sum_{y^0\in\{0,1\}}\pi_{0j}(y^0)\right)
\left(\sum_{y^1\in\{0,1\}}\pi_{1j}(y^1)\right)
=1.
\]
Thus \(G_\theta\) is a probability law. The same atom formula gives
\[
G_\theta(X=j)=\frac1d
\]
and
\[
\begin{aligned}
&G_\theta(X=j,A=a,Y^0=y^0,Y^1=y^1)\\
&\qquad=\frac1{2d}\pi_{0j}(y^0)\pi_{1j}(y^1)
=\frac12G_\theta(X=j,Y^0=y^0,Y^1=y^1).
\end{aligned}
\]
These are the conditions in
\cref{obj:ass:uniform-covariates,obj:ass:fair-randomization}.
The atom formula assigns all mass to \(Y=Y_A^{\mathrm{pot}}\), supplying
\cref{obj:ass:consistency}. Finally, summing over the other coordinates yields
\[
G_\theta(X=j,Y^a=1)
=\frac1d\,\frac{1+(2a-1)\theta_j}{2},
\]
and hence
\[
\mu_{0j}(G_\theta)=\frac{1-\theta_j}{2},
\qquad
\mu_{1j}(G_\theta)=\frac{1+\theta_j}{2}.
\]
Both means lie in \([1/4,3/4]\), as required by
\cref{obj:ass:interior-means}. All four conditions of
\cref{obj:def:causal-model} therefore hold.
% lean: CausalSmith.Stat.LdpOptvalueUniformFrontier.symmetricLaw_causalModel

\item \emph{Symmetric value and the ancillary assignment.}
The arm means just computed satisfy
\[
\tau_j(G_\theta)=\theta_j,
\qquad
\frac{\mu_{0j}(G_\theta)+\mu_{1j}(G_\theta)}2=\frac12.
\]
Applying the armwise maximum decomposition gives
\[
V(G_\theta)
=\frac1d\sum_{j=1}^d\frac12+\frac{F(\theta)}2
=\frac12+\frac{F(\theta)}2.
\]
% lean: CausalSmith.Stat.LdpOptvalueUniformFrontier.symmetricLaw_value

Define the observed-record reconstruction map on the whole paired
alphabet and both assignment values by
\[
\operatorname{rec}(j,s,a)
:=\left(j,a,\frac{1+s(2a-1)}2\right),
\qquad
j\in[d],\quad s\in\{-1,1\},\quad a\in\{0,1\}.
\]
Its outcome coordinate is binary. Summing the full-data atom masses over
the unused potential outcome gives
\[
\begin{aligned}
G_\theta(X=j,S=s,A=a)
&=\frac1{2d}\,
\pi_{aj}\!\left(\frac{1+s(2a-1)}2\right)\\
&=\frac1{2d}\,\frac{1+s\theta_j}{2}
=\frac{P_\theta(j,s)}2.
\end{aligned}
\]
Indeed, for \(a=0\), the reconstructed outcome is \((1-s)/2\), so the
two choices of \(s\) select the failure and success masses of
\(\operatorname{Bernoulli}((1-\theta_j)/2)\). For \(a=1\), it is
\((1+s)/2\), selecting the success and failure masses of
\(\operatorname{Bernoulli}((1+\theta_j)/2)\). In both cases the selected
mass is \((1+s\theta_j)/2\). Thus the joint signed-input and assignment
law is
\[
\mathcal L_{G_\theta}(X,S,A)
=P_\theta\otimes\operatorname{Bernoulli}(1/2).
\]
% lean: CausalSmith.Stat.LdpOptvalueUniformFrontier.symmetricLaw_ancillary

For every binary observed record, \((2A-1)^2=1\), and therefore
\[
S(2A-1)
=(2A-1)(2Y-1)(2A-1)=2Y-1.
\]
This identity holds pointwise and hence \(G_\theta\)-almost surely.
% lean: CausalSmith.Stat.LdpOptvalueUniformFrontier.outcome_sign_reconstruction

\item \emph{Class preservation by averaging.}
Fix \(\varepsilon\in\mathbb R\) and an original-record protocol \(Q\).
Using the same seed law and message spaces, define
\[
\mathsf K_i(\cdot\mid(j,s),h,r)
:=\frac12\sum_{a\in\{0,1\}}
Q_i(\cdot\mid\operatorname{rec}(j,s,a),h,r).
\]
Each kernel is a finite mixture of measurable probability kernels, so
this defines a protocol.

Suppose \(Q\) satisfies \cref{obj:ass:full-record-privacy}. For every
\(i,h,r\), measurable message event \(E\), and paired inputs
\((j,s),(j',s')\), apply that inequality separately with assignment zero
and assignment one:
\[
Q_i(E\mid\operatorname{rec}(j,s,a),h,r)
\le e^\varepsilon
Q_i(E\mid\operatorname{rec}(j',s',a),h,r),
\qquad a\in\{0,1\}.
\]
Adding the two inequalities with weights \(1/2\) gives
\[
\mathsf K_i(E\mid(j,s),h,r)
\le e^\varepsilon\mathsf K_i(E\mid(j',s'),h,r).
\]
This proves sequential membership as specified by
\cref{obj:def:sequential-class}. If \(Q\) also satisfies
\cref{obj:ass:noninteraction}, then for all histories \(h,h'\),
\[
\begin{aligned}
\mathsf K_i(\cdot\mid(j,s),h,r)
&=\frac12\sum_a
Q_i(\cdot\mid\operatorname{rec}(j,s,a),h,r)\\
&=\frac12\sum_a
Q_i(\cdot\mid\operatorname{rec}(j,s,a),h',r)
=\mathsf K_i(\cdot\mid(j,s),h',r).
\end{aligned}
\]
Thus \cref{obj:def:noninteractive-class} is preserved as well. These
arguments apply at every declared history and seed, for every real
\(\varepsilon\).
% lean: CausalSmith.Stat.LdpOptvalueUniformFrontier.signed_protocol_classes

\item \emph{Reconstruction of the complete sampled input.}
For deterministic vectors, define
\[
\begin{gathered}
\mathbf v=((j_i,s_i))_{i=1}^n
\in([d]\times\{-1,1\})^n,
\qquad
\mathbf a=(a_i)_{i=1}^n\in\{0,1\}^n,\\
\operatorname{rec}^{(n)}(\mathbf v,\mathbf a)
:=\bigl(\operatorname{rec}(j_i,s_i,a_i)\bigr)_{i=1}^n.
\end{gathered}
\]
The single-participant ancillary identity and
\cref{obj:ass:iid-people} imply
\[
\mathcal L_{G_\theta}\!\left(
((X_i,S_i))_{i=1}^n,(A_i)_{i=1}^n\right)
=P_\theta^{\otimes n}
 \otimes\operatorname{Bernoulli}(1/2)^{\otimes n}.
\]
To see the factorization directly, the probability of a specified pair
of vectors is
\[
\prod_{i=1}^n\frac{P_\theta(j_i,s_i)}2
=\left(\prod_{i=1}^nP_\theta(j_i,s_i)\right)2^{-n}.
\]
Reconstruction recovers every observed record, so pushing this product
law through \(\operatorname{rec}^{(n)}\) gives
\[
\bigl(P_\theta^{\otimes n}
 \otimes\operatorname{Bernoulli}(1/2)^{\otimes n}\bigr)
 \circ(\operatorname{rec}^{(n)})^{-1}
=(G_\theta)_O^{\otimes n}.
\]

Write \(\lambda=\mathcal L(R)\) for the protocol's seed law. Including
the independent randomness from
\cref{obj:ass:independent-randomness}, define
\[
\Lambda_\theta
:=P_\theta^{\otimes n}
 \otimes\operatorname{Bernoulli}(1/2)^{\otimes n}
 \otimes\lambda\otimes\operatorname{Uniform}[0,1]
\]
on tuples \((\mathbf v,\mathbf a,r,u)\). Its pushforward satisfies
\[
\mathcal L_{G_\theta}(\mathbf O,R,U)
=\Lambda_\theta\circ
\bigl[(\mathbf v,\mathbf a,r,u)
\mapsto(\operatorname{rec}^{(n)}(\mathbf v,\mathbf a),r,u)\bigr]^{-1}.
\]
Thus running \(Q\) on the sampled causal records has the same joint
message and randomizer law as running it on these reconstructed
product-law inputs.
% lean: CausalSmith.Stat.LdpOptvalueUniformFrontier.symmetricLaw_reconstructed_decisionLaw

\item \emph{Averaging through the sequential recursion.}
For \(0\le\ell\le n\), define the prefix message space
\[
\mathcal Z_{\le\ell}:=\prod_{i=1}^{\ell}\mathcal Z_i,
\qquad
\mathcal Z_{\le0}:=\{\langle\rangle\}.
\]
For a fixed observed-input vector
\(\mathbf o\in([d]\times\{0,1\}^2)^n\), paired-input vector
\(\mathbf v\in([d]\times\{-1,1\})^n\), and seed \(r\in\mathcal R\),
let \(\Pi_\ell^Q(\cdot\mid\mathbf o,r)\) and
\(\Pi_\ell^{\mathsf K}(\cdot\mid\mathbf v,r)\) be the prefix probability
laws defined by
\[
\Pi_0^Q=\Pi_0^{\mathsf K}=\delta_{\langle\rangle}
\]
and, for \(0\le\ell<n\),
\[
\begin{aligned}
\Pi_{\ell+1}^Q(dh,dz\mid\mathbf o,r)
&=\Pi_\ell^Q(dh\mid\mathbf o,r)\,
 Q_{\ell+1}(dz\mid o_{\ell+1},h,r),\\
\Pi_{\ell+1}^{\mathsf K}(dh,dz\mid\mathbf v,r)
&=\Pi_\ell^{\mathsf K}(dh\mid\mathbf v,r)\,
 \mathsf K_{\ell+1}(dz\mid v_{\ell+1},h,r),
\end{aligned}
\]
where \(h\in\mathcal Z_{\le\ell}\), \(z\in\mathcal Z_{\ell+1}\), and
\((h,z)\) denotes the appended history.

We prove by induction on \(\ell\) that, for every measurable
\(\varphi:\mathcal Z_{\le\ell}\to[0,\infty]\),
\[
\begin{aligned}
&2^{-n}\sum_{\mathbf a\in\{0,1\}^n}
 \int\varphi(h)\,
 \Pi_\ell^Q(dh\mid
 \operatorname{rec}^{(n)}(\mathbf v,\mathbf a),r)\\
&\qquad=\int\varphi(h)\,
 \Pi_\ell^{\mathsf K}(dh\mid\mathbf v,r).
\end{aligned}
\]
For \(\ell=0\), both prefix laws are concentrated on the empty history,
and \(2^{-n}\sum_{\mathbf a\in\{0,1\}^n}1=1\).

For the induction step, take measurable
\(\varphi:\mathcal Z_{\le\ell+1}\to[0,\infty]\) and define
\[
\varphi_{\mathrm{av},\ell}(h)
:=\int\varphi(h,z)\,
 \mathsf K_{\ell+1}(dz\mid v_{\ell+1},h,r),
\qquad h\in\mathcal Z_{\le\ell}.
\]
This is a measurable nonnegative function. For \(i\in[n]\) and
\(a\in\{0,1\}\), define the updated bit vector by
\[
(\mathbf a^{\,i\leftarrow a})_{i'}
:=
\begin{cases}
a,&i'=i,\\
a_{i'},&i'\ne i,
\end{cases}
\qquad i'\in[n].
\]
The uniform distribution on bit vectors is invariant under flipping
coordinate \(i\). Pairing each vector with its flipped vector shows that
averaging any nonnegative function over all bit vectors equals
averaging its two versions with coordinate \(i\) set to zero and one.

Apply this identity at \(i=\ell+1\). Updating that coordinate leaves the
first \(\ell\) inputs unchanged, hence leaves the preceding prefix law
unchanged. The recursion and the definition of \(\mathsf K_{\ell+1}\)
therefore give, for every \(\mathbf a\),
\[
\begin{aligned}
&\frac12\sum_{a\in\{0,1\}}
 \int\varphi(h,z)\,
 \Pi_{\ell+1}^Q
 \bigl(dh,dz\mid
 \operatorname{rec}^{(n)}
   (\mathbf v,\mathbf a^{\,\ell+1\leftarrow a}),r\bigr)\\
&\qquad=
 \int\left[
 \frac12\sum_{a\in\{0,1\}}
 \int\varphi(h,z)\,
 Q_{\ell+1}(dz\mid
 \operatorname{rec}(j_{\ell+1},s_{\ell+1},a),h,r)
 \right]\\
&\hspace{6em}\cdot
 \Pi_\ell^Q(dh\mid
 \operatorname{rec}^{(n)}(\mathbf v,\mathbf a),r)\\
&\qquad=
 \int\varphi_{\mathrm{av},\ell}(h)\,
 \Pi_\ell^Q(dh\mid
 \operatorname{rec}^{(n)}(\mathbf v,\mathbf a),r).
\end{aligned}
\]
Averaging this equality over \(\mathbf a\), applying the induction
hypothesis to \(\varphi_{\mathrm{av},\ell}\), and using the paired
recursion proves
\[
\begin{aligned}
&2^{-n}\sum_{\mathbf a}
 \int\varphi\,
 \Pi_{\ell+1}^Q
 (\,\cdot\mid\operatorname{rec}^{(n)}(\mathbf v,\mathbf a),r)\\
&\qquad=
 \int\varphi_{\mathrm{av},\ell}(h)\,
 \Pi_\ell^{\mathsf K}(dh\mid\mathbf v,r)
 =\int\varphi\,\Pi_{\ell+1}^{\mathsf K}
 (\,\cdot\mid\mathbf v,r).
\end{aligned}
\]
At \(\ell=n\), indicator functions of measurable transcript events give
the equality of probability laws
\[
2^{-n}\sum_{\mathbf a\in\{0,1\}^n}
\Pi_n^Q
(\,\cdot\mid\operatorname{rec}^{(n)}(\mathbf v,\mathbf a),r)
=\Pi_n^{\mathsf K}(\,\cdot\mid\mathbf v,r).
\]
Integrating against \(P_\theta^{\otimes n}\), \(\lambda\), and the uniform
analyst randomizer, and using the identity expressing
\(\mathcal L_{G_\theta}(\mathbf O,R,U)\) as the reconstruction
pushforward of \(\Lambda_\theta\), proves
\[
\begin{aligned}
&\mathcal L_{G_\theta,Q}(\mathbf Z,R,U)\\
&\qquad=
P_\theta^{\otimes n}\otimes\lambda
 \otimes\operatorname{Uniform}[0,1]
\ \text{followed by }\ 
\mathbf Z\sim\Pi_n^{\mathsf K}(\cdot\mid\mathbf v,r).
\end{aligned}
\]
Equivalently, for every measurable
\(\mathcal E\subseteq\mathcal Z_{\le n}\times\mathcal R\times[0,1]\),
the common probability is
\[
\int\!\int\!\int\!\int
\mathbf1_{\mathcal E}(\mathbf z,r,u)\,
\Pi_n^{\mathsf K}(d\mathbf z\mid\mathbf v,r)\,
\operatorname{Uniform}[0,1](du)\,
\lambda(dr)\,P_\theta^{\otimes n}(d\mathbf v).
\]
Projecting onto \((R,\mathbf Z)\) gives the first asserted joint
seed-transcript equality. The induction begins at zero and therefore
also covers \(n=0\).
% lean: CausalSmith.Stat.LdpOptvalueUniformFrontier.averagedProtocol_reconstructed_decisionLaw

\item \emph{Class preservation by signed preprocessing.}
For the converse direction, fix a paired-alphabet protocol \(\mathsf K\).
Define the deterministic paired-record map by
\[
\operatorname{pair}(j,a,y)
:=\bigl(j,(2a-1)(2y-1)\bigr),
\qquad
(j,a,y)\in[d]\times\{0,1\}^2,
\]
and define \(Q\), with the same seed law and message spaces, by
\[
Q_i(\cdot\mid(j,a,y),h,r)
:=\mathsf K_i(\cdot\mid\operatorname{pair}(j,a,y),h,r).
\]
Composition with this measurable finite-alphabet map preserves the
probability-kernel property.

If \(\mathsf K\) satisfies the paired privacy inequality, then for every
two observed records \(o,o'\), every history and seed, and every
measurable message event \(E\),
\[
\begin{aligned}
Q_i(E\mid o,h,r)
&=\mathsf K_i(E\mid\operatorname{pair}(o),h,r)\\
&\le e^\varepsilon
 \mathsf K_i(E\mid\operatorname{pair}(o'),h,r)
=e^\varepsilon Q_i(E\mid o',h,r).
\end{aligned}
\]
Thus \(Q\) satisfies \cref{obj:ass:full-record-privacy} and belongs to
the sequential class whenever \(\mathsf K\) does. If \(\mathsf K\) is
independent of history, then
\[
Q_i(\cdot\mid o,h,r)
=\mathsf K_i(\cdot\mid\operatorname{pair}(o),h,r)
=\mathsf K_i(\cdot\mid\operatorname{pair}(o),h',r)
=Q_i(\cdot\mid o,h',r).
\]
This supplies \cref{obj:ass:noninteraction} and hence the noninteractive
membership in \cref{obj:def:noninteractive-class}.
% lean: CausalSmith.Stat.LdpOptvalueUniformFrontier.signed_protocol_classes

\item \emph{The converse transcript identity.}
For a deterministic observed vector \(\mathbf o\), define its
coordinatewise paired image by
\[
\operatorname{pair}^{(n)}(\mathbf o)
:=\bigl(\operatorname{pair}(o_i)\bigr)_{i=1}^n.
\]
The previously established identities
\(\tau(G_\theta)=\theta\) and
\(\mathcal L_{G_\theta}(X,S)=P_\theta\) imply, by coordinatewise
pushforward of the independent observed sample,
\[
\begin{aligned}
&\bigl((G_\theta)_O^{\otimes n}
 \otimes\lambda\otimes\operatorname{Uniform}[0,1]\bigr)
 \circ
 \bigl[(\mathbf o,r,u)
 \mapsto(\operatorname{pair}^{(n)}(\mathbf o),r,u)\bigr]^{-1}\\
&\qquad=
P_\theta^{\otimes n}\otimes\lambda
 \otimes\operatorname{Uniform}[0,1].
\end{aligned}
\]
Here \cref{obj:ass:iid-people,obj:ass:independent-randomness} identify the
measure on the left before pushforward with the sampling scheme's
joint input law.

For each fixed \(\mathbf o,r\), induction through the prefix recursion
gives
\[
\Pi_\ell^Q(\cdot\mid\mathbf o,r)
=\Pi_\ell^{\mathsf K}
(\cdot\mid\operatorname{pair}^{(n)}(\mathbf o),r),
\qquad 0\le\ell\le n.
\]
The zero-length laws are the same point mass. At a successor step,
the induction hypothesis identifies the preceding prefix laws, and
the defining equality
\[
Q_{\ell+1}(dz\mid o_{\ell+1},h,r)
=\mathsf K_{\ell+1}
(dz\mid\operatorname{pair}(o_{\ell+1}),h,r)
\]
identifies their extension kernels. Thus their successor prefix laws
are equal as well. In particular, their full fixed-input transcript
laws agree at \(\ell=n\).

Integrating this fixed-input equality and using the displayed input
pushforward proves that the joint law of \((\mathbf Z,R,U)\) under
\(G_\theta,Q\) equals the law obtained from independent \(P_\theta\)
inputs, the seed law \(\lambda\), and \(\mathsf K\). Projecting onto
\((R,\mathbf Z)\) gives the asserted converse equality. The same
zero-length base case covers \(n=0\).
% lean: CausalSmith.Stat.LdpOptvalueUniformFrontier.pulledProtocol_decisionLaw
\end{enumerate}
\end{proof}

\begin{proof}[Proof of \cref{obj:thm:uniform-private-value-frontiers}]
We first establish uniform finite-resource bounds, then prove the numerical resource conclusions and transfer them to the minimax quantities. Throughout the finite-resource argument, fix an allowed tuple and a sampling scheme satisfying \cref{obj:ass:iid-people,obj:ass:independent-randomness}, and write
\[
t=n\varepsilon^2,\qquad L=\log(ed),\qquad
z_*=\frac{d^2L}{t},\qquad w=\log(e+z_*).
\]
Thus \(t>0\), and
\[
1\le L=1+\log d\le d.
\]

\begin{enumerate}
\item
Invoke the published Bernstein limit
\[
2k_{\mathrm{half}} e_{2k_{\mathrm{half}}}\longrightarrow\beta_*,
\qquad k_{\mathrm{half}}\in\mathbb N_{>0},
\qquad 0.278<\beta_*<0.286,
\]
as recorded in \citep[Section 3.1, p.~8]{CaiLow2011Nonsmooth}. Here
\[
e_K=\inf_{\deg p\le K}\sup_{|x|\le1}\bigl|p(x)-|x|\bigr|.
\]
Choose an integer \(N_{\mathrm{tail}}\) such that
\[
k_{\mathrm{half}}\ge N_{\mathrm{tail}}
\quad\Longrightarrow\quad
(2k_{\mathrm{half}})e_{2k_{\mathrm{half}}}\ge\frac1{10},
\]
and define
\[
N_{\mathrm B}=\max\{N_{\mathrm{tail}},1\},
\qquad
b_*=\min\left\{\frac1{10},e_{2N_{\mathrm B}}\right\}.
\]
The tail inequality gives \(e_{2N_{\mathrm B}}>0\), hence \(b_*>0\).
The approximation errors decrease with degree, because the admissible polynomial spaces increase. Therefore, for every positive even integer \(K\),
\[
e_K\ge\frac{b_*}{K}.
\]
Indeed, write \(K=2k_{\mathrm{half}}\). If \(k_{\mathrm{half}}\ge N_{\mathrm B}\), then \(Ke_K\ge1/10\ge b_*\). If \(k_{\mathrm{half}}<N_{\mathrm B}\), then
\[
Ke_K\ge e_K\ge e_{2N_{\mathrm B}}\ge b_*.
\]

For such a \(K\) and an amplitude \(0<\alpha\le1/2\), take the measures supplied by \cref{obj:lem:cai-low-moment-duality} and define
\[
\nu_\ell=(x\mapsto\alpha x)_{\#}\xi_\ell,
\qquad \ell\in\{0,1\}.
\]
These are probability measures supported on \([-\alpha,\alpha]\). Scaling their moments gives
\[
\int u^v\,\nu_0(du)=\int u^v\,\nu_1(du)
\quad(0\le v\le K),
\qquad
\int |u|\,\nu_1(du)-\int |u|\,\nu_0(du)
=2\alpha e_K.
\]
Define the coordinate priors, induced causal priors, welfare centers, and separation by
\[
\Pi_\ell=\nu_\ell^{\otimes d},\qquad
\Lambda_\ell=(\theta\mapsto G_\theta)_{\#}\Pi_\ell,
\qquad
v_\ell=\frac12+\frac12\int|u|\,\nu_\ell(du),
\qquad
\Delta=v_1-v_0=\alpha e_K.
\]
By \cref{obj:prop:signed-causal-bridge}, the priors \(\Lambda_\ell\) are supported on \(\mathcal V_d\), and
\[
V(G_\theta)=\frac12+\frac1{2d}\sum_{j=1}^d|\theta_j|.
\]
Independence of the prior coordinates consequently gives
\[
\mathbb E_{\Lambda_\ell}V(P)=v_\ell,
\qquad
\operatorname{Var}_{\Lambda_\ell}(V(P))
=\frac1{4d}\operatorname{Var}_{\nu_\ell}(|u|)
\le\frac{\alpha^2}{4d}.
\]
Here the last inequality follows from \(|u|\le\alpha\) almost surely. Chebyshev's inequality and the approximation lower bound yield
\[
\Delta\ge\frac{b_*\alpha}{K}>0,
\qquad
\Lambda_\ell\!\left\{|V(P)-v_\ell|>\frac{\Delta}{8}\right\}
\le\frac{16\alpha^2}{d\Delta^2}
\le\frac{16K^2}{b_*^2d}.
\]
% lean: CausalSmith.Stat.LdpOptvalueUniformFrontier.causal_dual_product_priors_target_concentration_of_gate

\item
Use the converse tuning in \cref{obj:def:frontier-handle}:
\[
(\alpha,k)=
\begin{cases}
\left(\dfrac{\sqrt{z_*}}{100},\,2\lceil16L\rceil+1\right),
&z_*\le1,\\[3pt]
\left(\dfrac1{100},\,2\lceil16L/w\rceil+1\right),
&z_*>1,
\end{cases}
\qquad K=k-1.
\]
Both branches have \(0<\alpha\le1/2\) and a positive even \(K\).
The ceiling inequalities give
\[
K\le
\begin{cases}
34L,&z_*\le1,\\
34\max\{1,L/w\},&z_*>1.
\end{cases}
\]
Since \(w\ge1\) and \(L\ge1\), both bounds imply \(K\le34L\).

We also have the calibrated separation inequality
\[
\frac{\alpha}{K}\ge\frac1{3400}\sqrt{\rho(t,d)}.
\]
For \(z_*\le1\), the dense formula gives
\[
\sqrt{\rho(t,d)}=\frac{\sqrt{z_*}}{L},
\qquad
\frac{\alpha}{K}
\ge\frac{\sqrt{z_*}}{3400L}.
\]
For \(z_*>1\), put
\[
u_{\mathrm{rate}}=\min\{1,w/L\}=\sqrt{\rho(t,d)}.
\]
Then \(u_{\mathrm{rate}}\le1\) and
\(u_{\mathrm{rate}}L/w\le1\). Hence
\[
u_{\mathrm{rate}}K
\le u_{\mathrm{rate}}(32L/w+2)\le34,
\qquad
\frac{\alpha}{K}\ge\frac{u_{\mathrm{rate}}}{3400}.
\]

Because
\[
\frac{L^2}{d}\longrightarrow0\qquad(d\longrightarrow\infty),
\]
choose a fixed integer \(D_{\mathrm{con}}\ge2\) such that
\[
d>D_{\mathrm{con}}
\quad\Longrightarrow\quad
\frac{16(34L)^2}{b_*^2d}\le\frac1{100}.
\]
For these dimensions, the priors constructed in the preceding step satisfy
\[
\Delta\ge\frac{b_*}{3400}\sqrt{\rho(t,d)},
\qquad
\Lambda_\ell\!\left\{|V(P)-v_\ell|>\frac{\Delta}{8}\right\}
\le\frac1{100}
\quad(\ell=0,1).
\]
The cutoff depends on the universal approximation constant and is uniform over the allowed resource tuples.
% lean: CausalSmith.Stat.LdpOptvalueUniformFrontier.causal_concentrated_resource_priors_of_gate

\item
Fix any \(Q\in\mathfrak Q_{\mathrm{SI}}(n,d,\varepsilon)\). The measurable-density requirement in \cref{obj:lem:adaptive-moment-contraction} is supplied by \cref{obj:lem:measurable-kernel-radon-nikodym}. Define
\[
\beta=\frac{e^\varepsilon-1}{2d},
\qquad
\gamma_{\mathrm{con}}
=\frac{\alpha\sqrt{et}}{d\sqrt{k}},
\qquad
w_{\mathrm{con}}=
\begin{cases}
1,&z_*\le1,\\
w,&z_*>1.
\end{cases}
\]
The matched moments through \(K=k-1\) give, when \(k\le n\),
\[
\operatorname{TV}(\mathsf M_0^Q,\mathsf M_1^Q)
\le d\alpha^k\sqrt{\binom nk}\,\beta^k
\le d\gamma_{\mathrm{con}}^k.
\]
To justify the second inequality, the budget range gives
\[
e^\varepsilon-1\le2\varepsilon,
\qquad \beta\le\frac{\varepsilon}{d}.
\]
Moreover, the exponential series implies \(k^k/k!\le e^k\), so
\[
\binom nk\le\frac{n^k}{k!}\le\left(\frac{en}{k}\right)^k,
\qquad
\sqrt{\binom nk}\le\left(\sqrt{\frac{en}{k}}\right)^k.
\]
For \(k>n\), \cref{obj:lem:adaptive-moment-contraction} instead gives
\(\mathsf M_0^Q=\mathsf M_1^Q\).

It remains to bound the geometric expression uniformly. In both branches,
\[
kw_{\mathrm{con}}\ge32L,
\qquad
\gamma_{\mathrm{con}}^2=\frac{\alpha^2eL}{kz_*}.
\]
For \(z_*\le1\), the amplitude formula and \(e<3\) give
\[
\alpha^2e\,w_{\mathrm{con}}e^{w_{\mathrm{con}}/2}
=\frac{z_*}{10000}e\,e^{1/2}
\le32z_*.
\]
For \(z_*>1\), use
\[
e^w=e+z_*,
\qquad
w\le2e^{w/2},
\]
the latter following from \(1+w/2\le e^{w/2}\). Thus
\[
we^{w/2}\le2(e+z_*)\le8z_*,
\qquad
\alpha^2e\,we^{w/2}\le32z_*.
\]
Combining either budget inequality with \(kw_{\mathrm{con}}\ge32L\) yields
\[
\gamma_{\mathrm{con}}^2
\le e^{-w_{\mathrm{con}}/2},
\qquad
\gamma_{\mathrm{con}}\le e^{-w_{\mathrm{con}}/4}.
\]
Since \(d\le e^L\) and \(L\ge1\),
\[
d\gamma_{\mathrm{con}}^k
\le de^{-kw_{\mathrm{con}}/4}
\le de^{-8L}
\le e^{-7L}
\le e^{-7}
\le\frac1{100}.
\]
For the last inequality, \(e\ge2\) gives \(e^7\ge128\).
Consequently, for every sequential protocol,
\[
\operatorname{TV}(\mathsf M_0^Q,\mathsf M_1^Q)\le\frac1{100}.
\]
The sampling and randomness assumptions identify these mixtures with the joint seed-transcript mixtures under the fixed sampling scheme; the independent analyst randomizer supplies the decision mixtures required by \cref{obj:lem:separated-value-mixtures}.
% lean: CausalSmith.Stat.LdpOptvalueUniformFrontier.causal_lawspace_resource_contraction

\item
Suppose first that \(d>D_{\mathrm{con}}\). Apply \cref{obj:lem:separated-value-mixtures} with
\[
\eta_0=\omega=\frac1{100}.
\]
The support, concentration, and transcript-proximity hypotheses were established above. For every estimator and every globally honest connected interval,
\[
\sup_{P\in\mathcal V_d}
\mathbb E_{P,Q}(T-V(P))^2
\ge\frac{9\Delta^2}{128}\frac{97}{100},
\]
\[
\sup_{P\in\mathcal V_d}
\mathbb E_{P,Q}\operatorname{length}I
\ge\frac{3\Delta}{4}\frac{77}{100}.
\]
Define
\[
s_{\mathrm{con}}=\frac{b_*}{3400},
\qquad
c_{\mathrm{large}}
=\min\left\{\frac{s_{\mathrm{con}}^2}{100},
                 \frac{s_{\mathrm{con}}}{100}\right\}>0.
\]
The separation bound gives
\[
\Delta^2\ge s_{\mathrm{con}}^2\rho(t,d),
\qquad
\Delta\ge s_{\mathrm{con}}\sqrt{\rho(t,d)}.
\]
Since \(9\cdot97/12800\ge1/100\) and
\(3\cdot77/400\ge1/100\), the preceding decision bounds imply
\[
\sup_{P\in\mathcal V_d}\mathbb E_{P,Q}(T-V(P))^2
\ge c_{\mathrm{large}}\rho(t,d),
\qquad
\sup_{P\in\mathcal V_d}\mathbb E_{P,Q}\operatorname{length}I
\ge c_{\mathrm{large}}\sqrt{\rho(t,d)}.
\]
% lean: CausalSmith.Stat.LdpOptvalueUniformFrontier.frontier_lower_of_gate

\item
For the bounded dimensions and the later resource analysis, we need
\[
\min\{1,t^{-1}\}\le\rho(t,d)
\le4\min\{1,d^2/t\},
\qquad
\frac14\min\{t,d\}\le t\rho(t,d).
\]
Here is their branchwise derivation. In the dense branch,
\[
\rho(t,d)=\frac{d^2}{tL},
\qquad
\frac1t\le\rho(t,d)\le\min\{1,d^2/t\},
\qquad
t\rho(t,d)=\frac{d^2}{L}\ge d,
\]
using \(1\le L\le d\).

For the nondense branch, \(z_*>1\) and
\[
\rho(t,d)=\min\{1,(w/L)^2\}.
\]
The inequalities \(\log d\ge1-d^{-1}\) and \(e<3\) give
\[
dL=d(1+\log d)\ge2d-1\ge3>e.
\]
Consequently,
\[
\log(e+d^2L)\ge L,
\qquad
\log(e+dL)\ge\frac L2.
\]
For the second inequality, the logarithm on the left is at least both
\(1\) and \(\log d\).

The function
\[
\Phi_{\mathrm{res}}(t')
=t'\left[\log\left(e+\frac{d^2L}{t'}\right)\right]^2,
\qquad t'\in(0,\infty),
\]
is increasing. Indeed, for every \(v_{\mathrm{shift}}\ge0\),
\[
\log(e+v_{\mathrm{shift}})
\ge2-\frac e{e+v_{\mathrm{shift}}},
\]
by \(\log x\ge1-x^{-1}\) applied to
\(x=(e+v_{\mathrm{shift}})/e\). Therefore
\[
\log(e+v_{\mathrm{shift}})
-\frac{2v_{\mathrm{shift}}}{e+v_{\mathrm{shift}}}
\ge\frac e{e+v_{\mathrm{shift}}}\ge0,
\]
and differentiation gives
\[
\Phi_{\mathrm{res}}'(t')
=
\log\left(e+\frac{d^2L}{t'}\right)
\left[
\log\left(e+\frac{d^2L}{t'}\right)
-\frac{2d^2L/t'}{e+d^2L/t'}
\right]\ge0.
\]

If \(t\le1\), then \(w\ge\log(e+d^2L)\ge L\), so \(\rho(t,d)=1\).
If \(t\ge1\), monotonicity gives
\[
tw^2\ge[\log(e+d^2L)]^2\ge L^2,
\]
which proves \(\rho(t,d)\ge\min\{1,t^{-1}\}\).
If \(t\le d\), then \(w\ge\log(e+dL)\ge L/2\), so
\[
t(w/L)^2\ge t/4.
\]
If \(t\ge d\), monotonicity instead gives
\[
tw^2\ge d[\log(e+dL)]^2\ge dL^2/4.
\]
Since \(t\ge\min\{t,d\}/4\), taking the minimum with \(t\) proves the asserted lower bound for \(t\rho\).

Finally, for \(z_*\ge1\),
\[
\log(e+z_*)
\le\log(e+1)+\log z_*
\le2+2(\sqrt{z_*}-1)=2\sqrt{z_*}.
\]
Here \(\log(e+1)\le2\), and
\(\log\sqrt{z_*}\le\sqrt{z_*}-1\).
Hence
\[
(w/L)^2\le\frac{4z_*}{L^2}
=\frac{4d^2}{tL}\le\frac{4d^2}{t}.
\]
Together with \(\rho\le1\), this proves the upper comparison.
% lean: CausalSmith.Stat.LdpOptvalueUniformFrontier.rho_elementary_comparisons

\item
We next obtain a per-decision two-point bound for every dimension. Define
\[
\alpha_{\mathrm{pt}}
=\frac18\min\{1,t^{-1/2}\},
\qquad
\theta^{(0)}_j=0,\qquad
\theta^{(1)}_j=\alpha_{\mathrm{pt}}
\quad(j\in[d]).
\]
These vectors belong to \([-1/2,1/2]^d\), and
\[
0<\alpha_{\mathrm{pt}}\le\frac18,\qquad
\alpha_{\mathrm{pt}}\sqrt t\le\frac18.
\]
By \cref{obj:prop:signed-causal-bridge},
\[
G_{\theta^{(0)}},G_{\theta^{(1)}}\in\mathcal V_d,
\qquad
V(G_{\theta^{(0)}})=\frac12,\qquad
V(G_{\theta^{(1)}})=\frac12+\frac{\alpha_{\mathrm{pt}}}{2}.
\]
Apply \cref{obj:lem:adaptive-moment-contraction} with matching order \(1\) and coordinate measures \(\delta_0,\delta_{\alpha_{\mathrm{pt}}}\); their zeroth moments agree. For their seed-transcript laws,
\[
\operatorname{TV}
\le d\alpha_{\mathrm{pt}}\sqrt n\,
       \frac{e^\varepsilon-1}{2d}
\le\alpha_{\mathrm{pt}}\sqrt t
\le\frac18.
\]
The point priors have zero target-concentration error. Thus \cref{obj:lem:separated-value-mixtures}, with centers separated by
\(\alpha_{\mathrm{pt}}/2\), gives
\[
\sup_{P\in\mathcal V_d}\mathbb E_{P,Q}(T-V(P))^2
\ge\frac{63\alpha_{\mathrm{pt}}^2}{4096},
\qquad
\sup_{P\in\mathcal V_d}\mathbb E_{P,Q}\operatorname{length}I
\ge\frac{81\alpha_{\mathrm{pt}}}{320}.
\]
Since
\[
\alpha_{\mathrm{pt}}^2=\frac1{64}\min\{1,t^{-1}\},
\]
these imply
\[
\sup_{P\in\mathcal V_d}\mathbb E_{P,Q}(T-V(P))^2
\ge\frac1{8192}\min\{1,t^{-1}\},
\qquad
\sup_{P\in\mathcal V_d}\mathbb E_{P,Q}\operatorname{length}I
\ge\frac1{8192}\min\{1,t^{-1/2}\}.
\]

For \(d\le D_{\mathrm{con}}\), the elementary upper comparison gives
\[
\rho(t,d)\le4D_{\mathrm{con}}^2\min\{1,t^{-1}\}.
\]
To see the comparison inside the minimum, use
\(\min\{1,d^2/t\}\le1\) when \(t\le1\), and
\(\min\{1,d^2/t\}\le d^2/t\) when \(t>1\).
Taking square roots and enlarging the factor gives
\[
\sqrt{\rho(t,d)}
\le2D_{\mathrm{con}}\min\{1,t^{-1/2}\}
\le4D_{\mathrm{con}}^2\min\{1,t^{-1/2}\}.
\]
Define
\[
c_{\mathrm{small}}=\frac1{32768D_{\mathrm{con}}^2},
\qquad
c_{\mathrm{lo}}=\min\{c_{\mathrm{large}},c_{\mathrm{small}}\}>0.
\]
For \(d\le D_{\mathrm{con}}\), the two-point bounds give the frontier lower bounds with \(c_{\mathrm{small}}\); for \(d>D_{\mathrm{con}}\), the product-prior bounds give them with \(c_{\mathrm{large}}\). Thus, for every sequential protocol and every admissible decision,
\[
\sup_{P\in\mathcal V_d}\mathbb E_{P,Q}(T-V(P))^2
\ge c_{\mathrm{lo}}\rho(t,d),
\qquad
\sup_{P\in\mathcal V_d}\mathbb E_{P,Q}\operatorname{length}I
\ge c_{\mathrm{lo}}\sqrt{\rho(t,d)}.
\]
Every noninteractive protocol is sequential by
\cref{obj:def:noninteractive-class,obj:def:sequential-class}.
Taking the decision and protocol infima therefore gives, for either class,
\[
\mathfrak R_{\mathcal C}\ge c_{\mathrm{lo}}\rho(t,d),
\qquad
\mathfrak H_{\mathcal C}\ge c_{\mathrm{lo}}\sqrt{\rho(t,d)}.
\]
The risk lower bounds remain valid for extended nonnegative expectations.
% lean: CausalSmith.Stat.LdpOptvalueUniformFrontier.frontier_minimax_lower_of_gate

\item
For the upper bounds, use the construction \(\mathscr H\) of
\cref{obj:def:frontier-handle}, with its singleton public seed and transcript-based output \(\widehat V\). Its channels are fixed by the public resources. For \(n=2\) the messages are constant. For \(n>2\), use the baseline, pilot, and evaluation blocks specified there, with
\[
m=\lfloor n/3\rfloor,\qquad m_0=n-2m,\qquad
\delta=\tanh(\varepsilon/2)=\frac{e^\varepsilon-1}{e^\varepsilon+1},
\qquad b=\frac d\delta,\qquad \sigma^2=\frac{b^2}{m}.
\]
The block sizes satisfy
\[
m\ge n/5>0,\qquad m_0\ge n/3>0,\qquad 2m\le n.
\]
For the first inequality, write \(n=3m+r\), \(r\in\{0,1,2\}\), and use \(m\ge1\), giving \(n\le5m\). The second follows from \(m\le n/3\).

The binary channel probabilities are between \((1-\delta)/2\) and
\((1+\delta)/2\), and
\[
\frac{1+\delta}{1-\delta}=e^\varepsilon.
\]
Thus their ratio across any two original records is at most
\(e^\varepsilon\). The vector channel satisfies
\[
2^{-d}(1-\delta)
\le Q^{\mathrm{vec}}(z\mid O=(j,A,Y))
\le 2^{-d}(1+\delta)
\qquad\bigl(z\in\{-1,1\}^d\bigr).
\]
Its probability normalization and noninteractive privacy are also supplied by
\cref{obj:lem:private-moment-and-dependence}. Summing the pointwise privacy inequalities over message sets proves full-record privacy. The kernels are independent of message history, so the entire construction belongs to the noninteractive class.

The sampling and randomness assumptions imply the product transcript law
\[
\mathcal L_P(\mathbf Z)
=\bigotimes_{i=1}^n
\left[\int Q_i(\,\cdot\mid o)\,P_O(do)\right].
\]
In particular, the disjoint blocks are independent and each vector block has the law used in \cref{obj:lem:finite-private-polynomial-calibration}.

The privacy attenuation obeys
\[
\delta\ge\varepsilon/4>0.
\]
Indeed, \(e^\varepsilon\ge1+\varepsilon\) and \(e^\varepsilon<3\) imply
\(4(e^\varepsilon-1)\ge\varepsilon(e^\varepsilon+1)\).
Consequently,
\[
\sigma^2=\frac{d^2}{m\delta^2}\le\frac{80d^2}{t}.
\]
For the baseline messages,
\[
\mathbb E_P\mathsf H_i=\delta(2B(P)-1),
\qquad
\operatorname{Var}_P(\mathsf H_i)\le1.
\]
Independence therefore gives
\[
\mathbb E_P\widehat B=B(P),
\qquad
\mathbb E_P(\widehat B-B(P))^2
\le\frac1{4m_0\delta^2}\le\frac{12}{t}.
\]

The interior arm means imply
\[
V(P)=\frac1d\sum_{j=1}^d
       \max\{\mu_{0j}(P),\mu_{1j}(P)\}\in[1/4,3/4].
\]
For \(n>2\), the decomposition in \cref{obj:prop:signed-causal-bridge} and clipping to this interval give
\[
(\widehat V-V(P))^2
\le2(\widehat B-B(P))^2
   +\frac12(\widehat F-F(\tau(P)))^2.
\]
This follows by first contracting the distance through clipping and then applying
\((x+y/2)^2\le2x^2+y^2/2\).
Hence
\[
\mathbb E_P(\widehat V-V(P))^2
\le\frac{24}{t}
+\frac12\mathbb E_P(\widehat F-F(\tau(P)))^2.
\]
All branches also satisfy the pointwise bound
\[
(\widehat V-V(P))^2\le\frac14.
\]
% lean: CausalSmith.Stat.LdpOptvalueUniformFrontier.frontier_product_risk_le_components

\item
We prove the risk bound with the fixed constant
\[
c_{\mathrm{up}}=10^{16}
\]
in the following order.

If \(n=2\), then \(\widehat V=1/2\) and \(t\le2\). Therefore
\[
(\widehat V-V(P))^2\le\frac1{16},
\qquad
\rho(t,d)\ge\min\{1,t^{-1}\}\ge\frac12,
\]
which proves the bound. For \(n>2\), if \(\rho(t,d)=1\), clipping gives
\[
\mathbb E_P(\widehat V-V(P))^2\le\frac14
\le c_{\mathrm{up}}\rho(t,d).
\]

Next suppose the selected branch is the fallback \(\widehat F=1/4\).
This branch has \(t<d^2L\) and, with
\[
H=\log(e+\sigma^2L),
\]
it satisfies \(L/(1024H)<4\). The noise bound gives
\[
H\le\log\!\bigl(80(e+z_*)\bigr)
=\log80+w\le6w,
\]
because \(\log80\le5\) and \(w\ge1\). Thus
\[
L<4096H\le24576w,
\qquad
\rho(t,d)=\bigl(\min\{1,w/L\}\bigr)^2
\ge\frac1{24576^2}.
\]
Clipping again yields
\[
\mathbb E_P(\widehat V-V(P))^2
\le\frac14
\le c_{\mathrm{up}}\rho(t,d).
\]

Next, when \(L<4096\) and \(t<d^2L\), use \(w\ge1\) to obtain
\[
\rho(t,d)\ge\frac1{4096^2}.
\]
The same clipping argument proves the bound.

If \(L<4096\) in the remaining case, then \(t\ge d^2L\), and the construction uses the absolute evaluation means. The private-message moments in
\cref{obj:lem:private-moment-and-dependence} imply
\[
\mathbb E_P\overline W_j=\tau_j(P),
\qquad
\mathbb E_P(\overline W_j-\tau_j(P))^2\le\sigma^2.
\]
The absolute-value inequality and the squared-average inequality give
\[
\left[
\frac1d\sum_{j=1}^d|\overline W_j|
-\frac1d\sum_{j=1}^d|\tau_j(P)|
\right]^2
\le\frac1d\sum_{j=1}^d
       (\overline W_j-\tau_j(P))^2.
\]
Hence
\[
\mathbb E_P(\widehat F-F(\tau(P)))^2\le\sigma^2
\]
and
\[
\mathbb E_P(\widehat V-V(P))^2
\le\frac{24}{t}+\frac{\sigma^2}{2}
\le\frac{64d^2}{t}
=64L\rho(t,d)
\le64\cdot4096\,\rho(t,d)
\le c_{\mathrm{up}}\rho(t,d).
\]
% lean: CausalSmith.Stat.LdpOptvalueUniformFrontier.frontier_upper_of_gate

\item
In the remaining dense case, \(L\ge4096\) and \(t\ge d^2L\).
The construction uses
\[
D=2\left\lfloor\frac{L}{2048}\right\rfloor,
\qquad T_*=4\sigma\sqrt L,\qquad
a_{\mathrm{rad}}=8\sigma\sqrt L.
\]
This is an even degree, and its moments fit the evaluation block:
\[
2D\le\frac{L}{512}
\le\frac{d^2L}{5}
\le\frac n5\le m.
\]
Here \(t\le n\) follows from \(\varepsilon\le1\).

The concentration and approximation hypotheses of
\cref{obj:lem:finite-private-polynomial-calibration} are supplied by
\cref{obj:lem:bounded-mean-concentration,obj:lem:cai-low-chebyshev-approximation}.
Its hybrid conclusion, applied to the independent pilot and evaluation blocks, gives
\[
\mathbb E_P(\widehat F-F(\tau(P)))^2
\le500000000\,\frac{\sigma^2}{L}
\le40000000000\,\rho(t,d).
\]
Since \(L\le d^2\), the dense formula also gives \(t^{-1}\le\rho(t,d)\).
Therefore
\[
\mathbb E_P(\widehat V-V(P))^2
\le\left(24+\frac{40000000000}{2}\right)\rho(t,d)
\le c_{\mathrm{up}}\rho(t,d).
\]
% lean: CausalSmith.Stat.LdpOptvalueUniformFrontier.frontierEstimator_hybrid_dense_rate

\item
The final case has \(L\ge4096\), \(t<d^2L\), \(\rho(t,d)\ne1\), and a branch distinct from the fallback. With \(H=\log(e+\sigma^2L)\), this means
\[
\frac{L}{1024H}\ge4,
\qquad
D=2\left\lfloor\frac{L}{2048H}\right\rfloor,
\qquad a_{\mathrm{rad}}=\frac12.
\]
Since \(H\ge1\), the floor inequalities give
\[
D\ge2,\qquad D\ \text{even},\qquad
\frac{L}{2048H}\le D\le\frac{L}{1024H}.
\]

We verify \(2D\le m\). Since \(0<\delta\le1\),
\[
\sigma^2\ge\frac{d^2}{m}.
\]
If \(m<2D\), the degree budget would imply
\[
m<2D\le\frac{L}{512H}\le\frac{L}{512},
\qquad
\sigma^2L>512d^2.
\]
Thus \(e+\sigma^2L\ge ed\) and \(H\ge L\). But then
\(D\le L/(1024H)\le1/1024\), contradicting \(D\ge2\).
This proves the moment-fit condition.

The global-polynomial conclusion of
\cref{obj:lem:finite-private-polynomial-calibration} now gives
\[
\mathbb E_P(\widehat F-F(\tau(P)))^2
\le\frac{e^{9DH}}{2d}+\frac1{(D+1)^2}.
\]
For completeness, the exponential dimension estimate used to calibrate this bound is
\[
L^2e^{L/100}\le4d.
\]
Indeed, \(L/4\le e^{L/4}\) gives \(L^2\le16e^{L/2}\), while
\(L\ge4096\) gives \(4\le e^{49L/100-1}\). Therefore
\[
L^2e^{L/100}
\le16e^{51L/100}
\le4e^{L-1}=4d.
\]
Since \(DH\le L/1024\), it follows that
\[
\frac{e^{9DH}}{2d}
\le\frac{e^{L/100}}{2d}\le\frac2{L^2}.
\]
As above, \(H\le6w\). The lower rounded-degree bound gives
\[
\frac1{D+1}\le\frac{2048H}{L}
\le12288\,\frac wL.
\]
The nonsaturation condition and the nondense formula imply
\[
\frac wL<1,\qquad \rho(t,d)=(w/L)^2.
\]
Using \(w\ge1\), we obtain
\[
\frac{e^{9DH}}{2d}\le2\rho(t,d),
\qquad
\frac1{(D+1)^2}\le150994944\,\rho(t,d).
\]
Also \(t^{-1}\le\rho(t,d)\): otherwise
\(\min\{1,t^{-1}\}=1\), and the elementary lower bound together with
\(\rho\le1\) would force saturation. Consequently,
\[
\mathbb E_P(\widehat V-V(P))^2
\le\left[24+\frac12(2+150994944)\right]\rho(t,d)
\le c_{\mathrm{up}}\rho(t,d).
\]
All cases have therefore established
\[
\sup_{P\in\mathcal V_d}
\mathbb E_P(\widehat V-V(P))^2
\le c_{\mathrm{up}}\rho(t,d).
\]
% lean: CausalSmith.Stat.LdpOptvalueUniformFrontier.frontier_global_resource_bound

\item
Use the interval in \cref{obj:def:frontier-handle}:
\[
q_*=\sqrt{10^{17}\rho(t,d)},
\qquad
I=
\left[
\max\{1/4,\widehat V-q_*\},
\min\{3/4,\widehat V+q_*\}
\right].
\]
The rate is positive by the elementary lower comparison, so \(q_*>0\).
Because \(\widehat V\in[1/4,3/4]\), these endpoints are ordered, and the interval is connected and closed. All decisions are Borel functions of the finite transcript.

For \(V(P)\in[1/4,3/4]\), failure of coverage implies
\(|\widehat V-V(P)|>q_*\). Markov's inequality therefore gives
\[
\Pr_P\{V(P)\notin I\}
\le\Pr_P\{(\widehat V-V(P))^2\ge q_*^2\}
\le\frac{10^{16}\rho(t,d)}{10^{17}\rho(t,d)}
=\frac1{10}.
\]
Clipping the endpoints gives
\[
\operatorname{length}I\le2q_*
=\sqrt{40c_{\mathrm{up}}}\,\sqrt{\rho(t,d)}.
\]
Thus the same noninteractive protocol, transcript-only estimator, and interval simultaneously satisfy
\[
\mathbb E_P(\widehat V-V(P))^2\le c_{\mathrm{up}}\rho(t,d),
\qquad
\Pr_P\{V(P)\in I\}\ge0.90,
\qquad
\mathbb E_P\operatorname{length}I
\le\sqrt{40c_{\mathrm{up}}}\,\sqrt{\rho(t,d)}.
\]
These hold for every \(P\in\mathcal V_d\), including the \(n=2\) branch.
% lean: CausalSmith.Stat.LdpOptvalueUniformFrontier.frontier_upper_of_gate

\item
Choose the common constants
\[
c=\min\{c_{\mathrm{lo}},1\},
\qquad
C_0=\max\{1,c_{\mathrm{up}},\sqrt{40c_{\mathrm{up}}}\}.
\]
Then
\[
0<c\le C_0<\infty,\qquad
c\le c_{\mathrm{lo}},\qquad
c_{\mathrm{up}}\le C_0,\qquad
\sqrt{40c_{\mathrm{up}}}\le C_0.
\]
The displayed attainment bounds consequently hold with \(C_0\).
The estimate is a Borel function of \(\mathbf Z\) alone, since the public seed is a singleton and the construction ignores analyst randomization.

The attaining protocol is admissible for both classes. Using it and its estimate as candidates in \cref{obj:def:risk}, and its globally honest interval as a candidate in \cref{obj:def:honest-length}, gives
\[
\mathfrak R_{\mathcal C}
\le\sup_{P\in\mathcal V_d}\mathbb E_P(\widehat V-V(P))^2
\le C_0\rho(t,d),
\]
\[
\mathfrak H_{\mathcal C}
\le\sup_{P\in\mathcal V_d}\mathbb E_P\operatorname{length}I
\le C_0\sqrt{\rho(t,d)}.
\]
Together with the lower bounds, this proves
\[
c\rho(t,d)\le\mathfrak R_{\mathcal C}\le C_0\rho(t,d),
\qquad
c\sqrt{\rho(t,d)}
\le\mathfrak H_{\mathcal C}\le C_0\sqrt{\rho(t,d)}.
\]
The attaining candidates are precisely the resource-selected construction, so its three guarantees hold with this same constant.
% lean: CausalSmith.Stat.LdpOptvalueUniformFrontier.causalAttainment_minimax_upper

\item
Protocol-class inclusion gives
\[
\mathfrak R_{\mathrm{SI}}\le\mathfrak R_{\mathrm{NI}},
\qquad
\mathfrak H_{\mathrm{SI}}\le\mathfrak H_{\mathrm{NI}}.
\]
Both denominators are positive by the lower bounds, and all four quantities are finite by the upper bounds. Hence
\[
1\le\frac{\mathfrak R_{\mathrm{NI}}}{\mathfrak R_{\mathrm{SI}}}
\le\frac{C_0\rho(t,d)}{c\rho(t,d)}
=\frac{C_0}{c},
\qquad
1\le\frac{\mathfrak H_{\mathrm{NI}}}{\mathfrak H_{\mathrm{SI}}}
\le\frac{C_0\sqrt{\rho(t,d)}}{c\sqrt{\rho(t,d)}}
=\frac{C_0}{c}.
\]
% lean: CausalSmith.Stat.LdpOptvalueUniformFrontier.frontier_ratio_comparison

\item
Squaring the nonnegative length bounds gives
\[
c^2\rho(t,d)\le\mathfrak H_{\mathcal C}^2
\le C_0^2\rho(t,d).
\]
The risk bounds then yield the two comparison chains
\[
\frac{c^2}{C_0}\mathfrak R_{\mathcal C}
\le\frac{c^2}{C_0}\,C_0\rho(t,d)
=c^2\rho(t,d)
\le\mathfrak H_{\mathcal C}^2,
\]
\[
\mathfrak H_{\mathcal C}^2
\le C_0^2\rho(t,d)
=\frac{C_0^2}{c}\,c\rho(t,d)
\le\frac{C_0^2}{c}\mathfrak R_{\mathcal C}.
\]
% lean: CausalSmith.Stat.LdpOptvalueUniformFrontier.frontier_square_comparison

\item
We now establish the resource conclusions. For any \(d\ge2\) and
\(0<t<d^2L\), the nondense expression has a positive logarithmic ratio, and therefore
\[
\begin{aligned}
\rho(t,d)=1
&\iff \log(e+d^2L/t)\ge L\\
&\iff e+d^2L/t\ge ed\\
&\iff t\le\frac{d^2L}{ed-e}.
\end{aligned}
\]
The final equivalence uses \(t>0\) and \(ed-e>0\).
% lean: CausalSmith.Stat.LdpOptvalueUniformFrontier.rho_saturation_iff

\item
Fix any sequence of allowed tuples and define
\[
t_k=n_k\varepsilon_k^2,\qquad
L_k=\log(ed_k),\qquad
\rho_k=\rho(t_k,d_k),\qquad
g_k=\frac{\log(1+d_k^2/t_k)}{L_k}.
\]
All \(t_k\) are positive, and \(\rho_k,g_k\ge0\).

Two numerical inequalities will give the consistency equivalence. For arbitrary \(d\ge2\), \(t>0\), define
\[
g(t,d)=\frac{\log(1+d^2/t)}{L},
\qquad
e_{\mathrm{corr}}(d)=\frac{\log(e+L)}{L}.
\]
For \(x_{\mathrm{res}}=d^2/t\ge0\), the inequalities
\[
1+x_{\mathrm{res}}\le e+Lx_{\mathrm{res}}
\le(1+x_{\mathrm{res}})(e+L)
\]
imply
\[
\log(1+x_{\mathrm{res}})
\le\log(e+Lx_{\mathrm{res}})
\le\log(1+x_{\mathrm{res}})+\log(e+L).
\]
It follows that
\[
\rho(t,d)<1\quad\Longrightarrow\quad
g(t,d)^2\le\rho(t,d).
\]
In the nondense branch, this follows by dividing the first logarithmic inequality by \(L\), because nonsaturation selects the untruncated ratio. In the dense branch,
\[
x_{\mathrm{res}}L\le1,\qquad
0\le g(t,d)\le\frac{x_{\mathrm{res}}}{L}\le1,
\]
using \(\log(1+x_{\mathrm{res}})\le x_{\mathrm{res}}\); hence
\[
g(t,d)^2\le\frac{x_{\mathrm{res}}}{L}=\rho(t,d).
\]
For both branches we also have
\[
\rho(t,d)
\le\bigl[g(t,d)+e_{\mathrm{corr}}(d)\bigr]^2+\frac1{L^2}.
\]
In the nondense branch use the upper logarithmic inequality and the nonnegative truncated ratio. In the dense branch use
\(\rho(t,d)\le L^{-2}\).
Finally,
\[
e_{\mathrm{corr}}(d)\longrightarrow0,
\qquad
L^{-1}\longrightarrow0
\qquad(d\longrightarrow\infty),
\]
by \(L\to\infty\) and \(\log(e+L)/L\to0\).

Suppose \(\rho_k\to0\). The elementary lower comparison forces \(t_k\to\infty\): for any real threshold \(M_{\mathrm{test}}\), if
\(t_k<M_{\mathrm{test}}\), then
\[
\rho_k\ge
\min\{1,t_k^{-1}\}
\ge\frac1{\max\{M_{\mathrm{test}},1\}}>0.
\]
Such indices are eventually excluded by convergence to zero.
Eventually \(\rho_k<1\), and the first numerical inequality gives
\[
0\le g_k\le\sqrt{\rho_k}\longrightarrow0.
\]

Conversely, suppose \(t_k\to\infty\) and \(g_k\to0\). Fix a tolerance
\(\zeta>0\), and define
\[
q_{\mathrm{cons}}=\frac{\sqrt{\zeta}}4.
\]
Choose an integer \(D_{\mathrm{cons}}\) such that, for every
\(d\ge D_{\mathrm{cons}}\),
\[
e_{\mathrm{corr}}(d)<q_{\mathrm{cons}},
\qquad
\frac1{\log(ed)}<q_{\mathrm{cons}}.
\]
Eventually,
\[
g_k<q_{\mathrm{cons}},
\qquad
t_k\ge\frac{4D_{\mathrm{cons}}^2}{\zeta}+1.
\]
If \(d_k\ge D_{\mathrm{cons}}\), the numerical upper bound gives
\[
\rho_k<(2q_{\mathrm{cons}})^2+q_{\mathrm{cons}}^2
=\frac{5\zeta}{16}<\zeta.
\]
If \(d_k<D_{\mathrm{cons}}\), the elementary comparison gives
\[
\rho_k\le\frac{4d_k^2}{t_k}
\le\frac{4D_{\mathrm{cons}}^2}{t_k}<\zeta.
\]
Thus \(\rho_k\to0\), establishing the criterion for arbitrary allowed sequences.
% lean: CausalSmith.Stat.LdpOptvalueUniformFrontier.rho_consistency_iff

\item
We next prove the private-parametric criterion. Write eventual comparability as the existence of constants
\[
0<c_{\mathrm{cmp}}\le C_{\mathrm{cmp}}<\infty,
\qquad
\frac{c_{\mathrm{cmp}}}{t_k}\le\rho_k
\le\frac{C_{\mathrm{cmp}}}{t_k}
\quad\text{eventually}.
\]
If these inequalities hold, then \(\rho_k\le1\) implies
\(t_k\ge c_{\mathrm{cmp}}\) eventually. The elementary scaled lower bound also gives
\[
\frac14\min\{t_k,d_k\}\le t_k\rho_k\le C_{\mathrm{cmp}},
\]
so \(\min\{t_k,d_k\}\le4C_{\mathrm{cmp}}\) eventually.

Conversely, suppose that for some \(a>0\) and \(M\in\mathbb R\),
\[
t_k\ge a,\qquad \min\{t_k,d_k\}\le M
\quad\text{eventually}.
\]
Define
\[
M_+=\max\{1,M\},\qquad
c_{\mathrm{cmp}}=\min\{1,a\},\qquad
C_{\mathrm{cmp}}=4M_+^2.
\]
The elementary lower comparison gives
\[
\frac{c_{\mathrm{cmp}}}{t_k}
\le\min\{1,t_k^{-1}\}\le\rho_k.
\]
For the upper bound, if \(t_k\le d_k\), then \(t_k\le M_+\), and
\[
\rho_k\le1\le\frac{4M_+^2}{t_k}.
\]
If \(d_k\le t_k\), then \(d_k\le M_+\), and
\[
\rho_k\le\frac{4d_k^2}{t_k}
\le\frac{4M_+^2}{t_k}.
\]
This proves the equivalence.

When \(t_k\to\infty\), an eventual positive lower bound on \(t_k\) is automatic. If \(\min\{t_k,d_k\}\le M\) eventually, then eventually \(t_k>M\), forcing \(d_k\le M\). Enlarging \(M\) to an integer yields the stated eventual dimension bound. Conversely, an eventual integer bound on \(d_k\) bounds the minimum.
% lean: CausalSmith.Stat.LdpOptvalueUniformFrontier.rho_parametric_comparable_iff

\item
Fix resource-ratio endpoints
\[
0<a\le b_{\mathrm{ratio}}<\infty.
\]
Choose \(D_0\) sufficiently large that, for every \(d\ge D_0\), with \(L=\log(ed)\),
\[
L\ge b_{\mathrm{ratio}}+1,\qquad
2\log b_{\mathrm{ratio}}\le\log L,\qquad
\log(e+1/a)\le\log L.
\]
These thresholds exist because \(L,\log L\to\infty\).
For \(d\ge\max\{D_0,2\}\) and \(a\le t/d^2\le b_{\mathrm{ratio}}\), define
\[
u_{\mathrm{elbow}}=\frac{t}{d^2}\in[a,b_{\mathrm{ratio}}].
\]
Then \(t<d^2L\), and
\[
\log(e+L/u_{\mathrm{elbow}})
\ge\log(L/b_{\mathrm{ratio}})
=\log L-\log b_{\mathrm{ratio}}\ge\frac12\log L.
\]
Also,
\[
e+\frac{L}{u_{\mathrm{elbow}}}
\le L(e+1/a),
\]
because \(L\ge1\). Hence
\[
\log(e+L/u_{\mathrm{elbow}})
\le\log L+\log(e+1/a)\le2\log L.
\]
Since \(0\le\log L/L\le1\), truncation at one preserves the bounds
\[
\frac12\frac{\log L}{L}
\le\min\left\{1,\frac{\log(e+L/u_{\mathrm{elbow}})}L\right\}
\le2\frac{\log L}{L}.
\]
Squaring gives the uniform comparison
\[
\frac14\left(\frac{\log L}{L}\right)^2
\le\rho(t,d)\le
4\left(\frac{\log L}{L}\right)^2.
\]
This proves the final numerical resource assertion, with constants \(1/4\) and \(4\).
% lean: CausalSmith.Stat.LdpOptvalueUniformFrontier.rate_resource_conclusions

\item
Finally, fix a family of sampling schemes satisfying the two stated assumptions at each \(n,d\), and fix a protocol class \(\mathcal C\).
Along an allowed sequence define
\[
\mathfrak R_k=\mathfrak R_{\mathcal C}(n_k,d_k,\varepsilon_k),
\qquad
\mathfrak H_k=\mathfrak H_{\mathcal C}(n_k,d_k,\varepsilon_k).
\]
The finite upper bounds make both quantities finite, and their real-valued bounds are
\[
c\rho_k\le\mathfrak R_k\le C_0\rho_k,
\qquad
c\sqrt{\rho_k}\le\mathfrak H_k\le C_0\sqrt{\rho_k}.
\]
The first sandwich gives
\[
\mathfrak R_k\to0\iff\rho_k\to0.
\]
The second, together with continuity of the square root and
\((\sqrt{\rho_k})^2=\rho_k\), gives
\[
\mathfrak H_k\to0\iff\rho_k\to0.
\]
Combining these with the numerical consistency criterion proves both consistency statements.

The same sandwiches transfer eventual comparability:
\[
\mathfrak R_k\asymp t_k^{-1}
\iff\rho_k\asymp t_k^{-1},
\qquad
\mathfrak H_k\asymp t_k^{-1/2}
\iff\rho_k\asymp t_k^{-1}.
\]
For risk, one combines either comparison with
\(\rho_k\le\mathfrak R_k/c\) and
\(\mathfrak R_k\le C_0\rho_k\).
For length, one first uses the comparison with \(\sqrt{\rho_k}\) and then squares, or takes square roots, of the nonnegative inequalities. The resource and bounded-dimension criteria therefore transfer exactly as stated.

To retain common constants for the final risk and length comparisons, define
\[
\chi_d=\frac{\log\log(ed)}{\log(ed)},\qquad
c_{\mathrm{seq}}=\min\{c/4,c^2/4\},\qquad
C_{\mathrm{seq}}=\max\{4C_0,4C_0^2\}.
\]
The numerical comparison in the preceding step and the finite bounds imply
\[
c_{\mathrm{seq}}\chi_d^2
\le\mathfrak R_{\mathcal C}(n,d,\varepsilon)
\le C_{\mathrm{seq}}\chi_d^2,
\]
and
\[
\sqrt{c_{\mathrm{seq}}}\,|\chi_d|
\le\mathfrak H_{\mathcal C}(n,d,\varepsilon)
\le\sqrt{C_{\mathrm{seq}}}\,|\chi_d|.
\]
Indeed, taking square roots of the numerical rate bounds gives
\[
\frac c2|\chi_d|
\le c\sqrt{\rho(t,d)}
\le\mathfrak H_{\mathcal C},
\qquad
\mathfrak H_{\mathcal C}
\le C_0\sqrt{\rho(t,d)}
\le2C_0|\chi_d|,
\]
and the definitions of \(c_{\mathrm{seq}},C_{\mathrm{seq}}\) encompass these factors as well as the risk factors.
For \(d\ge2\), \(\chi_d>0\).
Along a sequence with \(d_k\to\infty\) and
\(t_k/d_k^2\in[a,b_{\mathrm{ratio}}]\), the dimension threshold is eventually met.
The displayed bounds prove the two final asymptotic orders.
% lean: CausalSmith.Stat.LdpOptvalueUniformFrontier.value_sequence_conclusions_of_bounds
\end{enumerate}
\end{proof}

\begin{proof}[Proof of \cref{obj:thm:two-cell-calibration}]
\begin{enumerate}
\item
By \cref{obj:lem:measurable-kernel-radon-nikodym}, the measurable-density hypothesis of \cref{obj:lem:dimension-free-private-two-point} is satisfied. The latter supplies a universal constant \(c_1>0\) such that, for every allowed two-cell resource tuple and sampling scheme, and every \(\mathcal C\in\{\mathrm{NI},\mathrm{SI}\}\),
\[
\mathfrak R_{\mathcal C}(n,2,\varepsilon)
\ge c_1\min\left\{1,\frac1{n\varepsilon^2}\right\},
\qquad
\mathfrak H_{\mathcal C}(n,2,\varepsilon)
\ge c_1\min\left\{1,\frac1{\sqrt n\,\varepsilon}\right\}.
\]
Set
\[
c=\min\{c_1,1\},
\qquad C_0=1000.
\]
Then \(0<c\le1\le C_0<\infty\).

Fix \(n\ge2\), \(0<\varepsilon\le1\), and a sampling scheme satisfying
\cref{obj:ass:iid-people,obj:ass:independent-randomness}. Define the effective private resource by
\[
t=n\varepsilon^2>0.
\]
All expectations and probabilities below use the decision law induced by this sampling scheme and the specified protocol. We first construct simultaneous risk and interval witnesses, splitting according to \(t\le1\) or \(t>1\).
% lean: CausalSmith.Stat.LdpOptvalueUniformFrontier.two_cell_calibration

\item
Suppose \(t\le1\). Since
\[
(\sqrt n\,\varepsilon)^2=t,
\qquad 0<\sqrt n\,\varepsilon\le1,
\]
the two rate factors satisfy
\[
\min\left\{1,\frac1{n\varepsilon^2}\right\}
=
\min\left\{1,\frac1{\sqrt n\,\varepsilon}\right\}
=1.
\]
Use a singleton public seed and let every participant send the same fixed message. Set
\[
T(\mathbf Z,R,U)=\frac12,
\qquad
I(\mathbf Z,R,U)=\left[\frac14,\frac34\right].
\]
The message law is constant in the original record and history, so it satisfies full-record privacy and noninteraction. The estimate and interval endpoints are constant Borel functions.

For every \(P\in\mathcal V_2\), the interior means in
\cref{obj:def:causal-model} give
\[
\frac14\le
\max\{\mu_{0j}(P),\mu_{1j}(P)\}
\le\frac34
\qquad(j\in[2]).
\]
Averaging these inequalities yields
\[
\frac14\le
V(P)=\frac12\sum_{j=1}^{2}
\max\{\mu_{0j}(P),\mu_{1j}(P)\}
\le\frac34.
\]
Consequently,
\[
\mathbb E[(T-V(P))^2]
=\left(\frac12-V(P)\right)^2
\le\frac1{16},
\qquad
\mathbb P\{V(P)\in I\}=1,
\qquad
\mathbb E[\operatorname{length}I]=\frac12.
\]
These bounds attain both rate factors with comparison constant \(1\). Enlarging that constant to \(1000\) preserves the same witnesses and coverage.
% lean: CausalSmith.Stat.LdpOptvalueUniformFrontier.two_cell_saturated_attainment
% lean: CausalSmith.Stat.LdpOptvalueUniformFrontier.causalAttainment_mono_constant

\item
Suppose \(t>1\). We construct the three transcript statistics used in this branch. Define the binary and vector block sizes, privacy attenuation, and scaled-message magnitude by
\[
m_{\mathrm{bin}}=\left\lfloor\frac n2\right\rfloor,
\qquad
m_{\mathrm{vec}}=n-m_{\mathrm{bin}},
\qquad
\delta=\tanh(\varepsilon/2)
=\frac{e^\varepsilon-1}{e^\varepsilon+1},
\qquad
b=\frac2\delta.
\]
Here \(m_{\mathrm{bin}}\ge1\), \(m_{\mathrm{vec}}\ge1\), and \(0<\delta<1\).
For later calibration, the elementary inequalities
\[
e^\varepsilon-1\ge\varepsilon,
\qquad e^\varepsilon\le e<3
\]
also give
\[
\delta
=\frac{e^\varepsilon-1}{e^\varepsilon+1}
\ge\frac{\varepsilon}{4}.
\]

Use a singleton public seed. Each participant's message belongs to the finite alphabet
\(\{-1,1\}\times\{-1,1\}^{2}\). For an original record
\(O=(j,A,Y)\), with \(j\in[2]\) and \(A,Y\in\{0,1\}\), define, for
\(h\in\{-1,1\}\), \(z\in\{-1,1\}^{2}\), and \(1\le i\le n\),
\[
Q_i\bigl(\{(h,z)\}\mid O=(j,A,Y)\bigr)
=
\begin{cases}
\displaystyle
\frac{1+h\delta(2Y-1)}2\,
\mathbf1_{\{z=(-1,-1)\}},
& i\le m_{\mathrm{bin}},\\[6pt]
\displaystyle
\frac14\bigl(1+\delta(2A-1)(2Y-1)z_j\bigr)\,
\mathbf1_{\{h=-1\}},
& i>m_{\mathrm{bin}}.
\end{cases}
\]
These kernels are defined on every original record and are independent of message history. Their masses are nonnegative and sum to one: the first branch sums over the two binary signs, while in the second branch
\[
\sum_{z\in\{-1,1\}^{2}}z_j=0.
\]

For the binary branch, its two possible nonzero atom masses lie between
\((1-\delta)/2\) and \((1+\delta)/2\). The identity
\[
1+\delta=e^\varepsilon(1-\delta)
\]
therefore gives the atomwise privacy inequality for every pair of original records. The vector branch is precisely the \(d=2\) channel in
\cref{obj:lem:private-moment-and-dependence}, which supplies its full-record privacy. The fixed extra message coordinates preserve these inequalities. Summing over atoms gives
\[
Q_i(E\mid o)\le e^\varepsilon Q_i(E\mid o')
\]
for every message event \(E\) and every pair of original records \(o,o'\). Thus
\[
Q\in\mathfrak Q_{\mathrm{NI}}(n,2,\varepsilon).
\]

Write the transcript and its scaled vector coordinates as
\[
\mathbf Z
=\bigl((\mathsf H_i^{\mathrm{cal}},Z_i^{\mathrm{cal}})\bigr)_{i=1}^{n},
\qquad
W_{ij}^{\mathrm{cal}}
=bZ_{m_{\mathrm{bin}}+i,j}^{\mathrm{cal}}
\quad
(1\le i\le m_{\mathrm{vec}},\ j\in[2]).
\]
Define the baseline and contrast estimates by
\[
\widehat B_{\mathrm{cal}}
=\frac12+
\frac1{2\delta m_{\mathrm{bin}}}
\sum_{i=1}^{m_{\mathrm{bin}}}\mathsf H_i^{\mathrm{cal}},
\qquad
\widehat\tau_j
=\frac1{m_{\mathrm{vec}}}
\sum_{i=1}^{m_{\mathrm{vec}}}W_{ij}^{\mathrm{cal}}
\quad(j\in[2]).
\]

Conditional on the records, the transcript atom probability is the product of its participant-specific atom probabilities. Integrating against the product record law from
\cref{obj:ass:iid-people}, and using
\cref{obj:ass:independent-randomness}, gives the transcript identity
\[
\mathbb P\{\mathbf Z=((h_i,z_i))_{i=1}^{n}\}
=
\prod_{i=1}^{n}
\int Q_i(\{(h_i,z_i)\}\mid o)\,P_O(do).
\]
Thus distinct message rows are independent under the sampling decision law.

For \(i\le m_{\mathrm{bin}}\), direct summation over the binary message gives
\[
\mathbb E[\mathsf H_i^{\mathrm{cal}}\mid O_i]
=\delta(2Y_i-1),
\qquad
\mathbb E[\mathsf H_i^{\mathrm{cal}}]
=\delta(2B(P)-1),
\qquad
\mathbb E[(\mathsf H_i^{\mathrm{cal}})^2]=1.
\]
It follows that
\[
\mathbb E[\widehat B_{\mathrm{cal}}]=B(P),
\qquad
\operatorname{Var}\left(
\sum_{i=1}^{m_{\mathrm{bin}}}\mathsf H_i^{\mathrm{cal}}
\right)
=
\sum_{i=1}^{m_{\mathrm{bin}}}
\operatorname{Var}(\mathsf H_i^{\mathrm{cal}})
\le m_{\mathrm{bin}}.
\]
Scaling this centered sum proves
\[
\mathbb E[(\widehat B_{\mathrm{cal}}-B(P))^2]
\le\frac1{4m_{\mathrm{bin}}\delta^2}.
\]

For the vector block, \cref{obj:lem:private-moment-and-dependence} supplies
\[
\mathbb E[W_{ij}^{\mathrm{cal}}]=\tau_j(P),
\qquad
\mathbb E[(W_{ij}^{\mathrm{cal}})^2]=b^2.
\]
The hypotheses of that result hold with block size \(m_{\mathrm{vec}}\), since
\(0<m_{\mathrm{vec}}\le n\), \(d=2\), and \(P\in\mathcal V_2\).
Expanding the squared column sum and using independence of distinct rows yields, for each \(j\in[2]\),
\[
\begin{aligned}
\mathbb E[\widehat\tau_j]
&=\tau_j(P),\\
\mathbb E[\widehat\tau_j^2]
&=\frac{
m_{\mathrm{vec}}b^2+
m_{\mathrm{vec}}(m_{\mathrm{vec}}-1)\tau_j(P)^2
}{m_{\mathrm{vec}}^2}
=\tau_j(P)^2+
\frac{b^2-\tau_j(P)^2}{m_{\mathrm{vec}}}.
\end{aligned}
\]
Subtracting the squared mean therefore gives
\[
\mathbb E[(\widehat\tau_j-\tau_j(P))^2]
=
\frac{b^2-\tau_j(P)^2}{m_{\mathrm{vec}}}
\le\frac{b^2}{m_{\mathrm{vec}}}
=\frac4{m_{\mathrm{vec}}\delta^2}.
\]
All three statistics are Borel functions of the finite transcript alone. These are the required component bounds, simultaneously for every \(P\in\mathcal V_2\).
% lean: CausalSmith.Stat.LdpOptvalueUniformFrontier.twoCellComponents_of_sampling

\item
In the present branch,
\[
1<\sqrt n\,\varepsilon,
\qquad
\min\left\{1,\frac1{n\varepsilon^2}\right\}
=\frac1t,
\qquad
\min\left\{1,\frac1{\sqrt n\,\varepsilon}\right\}
=\frac1{\sqrt n\,\varepsilon}.
\]
Define the unprojected and projected welfare estimates by
\[
\widetilde V_{\mathrm{cal}}
=\widehat B_{\mathrm{cal}}
+\frac{|\widehat\tau_1|+|\widehat\tau_2|}{4},
\qquad
T(\mathbf Z,R,U)
=\max\left\{\frac14,
\min\left\{\frac34,\widetilde V_{\mathrm{cal}}\right\}\right\}.
\]
This \(T\) is a Borel function of \(\mathbf Z\) alone.
By \cref{obj:prop:signed-causal-bridge},
\[
V(P)=B(P)+\frac{|\tau_1(P)|+|\tau_2(P)|}{4}.
\]
The interior means also give \(V(P)\in[1/4,3/4]\), as established in the preceding branch.

Projection toward this interval reduces squared error:
\[
(T-V(P))^2
\le(\widetilde V_{\mathrm{cal}}-V(P))^2.
\]
Indeed, if \(\widetilde V_{\mathrm{cal}}\le1/4\), then
\(T=1/4\le V(P)\); if \(\widetilde V_{\mathrm{cal}}\ge3/4\), then
\(V(P)\le T=3/4\); and inside the interval \(T=\widetilde V_{\mathrm{cal}}\).

The reverse triangle inequality gives
\[
\bigl(|\widehat\tau_j|-|\tau_j(P)|\bigr)^2
\le(\widehat\tau_j-\tau_j(P))^2
\qquad(j\in[2]).
\]
Applying the squared-sum inequality first to the baseline and contrast contributions and then to the two contrast contributions proves the pointwise bound
\[
(T-V(P))^2
\le
2(\widehat B_{\mathrm{cal}}-B(P))^2
+\frac{
(\widehat\tau_1-\tau_1(P))^2+
(\widehat\tau_2-\tau_2(P))^2
}{4}.
\]
Integrating and substituting the component bounds gives
\[
\mathbb E[(T-V(P))^2]
\le
\frac1{2m_{\mathrm{bin}}\delta^2}
+\frac2{m_{\mathrm{vec}}\delta^2}.
\]

The half split satisfies
\[
m_{\mathrm{bin}}\ge\frac n3,
\qquad
m_{\mathrm{vec}}\ge\frac n2.
\]
For the first inequality, even \(n\) gives \(m_{\mathrm{bin}}=n/2\); odd \(n\ge3\) gives
\(m_{\mathrm{bin}}=(n-1)/2\ge n/3\).
The second follows from \(m_{\mathrm{bin}}\le n/2\).
Together with \(\delta\ge\varepsilon/4\), these bounds yield the numerical calibration
\[
\begin{aligned}
\mathbb E[(T-V(P))^2]
&\le
\frac1{2(n/3)\delta^2}
+\frac2{(n/2)\delta^2}\\
&=\frac{88}{n(4\delta)^2}
\le\frac{88}{n\varepsilon^2}
=\frac{88}{t}
\le\frac{1000}{t}.
\end{aligned}
\]
% lean: CausalSmith.Stat.LdpOptvalueUniformFrontier.two_cell_risk_of_component_mse

\item
Using the same protocol and estimator, define the deterministic interval radius and interval by
\[
q_{\mathrm{cal}}
=\sqrt{\frac{10000}{t}}
=\frac{100}{\sqrt n\,\varepsilon}>0,
\]
\[
I(\mathbf Z,R,U)
=
\left[
\max\left\{\frac14,T-q_{\mathrm{cal}}\right\},
\min\left\{\frac34,T+q_{\mathrm{cal}}\right\}
\right].
\]
Both endpoints are Borel functions of the transcript. Since
\(T\in[1/4,3/4]\) and \(q_{\mathrm{cal}}\ge0\),
\[
\frac14
\le\max\left\{\frac14,T-q_{\mathrm{cal}}\right\}
\le T
\le\min\left\{\frac34,T+q_{\mathrm{cal}}\right\}
\le\frac34.
\]
Thus the endpoints are ordered and the interval is connected, closed, and contained in the required domain.

For \(V(P)\in[1/4,3/4]\), the inequality
\(|T-V(P)|\le q_{\mathrm{cal}}\) implies \(V(P)\in I\). Hence
\[
\{V(P)\notin I\}
\subseteq
\{(T-V(P))^2\ge q_{\mathrm{cal}}^2\}.
\]
Markov's inequality and the risk bound give
\[
\begin{aligned}
q_{\mathrm{cal}}^2\,
\mathbb P\{(T-V(P))^2\ge q_{\mathrm{cal}}^2\}
&\le\mathbb E[(T-V(P))^2]\\
&\le\frac{88}{t}
\le\frac{1000}{t}
=\frac{q_{\mathrm{cal}}^2}{10}.
\end{aligned}
\]
Division by \(q_{\mathrm{cal}}^2>0\) and passage to the coverage complement therefore yield
\[
\mathbb P\{V(P)\in I\}\ge1-\frac1{10}=0.90.
\]

Finally, the upper endpoint is at most \(T+q_{\mathrm{cal}}\), and the lower endpoint is at least \(T-q_{\mathrm{cal}}\). Consequently,
\[
\operatorname{length}I
\le2q_{\mathrm{cal}}
=\frac{200}{\sqrt n\,\varepsilon},
\]
and, under the probability decision law,
\[
\mathbb E[\operatorname{length}I]
\le\frac{200}{\sqrt n\,\varepsilon}
\le\frac{1000}{\sqrt n\,\varepsilon}.
\]
This completes the simultaneous attainment bounds in the branch \(t>1\).
% lean: CausalSmith.Stat.LdpOptvalueUniformFrontier.twoCellInterval_coverage_of_risk
% lean: CausalSmith.Stat.LdpOptvalueUniformFrontier.twoCellInterval_expectedLength_le
% lean: CausalSmith.Stat.LdpOptvalueUniformFrontier.two_cell_unsaturated_attainment_of_components

\item
The two branches provide, for every allowed tuple and sampling scheme, a total noninteractive protocol on the original records and transcript-based decisions satisfying
\[
\sup_{P\in\mathcal V_2}\mathbb E[(T-V(P))^2]
\le1000\min\left\{1,\frac1{n\varepsilon^2}\right\},
\]
\[
\inf_{P\in\mathcal V_2}\mathbb P\{V(P)\in I\}\ge0.90,
\qquad
\sup_{P\in\mathcal V_2}\mathbb E[\operatorname{length}I]
\le1000\min\left\{1,\frac1{\sqrt n\,\varepsilon}\right\}.
\]
The definitions in
\cref{obj:def:noninteractive-class,obj:def:sequential-class} give
\[
\mathfrak Q_{\mathrm{NI}}(n,2,\varepsilon)
\subseteq
\mathfrak Q_{\mathrm{SI}}(n,2,\varepsilon).
\]
Thus these same witnesses are admissible for either protocol class. Taking the infima in
\cref{obj:def:risk,obj:def:honest-length} proves, for every
\(\mathcal C\in\{\mathrm{NI},\mathrm{SI}\}\),
\[
\mathfrak R_{\mathcal C}(n,2,\varepsilon)
\le1000\min\left\{1,\frac1{n\varepsilon^2}\right\},
\qquad
\mathfrak H_{\mathcal C}(n,2,\varepsilon)
\le1000\min\left\{1,\frac1{\sqrt n\,\varepsilon}\right\}.
\]
% lean: CausalSmith.Stat.LdpOptvalueUniformFrontier.causalAttainment_minimax_upper

\item
Apply the lower bounds obtained at the start to the fixed sampling scheme and either protocol class. Since \(c\le c_1\) and both rate factors are nonnegative,
\[
c\min\left\{1,\frac1{n\varepsilon^2}\right\}
\le
c_1\min\left\{1,\frac1{n\varepsilon^2}\right\}
\le\mathfrak R_{\mathcal C}(n,2,\varepsilon),
\]
\[
c\min\left\{1,\frac1{\sqrt n\,\varepsilon}\right\}
\le
c_1\min\left\{1,\frac1{\sqrt n\,\varepsilon}\right\}
\le\mathfrak H_{\mathcal C}(n,2,\varepsilon).
\]
Together with the upper bounds and the simultaneous witnesses, these prove all assertions with the chosen universal constants \(c\) and \(C_0\).
% lean: CausalSmith.Stat.LdpOptvalueUniformFrontier.two_cell_calibration
\end{enumerate}
\end{proof}

\begin{proof}[Proof of \cref{obj:thm:paired-uniform-tv-frontier}]
Throughout, resources are admissible:
\[
n,d\in\mathbb N,\qquad n,d\ge2,\qquad 0<\varepsilon\le1.
\]
We use the effective resource, logarithmic dimension, and converse tuning quantities
\[
t=n\varepsilon^2>0,\qquad L=\log(ed),\qquad
z_*=\frac{d^2L}{t},\qquad w=\log(e+z_*).
\]
In particular, \(1\le L\le d\). The rates in \cref{obj:def:frontier-handle} are independent of the protocol-class label:
\[
r_{\mathrm{NI}}=r_{\mathrm{SI}}=\rho(t,d),
\qquad
h_{\mathrm{NI}}=h_{\mathrm{SI}}=\sqrt{\rho(t,d)}.
\]

\begin{enumerate}
\item \textit{A uniform approximation constant.}
Invoke the published Bernstein limit
\[
K e_K\longrightarrow\beta_*
\quad\text{along positive even integers }K,
\qquad 0.278<\beta_*<0.286,
\]
as stated in \citep[Section 3.1, definition of the Bernstein constant and the immediately following paragraph, p.~8]{CaiLow2011Nonsmooth}; see also \citep{Bernstein1913}.
Choose an integer \(N_{\mathrm B}\ge1\) such that
\[
K e_K\ge\frac1{10}
\qquad\text{for every even }K\ge2N_{\mathrm B},
\]
and define
\[
b_{\mathrm{app}}
=\min\left\{\frac1{10},e_{2N_{\mathrm B}}\right\}>0.
\]
The positivity follows from the preceding tail bound at \(K=2N_{\mathrm B}\).
Approximation error decreases as the allowed degree increases. Hence, for every positive even \(K\),
\[
\frac{b_{\mathrm{app}}}{K}\le e_K.
\]
Indeed, the tail bound proves this when \(K\ge2N_{\mathrm B}\); otherwise,
\[
b_{\mathrm{app}}\le e_{2N_{\mathrm B}}\le e_K\le K e_K.
\]
% lean: CausalSmith.Stat.LdpOptvalueUniformFrontier.bernstein_uniform_positive_of_gate

\item \textit{Paired product priors.}
Take the converse amplitude and matching degree from \cref{obj:def:frontier-handle}:
\[
(\alpha,k)=
\begin{cases}
\left(\dfrac{\sqrt{z_*}}{100},\,2\lceil16L\rceil+1\right),&z_*\le1,\\[4pt]
\left(\dfrac1{100},\,2\lceil16L/w\rceil+1\right),&z_*>1,
\end{cases}
\qquad K=k-1.
\]
These choices satisfy \(0<\alpha\le1/2\), and \(K\) is a positive even integer.
By \cref{obj:lem:cai-low-moment-duality}, there are symmetric probability measures \(\xi_0,\xi_1\), supported on \([-1,1]\), whose moments agree through degree \(K\) and whose absolute moments differ by \(2e_K\). Define the scaled coordinate priors and their product priors by
\[
\nu_\ell=(x\mapsto\alpha x)_\#\xi_\ell,
\qquad
\Pi_\ell=\nu_\ell^{\otimes d},
\qquad \ell\in\{0,1\}.
\]
They are supported on \([-\alpha,\alpha]^d\subseteq[-1/2,1/2]^d\), and
\[
\int u^q\,d\nu_1(u)=\int u^q\,d\nu_0(u)
\qquad(0\le q\le K).
\]
Using the paired-target identity in \cref{obj:prop:signed-causal-bridge}, define their total-variation centers and separation by
\[
v_\ell^{\mathrm{TV}}=\frac12\int |u|\,d\nu_\ell(u),
\qquad
\Delta_{\mathrm{TV}}
=v_1^{\mathrm{TV}}-v_0^{\mathrm{TV}}
=\alpha e_K>0.
\]
Independence of the coordinates gives
\[
\mathbb E_{\Pi_\ell}D_{\mathrm{TV}}(P_\theta)=v_\ell^{\mathrm{TV}},
\qquad
\operatorname{Var}_{\Pi_\ell}\!\left(D_{\mathrm{TV}}(P_\theta)\right)
\le\frac{\alpha^2}{4d}.
\]
Consequently, Chebyshev's inequality and the approximation bound give
\[
\Pi_\ell\left\{
\left|D_{\mathrm{TV}}(P_\theta)-v_\ell^{\mathrm{TV}}\right|
\ge\frac{\Delta_{\mathrm{TV}}}{8}
\right\}
\le
\frac{16\alpha^2}{d\Delta_{\mathrm{TV}}^2}
\le\frac{16K^2}{b_{\mathrm{app}}^2d}.
\]

The ceiling inequalities also give the required rate calibration. If \(z_*\le1\), then
\[
K\le34L,\qquad
\sqrt{\rho(t,d)}=\frac{\sqrt{z_*}}{L},
\qquad
\frac{\alpha}{K}\ge\frac1{3400}\sqrt{\rho(t,d)}.
\]
If \(z_*>1\), then
\[
K\le34\max\{1,L/w\},
\qquad
\sqrt{\rho(t,d)}=\min\{1,w/L\},
\]
so again
\[
\frac{\alpha}{K}\ge\frac1{3400}\sqrt{\rho(t,d)}.
\]
Since \(w\ge1\), both branches satisfy \(K\le34L\). Define
\[
s_{\mathrm{app}}=\frac{b_{\mathrm{app}}}{3400}>0.
\]
Then
\[
\Delta_{\mathrm{TV}}\ge s_{\mathrm{app}}\sqrt{\rho(t,d)}.
\]
Finally, \(L^2/d\to0\) as \(d\to\infty\). Thus there is a universal integer \(D_{\mathrm{cut}}\ge2\) such that
\[
\frac{16\cdot34^2L^2}{b_{\mathrm{app}}^2d}\le\frac1{100}
\qquad(d>D_{\mathrm{cut}}).
\]
For those dimensions, each prior's target-concentration failure probability is at most \(1/100\).
% lean: CausalSmith.Stat.LdpOptvalueUniformFrontier.paired_lawspace_resource_priors_of_gate

\item \textit{Contraction of the paired transcript mixtures.}
Fix any sequential paired-input protocol \(\mathsf K\).
By \cref{obj:prop:signed-causal-bridge}, its pullback
\[
Q_i(E\mid(j,a,y),\eta,r)
=
\mathsf K_i\!\left(E\mid(j,(2a-1)(2y-1)),\eta,r\right)
\]
is sequentially private at the same budget, and its seed-transcript law under \(G_\theta\) equals that of \(\mathsf K\) under \(P_\theta\).
The density hypothesis for \cref{obj:lem:adaptive-moment-contraction} is supplied by \cref{obj:lem:measurable-kernel-radon-nikodym}.
Define the paired seed-transcript mixtures by
\[
\mathsf M_\ell^{\mathrm{TV}}
=
\int\mathcal L_{P_\theta,\mathsf K}(R,\mathbf Z^{\mathsf K})\,
d\Pi_\ell(\theta),
\qquad \ell\in\{0,1\},
\]
and the contraction scale by
\[
\beta=\frac{e^\varepsilon-1}{2d}.
\]
The moment agreement in the previous step allows \cref{obj:lem:adaptive-moment-contraction} to yield
\[
\operatorname{TV}(\mathsf M_0^{\mathrm{TV}},\mathsf M_1^{\mathrm{TV}})
\le
\min\left\{1,d\alpha^k\sqrt{\binom nk}\,\beta^k\right\}
\qquad(k\le n),
\]
whereas the two mixtures are equal when \(k>n\).

For \(0<\varepsilon\le1\), the exponential bound
\(e^\varepsilon-1\le2\varepsilon\) gives \(\beta\le\varepsilon/d\).
Moreover,
\[
\binom nk\le\frac{n^k}{k!}\le\left(\frac{en}{k}\right)^k,
\]
because \(k^k/k!\le e^k\), as follows from the exponential series.
Thus, when \(k\le n\),
\[
\operatorname{TV}(\mathsf M_0^{\mathrm{TV}},\mathsf M_1^{\mathrm{TV}})
\le
\min\left\{1,
d\left(\alpha\sqrt{\frac{en}{k}}\frac{\varepsilon}{d}\right)^k
\right\}.
\]

To evaluate this bound, define
\[
w_{\mathrm{con}}=
\begin{cases}
1,&z_*\le1,\\
w,&z_*>1,
\end{cases}
\qquad
B_{\mathrm{con}}=\alpha\sqrt{\frac{en}{k}}\frac{\varepsilon}{d}.
\]
The prescribed ceilings give
\[
32L\le k w_{\mathrm{con}},
\qquad
B_{\mathrm{con}}^2=\frac{\alpha^2eL}{kz_*}.
\]
In the branch \(z_*\le1\), the amplitude choice and \(e<3\) imply
\[
\alpha^2e\,w_{\mathrm{con}}e^{w_{\mathrm{con}}/2}
=\frac{z_*}{10^4}e^{3/2}\le32z_*.
\]
In the other branch, \(e^w=e+z_*\), and \(w\le2e^{w/2}\) gives
\[
w e^{w/2}\le2(e+z_*)\le8z_*.
\]
Hence the same bound
\[
\alpha^2e\,w_{\mathrm{con}}e^{w_{\mathrm{con}}/2}\le32z_*
\]
holds. Combining it with \(32L\le kw_{\mathrm{con}}\) yields
\[
B_{\mathrm{con}}^2\le e^{-w_{\mathrm{con}}/2},
\qquad
B_{\mathrm{con}}\le e^{-w_{\mathrm{con}}/4}.
\]
Therefore
\[
d B_{\mathrm{con}}^k
\le d e^{-8L}
\le e^{-7L}
\le e^{-7}\le\frac1{100}.
\]
For \(k>n\), exact equality of the mixtures gives the same conclusion:
\[
\operatorname{TV}(\mathsf M_0^{\mathrm{TV}},\mathsf M_1^{\mathrm{TV}})
\le\frac1{100}.
\]
% lean: CausalSmith.Stat.LdpOptvalueUniformFrontier.paired_lawspace_resource_contraction

\item \textit{Decision lower bounds in large dimension.}
We give the paired-model decision reduction explicitly.
For an estimator and a globally \(0.90\)-honest interval, define their worst-law risk and expected length by
\[
\begin{aligned}
\mathcal R_{\mathsf K}(\widehat D)
&=\sup_{\mathsf p\in\mathcal P_d^{\mathrm{pair}}}
\mathbb E_{\mathsf p,\mathsf K}
\left[(\widehat D-D_{\mathrm{TV}}(\mathsf p))^2\right],\\
\mathcal H_{\mathsf K}(J_{\mathrm{TV}})
&=\sup_{\mathsf p\in\mathcal P_d^{\mathrm{pair}}}
\mathbb E_{\mathsf p,\mathsf K}\operatorname{length}J_{\mathrm{TV}}.
\end{aligned}
\]
These definitions apply to every admissible Borel decision of \cref{obj:def:tv-private-decision-values}.
For this reduction, let \(\eta_0,\omega\ge0\) bound the two concentration failures and the seed-transcript total variation, respectively.
Let
\[
N_\ell^{\mathrm{TV}}
=\mathsf M_\ell^{\mathrm{TV}}\otimes\operatorname{Uniform}[0,1]
\]
be the joint seed-transcript-randomizer mixture.
Appending the independent analyst randomizer preserves the bound:
\[
\operatorname{TV}(N_0^{\mathrm{TV}},N_1^{\mathrm{TV}})\le\omega.
\]
Indeed, integrating an event indicator over the independent randomizer produces a measurable function taking values in \([0,1]\).

Threshold \(\widehat D\) at
\((v_0^{\mathrm{TV}}+v_1^{\mathrm{TV}})/2\).
On the target-concentration event, a wrong prior label entails an estimation error of magnitude at least \(3\Delta_{\mathrm{TV}}/8\).
If the worst-law risk is finite, Markov's inequality bounds the sum of testing errors by
\[
2\eta_0+
\frac{128\,\mathcal R_{\mathsf K}(\widehat D)}
{9\Delta_{\mathrm{TV}}^2}.
\]
The testing lower bound from total variation is \(1-\omega\). It follows that
\[
\mathcal R_{\mathsf K}(\widehat D)
\ge
\frac{9\Delta_{\mathrm{TV}}^2}{128}
(1-\omega-2\eta_0)_+.
\]
The same inequality holds immediately when the risk is infinite.

For the interval reduction, define
\[
\mathcal E_\ell^{\mathrm{TV}}
=
\left\{
J_{\mathrm{TV}}\cap
\left[v_\ell^{\mathrm{TV}}-\frac{\Delta_{\mathrm{TV}}}{8},
      v_\ell^{\mathrm{TV}}+\frac{\Delta_{\mathrm{TV}}}{8}\right]
\ne\varnothing
\right\}.
\]
Paired-model honesty and target concentration imply
\[
N_\ell^{\mathrm{TV}}(\mathcal E_\ell^{\mathrm{TV}})
\ge0.90-\eta_0.
\]
Transferring the second event to the first mixture and applying the intersection bound gives
\[
N_0^{\mathrm{TV}}
(\mathcal E_0^{\mathrm{TV}}\cap\mathcal E_1^{\mathrm{TV}})
\ge0.80-2\eta_0-\omega.
\]
On this intersection, ordered endpoints and connectedness force length at least \(3\Delta_{\mathrm{TV}}/4\). Therefore
\[
\mathcal H_{\mathsf K}(J_{\mathrm{TV}})
\ge
\frac{3\Delta_{\mathrm{TV}}}{4}
(0.80-2\eta_0-\omega)_+.
\]

For \(d>D_{\mathrm{cut}}\), the preceding steps permit
\(\eta_0=\omega=1/100\). Define
\[
c_{\mathrm{large}}
=\min\left\{\frac{s_{\mathrm{app}}^2}{100},
            \frac{s_{\mathrm{app}}}{100}\right\}>0.
\]
Since \(\Delta_{\mathrm{TV}}\ge s_{\mathrm{app}}\sqrt{\rho(t,d)}\),
\[
\begin{aligned}
\mathcal R_{\mathsf K}(\widehat D)
&\ge\frac9{128}\frac{97}{100}\Delta_{\mathrm{TV}}^2
\ge c_{\mathrm{large}}\rho(t,d),\\
\mathcal H_{\mathsf K}(J_{\mathrm{TV}})
&\ge\frac34\frac{77}{100}\Delta_{\mathrm{TV}}
\ge c_{\mathrm{large}}\sqrt{\rho(t,d)}.
\end{aligned}
\]
% lean: CausalSmith.Stat.LdpOptvalueUniformFrontier.paired_converse_of_gate

\item \textit{Elementary rate comparisons and bounded dimension.}
We first establish the numerical comparisons used to complete the converse:
\[
\min\{1,t^{-1}\}\le\rho(t,d)\le4\min\{1,d^2/t\},
\qquad
\frac14\min\{t,d\}\le t\rho(t,d).
\]
In the dense branch \(t\ge d^2L\), these follow from
\[
\rho(t,d)=\frac{d^2}{tL},\qquad
\frac1t\le\rho(t,d)\le\min\{1,d^2/t\},
\qquad
t\rho(t,d)=\frac{d^2}{L}\ge d.
\]

For the nondense branch, define the auxiliary function
\[
\Psi(s_{\mathrm{res}})
=s_{\mathrm{res}}
\left[\log\left(e+\frac{d^2L}{s_{\mathrm{res}}}\right)\right]^2,
\qquad s_{\mathrm{res}}>0.
\]
Writing \(v_{\mathrm{res}}=d^2L/s_{\mathrm{res}}>0\), differentiation gives
\[
\Psi'(s_{\mathrm{res}})
=
\log(e+v_{\mathrm{res}})
\left[
\log(e+v_{\mathrm{res}})
-\frac{2v_{\mathrm{res}}}{e+v_{\mathrm{res}}}
\right]\ge0.
\]
For the last inequality, \(\log y\ge1-y^{-1}\), applied at
\(y=(e+v_{\mathrm{res}})/e\), gives
\[
\log(e+v_{\mathrm{res}})
\ge1+\frac{v_{\mathrm{res}}}{e+v_{\mathrm{res}}}
\ge\frac{2v_{\mathrm{res}}}{e+v_{\mathrm{res}}}.
\]
Also,
\[
\log(e+d^2L)\ge L,
\qquad
\log(e+dL)\ge L/2.
\]
The first follows from \(dL\ge e\); the second follows by combining
\(\log(e+dL)\ge1\) and \(\log(e+dL)\ge\log d\).
Thus, if \(t\le1\), then \(w\ge L\) and \(\rho(t,d)=1\).
If \(t\ge1\), monotonicity of \(\Psi\) gives
\(tw^2\ge L^2\), establishing the lower comparison with \(\min\{1,t^{-1}\}\).
Likewise, if \(t\le d\), then \(w\ge L/2\); if \(t\ge d\), monotonicity gives
\(tw^2\ge dL^2/4\). Applying the cap at one proves the comparison with \(\min\{t,d\}/4\).

For the upper comparison, \(z_*>1\) in this branch and
\[
w\le\log(e+1)+\log z_*
\le2+2(\sqrt{z_*}-1)=2\sqrt{z_*}.
\]
Hence
\[
\frac{w^2}{L^2}\le\frac{4z_*}{L^2}
=\frac{4d^2}{tL}\le\frac{4d^2}{t}.
\]
Together with \(\rho(t,d)\le1\), this proves all three comparisons.

Now define the two-point amplitude
\[
\alpha_{\mathrm{two}}
=\frac18\min\{1,t^{-1/2}\}\in(0,1/2].
\]
Consider the paired laws with constant contrast zero and constant contrast
\(\alpha_{\mathrm{two}}\). Their targets are zero and
\(\alpha_{\mathrm{two}}/2\).
Applying \cref{obj:lem:adaptive-moment-contraction} at matching order one through the bridge in \cref{obj:prop:signed-causal-bridge} gives
\[
\operatorname{TV}\!\left(
\mathcal L_{P_0,\mathsf K}(R,\mathbf Z^{\mathsf K}),
\mathcal L_{P_{\alpha_{\mathrm{two}}\mathbf1},\mathsf K}
(R,\mathbf Z^{\mathsf K})
\right)
\le
\alpha_{\mathrm{two}}\sqrt n\,\varepsilon
\le\frac18.
\]
Here \(P_{\alpha_{\mathrm{two}}\mathbf1}\) denotes the law with
\(\theta_j=\alpha_{\mathrm{two}}\) for every \(j\).
The preceding decision reduction, with deterministic targets and
\(\omega=1/8\), therefore yields
\[
\mathcal R_{\mathsf K}(\widehat D)
\ge\frac{63\alpha_{\mathrm{two}}^2}{4096}
\ge\frac1{8192}\min\{1,t^{-1}\},
\]
\[
\mathcal H_{\mathsf K}(J_{\mathrm{TV}})
\ge\frac{81\alpha_{\mathrm{two}}}{320}
\ge\frac1{8192}\min\{1,t^{-1/2}\}.
\]
The calibrations use
\[
\alpha_{\mathrm{two}}^2
=\frac1{64}\min\{1,t^{-1}\}.
\]

If \(d\le D_{\mathrm{cut}}\), the elementary upper comparison implies
\[
\rho(t,d)
\le4D_{\mathrm{cut}}^2\min\{1,t^{-1}\},
\]
and taking square roots gives
\[
\sqrt{\rho(t,d)}
\le2D_{\mathrm{cut}}\min\{1,t^{-1/2}\}
\le4D_{\mathrm{cut}}^2\min\{1,t^{-1/2}\}.
\]
Define
\[
c_{\mathrm{small}}=\frac1{32768D_{\mathrm{cut}}^2},
\qquad
c_{\mathrm{lower}}=\min\{c_{\mathrm{large}},c_{\mathrm{small}}\}>0.
\]
The cases \(d\le D_{\mathrm{cut}}\) and \(d>D_{\mathrm{cut}}\) now give, for every sequential paired protocol and every admissible decision,
\[
\mathcal R_{\mathsf K}(\widehat D)\ge
c_{\mathrm{lower}}\rho(t,d),
\qquad
\mathcal H_{\mathsf K}(J_{\mathrm{TV}})\ge
c_{\mathrm{lower}}\sqrt{\rho(t,d)}.
\]
Every noninteractive protocol is sequential. Taking the decision and protocol infima in \cref{obj:def:tv-private-decision-values} therefore proves these lower bounds for both classes.
% lean: CausalSmith.Stat.LdpOptvalueUniformFrontier.paired_minimax_lower_of_gate

\item \textit{The concrete upper-risk construction.}
Choose the universal constants
\[
c=\min\{c_{\mathrm{lower}},1\},
\qquad C_0=10^{16}.
\]
Then \(0<c\le C_0\).
Use the paired procedure specified in \cref{obj:synth_5}, with a singleton public seed. Its channel quantities are
\[
m=\lfloor n/3\rfloor,\qquad
\delta=\tanh(\varepsilon/2),\qquad b=\frac d\delta.
\]
Each of the first \(2m\) participants releases a vector with probabilities
\[
\mathsf K_i(z'\mid(j,s))
=2^{-d}(1+\delta s z'_j),
\qquad z'\in\{-1,1\}^d;
\]
the remaining participants release \((-1,\ldots,-1)\).
Symmetry of the sign-vector alphabet makes the active probabilities sum to one, and
\[
\frac{1+\delta}{1-\delta}=e^\varepsilon
\]
bounds the likelihood ratio at every active message. The constant releases satisfy the privacy inequality as well. All kernels are independent of history, so this is a pure \(\varepsilon\)-private noninteractive protocol with the declared sign-vector message spaces.

For \(n>2\), define the pilot and evaluation matrices, means, and calibration quantities by
\[
W^{\mathrm{pil}}_{ij}=bZ_{i,j},\qquad
W^{\mathrm{ev}}_{ij}=bZ_{m+i,j},
\qquad 1\le i\le m,\quad j\in[d],
\]
\[
M_j=\frac1m\sum_iW^{\mathrm{pil}}_{ij},\qquad
\overline W_j=\frac1m\sum_iW^{\mathrm{ev}}_{ij},
\qquad
\sigma^2=\frac{b^2}{m},\qquad \sigma=\sqrt{\sigma^2},
\qquad H=\log(e+\sigma^2L).
\]
The statistic \(\widehat F\) is the ordered branch-selected norm estimate in
\cref{obj:synth_5}, using the private moments and polynomial coefficients specified there and in \cref{obj:def:private-moments,obj:def:frontier-handle}.
Set
\[
\widehat D=
\begin{cases}
1/8,&n=2,\\
\max\{0,\min\{1/4,\widehat F/2\}\},&n>2.
\end{cases}
\]
The decisions are measurable functions of the finite transcript.

The cube target belongs to \([0,1/4]\), so
\[
(\widehat D-D_{\mathrm{TV}}(P_\theta))^2\le\frac1{16}.
\]
For \(n>2\), clipping also gives the sharper comparison
\[
\mathbb E_{P_\theta,\mathsf K}
(\widehat D-D_{\mathrm{TV}}(P_\theta))^2
\le\frac14
\mathbb E_{P_\theta,\mathsf K}(\widehat F-F(\theta))^2.
\]
To apply the existing moment and polynomial calibrations, use \(G_\theta\) from
\cref{obj:def:symmetric-law}. Its contrasts satisfy
\[
\tau_j(G_\theta)
=\frac{1+\theta_j}{2}-\frac{1-\theta_j}{2}
=\theta_j.
\]
The bridge in \cref{obj:prop:signed-causal-bridge} identifies the active paired blocks with independent vector-channel blocks under this law. Thus
\cref{obj:lem:private-moment-and-dependence,obj:lem:finite-private-polynomial-calibration}
apply with contrast vector \(\theta\); their concentration and approximation hypotheses are supplied by
\cref{obj:lem:bounded-mean-concentration,obj:lem:cai-low-chebyshev-approximation}.

We verify the resource calibration used in the risk bounds. For \(n\ge3\),
\[
m\ge n/5,\qquad
\delta=\frac{e^\varepsilon-1}{e^\varepsilon+1}\ge\varepsilon/4,
\qquad
\sigma^2\le\frac{80d^2}{t}.
\]
The attenuation bound follows from \(e^\varepsilon\ge1+\varepsilon\) and
\(e^\varepsilon+1<4\).
Consequently,
\[
1\le H\le\log(80(e+z_*))
\le5+w\le6w.
\]

We now follow the risk case split.

If \(n=2\), then \(t\le2\), and the elementary lower comparison gives
\(\rho(t,d)\ge1/2\). The constant estimate satisfies
\[
\mathbb E(\widehat D-D_{\mathrm{TV}}(P_\theta))^2
\le\frac1{64}\le C_0\rho(t,d).
\]

Suppose \(n>2\). If \(\rho(t,d)=1\), the uniform \(1/16\) loss bound gives the desired result.
Next consider the constant-fallback branch. Its conditions include
\[
t<d^2L,\qquad \frac{L}{1024H}<4.
\]
Since \(H\le6w\),
\[
L<4096H\le24576w,
\qquad
\rho(t,d)\ge\left(\frac1{24576}\right)^2.
\]
Again,
\[
\mathbb E(\widehat D-D_{\mathrm{TV}}(P_\theta))^2
\le\frac1{16}\le C_0\rho(t,d).
\]

For the remaining branches, suppose first that \(L<4096\).
In the dense case \(t\ge d^2L\), the absolute-mean statistic satisfies
\[
\mathbb E\left[
\left(\frac1d\sum_j|\overline W_j|-F(\theta)\right)^2
\right]
\le\frac1d\sum_j\mathbb E(\overline W_j-\theta_j)^2
\le\sigma^2.
\]
Here the first inequality uses the absolute-value Lipschitz inequality and the quadratic mean inequality; the second uses the moment bounds and independence across participants. The numerical calibration is
\[
\frac{24}{t}+\frac{\sigma^2}{2}
\le\frac{64d^2}{t}
=64L\rho(t,d)
\le64\cdot4096\,\rho(t,d).
\]
Since the TV risk is at most \(\sigma^2/4\), this proves its \(C_0\rho(t,d)\) bound.
In the nondense case, \(w\ge1\) implies
\[
\rho(t,d)\ge\left(\frac1{4096}\right)^2,
\]
so the uniform loss bound suffices.

Finally, suppose \(L\ge4096\).
In the dense case, the prescribed hybrid degree and radius are
\[
D=2\lfloor L/2048\rfloor,\qquad
T_*=4\sigma\sqrt L,\qquad a_{\mathrm{rad}}=8\sigma\sqrt L.
\]
Its degree fits the block:
\[
2D\le L/512\le m.
\]
Indeed, \(n\ge t\ge d^2L\ge4L\) and \(m\ge n/5\).
The hybrid conclusion of \cref{obj:lem:finite-private-polynomial-calibration} gives
\[
\mathbb E(\widehat F-F(\theta))^2
\le500000000\,\frac{\sigma^2}{L}
\le40000000000\,\rho(t,d).
\]
Also \(t^{-1}\le\rho(t,d)\), and hence
\[
\frac{24}{t}
+\frac12\left(500000000\,\frac{\sigma^2}{L}\right)
\le(24+20000000000)\rho(t,d)
\le C_0\rho(t,d).
\]
The clipping comparison proves the claimed TV risk bound.

In the nondense case remaining after the saturation and fallback cases,
\[
\rho(t,d)=(w/L)^2,\qquad
a_{\mathrm{rad}}=\frac12,\qquad
D=2\left\lfloor\frac{L}{2048H}\right\rfloor,
\]
and floor rounding gives
\[
2\le D,\qquad
\frac{L}{2048H}\le D\le\frac{L}{1024H}.
\]
The moment-fit condition \(2D\le m\) also holds. To see this, if \(m<2D\), the degree budget and \(H\ge1\) imply \(m<L/512\).
Since \(\delta<1\),
\[
\sigma^2\ge\frac{d^2}{m},
\qquad
e+\sigma^2L>512d^2\ge ed,
\]
so \(H\ge L\), contradicting \(1024DH\le L\) and \(D\ge2\).

The global-polynomial conclusion of
\cref{obj:lem:finite-private-polynomial-calibration} now yields
\[
\mathbb E(\widehat F-F(\theta))^2
\le\frac{e^{9DH}}{2d}+\frac1{(D+1)^2}.
\]
The degree budget gives \(9DH\le L/100\).
For \(L\ge4096\), the elementary exponential inequalities
\[
L^2\le16e^{L/2},
\qquad
4\le e^{49L/100-1}
\]
imply
\[
L^2e^{L/100}\le4e^{L-1}=4d.
\]
Therefore
\[
\frac{e^{9DH}}{2d}\le\frac2{L^2}\le2\rho(t,d),
\qquad
\frac1{D+1}\le\frac{2048H}{L}
\le\frac{12288w}{L},
\]
and hence
\[
\frac1{(D+1)^2}\le150994944\,\rho(t,d).
\]
The elementary lower comparison, together with \(\rho(t,d)<1\), gives
\(t^{-1}\le\rho(t,d)\). Thus
\[
\frac{24}{t}
+\frac12\left(\frac{e^{9DH}}{2d}+\frac1{(D+1)^2}\right)
\le75497500\,\rho(t,d)
\le C_0\rho(t,d).
\]
Clipping again completes the risk bound. Every admissible tuple has therefore been covered:
\[
\sup_{\mathsf p\in\mathcal P_d^{\mathrm{pair}}}
\mathbb E_{\mathsf p,\mathsf K}
(\widehat D-D_{\mathrm{TV}}(\mathsf p))^2
\le C_0\rho(t,d).
\]
% lean: CausalSmith.Stat.LdpOptvalueUniformFrontier.tvFrontierEstimator_uniform_rate

\item \textit{Intervals and the common minimax bounds.}
Define
\[
q_*=\sqrt{10C_0\rho(t,d)}=\sqrt{10^{17}\rho(t,d)},
\qquad
J_{\mathrm{TV}}
=
\left[
\max\{0,\widehat D-q_*\},
\min\{1/4,\widehat D+q_*\}
\right].
\]
The elementary lower comparison gives \(\rho(t,d)>0\).
The interval contains its center, has ordered Borel endpoints, and is connected, closed, and contained in \([0,1/4]\).
Since the target is also in this range, failure of coverage entails
\(|\widehat D-D_{\mathrm{TV}}(\mathsf p)|>q_*\). Markov's inequality gives
\[
\Pr_{\mathsf p,\mathsf K}
\{D_{\mathrm{TV}}(\mathsf p)\notin J_{\mathrm{TV}}\}
\le\frac{C_0\rho(t,d)}{q_*^2}
=\frac1{10}.
\]
Furthermore,
\[
\operatorname{length}J_{\mathrm{TV}}
\le2q_*
=2\sqrt{10C_0}\sqrt{\rho(t,d)}
\le C_0\sqrt{\rho(t,d)}.
\]
This proves all three concrete attainment guarantees.
The same construction supplies the general attainment witness: its message spaces are exactly \(\{-1,1\}^d\), and its estimate depends on the transcript alone.

The construction belongs to both protocol classes. Its interval, being contained in \([0,1/4]\), is also an admissible interval in the optimization over \([0,1]\).
Taking the corresponding candidates in the minimax infima gives, for either \(\mathcal C\),
\[
\mathfrak R^{\mathrm{TV}}_{\mathcal C}
(\mathcal P_d^{\mathrm{pair}};n,d,\varepsilon)
\le C_0\rho(t,d),
\qquad
\mathfrak H^{\mathrm{TV}}_{\mathcal C}
(\mathcal P_d^{\mathrm{pair}};n,d,\varepsilon)
\le C_0\sqrt{\rho(t,d)}.
\]
Since \(c\le c_{\mathrm{lower}}\) and both rates are nonnegative, the lower bounds proved above give
\[
\begin{aligned}
c\rho(t,d)
&\le\mathfrak R^{\mathrm{TV}}_{\mathcal C}
(\mathcal P_d^{\mathrm{pair}};n,d,\varepsilon)
\le C_0\rho(t,d),\\
c\sqrt{\rho(t,d)}
&\le\mathfrak H^{\mathrm{TV}}_{\mathcal C}
(\mathcal P_d^{\mathrm{pair}};n,d,\varepsilon)
\le C_0\sqrt{\rho(t,d)}.
\end{aligned}
\]
% lean: CausalSmith.Stat.LdpOptvalueUniformFrontier.tvAttainment_minimax_upper

\item \textit{Identification of the paired functional.}
For \(d\ge2\) and \(\theta\in[-1/2,1/2]^d\), the masses
\[
P_\theta(j,s)=\frac{1+s\theta_j}{2d}
\]
are nonnegative, and their sum is one because each pair has total mass \(1/d\).
Taking zero contrast gives the uniform probability law.
For two probability laws on a finite alphabet, total variation equals half the sum of the absolute mass differences. Therefore
\[
D_{\mathrm{TV}}(P_\theta)
=\operatorname{TV}(P_\theta,\mathsf U_{2d}).
\]
For each coordinate,
\[
\sum_{s\in\{-1,1\}}
\left|\frac{1+s\theta_j}{2d}-\frac1{2d}\right|
=
\frac{|\theta_j|}{d}.
\]
Summing and multiplying by one half gives
\[
D_{\mathrm{TV}}(P_\theta)
=\frac1{2d}\sum_{j=1}^d|\theta_j|
=\frac12F(\theta),
\qquad
0\le D_{\mathrm{TV}}(P_\theta)\le\frac14.
\]
% lean: CausalSmith.Stat.LdpOptvalueUniformFrontier.tvFromUniform_pairedLaw

\item \textit{Attainment and protocol-class ratios.}
The witnesses constructed above establish the existential attainment assertions for either class, because the class-selected rates agree. They also establish the concrete guarantees for the procedure of
\cref{obj:synth_5,obj:def:frontier-handle,obj:thm:uniform-private-value-frontiers}.

The protocol inclusion in \cref{obj:def:tv-private-decision-values} gives
\[
\mathfrak R^{\mathrm{TV}}_{\mathrm{SI}}
(\mathcal P_d^{\mathrm{pair}};n,d,\varepsilon)
\le
\mathfrak R^{\mathrm{TV}}_{\mathrm{NI}}
(\mathcal P_d^{\mathrm{pair}};n,d,\varepsilon),
\]
and the corresponding inequality for honest length.
The common bounds make all four values finite and strictly positive.
For risk, they also give
\[
\mathfrak R^{\mathrm{TV}}_{\mathrm{NI}}
(\mathcal P_d^{\mathrm{pair}};n,d,\varepsilon)
\le C_0\rho(t,d)
=\frac{C_0}{c}\,c\rho(t,d)
\le\frac{C_0}{c}
\mathfrak R^{\mathrm{TV}}_{\mathrm{SI}}
(\mathcal P_d^{\mathrm{pair}};n,d,\varepsilon).
\]
Dividing by the positive SI value yields the asserted risk ratio bounds.
Replacing \(\rho(t,d)\) by \(\sqrt{\rho(t,d)}\) proves the same ratio bounds for honest length.
% lean: CausalSmith.Stat.LdpOptvalueUniformFrontier.frontier_ratio_comparison

\item \textit{Saturation and consistency of the rate.}
For \(0<t<d^2L\), positivity of \(w/L\) gives
\[
\begin{aligned}
\rho(t,d)=1
&\Longleftrightarrow w\ge L\\
&\Longleftrightarrow e+\frac{d^2L}{t}\ge ed\\
&\Longleftrightarrow
t\le\frac{d^2L}{ed-e}.
\end{aligned}
\]
The last denominator is positive for every \(d\ge2\).

For an admissible resource sequence, define
\[
t_k=n_k\varepsilon_k^2,\qquad
L_k=\log(ed_k),\qquad
g_k=\frac{\log(1+d_k^2/t_k)}{L_k}.
\]
If \(\rho(t_k,d_k)\to0\), the lower comparison
\(\min\{1,t_k^{-1}\}\le\rho(t_k,d_k)\) forces \(t_k\to\infty\).
Also, whenever \(\rho(t_k,d_k)<1\),
\[
g_k^2\le\rho(t_k,d_k).
\]
In the nondense branch this follows from
\[
\log(1+d_k^2/t_k)
\le\log(e+L_kd_k^2/t_k)
\]
and removal of the cap when the rate is below one.
In the dense branch, using \(\log(1+x)\le x\) and
\(d_k^2/t_k\le1/L_k\) gives the same inequality.
Thus \(g_k\to0\).

Conversely, for every \(d\ge2\) and \(t>0\),
\[
\rho(t,d)\le
\left[
\frac{\log(1+d^2/t)}{L}
+\frac{\log(e+L)}{L}
\right]^2+\frac1{L^2}.
\]
For the nondense branch, this follows from
\[
e+L(d^2/t)\le(1+d^2/t)(e+L);
\]
for the dense branch, \(\rho(t,d)\le L^{-2}\).
Both correction terms tend to zero as \(d\to\infty\).

This bound also handles sequences with oscillating dimensions.
Given any positive tolerance, choose a dimension threshold so that both corrections are sufficiently small above that threshold.
Then \(g_k\to0\) makes the displayed upper bound small for all sufficiently large \(k\) whose dimensions exceed the threshold.
Below it, the elementary upper comparison gives
\[
\rho(t_k,d_k)\le\frac{4d_k^2}{t_k},
\]
which tends uniformly to zero as \(t_k\to\infty\).
Hence
\[
\rho(t_k,d_k)\to0
\quad\Longleftrightarrow\quad
t_k\to\infty\ \text{and}\ g_k\to0.
\]
% lean: CausalSmith.Stat.LdpOptvalueUniformFrontier.rho_consistency_iff

\item \textit{Parametric comparison and the quadratic-resource regime.}
Suppose that fixed positive finite constants
\(c_{\mathrm{cmp}}\le C_{\mathrm{cmp}}\) satisfy, eventually,
\[
\frac{c_{\mathrm{cmp}}}{t_k}
\le\rho(t_k,d_k)
\le\frac{C_{\mathrm{cmp}}}{t_k}.
\]
Since \(\rho(t_k,d_k)\le1\) and
\(t_k\rho(t_k,d_k)\ge\min\{t_k,d_k\}/4\), these imply
\[
t_k\ge c_{\mathrm{cmp}},
\qquad
\min\{t_k,d_k\}\le4C_{\mathrm{cmp}}
\]
eventually.

Conversely, suppose that there are \(a_{\mathrm{par}}>0\) and
\(M_{\mathrm{par}}\in\mathbb R\) such that, eventually,
\[
t_k\ge a_{\mathrm{par}},
\qquad
\min\{t_k,d_k\}\le M_{\mathrm{par}}.
\]
Define
\[
B_{\mathrm{par}}=\max\{1,M_{\mathrm{par}}\}.
\]
The elementary lower comparison gives
\[
\frac{\min\{1,a_{\mathrm{par}}\}}{t_k}
\le\min\{1,t_k^{-1}\}\le\rho(t_k,d_k).
\]
For the upper comparison, split according to \(t_k\le d_k\) or \(d_k\le t_k\).
In the first case \(t_k\le B_{\mathrm{par}}\), so
\(\rho(t_k,d_k)\le1\le4B_{\mathrm{par}}^2/t_k\).
In the second case \(d_k\le B_{\mathrm{par}}\), and
\[
\rho(t_k,d_k)\le\frac{4d_k^2}{t_k}
\le\frac{4B_{\mathrm{par}}^2}{t_k}.
\]
This proves the asserted general parametric equivalence.
When \(t_k\to\infty\), the positive lower-resource condition holds eventually. Moreover,
\[
\min\{t_k,d_k\}\le M_{\mathrm{par}}
\]
eventually forces \(d_k\le M_{\mathrm{par}}\) eventually, since
\(t_k>M_{\mathrm{par}}\) eventually. Enlarging to an integer bound proves the bounded-dimension characterization; its converse follows from
\(\min\{t_k,d_k\}\le d_k\).

For the quadratic-resource regime, fix
\(0<a_{\mathrm{el}}\le b_{\mathrm{el}}\).
Choose an integer \(D_0\) large enough that, whenever \(d\ge D_0\),
\[
L\ge b_{\mathrm{el}}+1,\qquad
\log L\ge2\log b_{\mathrm{el}},\qquad
\log L\ge\log(e+a_{\mathrm{el}}^{-1}).
\]
For every \(t\) in the specified band, define
\[
u_{\mathrm{el}}=\frac{t}{d^2}\in[a_{\mathrm{el}},b_{\mathrm{el}}].
\]
Then \(t<d^2L\), and
\[
\log(e+L/u_{\mathrm{el}})
\ge\log(L/b_{\mathrm{el}})
\ge\frac12\log L,
\]
\[
\log(e+L/u_{\mathrm{el}})
\le\log\!\left(L(e+a_{\mathrm{el}}^{-1})\right)
\le2\log L.
\]
Since \(0\le(\log L)/L\le1\), applying the cap and squaring gives
\[
\frac14\left(\frac{\log L}{L}\right)^2
\le\rho(t,d)
\le4\left(\frac{\log L}{L}\right)^2.
\]
These bounds supply the required constants uniformly over the band.
% lean: CausalSmith.Stat.LdpOptvalueUniformFrontier.rate_resource_conclusions

\item \textit{Transfer to the actual minimax values.}
The upper bounds already proved make the minimax values finite, so their real values satisfy the same four inequalities. For a fixed class and an admissible sequence, define
\[
\begin{aligned}
\mathcal R_{\mathcal C,k}^{\mathrm{TV}}
&=\mathfrak R^{\mathrm{TV}}_{\mathcal C}
(\mathcal P_{d_k}^{\mathrm{pair}};n_k,d_k,\varepsilon_k),\\
\mathcal H_{\mathcal C,k}^{\mathrm{TV}}
&=\mathfrak H^{\mathrm{TV}}_{\mathcal C}
(\mathcal P_{d_k}^{\mathrm{pair}};n_k,d_k,\varepsilon_k).
\end{aligned}
\]
The common constants give
\[
c\rho(t_k,d_k)\le\mathcal R_{\mathcal C,k}^{\mathrm{TV}}
\le C_0\rho(t_k,d_k),
\]
\[
c\sqrt{\rho(t_k,d_k)}
\le\mathcal H_{\mathcal C,k}^{\mathrm{TV}}
\le C_0\sqrt{\rho(t_k,d_k)}.
\]
Squeezing nonnegative sequences therefore yields
\[
\mathcal R_{\mathcal C,k}^{\mathrm{TV}}\to0
\Longleftrightarrow\rho(t_k,d_k)\to0,
\]
\[
\mathcal H_{\mathcal C,k}^{\mathrm{TV}}\to0
\Longleftrightarrow\sqrt{\rho(t_k,d_k)}\to0
\Longleftrightarrow\rho(t_k,d_k)\to0.
\]
Combining these equivalences with the rate consistency criterion proves both minimax consistency assertions.

Two-sided comparison is symmetric and transitive: reversing
\(c\rho\le\mathcal R\le C_0\rho\) gives
\(\mathcal R/C_0\le\rho\le\mathcal R/c\).
Consequently,
\[
\mathcal R_{\mathcal C,k}^{\mathrm{TV}}\asymp t_k^{-1}
\Longleftrightarrow
\rho(t_k,d_k)\asymp t_k^{-1}.
\]
The general and diverging-resource characterizations thus transfer directly.

Finally, in the band from the preceding step, the common bounds imply
\[
\frac c4\left(\frac{\log L}{L}\right)^2
\le
\mathfrak R^{\mathrm{TV}}_{\mathcal C}
(\mathcal P_d^{\mathrm{pair}};n,d,\varepsilon)
\le
4C_0\left(\frac{\log L}{L}\right)^2,
\]
and
\[
\frac c2\left|\frac{\log L}{L}\right|
\le
\mathfrak H^{\mathrm{TV}}_{\mathcal C}
(\mathcal P_d^{\mathrm{pair}};n,d,\varepsilon)
\le
2C_0\left|\frac{\log L}{L}\right|.
\]
To give the risk and length bounds the coupled constants in the statement, define
\[
c_{\mathrm{el}}=\min\{c/4,c^2/4\},
\qquad
C_{\mathrm{el}}=\max\{4C_0,4C_0^2\}.
\]
Then \(0<c_{\mathrm{el}}\le C_{\mathrm{el}}<\infty\), and
\[
\sqrt{c_{\mathrm{el}}}\le c/2,
\qquad
2C_0\le\sqrt{C_{\mathrm{el}}}.
\]
The displayed risk and length inequalities therefore imply exactly the final two bounds, with \(c_{\mathrm{el}},C_{\mathrm{el}}\) as the band-dependent constants. This completes all the sequence conclusions.
% lean: CausalSmith.Stat.LdpOptvalueUniformFrontier.value_sequence_conclusions_of_bounds
\end{enumerate}
\end{proof}

\begin{proof}[Proof of \cref{obj:thm:full-simplex-tv-converse}]
The lower bounds on the full simplex follow from the paired-family lower bounds and model inclusion. We record both steps to make the interval scope explicit.

For every \(\theta\in[-1/2,1/2]^d\),
\[
P_\theta(j,s)=\frac{1+s\theta_j}{2d}\ge0,
\qquad
\sum_{s\in\{-1,1\}}P_\theta(j,s)=\frac1d.
\]
Summing over \(j\) shows that \(P_\theta\) is a probability law on \(\mathcal A_d\), hence
\[
\mathcal P_d^{\mathrm{pair}}\subseteq\Delta_{2d}.
\]
This is the inclusion in \cref{obj:def:tv-corollary-models}. % lean: CausalSmith.Stat.LdpOptvalueUniformFrontier.pairedFamily_subset_simplex

Fix \(\mathcal C\in\{\mathrm{NI},\mathrm{SI}\}\). By the converse part of \cref{obj:thm:paired-uniform-tv-frontier}, there is a universal constant \(c>0\), common to both protocol classes and decision criteria, such that
\[
\mathfrak R^{\mathrm{TV}}_{\mathcal C}
 (\mathcal P_d^{\mathrm{pair}};n,d,\varepsilon)
\ge c\,r_{\mathcal C}(n,d,\varepsilon)
\]
and
\[
\mathfrak H^{\mathrm{TV}}_{\mathcal C}
 (\mathcal P_d^{\mathrm{pair}};n,d,\varepsilon)
\ge c\,h_{\mathcal C}(n,d,\varepsilon).
\]
These are the paired lower bounds established under the same cited Bernstein premise stated in both theorem environments.

For squared-error risk, any admissible protocol-estimator pair on \(\Delta_{2d}\) is also admissible on the paired subfamily, and
\[
\sup_{\mathsf p\in\Delta_{2d}}
 \mathbb E_{\mathsf p,\mathsf K}
 \bigl[(\widehat D-D_{\mathrm{TV}}(\mathsf p))^2\bigr]
\ge
\sup_{\mathsf p\in\mathcal P_d^{\mathrm{pair}}}
 \mathbb E_{\mathsf p,\mathsf K}
 \bigl[(\widehat D-D_{\mathrm{TV}}(\mathsf p))^2\bigr].
\]
Taking the infimum over the common protocol and estimator class gives
\[
\mathfrak R^{\mathrm{TV}}_{\mathcal C}
 (\Delta_{2d};n,d,\varepsilon)
\ge
\mathfrak R^{\mathrm{TV}}_{\mathcal C}
 (\mathcal P_d^{\mathrm{pair}};n,d,\varepsilon)
\ge c\,r_{\mathcal C}(n,d,\varepsilon).
\]

For interval length, a connected closed interval in \([0,1]\) that has coverage at least \(0.90\) at every law in \(\Delta_{2d}\) has the same coverage on \(\mathcal P_d^{\mathrm{pair}}\). Its Borel endpoints, connectedness, and range therefore satisfy the paired decision requirements. Moreover,
\[
\sup_{\mathsf p\in\Delta_{2d}}
 \mathbb E_{\mathsf p,\mathsf K}\operatorname{length}J_{\mathrm{TV}}
\ge
\sup_{\mathsf p\in\mathcal P_d^{\mathrm{pair}}}
 \mathbb E_{\mathsf p,\mathsf K}\operatorname{length}J_{\mathrm{TV}}.
\]
Taking the infimum over globally simplex-honest protocol-interval pairs yields
\[
\mathfrak H^{\mathrm{TV}}_{\mathcal C}
 (\Delta_{2d};n,d,\varepsilon)
\ge
\mathfrak H^{\mathrm{TV}}_{\mathcal C}
 (\mathcal P_d^{\mathrm{pair}};n,d,\varepsilon)
\ge c\,h_{\mathcal C}(n,d,\varepsilon).
\]
Both inequalities use the same universal positive constant. % lean: CausalSmith.Stat.LdpOptvalueUniformFrontier.full_simplex_tv_converse
\end{proof}

\begin{proof}[Proof of \cref{obj:lem:bounded-mean-concentration}]
\begin{enumerate}
\item
Since \(m\ge1\) and a probability space is nonempty, the pointwise bound for \(H_1\) gives
\[
-b\le H_1\le b
\quad\Longrightarrow\quad
-b\le b
\quad\Longrightarrow\quad
0\le b.
\]
If \(b=0\), both asserted bounds follow from
\[
\Pr\left\{m^{-1}\sum_{i=1}^m(H_i-\mathbb EH_i)\ge x\right\}\le1,
\qquad
\Pr\left\{m^{-1}\sum_{i=1}^m(H_i-\mathbb EH_i)\le-x\right\}\le1.
\]
Henceforth suppose \(b>0\), and define the variance proxy
\[
c_b:=\left(\frac{|b-(-b)|}{2}\right)^2=b^2>0.
\]
% lean: CausalSmith.Stat.LdpOptvalueUniformFrontier.boundedMeanConcentration_proved

\item
We first establish an upper-tail bound for any mutually independent, measurable real random variables \(H_1^\star,\ldots,H_m^\star\) on the given probability space satisfying
\[
-b\le H_i^\star\le b,
\qquad i=1,\ldots,m.
\]
Define their centered versions and interval endpoints by
\[
H_i^\circ:=H_i^\star-\mathbb EH_i^\star,
\qquad
\ell_i^\circ:=-b-\mathbb EH_i^\star,
\qquad
u_i^\circ:=b-\mathbb EH_i^\star.
\]
Boundedness ensures integrability, and therefore
\[
\mathbb EH_i^\circ=0,
\qquad
\ell_i^\circ\le H_i^\circ\le u_i^\circ,
\qquad
u_i^\circ-\ell_i^\circ=2b.
\]
The centered variables remain mutually independent because each is a measurable function of its corresponding original variable.

For each \(i\), define the moment-generating function and its logarithm on all of \(\mathbb R\) by
\[
\Phi_i(\upsilon):=\mathbb E\exp(\upsilon H_i^\circ),
\qquad
\psi_i(\upsilon):=\log\Phi_i(\upsilon),
\qquad \upsilon\in\mathbb R.
\]
Boundedness makes \(\Phi_i\) finite and strictly positive and permits differentiation under the expectation. Define expectation under the exponential tilt by
\[
\mathbb E_{\upsilon}^{(i)}[\cdot]
:=
\frac{\mathbb E[\cdot\,\exp(\upsilon H_i^\circ)]}{\Phi_i(\upsilon)},
\qquad
\bar h_{i,\upsilon}:=\mathbb E_{\upsilon}^{(i)}H_i^\circ.
\]
This is expectation under a probability measure absolutely continuous with respect to the original measure. Direct differentiation gives
\[
\psi_i(0)=0,
\qquad
\psi_i'(0)=\mathbb EH_i^\circ=0,
\qquad
\psi_i''(\upsilon)
=
\mathbb E_{\upsilon}^{(i)}[(H_i^\circ)^2]
-\bar h_{i,\upsilon}^{\,2}.
\]

The interval bound persists under the tilt. Consequently,
\[
\begin{aligned}
0
&\le
\mathbb E_{\upsilon}^{(i)}
\bigl[(u_i^\circ-H_i^\circ)(H_i^\circ-\ell_i^\circ)\bigr]\\
&=
-\mathbb E_{\upsilon}^{(i)}[(H_i^\circ)^2]
+(u_i^\circ+\ell_i^\circ)\bar h_{i,\upsilon}
-u_i^\circ\ell_i^\circ.
\end{aligned}
\]
Subtracting the squared tilted mean yields
\[
\begin{aligned}
\psi_i''(\upsilon)
&\le
(u_i^\circ-\bar h_{i,\upsilon})
(\bar h_{i,\upsilon}-\ell_i^\circ)\\
&=
\left(\frac{u_i^\circ-\ell_i^\circ}{2}\right)^2
-
\left(\bar h_{i,\upsilon}
-\frac{u_i^\circ+\ell_i^\circ}{2}\right)^2
\le c_b.
\end{aligned}
\]
For \(\upsilon>0\), Taylor's theorem supplies an intermediate point
\[
\upsilon_{i,\mathrm{mid}}\in(0,\upsilon),
\qquad
\psi_i(\upsilon)
=
\frac{\upsilon^2}{2}\psi_i''(\upsilon_{i,\mathrm{mid}})
\le\frac{c_b\upsilon^2}{2}.
\]
Thus \(\Phi_i(\upsilon)\le\exp(c_b\upsilon^2/2)\) for positive \(\upsilon\); at zero both sides equal one. For negative \(\upsilon\), apply the same positive-parameter calculation to the centered variable \(-H_i^\circ\), whose interval has the same width:
\[
-u_i^\circ\le-H_i^\circ\le-\ell_i^\circ,
\qquad
\mathbb E[-H_i^\circ]=0.
\]
It follows that
\[
\Phi_i(\upsilon)
=
\mathbb E\exp\bigl((-\upsilon)(-H_i^\circ)\bigr)
\le
\exp\left(\frac{c_b(-\upsilon)^2}{2}\right).
\]
We have therefore established
\[
\mathbb E\exp(\upsilon H_i^\circ)
\le\exp\left(\frac{c_b\upsilon^2}{2}\right),
\qquad \upsilon\in\mathbb R,\quad i=1,\ldots,m.
\]
% lean: ProbabilityTheory.hasSubgaussianMGF_of_mem_Icc

\item
Define the centered sum
\[
S^\star:=\sum_{i=1}^m H_i^\circ.
\]
Independence gives the factorization
\[
\begin{aligned}
\mathbb E\exp(\upsilon S^\star)
&=\prod_{i=1}^m\mathbb E\exp(\upsilon H_i^\circ)\\
&\le
\prod_{i=1}^m\exp\left(\frac{c_b\upsilon^2}{2}\right)
=
\exp\left(\frac{mc_b\upsilon^2}{2}\right),
\qquad \upsilon\in\mathbb R.
\end{aligned}
\]
Indeed, adjoining one centered variable to a partial sum multiplies their moment-generating functions, since that variable is independent of the preceding sum; induction gives the displayed formula.

Since \(mx\ge0\), for every \(\upsilon\ge0\) exponential Markov's inequality gives
\[
\begin{aligned}
\Pr\{S^\star\ge mx\}
&\le
\exp(-\upsilon mx)\,\mathbb E\exp(\upsilon S^\star)\\
&\le
\exp\left(-\upsilon mx+\frac{mc_b\upsilon^2}{2}\right).
\end{aligned}
\]
Choose
\[
\upsilon_{\mathrm{opt}}:=\frac{mx}{mc_b}=\frac{x}{c_b}\ge0.
\]
This choice is defined also when \(x=0\), and substitution yields
\[
\Pr\{S^\star\ge mx\}
\le
\exp\left(-\frac{(mx)^2}{2mc_b}\right).
\]
% lean: ProbabilityTheory.HasSubgaussianMGF.measure_sum_ge_le_of_iIndepFun

\item
Because \(m>0\), multiplying by \(m\) gives the event identity
\[
\left\{
x\le m^{-1}\sum_{i=1}^m
(H_i^\star-\mathbb EH_i^\star)
\right\}
=
\{mx\le S^\star\}.
\]
Using \(c_b=b^2\), the exponent simplifies to
\[
-\frac{(mx)^2}{2mc_b}
=
-\frac{mx^2}{2b^2}.
\]
Hence the upper-tail calculation establishes
\[
\Pr\left\{
m^{-1}\sum_{i=1}^m(H_i^\star-\mathbb EH_i^\star)\ge x
\right\}
\le
\exp\left(-\frac{mx^2}{2b^2}\right).
\]
Taking \(H_i^\star=H_i\) proves the asserted upper-tail bound.
% lean: CausalSmith.Stat.LdpOptvalueUniformFrontier.boundedMeanConcentration_proved.upper

\item
For the lower tail, define
\[
H_i^-:=-H_i,
\qquad i=1,\ldots,m.
\]
These are mutually independent and measurable, and negating the original bounds gives
\[
-b\le H_i^-\le b.
\]
Linearity of expectation and finite summation gives
\[
\mathbb EH_i^-=-\mathbb EH_i,
\qquad
\sum_{i=1}^m(H_i^--\mathbb EH_i^-)
=
-\sum_{i=1}^m(H_i-\mathbb EH_i).
\]
Therefore,
\[
\left\{
m^{-1}\sum_{i=1}^m(H_i-\mathbb EH_i)\le-x
\right\}
=
\left\{
m^{-1}\sum_{i=1}^m(H_i^--\mathbb EH_i^-)\ge x
\right\}.
\]
Applying the upper-tail bound just established to \(H_1^-,\ldots,H_m^-\) yields
\[
\Pr\left\{
m^{-1}\sum_{i=1}^m(H_i-\mathbb EH_i)\le-x
\right\}
\le
\exp\left(-\frac{mx^2}{2b^2}\right),
\]
which completes the proof.
% lean: CausalSmith.Stat.LdpOptvalueUniformFrontier.boundedMeanConcentration_proved
\end{enumerate}
\end{proof}

\begin{proof}[Proof of \cref{obj:lem:cai-low-chebyshev-approximation}]
\begin{enumerate}
\item Fix an even integer \(D\ge2\), and set
\[
\ell_{\mathrm{tr}}:=D/2\in\mathbb N,\qquad
\ell_{\mathrm{tr}}\ge1,\qquad D=2\ell_{\mathrm{tr}}.
\]
For the truncation argument, extend the polynomial family to degree zero by
\[
q_0(x):=\frac2\pi.
\]
Thus, for every \(\ell\in\mathbb N_0\),
\[
q_{2\ell}(x)
=\frac2\pi+\frac4\pi\sum_{\ell'=1}^{\ell}
 \frac{(-1)^{\ell'+1}\mathcal T_{2\ell'}(x)}
 {4(\ell')^2-1},
\qquad q_D=q_{2\ell_{\mathrm{tr}}}.
\]
This is precisely the one-based truncation of the even Chebyshev modes; the sum is empty when \(\ell=0\).
% lean: CausalSmith.Stat.LdpOptvalueUniformFrontier.chebyshevAbsPoly_eq_library

\item Define the continuous, \(\pi\)-periodic function \(\varphi_{\sin}:\mathbb R\to\mathbb R\) by
\[
\varphi_{\sin}(\vartheta)=\sin\vartheta
\quad(0\le\vartheta\le\pi),\qquad
\varphi_{\sin}(\vartheta+\pi)=\varphi_{\sin}(\vartheta)
\quad(\vartheta\in\mathbb R).
\]
The endpoint values agree, so the periodic extension is continuous. With \(\mathrm i:=\sqrt{-1}\), define its Fourier coefficients by
\[
\widehat\varphi_{\sin}(\ell)
:=\frac1\pi\int_0^\pi
 \sin\vartheta\,e^{-2\mathrm i\ell\vartheta}\,d\vartheta,
\qquad \ell\in\mathbb Z.
\]
For each such \(\ell\), \(1-4\ell^2\ne0\), and the function
\[
\Phi_\ell(\vartheta)
:=e^{-2\mathrm i\ell\vartheta}
 \frac{-2\mathrm i\ell\sin\vartheta-\cos\vartheta}
 {1-4\ell^2},
\qquad \vartheta\in\mathbb R,
\]
satisfies
\[
\Phi_\ell'(\vartheta)
=e^{-2\mathrm i\ell\vartheta}\sin\vartheta.
\]
Since \(e^{-2\mathrm i\ell\pi}=1\), evaluation at the endpoints gives
\[
\widehat\varphi_{\sin}(\ell)
=\frac{\Phi_\ell(\pi)-\Phi_\ell(0)}{\pi}
=\frac{2}{\pi(1-4\ell^2)}
=-\frac{2}{\pi(4\ell^2-1)}.
\]
In particular, \(\widehat\varphi_{\sin}(0)=2/\pi\).
% lean: Causalean.Mathlib.Analysis.AbsoluteValueMomentPriorDuality.fourierCoeff_sinCircle

\item For every integer \(\ell\ge1\),
\[
\bigl|\widehat\varphi_{\sin}(\ell)\bigr|
=\frac{2}{\pi(4\ell^2-1)}
\le\frac1{\ell^2}.
\]
Indeed, \(\ell^2\ge1\) and \(\pi\ge2\) imply
\[
2\ell^2\le\pi(4\ell^2-1).
\]
The coefficients at \(\ell\) and \(-\ell\) agree. Comparison with the convergent series \(\sum_{\ell=1}^\infty\ell^{-2}\) therefore yields
\[
\sum_{\ell\in\mathbb Z}
 \bigl|\widehat\varphi_{\sin}(\ell)\bigr|<\infty.
\]
The Fourier reconstruction theorem for a continuous periodic function with absolutely summable coefficients now gives
\[
\varphi_{\sin}(\vartheta)
=\sum_{\ell\in\mathbb Z}
 \widehat\varphi_{\sin}(\ell)e^{2\mathrm i\ell\vartheta},
\qquad \vartheta\in\mathbb R.
\]
% lean: Causalean.Mathlib.Analysis.AbsoluteValueMomentPriorDuality.summable_fourierCoeff_sinCircle

\item For \(0\le\vartheta\le\pi\), periodicity and the sine addition identities give
\[
\varphi_{\sin}(\vartheta+\pi/2)
=
\begin{cases}
\cos\vartheta,&0\le\vartheta\le\pi/2,\\
-\cos\vartheta,&\pi/2<\vartheta\le\pi,
\end{cases}
=|\cos\vartheta|.
\]
The defining Chebyshev recurrence and the cosine addition formula also give
\[
\mathcal T_\ell(\cos\vartheta)=\cos(\ell\vartheta),
\qquad \ell\in\mathbb N_0,\quad \vartheta\in\mathbb R.
\]
Consequently, for \(\ell\ge1\), the paired Fourier terms at the shifted argument satisfy
\[
\begin{aligned}
&\widehat\varphi_{\sin}(\ell)
 e^{2\mathrm i\ell(\vartheta+\pi/2)}
+\widehat\varphi_{\sin}(-\ell)
 e^{-2\mathrm i\ell(\vartheta+\pi/2)}
\\
&\qquad
=-\frac{4\cos(2\ell\vartheta+\ell\pi)}
 {\pi(4\ell^2-1)}
=\frac4\pi
 \frac{(-1)^{\ell+1}\mathcal T_{2\ell}(\cos\vartheta)}
 {4\ell^2-1}.
\end{aligned}
\]
For \(x\in[-1,1]\), define
\[
\vartheta_x:=\arccos x\in[0,\pi],
\qquad \cos\vartheta_x=x.
\]
Pairing the positive and negative frequencies in the absolutely convergent Fourier series and taking real parts therefore yields
\[
|x|=\frac2\pi+\frac4\pi\sum_{\ell=1}^{\infty}
 \frac{(-1)^{\ell+1}\mathcal T_{2\ell}(x)}
 {4\ell^2-1}.
\]
% lean: Causalean.Mathlib.Analysis.AbsoluteValueMomentPriorDuality.hasSum_absChebTerm_real

\item Splitting this series at \(\ell_{\mathrm{tr}}\) gives
\[
|x|-q_D(x)
=\frac4\pi\sum_{\ell=\ell_{\mathrm{tr}}+1}^{\infty}
 \frac{(-1)^{\ell+1}\mathcal T_{2\ell}(x)}
 {4\ell^2-1}.
\]
The scalar majorant has an exact telescoping sum, because
\[
\frac1{4\ell^2-1}
=\frac12\left(\frac1{2\ell-1}-\frac1{2\ell+1}\right)
\qquad(\ell\ge1).
\]
In particular, for every integer \(\ell_{\mathrm{end}}\ge\ell_{\mathrm{tr}}\),
\[
\sum_{\ell=\ell_{\mathrm{tr}}+1}^{\ell_{\mathrm{end}}}
 \frac1{4\ell^2-1}
=\frac12\left(
 \frac1{2\ell_{\mathrm{tr}}+1}
 -\frac1{2\ell_{\mathrm{end}}+1}\right),
\]
including the empty sum when \(\ell_{\mathrm{end}}=\ell_{\mathrm{tr}}\). Passing to the limit gives
\[
\sum_{\ell=\ell_{\mathrm{tr}}+1}^{\infty}
 \frac1{4\ell^2-1}
=\frac1{2(2\ell_{\mathrm{tr}}+1)}.
\]
Moreover, the cosine identity implies
\[
|\mathcal T_{2\ell}(x)|
=|\cos(2\ell\vartheta_x)|\le1.
\]
Thus the residual series is absolutely convergent, and
\[
\begin{aligned}
\bigl|q_D(x)-|x|\bigr|
&\le\frac4\pi
 \sum_{\ell=\ell_{\mathrm{tr}}+1}^{\infty}
 \frac{|\mathcal T_{2\ell}(x)|}{4\ell^2-1}\\
&\le\frac4\pi
 \sum_{\ell=\ell_{\mathrm{tr}}+1}^{\infty}
 \frac1{4\ell^2-1}
=\frac2{\pi(2\ell_{\mathrm{tr}}+1)}
=\frac2{\pi(D+1)}.
\end{aligned}
\]
% lean: Causalean.Mathlib.Analysis.AbsoluteValueMomentPriorDuality.absChebPoly_error_le

\item For real polynomials
\[
\mathcal B(x)=\sum_{\ell\ge0}\beta_\ell x^\ell,
\qquad
\mathcal E(x)=\sum_{\ell\ge0}\gamma_\ell x^\ell,
\]
where the real coefficient sequences have finite support, define
\[
\|\mathcal B\|_{\mathrm{coef},1}
:=\sum_{\ell\ge0}|\beta_\ell|,
\qquad
\|\mathcal E\|_{\mathrm{coef},1}
:=\sum_{\ell\ge0}|\gamma_\ell|.
\]
Coefficientwise triangle inequalities give
\[
\|\mathcal B\pm\mathcal E\|_{\mathrm{coef},1}
\le
\|\mathcal B\|_{\mathrm{coef},1}
+\|\mathcal E\|_{\mathrm{coef},1}.
\]
For every real scalar \(\zeta\), the same definition gives
\[
\|0\|_{\mathrm{coef},1}=0,\qquad
\|\zeta\|_{\mathrm{coef},1}=|\zeta|,\qquad
\|\zeta\mathcal B\|_{\mathrm{coef},1}
=|\zeta|\|\mathcal B\|_{\mathrm{coef},1},\qquad
\|x\mathcal B\|_{\mathrm{coef},1}
=\|\mathcal B\|_{\mathrm{coef},1}.
\]
Here the last identity follows by shifting the coefficient indices. Finally, expanding the product into monomials and applying the finite-sum triangle inequality yields
\[
\begin{aligned}
\mathcal B(x)\mathcal E(x)
&=\sum_{\ell,\ell'\ge0}
 \beta_\ell\gamma_{\ell'}x^{\ell+\ell'},\\
\|\mathcal B\mathcal E\|_{\mathrm{coef},1}
&\le\sum_{\ell,\ell'\ge0}
 |\beta_\ell\gamma_{\ell'}|
=\|\mathcal B\|_{\mathrm{coef},1}
 \|\mathcal E\|_{\mathrm{coef},1}.
\end{aligned}
\]
All sums here have finitely many nonzero terms.
% lean: Causalean.Mathlib.Analysis.Approximation.Chebyshev.polynomialCoeffOneNorm_mul_le

\item Define
\[
\chi:=1+\sqrt2>0,\qquad \chi^2=2\chi+1.
\]
We claim that
\[
\|\mathcal T_\ell\|_{\mathrm{coef},1}\le\chi^\ell
\qquad(\ell\in\mathbb N_0).
\]
The two initial cases are
\[
\|\mathcal T_0\|_{\mathrm{coef},1}=1=\chi^0,
\qquad
\|\mathcal T_1\|_{\mathrm{coef},1}=1\le\chi.
\]
For the induction step, the Chebyshev recurrence and the norm identities just established give
\[
\begin{aligned}
\|\mathcal T_{\ell+2}\|_{\mathrm{coef},1}
&=\|2x\mathcal T_{\ell+1}-\mathcal T_\ell\|_{\mathrm{coef},1}\\
&\le2\|\mathcal T_{\ell+1}\|_{\mathrm{coef},1}
 +\|\mathcal T_\ell\|_{\mathrm{coef},1}\\
&\le2\chi^{\ell+1}+\chi^\ell
=\chi^\ell(2\chi+1)
=\chi^{\ell+2}.
\end{aligned}
\]
This proves the claim by two-step induction.
% lean: Causalean.Mathlib.Analysis.Approximation.Chebyshev.polynomialCoeffOneNorm_chebyshev_le

\item For every \(\ell\in\mathbb N_0\),
\[
4(\ell+1)^2-1\ge3,\qquad
\frac4\pi\le\frac43,
\]
where the second inequality uses \(\pi>3\). Hence
\[
\left|
 \frac{(4/\pi)(-1)^\ell}{4(\ell+1)^2-1}
\right|
=\frac{4/\pi}{4(\ell+1)^2-1}
\le\frac12.
\]
Also, \(2\le(3/2)^2\) implies \(\sqrt2\le3/2\), and therefore
\[
\chi^2=3+2\sqrt2\le6.
\]
The product bound and the Chebyshev norm estimate now imply
\[
\begin{aligned}
\left\|
 \frac{(4/\pi)(-1)^\ell}{4(\ell+1)^2-1}
 \mathcal T_{2(\ell+1)}
\right\|_{\mathrm{coef},1}
&\le\frac12\chi^{2(\ell+1)}\\
&\le\frac12\,6^{\ell+1}
=3\cdot6^\ell
\le3\cdot8^\ell.
\end{aligned}
\]
% lean: Causalean.Mathlib.Analysis.Approximation.Chebyshev.AbsoluteValue.absChebPoly_coeffOneNorm_le

\item We next prove, by induction on \(\ell\in\mathbb N_0\), that
\[
\|q_{2\ell}\|_{\mathrm{coef},1}\le8^\ell.
\]
At \(\ell=0\),
\[
\|q_0\|_{\mathrm{coef},1}=\frac2\pi\le1=8^0.
\]
The truncations satisfy
\[
q_{2(\ell+1)}
=q_{2\ell}
+\frac{(4/\pi)(-1)^\ell}{4(\ell+1)^2-1}
 \mathcal T_{2(\ell+1)}.
\]
Thus the induction hypothesis and the preceding mode bound give
\[
\begin{aligned}
\|q_{2(\ell+1)}\|_{\mathrm{coef},1}
&\le\|q_{2\ell}\|_{\mathrm{coef},1}
+\left\|
 \frac{(4/\pi)(-1)^\ell}{4(\ell+1)^2-1}
 \mathcal T_{2(\ell+1)}
 \right\|_{\mathrm{coef},1}\\
&\le8^\ell+3\cdot8^\ell
\le8^{\ell+1}.
\end{aligned}
\]
In particular,
\[
\|q_D\|_{\mathrm{coef},1}
\le8^{\ell_{\mathrm{tr}}}
=2^{3\ell_{\mathrm{tr}}}
=2^{3D/2}.
\]
% lean: Causalean.Mathlib.Analysis.Approximation.Chebyshev.AbsoluteValue.absChebPoly_coeffOneNorm_le

\item Fix an integer \(v\) with \(0\le v\le D\). If \(a_{Dv}=0\), then
\[
|a_{Dv}|=0\le2^{3D/2}.
\]
Otherwise, its absolute value is one of the nonnegative summands in the coefficient norm. Using the stated monomial expansion of \(q_D\), we obtain
\[
|a_{Dv}|
\le\sum_{\ell=0}^{D}|a_{D\ell}|
=\|q_D\|_{\mathrm{coef},1}
\le2^{3D/2}.
\]
Together with the approximation bound, this proves both conclusions.
% lean: Causalean.Mathlib.Analysis.Approximation.Chebyshev.AbsoluteValue.absChebPoly_halfDegree_coeff_le
\end{enumerate}
\end{proof}

\begin{proof}[Proof of \cref{obj:lem:cai-low-moment-duality}]
Fix a positive even integer \(K\).

\begin{enumerate}
\item Define the approximation space and target function by
\[
I_{\mathrm{appr}}=[-1,1],\qquad
\mathcal C_{\mathrm{appr}}=C(I_{\mathrm{appr}},\mathbb R),\qquad
f_{\mathrm{abs}}(x)=|x|,
\]
and equip \(\mathcal C_{\mathrm{appr}}\) with the uniform norm. Let
\[
\mathcal E_K
=\{p|_{I_{\mathrm{appr}}}:p\in\mathbb R[x],\ \deg p\le K\}.
\]
For every real polynomial \(p\),
\[
\|f_{\mathrm{abs}}-p|_{I_{\mathrm{appr}}}\|_\infty
=\sup_{|x|\le1}\bigl||x|-p(x)\bigr|.
\]
Thus \(e_K\) is the distance from \(f_{\mathrm{abs}}\) to \(\mathcal E_K\).

We first verify attainment of this distance. Define
\[
R_{\mathrm{ball}}=2\|f_{\mathrm{abs}}\|_\infty+1,\qquad
\mathcal B_K
=\{h\in\mathcal E_K:\|h\|_\infty\le R_{\mathrm{ball}}\}.
\]
The space \(\mathcal E_K\) is finite dimensional, so \(\mathcal B_K\) is compact. The continuous function
\(h\mapsto\|f_{\mathrm{abs}}-h\|_\infty\) therefore has a minimizer
\(h_{\mathrm{min}}\in\mathcal B_K\). Since \(0\in\mathcal B_K\),
\[
\|f_{\mathrm{abs}}-h_{\mathrm{min}}\|_\infty
\le\|f_{\mathrm{abs}}\|_\infty.
\]
For \(h\in\mathcal E_K\setminus\mathcal B_K\), the reverse triangle inequality gives
\[
\|f_{\mathrm{abs}}-h\|_\infty
\ge\|h\|_\infty-\|f_{\mathrm{abs}}\|_\infty
>\|f_{\mathrm{abs}}\|_\infty+1.
\]
Consequently \(h_{\mathrm{min}}\) minimizes over all of \(\mathcal E_K\), and
\[
\|f_{\mathrm{abs}}-h_{\mathrm{min}}\|_\infty=e_K.
\]
% lean: Causalean.Mathlib.Analysis.AbsoluteValueMomentPriorDuality.exists_bestPolynomialAbs

\item Let the quotient map and the target's quotient class be
\[
\pi_{\mathrm{quot}}:\mathcal C_{\mathrm{appr}}
\longrightarrow\mathcal C_{\mathrm{appr}}/\mathcal E_K,
\qquad
\bar f_{\mathrm{abs}}=\pi_{\mathrm{quot}}(f_{\mathrm{abs}}).
\]
The quotient norm is the distance to the subspace. The minimizer above gives its upper bound, while the definition of \(e_K\) gives its lower bound:
\[
\|\bar f_{\mathrm{abs}}\|
=\inf_{h\in\mathcal E_K}\|f_{\mathrm{abs}}-h\|_\infty
=e_K.
\]
By the norming-functional form of Hahn--Banach, there is a continuous real linear functional
\[
\gamma_K:\mathcal C_{\mathrm{appr}}/\mathcal E_K\longrightarrow\mathbb R,
\qquad
\|\gamma_K\|\le1,\qquad
\gamma_K(\bar f_{\mathrm{abs}})=\|\bar f_{\mathrm{abs}}\|.
\]
This also covers a zero quotient class, for which the zero functional suffices. Define
\[
\Lambda_K=\gamma_K\circ\pi_{\mathrm{quot}}.
\]
Since the quotient map is contractive and vanishes on \(\mathcal E_K\),
\[
|\Lambda_K(h)|\le\|h\|_\infty
\quad(h\in\mathcal C_{\mathrm{appr}}),\qquad
\Lambda_K(1)=0,\qquad
\Lambda_K(f_{\mathrm{abs}})=e_K.
\]
In particular, with
\[
h_\ell(x)=x^\ell
\qquad(x\in I_{\mathrm{appr}},\ 0\le\ell\le K),
\]
we have
\[
\Lambda_K(h_\ell)=0.
\]
% lean: Causalean.Mathlib.Analysis.AbsoluteValueMomentPriorDuality.exists_dual_abs_certificate

\item We construct two positive linear functionals of mass one half whose difference is \(\Lambda_K\). On
\(\mathcal C_{\mathrm{appr}}\times\mathbb R\), define the subspace, its linear functional, and a cone by
\[
\begin{aligned}
\mathcal D_K
&=\{(c_{\mathrm{const}}1,s_{\mathrm{aux}}):
       c_{\mathrm{const}},s_{\mathrm{aux}}\in\mathbb R\},\\
\beta_K(c_{\mathrm{const}}1,s_{\mathrm{aux}})
&=\frac{c_{\mathrm{const}}}{2}+s_{\mathrm{aux}},\\
\mathcal S_K
&=\{(h_++h_-,-\Lambda_K(h_-)):
       h_+,h_-\in\mathcal C_{\mathrm{appr}},\
       h_+\ge0,\ h_-\ge0\}.
\end{aligned}
\]
The representation of an element of \(\mathcal D_K\) is unique because
\(I_{\mathrm{appr}}\) is nonempty. The set \(\mathcal S_K\) is closed under addition and multiplication by nonnegative scalars.

Suppose
\[
(c_{\mathrm{const}}1,s_{\mathrm{aux}})
=(h_++h_-,-\Lambda_K(h_-))\in\mathcal D_K\cap\mathcal S_K.
\]
Then
\[
c_{\mathrm{const}}\ge0,\qquad
0\le h_-(x)\le c_{\mathrm{const}},
\]
so the function
\[
h_{\mathrm{cent}}=2h_--c_{\mathrm{const}}1
\]
satisfies
\[
\|h_{\mathrm{cent}}\|_\infty\le c_{\mathrm{const}},\qquad
2\Lambda_K(h_-)
=\Lambda_K(h_{\mathrm{cent}})
\le|\Lambda_K(h_{\mathrm{cent}})|
\le c_{\mathrm{const}}.
\]
It follows that
\[
\beta_K(c_{\mathrm{const}}1,s_{\mathrm{aux}})
=\frac{c_{\mathrm{const}}}{2}-\Lambda_K(h_-)\ge0.
\]
Moreover, for every
\((h_{\mathrm{test}},s_{\mathrm{test}})
\in\mathcal C_{\mathrm{appr}}\times\mathbb R\),
\[
(\|h_{\mathrm{test}}\|_\infty1,-s_{\mathrm{test}})
\in\mathcal D_K,
\]
and
\[
(\|h_{\mathrm{test}}\|_\infty1,-s_{\mathrm{test}})
+(h_{\mathrm{test}},s_{\mathrm{test}})
=(\|h_{\mathrm{test}}\|_\infty1+h_{\mathrm{test}},0)
\in\mathcal S_K,
\]
because the first coordinate is nonnegative. The ordered linear extension theorem therefore supplies a linear functional
\[
\Gamma_K:\mathcal C_{\mathrm{appr}}\times\mathbb R\longrightarrow\mathbb R,
\qquad
\Gamma_K|_{\mathcal D_K}=\beta_K,\qquad
\Gamma_K\ge0\ \text{on }\mathcal S_K.
\]
Define, for \(h\in\mathcal C_{\mathrm{appr}}\),
\[
\Phi_+(h)=\Gamma_K(h,0),\qquad
\Phi_-(h)=\Phi_+(h)-\Lambda_K(h).
\]
For \(h\ge0\), both \((h,0)\) and \((h,-\Lambda_K(h))\) belong to \(\mathcal S_K\). Since \(\Gamma_K(0,s)=s\) for every real \(s\),
\[
\Phi_+(h)\ge0,\qquad
\Phi_-(h)=\Gamma_K(h,-\Lambda_K(h))\ge0.
\]
Finally,
\[
\Phi_+(1)=\frac12,\qquad
\Phi_-(1)=\frac12,\qquad
\Phi_+(h)-\Phi_-(h)=\Lambda_K(h).
\]
% lean: Causalean.Mathlib.Analysis.AbsoluteValueMomentPriorDuality.exists_positive_half_split

\item Riesz representation on the compact interval gives finite positive Borel measures \(\mu_+,\mu_-\) on \(I_{\mathrm{appr}}\), characterized by
\[
\int_{I_{\mathrm{appr}}}h\,d\mu_+=\Phi_+(h),\qquad
\int_{I_{\mathrm{appr}}}h\,d\mu_-=\Phi_-(h)
\quad(h\in\mathcal C_{\mathrm{appr}}).
\]
Applying these identities to the constant function yields
\[
\mu_+(I_{\mathrm{appr}})=\mu_-(I_{\mathrm{appr}})=\frac12.
\]
Define the inclusion and reflection maps, and the symmetrized measures on \(\mathbb R\), by
\[
\iota:I_{\mathrm{appr}}\longrightarrow\mathbb R,\quad \iota(x)=x,
\qquad
r_{\mathrm{refl}}:\mathbb R\longrightarrow\mathbb R,\quad
r_{\mathrm{refl}}(x)=-x,
\]
\[
\zeta_\pm
=\frac12\left(\iota_\#\mu_\pm
 +(r_{\mathrm{refl}}\circ\iota)_\#\mu_\pm\right).
\]
Here \(\#\) denotes pushforward. Both pushforwards have mass one half and are supported on \(I_{\mathrm{appr}}\); reflection interchanges them. Hence
\[
\zeta_\pm(\mathbb R)=\frac12,\qquad
\zeta_\pm(\mathbb R\setminus I_{\mathrm{appr}})=0,\qquad
(r_{\mathrm{refl}})_\#\zeta_\pm=\zeta_\pm.
\]
For every continuous \(h_{\mathrm{test}}:\mathbb R\to\mathbb R\), the pushforward integral formula gives
\[
\int_{\mathbb R}h_{\mathrm{test}}\,d\zeta_\pm
=\frac12\left(
\int_{I_{\mathrm{appr}}}h_{\mathrm{test}}(x)\,d\mu_\pm(x)
+\int_{I_{\mathrm{appr}}}h_{\mathrm{test}}(-x)\,d\mu_\pm(x)
\right).
\]
All these integrals exist because the measures are finite and the functions are bounded on the relevant compact interval.
% lean: Causalean.Mathlib.Analysis.AbsoluteValueMomentPriorDuality.symmPush_mass
% lean: Causalean.Mathlib.Analysis.AbsoluteValueMomentPriorDuality.symmPush_supported
% lean: Causalean.Mathlib.Analysis.AbsoluteValueMomentPriorDuality.symmPush_symmetric
% lean: Causalean.Mathlib.Analysis.AbsoluteValueMomentPriorDuality.integral_symmPush

\item For \(0\le\ell\le K\), define the reflected monomial by
\[
h_{\ell,\mathrm{refl}}(x)=(-x)^\ell=(-1)^\ell h_\ell(x)
\qquad(x\in I_{\mathrm{appr}}).
\]
Linearity and polynomial annihilation give
\[
\Lambda_K(h_\ell)=0,\qquad
\Lambda_K(h_{\ell,\mathrm{refl}})
=(-1)^\ell\Lambda_K(h_\ell)=0.
\]
Combining the integral formula with the difference of the representing functionals therefore yields
\[
\begin{aligned}
\int_{\mathbb R}x^\ell\,d\zeta_+
-\int_{\mathbb R}x^\ell\,d\zeta_-
&=\frac12\left(
\Lambda_K(h_\ell)+\Lambda_K(h_{\ell,\mathrm{refl}})
\right)\\
&=0.
\end{aligned}
\]
For the absolute-value function, reflection preserves the function itself:
\[
|-x|=|x|.
\]
Consequently,
\[
\int_{\mathbb R}|x|\,d\zeta_+
-\int_{\mathbb R}|x|\,d\zeta_-
=\frac12\left(
\Lambda_K(f_{\mathrm{abs}})+\Lambda_K(f_{\mathrm{abs}})
\right)
=e_K.
\]
% lean: Causalean.Mathlib.Analysis.AbsoluteValueMomentPriorDuality.exists_absExtremalDecomposition

\item Normalize the two measures by defining
\[
\xi_0=2\zeta_-,\qquad \xi_1=2\zeta_+.
\]
Their masses are one, and multiplication by two preserves their support and reflection invariance:
\[
\xi_0(\mathbb R)=\xi_1(\mathbb R)=1,\qquad
\xi_0(\mathbb R\setminus[-1,1])
=\xi_1(\mathbb R\setminus[-1,1])=0,
\]
\[
(r_{\mathrm{refl}})_\#\xi_0=\xi_0,\qquad
(r_{\mathrm{refl}})_\#\xi_1=\xi_1.
\]
Scaling the moment identities gives, for every integer \(0\le\ell\le K\),
\[
\int x^\ell\,d\xi_1(x)-\int x^\ell\,d\xi_0(x)
=2\left(\int x^\ell\,d\zeta_+(x)-\int x^\ell\,d\zeta_-(x)\right)
=0.
\]
Likewise, scaling the absolute-moment identity gives
\[
\int |x|\,d\xi_1(x)-\int |x|\,d\xi_0(x)=2e_K.
\]
The approximation error used throughout agrees with the stated one because
\[
\inf_{\deg p\le K}\sup_{|x|\le1}\bigl||x|-p(x)\bigr|
=
\inf_{\deg p\le K}\sup_{|x|\le1}\bigl|p(x)-|x|\bigr|
=e_K.
\]
These are the required measures and identities.
% lean: Causalean.Mathlib.Analysis.AbsoluteValueMomentPriorDuality.AbsExtremalDecomposition.toMomentMatchedPriors
% lean: CausalSmith.Stat.LdpOptvalueUniformFrontier.bestAbsApproxError_eq_library
% lean: CausalSmith.Stat.LdpOptvalueUniformFrontier.caiLowMomentDuality_proved
\end{enumerate}
\end{proof}

\begin{proof}[Proof of \cref{obj:lem:measurable-kernel-radon-nikodym}]
\begin{enumerate}
\item Use the canonical kernel Radon--Nikodym derivative to define
\[
f:H\times Z\longrightarrow[0,\infty],
\qquad
f(h,z):=\frac{dK}{dL}(h,z).
\]
The kernel Radon--Nikodym construction gives joint measurability of this function for finite kernels with standard-Borel target. % lean: ProbabilityTheory.Kernel.measurable_rnDeriv

\item Fix \(h\in H\) and a measurable set \(E\subseteq Z\). The canonical kernel derivative agrees with the ordinary measure derivative:
\[
f(h,z)
=
\frac{dK(h,\cdot)}{dL(h,\cdot)}(z)
\qquad\text{for \(L(h,\cdot)\)-almost every \(z\)}.
\]
Consequently,
\[
\int_E f(h,z)\,L(h,dz)
=
\int_E
\frac{dK(h,\cdot)}{dL(h,\cdot)}(z)\,L(h,dz)
=
K(h,E).
\]
The first equality uses the almost-everywhere agreement, and the second is the measure Radon--Nikodym identity, applicable because \(K(h,\cdot)\ll L(h,\cdot)\) and both measures are finite. Reversing this equality gives the required identity for every \(h\) and \(E\). % lean: ProbabilityTheory.Kernel.setLIntegral_rnDeriv
\end{enumerate}
\end{proof}

\bibliographystyle{plainnat}

\bibliography{references}

\end{document}
